\documentclass[twoside,11pt]{article}

% ? Specify used packages
\usepackage{graphicx}        %  Use this one for final production.
% \usepackage[draft]{graphicx} %  Use this one for drafting.
% ? End of specify used packages

\usepackage{lscape}
\usepackage{amsmath}

\pagestyle{myheadings}

% -----------------------------------------------------------------------------
% ? Document identification
% Fixed part
\newcommand{\stardoccategory}  {Starlink System Note}
\newcommand{\stardocinitials}  {SSN}
\newcommand{\stardocsource}    {ssn\stardocnumber}

% Variable part - replace [xxx] as appropriate.
\newcommand{\stardocnumber}    {72.1}
\newcommand{\stardocauthors}   {Tim Jenness, John F.\ Lightfoot\\
                                Joint Astronomy Centre, Hilo, Hawaii}
\newcommand{\stardocdate}      {09 August 1999}
\newcommand{\stardoctitle}     {SURF Programming Interface}
\newcommand{\stardocabstract}  {This document describes the SURF programming
interface. Including the library API, instructions on how to build
a SURF task, the SURF file format and instructions on how to write
a SURF import task.}

% Environment for indenting and using a small font.
\newenvironment{myquote}{\begin{quote}\begin{small}}{\end{small}\end{quote}}

\newcommand{\mytext}[1]{{\small \tt #1}}

% set up some common package names

\newcommand{\Kappa}{\xref{{\sc{Kappa}}}{sun95}{}}

\newcommand{\Figaro}{\xref{{\sc{Figaro}}}{sun86}{}}

\newcommand{\gaia}{\xref{{\sc{Gaia}}}{sun214}{}}

\newcommand{\convert}{\xref{{\sc{Convert}}}{sun55}{}}

\newcommand{\fluxes}{\xref{{\sc{Fluxes}}}{sun213}{}}

\newcommand{\Iras}{\xref{{\sc{Iras90}}}{sun163}{}}

\newcommand{\ndf}{\xref{NDF}{sun33}{}}

\newcommand{\hds}{\xref{HDS}{sun92}{}}

\newcommand{\agi}{\xref{AGI}{sun48}{}}

\newcommand{\ESP}{\xref{ESP}{sun180}{}}

\newcommand{\specx}{\xref{{\sc{Specx}}}{sc8}{}}

\newcommand{\polpack}{\xref{{\sc{Polpack}}}{sun223}{}}

\newcommand{\ccdpack}{\xref{CCDPACK}{sun139}{}}

\newcommand{\cursa}{\xref{{\sc{Cursa}}}{sun190}{}}

\newcommand{\surfref}{\xref{{\scusoft}}{sun216}{}}

% Application tasks
\newcommand{\task}[1]{{\sf #1}}

% ADAM parameters
\newcommand{\param}[1]{{\tt #1}}


% Common tasks
\newcommand{\renoise}{\xref{\task{reduce\_noise}}{sun216}{REDUCE_NOISE}}

\newcommand{\rebin}{\xref{\task{rebin}}{sun216}{REBIN}}

\newcommand{\bolrebin}{\xref{\task{bolrebin}}{sun216}{BOLREBIN}}

\newcommand{\intrebin}{\xref{\task{intrebin}}{sun216}{INTREBIN}}

\newcommand{\calcsky}{\xref{\task{calcsky}}{sun216}{CALCSKY}}

\newcommand{\chgqual}{\xref{\task{change\_quality}}{sun216}{CHANGE_QUALITY}}

\newcommand{\chgflat}{\xref{\task{change\_flat}}{sun216}{CHANGE_FLAT}}

\newcommand{\chgpnt}{\xref{\task{change\_pointing}}{sun216}{CHANGE_POINTING}}

\newcommand{\chgdata}{\xref{\task{change\_data}}{sun216}{CHANGE_DATA}}

\newcommand{\resw}{\xref{\task{reduce\_switch}}{sun216}{REDUCE_SWITCH}}

\newcommand{\flatf}{\xref{\task{flatfield}}{sun216}{FLATFIELD}}

\newcommand{\skydip}{\xref{\task{skydip}}{sun216}{SKYDIP}}

\newcommand{\scuphot}{\xref{\task{scuphot}}{sun216}{SCUPHOT}}

\newcommand{\ext}{\xref{\task{extinction}}{sun216}{EXTINCTION}}

\newcommand{\scuquick}{\xref{\task{scuquick}}{sun216}{SCUQUICK}}

\newcommand{\scuhelp}{\xref{\task{scuhelp}}{sun216}{SCUHELP}}

\newcommand{\remsky}{\xref{\task{remsky}}{sun216}{REMSKY}}

\newcommand{\scuover}{\xref{\task{scuover}}{sun216}{SCUOVER}}

\newcommand{\extdata}{\xref{\task{extract\_data}}{sun216}{EXTRACT_DATA}}

\newcommand{\sculog}{\xref{\task{sculog}}{sun216}{SCULOG}}

\newcommand{\scucat}{\xref{\task{scucat}}{sun216}{SCUCAT}}

\newcommand{\photsum}{\xref{\task{photsum}}{sun216}{PHOTSUM}}

\newcommand{\mapsum}{\xref{\task{mapsum}}{sun216}{MAPSUM}}

\newcommand{\skysum}{\xref{\task{skysum}}{sun216}{SKYSUM}}

\newcommand{\pointsum}{\xref{\task{pointsum}}{sun216}{POINTSUM}}

\newcommand{\qdraw}{\xref{\task{qdraw}}{sun216}{QDRAW}}

\newcommand{\sigclip}{\xref{\task{sigclip}}{sun216}{SIGCLIP}}

\newcommand{\restore}{\xref{\task{restore}}{sun216}{RESTORE}}

\newcommand{\sdip}{\xref{\task{sdip}}{sun216}{SDIP}}

\newcommand{\scupa}{\xref{\task{scupa}}{sun216}{SCUPA}}

\newcommand{\obssum}{\xref{\task{obssum}}{sun216}{OBSSUM}}

\newcommand{\despike}{\xref{\task{despike}}{sun216}{DESPIKE}}

\newcommand{\despikeb}{\xref{\task{despike2}}{sun216}{DESPIKE2}}

\newcommand{\extflat}{\xref{\task{extract\_flat}}{sun216}{EXTRACT_FLAT}}

\newcommand{\scushift}{\xref{\task{scushift}}{sun216}{SCUSHIFT}}

\newcommand{\scuclip}{\xref{\task{scuclip}}{sun216}{SCUCLIP}}

\newcommand{\scanrlb}{\xref{\task{scan\_rlb}}{sun216}{SCAN_RLB}}

\newcommand{\dspbol}{\xref{\task{dspbol}}{sun216}{DSPBOL}}

\newcommand{\pltbol}{\xref{\task{pltbol}}{sun216}{PLTBOL}}

\newcommand{\rlinplot}{\xref{\task{rlinplot}}{sun216}{RLINPLOT}}

\newcommand{\scuplot}{\xref{\task{scuplot}}{sun216}{SCUPLOT}}

\newcommand{\remdbm}{\xref{\task{remdbm}}{sun216}{REMDBM}}

\newcommand{\scumakewt}{\xref{\task{scumakewt}}{sun216}{SCUMAKEWT}}

\newcommand{\chgnacent}{\xref{\task{change\_nacentre}}{sun216}{CHANGE_NACENTRE}}

\newcommand{\scunoise}{\xref{\task{scunoise}}{sun216}{SCUNOISE}}

\newcommand{\setbolwt}{\xref{\task{setbolwt}}{sun216}{SETBOLWT}}

\newcommand{\adddbm}{\xref{\task{add\_dbm}}{sun216}{ADD_DBM}}

\newcommand{\remip}{\xref{\task{remip}}{sun216}{REMIP}}

\newcommand{\scusetenv}{\xref{\task{scusetenv}}{sun216}{SCUSETENV}}


% Non surf tasks

\newcommand{\fillbad}{\xref{\task{fillbad}}{sun95}{FILLBAD}}

\newcommand{\ndftrace}{\xref{\task{ndftrace}}{sun95}{NDFTRACE}}

\newcommand{\glitch}{\xref{\task{glitch}}{sun95}{GLITCH}}

\newcommand{\display}{\xref{\task{display}}{sun95}{DISPLAY}}

\newcommand{\linplot}{\xref{\task{linplot}}{sun95}{LINPLOT}}

\newcommand{\mlinplot}{\xref{\task{mlinplot}}{sun95}{MLINPLOT}}

\newcommand{\drawsig}{\xref{\task{drawsig}}{sun95}{DRAWSIG}}

\newcommand{\centroid}{\xref{\task{centroid}}{sun95}{CENTROID}}

\newcommand{\setaxis}{\xref{\task{setaxis}}{sun95}{SETAXIS}}

\newcommand{\kstest}{\xref{\task{kstest}}{sun95}{KSTEST}}

\newcommand{\stats}{\xref{\task{stats}}{sun95}{STATS}}

\newcommand{\thresh}{\xref{\task{thresh}}{sun95}{THRESH}}

\newcommand{\setbb}{\xref{\task{setbb}}{sun95}{SETBB}}

\newcommand{\fitslist}{\xref{\task{fitslist}}{sun95}{FITSLIST}}

\newcommand{\fourier}{\xref{\task{fourier}}{sun95}{FOURIER}}

\newcommand{\hislist}{\xref{\task{hislist}}{sun95}{HISLIST}}

\newcommand{\ndffits}{\xref{\task{ndf2fits}}{sun55}{NDF2FITS}}

\newcommand{\cadd}{\xref{\task{cadd}}{sun95}{CADD}}

\newcommand{\hdstrace}{\xref{\task{hdstrace}}{sun102}{}}

\newcommand{\psmerge}{\xref{\task{psmerge}}{sun164}{}}

\newcommand{\wdfits}{\xref{\task{wdfits}}{sun86}{WDFITS}}

\newcommand{\delobj}{\xref{\task{delobj}}{sun86}{DELOBJ}}

\newcommand{\sclean}{\xref{\task{sclean}}{sun86}{SCLEAN}}

\newcommand{\setvar}{\xref{\task{setvar}}{sun95}{SETVAR}}

\newcommand{\wcsalign}{\xref{\task{wcsalign}}{sun95}{WCSALIGN}}

\newcommand{\ndfcopy}{\xref{\task{ndfcopy}}{sun95}{NDFCOPY}}

\newcommand{\manic}{\xref{\task{manic}}{sun95}{MANIC}}

\newcommand{\cdiv}{\xref{\task{cdiv}}{sun95}{CDIV}}

\newcommand{\sub}{\xref{\task{sub}}{sun95}{SUB}}

\newcommand{\cursor}{\xref{\task{cursor}}{sun95}{CURSOR}}

\newcommand{\gdset}{\xref{\task{gdset}}{sun95}{GDSET}}

\newcommand{\image}{\xref{\task{image}}{sun86}{IMAGE}}

\newcommand{\gaufit}{\xref{\task{gaufit}}{sun180}{GAUFIT}}

\newcommand{\makemos}{\xref{\task{makemos}}{sun139}{MAKEMOS}}

\newcommand{\polimp}{\xref{\task{polimp}}{sun223}{POLIMP}}

\newcommand{\polcal}{\xref{\task{polcal}}{sun223}{POLCAL}}

\newcommand{\polbin}{\xref{\task{polbin}}{sun223}{POLBIN}}

\newcommand{\polvec}{\xref{\task{polvec}}{sun223}{POLVEC}}

\newcommand{\polplot}{\xref{\task{polplot}}{sun223}{POLPLOT}}

\newcommand{\catselect}{\xref{\task{catselect}}{sun190}{}}

% General commands
\newcommand{\scusoft}          {{\sc Surf}}

\newcommand{\micron}           {$\mu$m}  


% ? End of document identification
% -----------------------------------------------------------------------------

% +
%  Name:
%     ssn72.tex
%
%  Purpose:
%     SURF (see SUN/216) programming interface documentation
%
%  Authors:
%     AJC: A.J.Chipperfield (Starlink, RAL)
%     BLY: M.J.Bly (Starlink, RAL)
%     PWD: Peter W. Draper (Starlink, Durham University)
%     TIMJ: Tim Jenness (JAC, Hawaii)
%
%  History:
%     $Log$
%     Revision 1.4  2000/05/11 19:43:47  timj
%     Include some minor mods made by Martin Bly for the Starlink release.
%
%     Revision 1.3  1999/11/06 02:18:33  timj
%     Fix use of amstex.
%
%     Revision 1.2  1999/11/05 23:12:56  timj
%     Release V1.5-0
%
%     Revision 1.1  1999/08/19 03:52:10  timj
%     First draft.
%     09-AUG-1999 (TIMJ):
%        Convert template to SURF API document. Write complete document.
%     13-AUG-1998 (PWD):
%        Converted for use with LaTeX2HTML version 98.2 and
%        Star2HTML version 1.3.
%     16-JUN-1997 (BLY):
%        Adapted for LaTeX2e.
%     17-JAN-1996 (AJC):
%        Original with hypertext macros, based on MDL plain originals.


%     {Add further history here}

%  Copyright:
%     Copyright (C) 1996-1998 Central Laboratory of the Research Councils
%     Copyright (C) 1999 Particle Physics and Astronomy Research Council

%
% -

\newcommand{\stardocname}{\stardocinitials /\stardocnumber}
\markboth{\stardocname}{\stardocname}
\setlength{\textwidth}{160mm}
\setlength{\textheight}{230mm}
\setlength{\topmargin}{-2mm}
\setlength{\oddsidemargin}{0mm}
\setlength{\evensidemargin}{0mm}
\setlength{\parindent}{0mm}
\setlength{\parskip}{\medskipamount}
\setlength{\unitlength}{1mm}

% -----------------------------------------------------------------------------
%  Hypertext definitions.
%  ======================
%  These are used by the LaTeX2HTML translator in conjunction with star2html.

%  Comment.sty: version 2.0, 19 June 1992
%  Selectively in/exclude pieces of text.
%
%  Author
%    Victor Eijkhout                                      <eijkhout@cs.utk.edu>
%    Department of Computer Science
%    University Tennessee at Knoxville
%    104 Ayres Hall
%    Knoxville, TN 37996
%    USA

%  Do not remove the %begin{latexonly} and %end{latexonly} lines (used by 
%  LaTeX2HTML to signify text it shouldn't process).
%begin{latexonly}
\makeatletter
\def\makeinnocent#1{\catcode`#1=12 }
\def\csarg#1#2{\expandafter#1\csname#2\endcsname}

\def\ThrowAwayComment#1{\begingroup
    \def\CurrentComment{#1}%
    \let\do\makeinnocent \dospecials
    \makeinnocent\^^L% and whatever other special cases
    \endlinechar`\^^M \catcode`\^^M=12 \xComment}
{\catcode`\^^M=12 \endlinechar=-1 %
 \gdef\xComment#1^^M{\def\test{#1}
      \csarg\ifx{PlainEnd\CurrentComment Test}\test
          \let\html@next\endgroup
      \else \csarg\ifx{LaLaEnd\CurrentComment Test}\test
            \edef\html@next{\endgroup\noexpand\end{\CurrentComment}}
      \else \let\html@next\xComment
      \fi \fi \html@next}
}
\makeatother

\def\includecomment
 #1{\expandafter\def\csname#1\endcsname{}%
    \expandafter\def\csname end#1\endcsname{}}
\def\excludecomment
 #1{\expandafter\def\csname#1\endcsname{\ThrowAwayComment{#1}}%
    {\escapechar=-1\relax
     \csarg\xdef{PlainEnd#1Test}{\string\\end#1}%
     \csarg\xdef{LaLaEnd#1Test}{\string\\end\string\{#1\string\}}%
    }}

%  Define environments that ignore their contents.
\excludecomment{comment}
\excludecomment{rawhtml}
\excludecomment{htmlonly}

%  Hypertext commands etc. This is a condensed version of the html.sty
%  file supplied with LaTeX2HTML by: Nikos Drakos <nikos@cbl.leeds.ac.uk> &
%  Jelle van Zeijl <jvzeijl@isou17.estec.esa.nl>. The LaTeX2HTML documentation
%  should be consulted about all commands (and the environments defined above)
%  except \xref and \xlabel which are Starlink specific.

\newcommand{\htmladdnormallinkfoot}[2]{#1\footnote{#2}}
\newcommand{\htmladdnormallink}[2]{#1}
\newcommand{\htmladdimg}[1]{}
\newcommand{\hyperref}[4]{#2\ref{#4}#3}
\newcommand{\htmlref}[2]{#1}
\newcommand{\htmlimage}[1]{}
\newcommand{\htmladdtonavigation}[1]{}

\newenvironment{latexonly}{}{}
\newcommand{\latex}[1]{#1}
\newcommand{\html}[1]{}
\newcommand{\latexhtml}[2]{#1}
\newcommand{\HTMLcode}[2][]{}

%  Starlink cross-references and labels.
\newcommand{\xref}[3]{#1}
\newcommand{\xlabel}[1]{}

%  LaTeX2HTML symbol.
\newcommand{\latextohtml}{\LaTeX2\texttt{HTML}}

%  Define command to re-centre underscore for Latex and leave as normal
%  for HTML (severe problems with \_ in tabbing environments and \_\_
%  generally otherwise).
\renewcommand{\_}{\texttt{\symbol{95}}}

% -----------------------------------------------------------------------------
%  Debugging.
%  =========
%  Remove % on the following to debug links in the HTML version using Latex.

% \newcommand{\hotlink}[2]{\fbox{\begin{tabular}[t]{@{}c@{}}#1\\\hline{\footnotesize #2}\end{tabular}}}
% \renewcommand{\htmladdnormallinkfoot}[2]{\hotlink{#1}{#2}}
% \renewcommand{\htmladdnormallink}[2]{\hotlink{#1}{#2}}
% \renewcommand{\hyperref}[4]{\hotlink{#1}{\S\ref{#4}}}
% \renewcommand{\htmlref}[2]{\hotlink{#1}{\S\ref{#2}}}
% \renewcommand{\xref}[3]{\hotlink{#1}{#2 -- #3}}
%end{latexonly}
% -----------------------------------------------------------------------------
% ? Document specific \newcommand or \newenvironment commands.

%+
%  Name:
%     SST.TEX

%  Purpose:
%     Define LaTeX commands for laying out Starlink routine descriptions.

%  Language:
%     LaTeX

%  Type of Module:
%     LaTeX data file.

%  Description:
%     This file defines LaTeX commands which allow routine documentation
%     produced by the SST application PROLAT to be processed by LaTeX and
%     by LaTeX2html. The contents of this file should be included in the
%     source prior to any statements that make of the sst commnds.

%  Notes:
%     The style file html.sty provided with LaTeX2html needs to be used.
%     This must be before this file.

%  Authors:
%     RFWS: R.F. Warren-Smith (STARLINK)
%     PDRAPER: P.W. Draper (Starlink - Durham University)

%  History:
%     10-SEP-1990 (RFWS):
%        Original version.
%     10-SEP-1990 (RFWS):
%        Added the implementation status section.
%     12-SEP-1990 (RFWS):
%        Added support for the usage section and adjusted various spacings.
%     8-DEC-1994 (PDRAPER):
%        Added support for simplified formatting using LaTeX2html.
%     {enter_further_changes_here}

%  Bugs:
%     {note_any_bugs_here}

%-

%  Define length variables.
\newlength{\sstbannerlength}
\newlength{\sstcaptionlength}
\newlength{\sstexampleslength}
\newlength{\sstexampleswidth}

%  Define a \tt font of the required size.
\latex{\newfont{\ssttt}{cmtt10 scaled 1095}}
\html{\newcommand{\ssttt}{\tt}}

%  Define a command to produce a routine header, including its name,
%  a purpose description and the rest of the routine's documentation.
\newcommand{\sstroutine}[3]{
   \goodbreak
   \rule{\textwidth}{0.5mm}
   \vspace{-7ex}
   \newline
   \settowidth{\sstbannerlength}{{\Large {\bf #1}}}
   \setlength{\sstcaptionlength}{\textwidth}
   \setlength{\sstexampleslength}{\textwidth}
   \addtolength{\sstbannerlength}{0.5em}
   \addtolength{\sstcaptionlength}{-1.0\sstbannerlength}
   \addtolength{\sstcaptionlength}{-5.0pt}
   \settowidth{\sstexampleswidth}{{\bf Examples:}}
   \addtolength{\sstexampleslength}{-\sstexampleswidth}
   \parbox[t]{\sstbannerlength}{\flushleft{\Large {\bf #1}}}
   \parbox[t]{\sstcaptionlength}{\center{\Large #2}}
   \begin{description}
      #3
   \end{description}
}

%  Format the description section.
\newcommand{\sstdescription}[1]{\item[Description:] #1}

%  Format the usage section.
\newcommand{\sstusage}[1]{\item[Usage:] \mbox{}
\\[1.3ex]{\raggedright \ssttt #1}}

%  Format the invocation section.
\newcommand{\sstinvocation}[1]{\item[Invocation:]\hspace{0.4em}{\tt #1}}

%  Format the arguments section.
\newcommand{\sstarguments}[1]{
   \item[Arguments:] \mbox{} \\
   \vspace{-3.5ex}
   \begin{description}
      #1
   \end{description}
}

%  Format the returned value section (for a function).
\newcommand{\sstreturnedvalue}[1]{
   \item[Returned Value:] \mbox{} \\
   \vspace{-3.5ex}
   \begin{description}
      #1
   \end{description}
}

%  Format the parameters section (for an application).
\newcommand{\sstparameters}[1]{
   \item[Parameters:] \mbox{} \\
   \vspace{-3.5ex}
   \begin{description}
      #1
   \end{description}
}

%  Format the examples section.
\newcommand{\sstexamples}[1]{
   \item[Examples:] \mbox{} \\
   \vspace{-3.5ex}
   \begin{description}
      #1
   \end{description}
}

%  Define the format of a subsection in a normal section.
\newcommand{\sstsubsection}[1]{ \item[{#1}] \mbox{} \\}

%  Define the format of a subsection in the examples section.
\newcommand{\sstexamplesubsection}[2]{\sloppy
\item[\parbox{\sstexampleslength}{\ssttt #1}] \mbox{} \vspace{1.0ex}
\\ #2 }

%  Format the notes section.
\newcommand{\sstnotes}[1]{\item[Notes:] \mbox{} \\[1.3ex] #1}

%  Provide a general-purpose format for additional (DIY) sections.
\newcommand{\sstdiytopic}[2]{\item[{\hspace{-0.35em}#1\hspace{-0.35em}:}]
\mbox{} \\[1.3ex] #2}

%  Format the implementation status section.
\newcommand{\sstimplementationstatus}[1]{
   \item[{Implementation Status:}] \mbox{} \\[1.3ex] #1}

%  Format the bugs section.
\newcommand{\sstbugs}[1]{\item[Bugs:] #1}

%  Format a list of items while in paragraph mode.
\newcommand{\sstitemlist}[1]{
  \mbox{} \\
  \vspace{-3.5ex}
  \begin{itemize}
     #1
  \end{itemize}
}

%  Define the format of an item.
\newcommand{\sstitem}{\item}

%% Now define html equivalents of those already set. These are used by
%  latex2html and are defined in the html.sty files.
\begin{htmlonly}

%  sstroutine.
   \newcommand{\sstroutine}[3]{
      \subsection{#1\xlabel{#1}-\label{#1}#2}
      \begin{description}
         #3
      \end{description}
   }

%  sstdescription
   \newcommand{\sstdescription}[1]{\item[Description:]
      \begin{description}
         #1
      \end{description}
      \\
   }

%  sstusage
   \newcommand{\sstusage}[1]{\item[Usage:]
      \begin{description}
         {\ssttt #1}
      \end{description}
      \\
   }

%  sstinvocation
   \newcommand{\sstinvocation}[1]{\item[Invocation:]
      \begin{description}
         {\ssttt #1}
      \end{description}
      \\
   }

%  sstarguments
   \newcommand{\sstarguments}[1]{
      \item[Arguments:] \\
      \begin{description}
         #1
      \end{description}
      \\
   }

%  sstreturnedvalue
   \newcommand{\sstreturnedvalue}[1]{
      \item[Returned Value:] \\
      \begin{description}
         #1
      \end{description}
      \\
   }

%  sstparameters
   \newcommand{\sstparameters}[1]{
      \item[Parameters:] \\
      \begin{description}
         #1
      \end{description}
      \\
   }

%  sstexamples
   \newcommand{\sstexamples}[1]{
      \item[Examples:] \\
      \begin{description}
         #1
      \end{description}
      \\
   }

%  sstsubsection
   \newcommand{\sstsubsection}[1]{\item[{#1}]}

%  sstexamplesubsection
   \newcommand{\sstexamplesubsection}[2]{\item[{\ssttt #1}] #2}

%  sstnotes
   \newcommand{\sstnotes}[1]{\item[Notes:] #1 }

%  sstdiytopic
   \newcommand{\sstdiytopic}[2]{\item[{#1}] #2 }

%  sstimplementationstatus
   \newcommand{\sstimplementationstatus}[1]{
      \item[Implementation Status:] #1
   }

%  sstitemlist
   \newcommand{\sstitemlist}[1]{
      \begin{itemize}
         #1
      \end{itemize}
      \\
   }
%  sstitem
   \newcommand{\sstitem}{\item}

\end{htmlonly}

%  End of "sst.tex" layout definitions.
%.





% ? End of document specific commands
% -----------------------------------------------------------------------------
%  Title Page.
%  ===========
\renewcommand{\thepage}{\roman{page}}
\begin{document}
\thispagestyle{empty}

%  Latex document header.
%  ======================
\begin{latexonly}
   CCLRC / \textsc{Rutherford Appleton Laboratory} \hfill \textbf{\stardocname}\\
   {\large Particle Physics \& Astronomy Research Council}\\
   {\large Starlink Project\\}
   {\large \stardoccategory\ \stardocnumber}
   \begin{flushright}
   \stardocauthors\\
   \stardocdate
   \end{flushright}
   \vspace{-4mm}
   \rule{\textwidth}{0.5mm}
   \vspace{5mm}
   \begin{center}
   {\Large\textbf{\stardoctitle}}
   \end{center}
   \vspace{5mm}

\begin{center}
\includegraphics[width=2.0in]{ssn72_logo.eps}
\end{center}

% ? Heading for abstract if used.
   \vspace{10mm}
   \begin{center}
      {\Large\textbf{Abstract}}
   \end{center}
% ? End of heading for abstract.
\end{latexonly}

%  HTML documentation header.
%  ==========================
\begin{htmlonly}
   \xlabel{}
   \begin{rawhtml} <H1 ALIGN=CENTER> \end{rawhtml}
      \stardoctitle
   \begin{rawhtml} </H1> <HR> \end{rawhtml}

\includegraphics[width=2.0in]{ssn72_logo.eps}

   \begin{rawhtml} <P> <I> \end{rawhtml}
   \stardoccategory\ \stardocnumber \\
   \stardocauthors \\
   \stardocdate
   \begin{rawhtml} </I> </P> <H3> \end{rawhtml}
      \htmladdnormallink{CCLRC}{http://www.cclrc.ac.uk} /
      \htmladdnormallink{Rutherford Appleton Laboratory}
                        {http://www.cclrc.ac.uk/ral} \\
      \htmladdnormallink{Particle Physics \& Astronomy Research Council}
                        {http://www.pparc.ac.uk} \\
   \begin{rawhtml} </H3> <H2> \end{rawhtml}
      \htmladdnormallink{Starlink Project}{http://www.starlink.rl.ac.uk/}
   \begin{rawhtml} </H2> \end{rawhtml}
   \htmladdnormallink{\htmladdimg{source.gif} Retrieve hardcopy}
      {http://www.starlink.rl.ac.uk/cgi-bin/hcserver?\stardocsource}\\

%  HTML document table of contents. 
%  ================================
%  Add table of contents header and a navigation button to return to this 
%  point in the document (this should always go before the abstract \section). 
  \label{stardoccontents}
  \begin{rawhtml} 
    <HR>
    <H2>Contents</H2>
  \end{rawhtml}
  \htmladdtonavigation{\htmlref{\htmladdimg{contents_motif.gif}}
        {stardoccontents}}

% ? New section for abstract if used.
  \section{\xlabel{abstract}Abstract}
% ? End of new section for abstract

\end{htmlonly}

% -----------------------------------------------------------------------------
% ? Document Abstract. (if used)
%  ==================
\stardocabstract
% ? End of document abstract
% -----------------------------------------------------------------------------
% ? Latex document Table of Contents (if used).
%  ===========================================
  \newpage
  \begin{latexonly}
    \setlength{\parskip}{0mm}
    \tableofcontents
    \setlength{\parskip}{\medskipamount}
    \markboth{\stardocname}{\stardocname}
  \end{latexonly}
% ? End of Latex document table of contents
% -----------------------------------------------------------------------------
\cleardoublepage
\renewcommand{\thepage}{\arabic{page}}
\setcounter{page}{1}

%%%%%%%%%%%%%%%% S T A R T   H E R E %%%%%%%%%%%%%%%%%%%%%%%%%%
% ? Main text

\section{Introduction}

The SCUBA User Reduction Facility (\surfref, SUN/33) provides data reduction 
software for submillimetre bolometer cameras. It was specifically designed for 
the Submillimetre Common-User Bolometer Array (SCUBA; Holland et al 1999,
MNRAS, 303, 659) but could be ported to similar bolometer cameras. This
document describes the programming interface used to write \scusoft\ tasks and 
import foreign data formats into the system.


\section{SURF File Format}

\scusoft\ uses that standard Starlink \ndf\ (SUN/33) file format. This 
is a hierarchical format based on \hds\ (SUN/92) and has native support for 
variances, quality flags, non-linear axes and arbritrary extensions.

\subsection{Raw data format\label{sec:raw}}

SCUBA raw data (i.e.\ data in the format expected by the \scusoft\ task
\resw) has the following general layout:

\begin{myquote}
\begin{verbatim}
STRUCT  <NDF>

   MORE           <EXT>           {structure}
      FIGARO         <FIGARO_EXT>    {structure}

      SCUCD          <SCUCD_ST>      {structure}
         
      SCUBA          <SCUBA_ST>      {structure}

      FITS(171)      <_CHAR*80>      'ACCEPT  = 'not used'           / ac...'
                                     ... 'WAVE_5  =                 ...','END'

   DATA_ARRAY(5,130,900)  <_REAL>   0.0102353,1.9114159E-6,0.0184784,
                                    ... 1.7097691E-5,0.01525085,1.841503E-5,0
\end{verbatim}
\end{myquote}
It is a very simple NDF format with a single \texttt{DATA\_ARRAY} and
four extensions. The extensions are described in section \ref{scu:ext}. 
The \texttt{DATA\_ARRAY} is three dimensional with the following meanings:

\begin{enumerate}
\item The size of the first dimension depends on the observing mode. For
jiggle and noise observations it has size 5 corresponding to the demodulated
data, variance on demodulated data, internal calibrator signal, variance on
calibrator signal and data quality. For SCAN/MAP observations it has size 4
corresponding to the demodulated data, variance on demodulated data, internal
calibrator signal and data quality (i.e.\ the calibrator variance is not
stored for scan maps). The data quality is converted from a REAL number to a
standard UBYTE quality (appendix \ref{sec:quality}): if the quality is odd,
this indicates the flatfield was set to zero so bit 1 is turned on (this is
also done using the \flatf\ task), if the quality is greater than the value
specified via the \param{SPIKE\_LEVEL} parameter (i.e.\ there were more spikes
in the individual sample than the specified threshold) bit 2 is turned on. For
SKYDIP observations it has a size of SCUBA\_\_N\_TEMPS (see Appendix
\ref{app:const}, usually 3) corresponding to the different temperatures
measured (AMBIENT, SKY and COLD).

\item The second dimension corresponds to the number of bolometers
observed. This should match the size of the \texttt{BOL\_CHAN} and
\texttt{BOL\_ADC} arrays in the \texttt{SCUBA} extension.

\item The third dimension is the time axis and is equal to the total number of 
samples (of length \texttt{EXP\_TIME}). This dimension must be greater than
any value stored in the \texttt{DEM\_PNTR} array. For a JIGGLE observation
the size can be calculated (assuming the observation was not aborted early)

\end{enumerate}

\subsection{SURF Format\label{sec:surfform}}

The \resw\ task converts the raw SCUBA format into a more NDF-like format
with proper variance and quality arrays (extracted by \resw\ from the SCUBA
data array), axes, history reporting, a title and units.

Here is an example of the \resw\ output for a Mars pointing observation:

\begin{myquote}
\begin{verbatim}
O10  <NDF>

   DATA_ARRAY(37,32)  <_REAL>     -0.003304771,0.0002364833,0.00320935,
                                  ... 0.0003040197,0.001364953,-0.007496601
   MORE           <EXT>           {structure}
      FIGARO         <FIGARO_EXT>    {structure}
         {structure is empty}

      SCUCD          <SCUCD_ST>      {structure}
         JIGL_X(16)     <_REAL>         5.3766,10.6914,10.6914,5.3766,0,
                                        ... -10.6914,-5.3766,0,5.3766,-5.3766
         JIGL_Y(16)     <_REAL>         -3.09,0,6.18,9.27,12.36,9.27,6.18,
                                        ... 6.18,0,-3.09,-6.18,-9.27,-9.27
         DEC1           <_REAL>         0
         DEC2           <_REAL>         0
         LST_STRT(2,1,2,1)  <_DOUBLE>   3.9425429572598,3.9440140603798,
                                        3.9469952223685,3.9454089275179
         RA1            <_REAL>         0
         RA2            <_REAL>         0

      SCUBA          <SCUBA_ST>      {structure}
         BOL_CALB(16,9)  <_REAL>        1.286,0.9311,0.9499,1.002,0.9585,
                                        ... 1.054,1.065,1.045,1,1,1,1,1,1,0
         BOL_DU3(16,9)  <_REAL>         57.08,45.9,36.15,24.64,14.55,3.62,
                                        ... -15.78,9.047,-54.65,0,0,0,0,0,0
         BOL_DU4(16,9)  <_REAL>         -28.14,-35.91,-41.46,-49.62,
                                        ... 63.89,77.76,71.16,0,0,0,0,0,0
         BOL_QUAL(16,9)  <_INTEGER>     0,0,0,0,0,0,0,0,0,0,0,0,0,0,0,0,0,
                                        ... 0,0,0,0,0,0,0,0,0,0,1,1,0,0,0,1,1
         BOL_TYPE(16,9)  <_CHAR*20>     'SHORT','SHORT','SHORT','SHORT',
                                        ... 'P...','P1350_DC','P1100_DC','BAD'
         BOL_ADC(37)    <_INTEGER>      7,7,7,7,7,7,7,7,7,7,7,7,7,8,8,8,8,
                                        ... 8,8,8,8,8,8,8,8,9,9,9,9,9,9,9,9,9
         BOL_CHAN(37)   <_INTEGER>      1,2,3,4,7,8,9,10,11,13,14,15,16,1,
                                        ... 12,13,14,15,16,1,2,3,4,5,6,7,8,9
         FLAT_ADC       <_INTEGER>      0
         FLAT_CHN       <_INTEGER>      0
         FLAT_IND       <_INTEGER>      0
         PHOT_BB        <_INTEGER>      0
         ISTART         <_INTEGER>      0
         NPIX           <_INTEGER>      0
         POINTER        <_INTEGER>      0
         DEM_PNTR(1,2,1)  <_INTEGER>    1,17

      FITS(171)      <_CHAR*80>      'ACCEPT  = 'PROMPT  '           / ac...'
                                     ... 'WAVE_5  =                 ...','END'
      REDS           <SURF_EXTENSION>   {structure}
         BEAM_WT(3)     <_REAL>         0,1,0

   HISTORY        <HISTORY>       {structure}
      CREATED        <_CHAR*24>      '1999-AUG-11 17:04:57.000'
      CURRENT_RECORD  <_INTEGER>     1
      RECORDS(10)    <HIST_REC>      {array of structures}

      Contents of RECORDS(1)
         DATE           <_CHAR*24>      '1999-AUG-11 17:04:57.000'
         COMMAND        <_CHAR*29>      'REDUCE_SWITCH   (SURF V1.4-3)'
         USER           <_CHAR*4>       'timj'
         HOST           <_CHAR*6>       'lapaki'
         DATASET        <_CHAR*33>      '/local/jcmt/scratch/timj/orac/o10'
         TEXT(3)        <_CHAR*72>      'Parameters: IN=@10 MSG_FILTER='N...'
                                        ... 'Software: /star/bin/surf/redu...'

   QUALITY        <QUALITY>       {structure}
      QUALITY(37,32)  <_UBYTE>       0,0,0,0,0,0,0,0,0,0,0,0,0,0,0,0,0,0,
                                     ... 0,0,0,0,0,0,0,0,0,0,0,0,0,0,0,0,0,0,0
      BADBITS        <_UBYTE>        7

   VARIANCE(37,32)  <_REAL>       1.0950339E-6,7.1230471E-7,6.6492021E-7,
                                  ... 4.2760187E-9,7.4043317E-5,4.8094225E-6
   AXIS(2)        <AXIS>          {array of structures}

   Contents of AXIS(1)
      DATA_ARRAY(37)  <_REAL>        1,2,3,4,5,6,7,8,9,10,11,12,13,14,15,
                                     ... 26,27,28,29,30,31,32,33,34,35,36,37
      LABEL          <_CHAR*9>       'Bolometer'

   LABEL          <_CHAR*8>       'Raw data'
   TITLE          <_CHAR*30>      'mars'
   UNITS          <_CHAR*5>       'Volts'

End of Trace.
\end{verbatim}
\end{myquote}

The data arrays are now 2 dimensional, corresponding to the bolometer number
and time axis. The axis information is written so that the integration number
is reflected on the Y axis. Additionally, a new extension, \texttt{REDS}, is
created to hold the relative beam weights; for maps these are always set to 0,1,0.

A photometry observation is slightly different from a map observation in that
the output data array is now 3 dimensional rather than 2:

\begin{myquote}
\begin{verbatim}
O107  <NDF>

   MORE           <EXT>           {structure}
      FIGARO         <FIGARO_EXT>    {structure}
         {structure is empty}

      SCUCD          <SCUCD_ST>      {structure}
         JIGL_X(9)      <_REAL>         0,2,2,0,-2,-2,-2,0,2
         JIGL_Y(9)      <_REAL>         0,0,2,2,2,0,-2,-2,-2
         DEC1           <_REAL>         0
         DEC2           <_REAL>         0
         LST_STRT(2,1,8,1)  <_DOUBLE>   1.8834042737017,1.8844116777459,
                                        ... 1.8981978240679,1.8971856029149
         RA1            <_REAL>         0
         RA2            <_REAL>         0

      SCUBA          <SCUBA_ST>      {structure}
         BOL_CALB(16,9)  <_REAL>        1.286,0.9311,0.9499,1.002,0.9585,
                                        ... 1.054,1.065,1.045,1,1,1,1,1,1,0
         BOL_DU3(16,9)  <_REAL>         57.08,45.9,36.15,24.64,14.55,3.62,
                                        ... -15.78,9.047,-54.65,0,0,0,0,0,0
         BOL_DU4(16,9)  <_REAL>         -28.14,-35.91,-41.46,-49.62,
                                        ... 63.89,77.76,71.16,0,0,0,0,0,0
         BOL_QUAL(16,9)  <_INTEGER>     0,0,0,0,0,0,0,0,0,0,0,0,0,0,0,0,0,
                                        ... 0,0,0,0,0,0,0,0,0,0,1,1,0,0,0,1,1
         BOL_TYPE(16,9)  <_CHAR*20>     'SHORT','SHORT','SHORT','SHORT',
                                        ... 'P...','P1350_DC','P1100_DC','BAD'
         BOL_ADC(130)   <_INTEGER>      1,1,1,1,1,1,1,1,1,1,1,1,1,1,1,1,2,
                                        ... 8,8,8,8,8,8,8,8,9,9,9,9,9,9,9,9,9
         BOL_CHAN(130)  <_INTEGER>      1,2,3,4,5,6,7,8,9,10,11,12,13,14,
                                        ... 12,13,14,15,16,1,2,3,4,5,6,7,8,9
         FLAT_ADC       <_INTEGER>      0
         FLAT_CHN       <_INTEGER>      0
         FLAT_IND       <_INTEGER>      0
         PHOT_BB(3,2)   <_INTEGER>      0,112,0,0,46,0
         ISTART         <_INTEGER>      0
         NPIX           <_INTEGER>      0
         POINTER        <_INTEGER>      0
         DEM_PNTR(1,8,1)  <_INTEGER>    1,10,19,28,37,46,55,64

      FITS(171)      <_CHAR*80>      'ACCEPT  = 'not used'           / ac...'
                                     ... 'WAVE_5  =                 ...','END'
      REDS           <SURF_EXTENSION>   {structure}
         BEAM_WT(3)     <_REAL>         0.5,1,0.5

   HISTORY        <HISTORY>       {structure}
      CREATED        <_CHAR*24>      '1999-AUG-11 17:08:24.000'
      CURRENT_RECORD  <_INTEGER>     1
      RECORDS(10)    <HIST_REC>      {array of structures}

      Contents of RECORDS(1)
         DATE           <_CHAR*24>      '1999-AUG-11 17:08:24.000'
         COMMAND        <_CHAR*29>      'REDUCE_SWITCH   (SURF V1.4-3)'
         USER           <_CHAR*4>       'timj'
         HOST           <_CHAR*6>       'lapaki'
         DATASET        <_CHAR*34>      '/local/jcmt/scratch/timj/orac/o107'
         TEXT(3)        <_CHAR*72>      'Parameters: IN=@107 MSG_FILTER='...'
                                        ... 'Software: /star/bin/surf/redu...'

   DATA_ARRAY(130,72,3)  <_REAL>   -0.0002527004,-0.0006291242,
                                   ... 1.1537035E-5,0.001769435,0.001728294
   QUALITY        <QUALITY>       {structure}
      QUALITY(130,72,3)  <_UBYTE>    0,0,0,0,0,0,0,0,0,0,0,0,0,0,0,0,0,0,
                                     ... 0,0,0,0,0,0,0,0,0,0,0,0,0,0,0,0,0,0,0
      BADBITS        <_UBYTE>        7

   VARIANCE(130,72,3)  <_REAL>    3.6919182E-6,3.6207405E-6,4.2371025E-6,
                                  ... 1.3244949E-8,0.0002660675,2.1535034E-5
   AXIS(3)        <AXIS>          {array of structures}

   Contents of AXIS(1)
      DATA_ARRAY(130)  <_REAL>       1,2,3,4,5,6,7,8,9,10,11,12,13,14,15,
                                     ... 122,123,124,125,126,127,128,129,130
      LABEL          <_CHAR*9>       'Bolometer'

   LABEL          <_CHAR*8>       'Raw data'
   TITLE          <_CHAR*30>      'hltau'
   UNITS          <_CHAR*5>       'Volts'

End of Trace.
\end{verbatim}
\end{myquote}

The 3rd dimension corresponds to beam position. In most cases only beam 2 is
used (the middle beam) but for two or three  bolometer chopping the 1st and
3rd beams will be accessed. In general, the 1st and 3rd beams are simply
scaled versions of the middle beam; the relative weighting of the beams is
stored in the newly created \texttt{BEAM\_WT} array in the \texttt{REDS}
extension. In the above example, the middle beam is -0.5 times the left or
right beams.



\subsection{Extensions\label{scu:ext}}

The information describing the observation can be found in the NDF
extensions. \scusoft\ expects to find the following extensions in
data files:

\subsubsection{FITS}

The FITS extension is an array of 80 character strings stored in a format
identical to a standard FITS header. This contains general information
about the observation but no array information. Listings of full SCUBA
FITS headers can be found in Appendix \ref{app:fits}. \scusoft\ only requires
a subset of the full SCUBA headers. For a jiggle map these are listed
in table \ref{tab:jigfits}.

\begin{table}
\begin{myquote}
\begin{tabular}{ll}
OBJECT & Name of object \\
RUN    & Run number of observation \\
MODE   & Type of observation (MAP, PHOTOM, SKYDIP etc) \\
SAM\_MODE & Sampling method (JIGGLE or RASTER) \\
SAM\_CRDS & Coordinate frame of jiggle or scan\\
SAM\_PA & Scan PA (not required for jiggle)\\
LAT    & Latitude of object (eg Declination)\\
LONG   & Longitude of object (eg Right Ascension)\\
CENT\_CRD & Coordinate frame of LONG/LAT (RB,RJ,GA etc)\\
UTDATE & Date of observation in YYYY:MM:DD format\\
UTSTART& UT time of start of observation (HH:MM:SS) \\
STSTART& Local sidereal time of start of observation (HH:MM:SS)\\
STEND  & LST of end of observation\\
EXP\_TIME&Exposure time (seconds) of each individual sample\\
MAP\_X  & Tangent plane X offset from tracking centre\\
MAP\_Y &  Tangent plane Y offset from tracking centre\\
N\_BOLS & Number of bolometers selected (should match data array)\\
STATE & State at end of observation (ABORT indicates early abort)\\
VERSION & Version of real-time software (see SURFLIB\_PROCESS\_BOLS)\\
JIGL\_CNT & Number of offsets in jiggle pattern\\
J\_PER\_S & Number of jiggles per switch\\
J\_REPEAT & No. of jiggle pattern repeats in switch\\
CNTR\_DU3 & Nasmyth  dU3 coord of instrument centre\\
CNTR\_DU4 & Nasmyth dU4 coord of instrument centre\\
LAT-OBS & Latitude of observatory (degrees)\\
LONG-OBS & East longitude of observatory (degrees)\\
TELESCOP & Telescope name\\
INSTRUME & Instrument name\\
CHOP\_THR & Chopper throw (arcsec)\\
CHOP\_PA & Chopper PA 0=in lat, 90=in long\\
CHOP\_FRQ & Chopper frequency (Hz)\\
CHOP\_CRD & Coordinate frame of chop\\
CHOP\_FUN & Chopper waveform (SQUARE, TRIPOS)\\
N\_SUBS   & Number of sub-instruments \\
SUB\_1    & Name of sub-instrument 1 (or 'not used')\\
SUB\_2    & Name of sub-instrument 2 (or 'not used')\\
SUB\_3    & Name of sub-instrument 3 (or 'not used')\\
SUB\_4    & Name of sub-instrument 4 (or 'not used')\\
SUB\_5    & Name of sub-instrument 5 (or 'not used')\\
FILT\_1   & Filter name for sub-instrument 1 (or 'not used')\\
FILT\_2   & Filter name for sub-instrument 2 (or 'not used')\\
FILT\_3   & Filter name for sub-instrument 3 (or 'not used')\\
FILT\_4   & Filter name for sub-instrument 4 (or 'not used')\\
FILT\_5   & Filter name for sub-instrument 5 (or 'not used')\\
WAVE\_1   & Wavelength of sub-instrument 1 (microns)\\
WAVE\_2   & Wavelength of sub-instrument 2 (microns)\\
WAVE\_3   & Wavelength of sub-instrument 3 (microns)\\
WAVE\_4   & Wavelength of sub-instrument 4 (microns)\\
WAVE\_5   & Wavelength of sub-instrument 5 (microns)\\
\end{tabular}
\end{myquote}
\caption{FITS keywords required by \scusoft\ for Jiggle mapping}
\label{tab:jigfits}
\end{table}

\subsubsection{SCUBA}

The \texttt{SCUBA} extension contains information on the instrument and the
structure of the data array. It has the following format (output of
\hdstrace):

\begin{myquote}
\begin{verbatim}
SCUBA          <SCUBA_ST>      {structure}
   BOL_CALB(16,9)  <_REAL>        1.286,0.9311,0.9499,1.002,0.9585,
                                  ... 1.054,1.065,1.045,1,1,1,1,1,1,0
   BOL_DU3(16,9)  <_REAL>         57.08,45.9,36.15,24.64,14.55,3.62,
                                  ... -15.78,9.047,-54.65,0,0,0,0,0,0
   BOL_DU4(16,9)  <_REAL>         -28.14,-35.91,-41.46,-49.62,
                                  ... 63.89,77.76,71.16,0,0,0,0,0,0
   BOL_QUAL(16,9)  <_INTEGER>     0,0,0,0,0,0,0,0,0,0,0,0,0,0,0,0,0,
                                  ... 0,0,0,0,0,0,0,0,0,0,1,1,0,0,0,1,1
   BOL_TYPE(16,9)  <_CHAR*20>     'SHORT','SHORT','SHORT','SHORT',
                                  ... 'P...','P1350_DC','P1100_DC','BAD'
   BOL_ADC(128)   <_INTEGER>      1,1,1,1,1,1,1,1,1,1,1,1,1,1,1,1,2,
                                  ... 8,8,8,8,8,8,8,8,9,9,9,9,9,9,9,9,9
   BOL_CHAN(128)  <_INTEGER>      1,2,3,4,5,6,7,8,9,10,11,12,13,14,
                                  ... 12,13,14,15,16,1,2,3,4,5,6,7,8,9
   FLAT_ADC       <_INTEGER>      0
   FLAT_CHN       <_INTEGER>      0
   FLAT_IND       <_INTEGER>      0
   PHOT_BB(3,2)   <_INTEGER>      0,112,0,0,46,0
   DEM_PNTR(2,4,3,1)  <_INTEGER>   1,17,49,33,65,81,113,97,129,145,
                                   ... 257,273,305,289,321,337,369,353
   ISTART         <_INTEGER>      0
   NPIX           <_INTEGER>      0
   POINTER        <_INTEGER>      0
\end{verbatim}
\end{myquote}

The first 5 entries are simply the flatfield information (see Appendix
\ref{app:ff} for an example flatfield file). The shape of the arrays
(in this case $16\times9$) corresponds to the number of channels on each A/D
card and the number of A/D cards. These dimensions must match the values of
the constants SCUBA\_\_NUM\_ADC and SCUBA\_\_NUM\_CHAN (Appendix
\ref{app:const}).

The full description of these entries follows:

\begin{description}
\item[BOL\_CALB]\mbox{} \\
These are the actual flatfield values and contain the response of each pixel
relative to a reference pixel (usually the centre pixel on each array, H7 and
C14). This correction is applied to the data using the \flatf\ task
(specifically the \texttt{sculib\_flatfield\_data} routine):
\begin{equation}
V_\mathrm{flatfielded}(B) = V_\mathrm{original}(B) \times \mathrm{BOL\_CALB}(B)
\end{equation}
where $B$ is the selected bolometer.

\item[BOL\_DU3]\mbox{} \\
These are the X coordinates of every bolometer. For SCUBA, these are arcsec X
(U3) offsets (i.e.\ $\mathrm{d}U3$) from the left-hand Nasmyth centre.

\item[BOL\_DU4]\mbox{} \\
These are the Y coordinate of every bolometer. For SCUBA, these are arcsec Y
(U4) offsets (i.e.\ $\mathrm{d}U4$) from the left-hand Nasmyth centre.

\item[BOL\_QUAL]\mbox{} \\
This is an array of quality flags. A 0 indicates that the bolometer is `good', 
a 1 indicates that the bolometer is `bad'. This array is used during
flatfielding (\texttt{sculib\_flatfield\_data})to set bit 1 in the NDF
quality array.

\item[BOL\_TYPE]\mbox{}\\
This array contains a textual description of the sub-instrument to which each
bolometer belongs. This information is used by the \ext\ task to select the
bolometers from the specified sub-instrument.

\item[BOL\_ADC]\mbox{}\\
This array contains the A/D card numbers corresponding to each of the
bolometers that were used for the observation. The order and size should match
the size of the second dimension of the main data array for raw SCUBA data and 
the first dimension for data that has been processed by \resw. The \ext\ task
modifies this array to reflect the sub-instrument selection.

\item[BOL\_CHAN]\mbox{}\\
This array contains the A/D channel numbers corresponding to each of the
bolometers that were used for the observation. The order and size should match
the size of the second dimension of the main data array for raw SCUBA data and 
the first dimension for data that has been processed by \resw. The \ext\ task
modifies this array to reflect the sub-instrument selection.

\item[FLAT\_ADC]\mbox{}\\
Not used by \scusoft.

\item[FLAT\_CHN]\mbox{}\\
Not used by \scusoft.

\item[FLAT\_IND]\mbox{}\\
Not used by \scusoft.

\item[PHOT\_BB]\mbox{}\\
This is only used for PHOTOM (and POLPHOT) observations. The first dimension
corresponds to the number of beam positions (3 for 2 position chopping with
nodding) and the second dimension corresponds to the number of sub-instruments 
used for the observation. The array contains the bolometer number (using the
same indexing scheme as for the BOL\_CHAN and BOL\_ADC arrays) that is on
source for each of the beam positions. For standard single bolometer
photometry only the middle beam is on source and in the example shown only
H7 and C14 are used for the photometry. The \scuphot\ task uses this array to
decide which bolometers to extract from the full dataset. The \ext\ task
modifies this array to reflect the sub-instrument selection.

This information is also available in string form in the \texttt{PHOT\_BBF}
FITS keyword. For non-PHOTOM based observations this structure is 0.


\item[DEM\_PNTR]\mbox{}\\

This array is used to convert from measurement, integration, exposure and
switch number to sample number (i.e.\ the final dimension in the
data array). The routines \texttt{SCULIB\_FIND\_SWITCH} and
\texttt{SCULIB\_FIND\_INT} are provided to simplify this (they also
return the position of the end of the switch or integration).
It has dimensions of N\_SWITCHES, N\_EXPOSURES, N\_INTEGRATIONS,
and N\_MEASUREMENTS.

\item[ISTART]\mbox{}\\
Not used by \scusoft.

\item[NPIX]\mbox{}\\
Not used by \scusoft.

\item[POINTER]\mbox{}\\
Not used by \scusoft.


\end{description}

\subsubsection{SCUCD}

The \texttt{SCUCD} extension contains information on telescope movements.
It is slightly different for JIGGLE (MAP and PHOTOM) and SCAN observations.
For a JIGGLE observation it has the following format:

\begin{myquote}
\begin{verbatim}
STRUCT.MORE.SCUCD  <SCUCD_ST>

   JIGL_X(64)     <_REAL>         -2.6883,-5.3457,-8.033999,-13.3797,
                                  ... 5.3457,2.6883,0,2.6883,0,2.6883,0,0
   JIGL_Y(64)     <_REAL>         1.545,3.09,1.545,1.545,3.09,4.635,6.18,
                                  ... -9.27,-7.725,-6.18,-4.635,-3.09,0
   DEC1           <_REAL>         0
   DEC2           <_REAL>         0
   LST_STRT(2,4,3,1)  <_DOUBLE>   5.6288711444511,5.630383297715,
                                  ... 5.6639924748021,5.6623806283313
   RA1            <_REAL>         0
   RA2            <_REAL>         0
\end{verbatim}
\end{myquote}

and for a SCAN observation:

\begin{myquote}
\begin{verbatim}
STRUCT.MORE.SCUCD  <SCUCD_ST>

   JIGL_X         <_REAL>         0
   JIGL_Y         <_REAL>         0
   DEC1(1,5,1,1)  <_REAL>         -17.06608,-17.04154,-17.10884,-17.07168,
                                  -17.11085
   DEC2(1,5,1,1)  <_REAL>         -17.03781,-17.09083,-17.04711,-17.11225,
                                  -17.09126
   LST_STRT(1,5,1,1)  <_DOUBLE>   5.521871430823,5.5231108938446,
                                  ... 5.5266851885263,5.5281428875176
   RA1(1,5,1,1)   <_REAL>         21.16225,21.16618,21.16226,21.1673,
                                  21.16577
   RA2(1,5,1,1)   <_REAL>         21.16466,21.16197,21.16755,21.1638,
                                  21.16747
   WPLATE         <_REAL>         0
\end{verbatim}
\end{myquote}

The full description of these entries follows:

\begin{description}

\item[JIGL\_X and JIGL\_Y]\mbox{} \\

These are the X and Y coordinates of the jiggle pattern in arcsecond
offsets. The coordinate frame is given in the \texttt{SAM\_CRDS} FITS keyword.
They are only used for JIGGLE observations.

\item[LST\_STRT]\mbox{} \\

This is the local sidereal time of the start of every switch expressed as
decimal radians. It has dimensions of N\_SWITCHES, N\_EXPOSURES,
N\_INTEGRATIONS, and N\_MEASUREMENTS. The current epoch of observation
(Modified Julian Date) is calculated from the UTSTART and UTDATE FITS
keywords in \texttt{SCULIB\_GET\_MJD}.

\item[RA1, DEC1, RA2 and DEC2]\mbox{} \\

This is the apparent RA and Dec of the start (RA1, DEC1) and end (RA2, DEC2)
of the scans in radians. The dimensions are N\_SWITCHES, N\_EXPOSURES,
N\_INTEGRATIONS, and N\_MEASUREMENTS. The current epoch of observation
(Modified Julian Date) is calculated from the UTSTART and UTDATE FITS
keywords in \texttt{SCULIB\_GET\_MJD}.

\item[WPLATE]\mbox{} \\

For polarimetry observations (POLPHOT or POLPMAP) this array contains the 
waveplate angle for each measurement (where a measurement is defined as a move 
of the waveplate). It is used by the \texttt{SURFLIB\_FILL\_WPLATE} routine.

\end{description}

\subsubsection{REDS}

The \texttt{REDS}\footnote{The name \texttt{REDS} is historical, at some time
in the distant past the SCUBA data reduction software was called
\textsc{Reds}, standing for `REDuction Scuba'} extension is created by SURF
as a general purpose repository of information that should be passed between
SURF tasks.

It can contain the following entries:

\begin{description}
\item[BEAM\_WT]\mbox{} \\
The relative weight of the beams. Only used for PHOTOM observations. This is 
created by \resw.
\item[SKY]\mbox{} \\
An NDF containing the estimate of the sky fluctuations. This is usually
created by \calcsky\ and is used by \remsky\ to remove the sky fluctuations.
The size of the NDF should match the size of the time axis in the main data
array.
\item[BOLWT]\mbox{} \\
Relative weight of each bolometer. Created by \setbolwt\ and used by \rebin.
\end{description}


\section{Aborted data files}

When an observation completes normally (either because the specified number of
measurements has been taken or a \texttt{TERMINATE} command has been issued)
the size of the final dimension of the data array and the contents and shape
of the \texttt{DEM\_PNTR} array agree.  In some cases an observation must be
aborted suddenly using the \texttt{ABORT} command. 
Aborted observation files do not always have the correct shape for 
\texttt{DEM\_PNTR} (e.g.\ the number of integrations specified by the second
dimension of the array does not agree with the number that were actually
observed). \scusoft\ overcomes this problem by looking in the \texttt{STATE}
FITS keyword; if it contains the string ``ABORTING'' the exposure, integration 
and measurement values are read from the FITS keywords \texttt{LAST\_EXP}, 
\texttt{LAST\_INT} and \texttt{LAST\_MEAS} respectively.



\section{\xlabel{quality}Quality flags\label{sec:quality}}

The \scusoft\ software conforms to the NDF standard concerning the processing
of quality or bad-pixel masks. Each method of setting a pixel bad is
associated with a bit in the quality masking flag (the NDF.QUALITY.BADBITS
component). The bad bits and their meaning in \scusoft\ are described in table
\ref{badbits}.

\begin{table}
\caption{Quality bits used by \scusoft}
\label{badbits}
\begin{center}
\begin{tabular}{ccl}
\hline\hline
Bit\# & Value & Meaning \\ \hline
0 & 1  &Infinity (eg division by zero) \\
1 & 2  &Set by \flatf.\\
2 & 4  &Set by \resw\ if the transputers detected more spikes\\
  &    & than specified by the \param{SPIKE\_LEVEL} parameter.\\
3 & 8  &Set by \chgqual.\\
4 & 16 &Set by despiking (\scuclip\ and \despike)\\ \hline\hline
\end{tabular}
\end{center}
\end{table}

In order to remove the effect of a particular bit (i.e.\ to ignore a
despiking), the \Kappa\ task \setbb\ can be used to change the bad-bits
mask in the NDF. Simply calculate the value related to the bits you are
interested in keeping and use this value in \setbb. Note that care must
be taken in deciding which bits are to be used for masking bad data. Bits zero
and one must always be set whereas the other three bits are
optional. \chgqual\ is the only task that acts on a file rather than producing
a processed copy and it is probably better if \chgqual\ is used directly if
you wish to manipulate the mask associated with bit three.

Additionally, regridded images also use quality flags. Bit 0 is used
to represent areas where no data were available and bit 1 is used to mask
data at the edge of the regridded area via the TRIM parameter.

\section{SURF Libraries}

\scusoft\ consists of the top level task routines (i.e.\ routines that
provide the structure of the task and do all disk I/O and most of the
interaction with the parameter system) and lower level
routines that do jobs that could be used by multiple tasks or other systems
(such as the real-time observing system) or provide the data reduction
algorithms. These lower level routines are stored in libraries and are
discussed in the following sections.

\subsection{SURFLIB}

\texttt{SURFLIB} is a library of data reduction algorithms and routines that
are directly relevant to \scusoft\ itself. All new subroutines required by
\scusoft\ are added to this library unless there is an obvious reason to put
it somewhere else. Contains routines for sky noise removal, instrumental
polarisation removal, baseline removal from scan maps etc).


\subsection{SCULIB}

This library was designed to provide the fundamental building blocks for
manipulating SCUBA data in the off-line and real-time system. Many of the
routines in this library are not used in \scusoft\ itself. In general,
new routines are added to \texttt{SURFLIB} rather than \texttt{SCULIB}
although \texttt{SCULIB} does contain some routines that are used solely in 
\scusoft.\footnote{\ldots for historical reasons.}


\section{Anatomy of a SURF task}

\scusoft\ tasks follow the same general format internally:

\begin{enumerate}
\item Request the messaging level from the parameter system
(\param{MSG\_LEVEL}).
\item Read in some information concerning the object name, observation mode
and run number and write this to the screen (at normal messaging level) to
keep the user informed.
\item Check the history component to make sure that the correct tasks have
been run prior to running the current task. This always contains a check to
see that the \resw\ task has been run on the data.
\item Check the dimensions of the data array, the \texttt{LST\_STRT} array and
the \texttt{DEM\_PNTR} arrays.
\item Retrieve information from NDF extensions (see e.g.\ 
\texttt{SCULIB\_GET\_BOL\_DESC}, \texttt{SCULIB\_GET\_JIGGLE}). [as required]
\item If required, construct an output name based on the input filename
(\texttt{SCULIB\_CONSTRUCT\_OUT}). The output file is usually propogated from
the input file using \texttt{NDF\_PROP}.
\item Do whatever has to be done.
\item Shut down NDF, free memory, free locators.

\end{enumerate}

\section{Writing an import task}

In order to import data into \scusoft\ that has not been taken with the
SCUBA VAX/VMS data acquisition system, it is necessary to convert the
original file format to a form that \scusoft\ can understand.
The file can be converted to the same format as written by the SCUBA
acquisition system (\S\ref{sec:raw}) or the format written by the
\resw\ task (\S\ref{sec:surfform}).

It is expected that the conversion of the data structure to the correct format 
will be fairly simple. The difficulty lies in generating the required FITS
information (with the correct FITS keywords) and the NDF extensions.

\subsection{Case Study: DREAM}

DREAM is a SCUBA observing mode where the chopping is done as part of the
jiggle pattern and the resulting signal is reconstructued by tiling the
bolometer patterns (le Poole and Greve, 1998, SPIE 3357). In this case
it is not necessary to subtract switches, despike or flatfield the data and
\scusoft\ is required simply to perform the regridding and coadding.
In this case, the \scusoft\ format was chosen to simulate the output of the
\flatf\ task (the HISTORY component of the NDF was updated to reflect this).

The required stages are:

\begin{enumerate}
\item Read in the import data file.
\item Create output NDF.
\item Convert data array to a 2-d array of dimensions (N\_BOLS $\times$
N\_SAMPLES). Also create Variance and Quality arrays.
\item Write AXIS information. This is simply bolometer number for the X axis
and integration number for the Y axis.
\item Create the \texttt{SCUCD}, \texttt{SCUBA} and \texttt{REDS} extensions.
\item Construct the \texttt{DEM\_PNTR} array.
\item Construct the \texttt{LST\_STRT} array.
\item Store the jiggle patterns in \texttt{JIGL\_X} and \texttt{JIGL\_Y}.
\item Construct \texttt{BOL\_CHAN} and \texttt{BOL\_ADC} arrays.
\item Store the flatfield information (essentially for the X and Y offsets of
each bolometer). This is done by reading in a flatfield file using
\texttt{SCULIB\_READBOLS}.
\item Create FITS header (see table \ref{tab:jigfits} for required keywords for
jiggle map).
\item Write fake history information to convince \scusoft\ that \resw\ and
\flatf\ have been run on the file.
\item Shut down.
\end{enumerate}


\appendix

\clearpage

\section*{\xlabel{glossary}Glossary\label{glossary}}
\addcontentsline{toc}{section}{Glossary}

\begin{description}

\item{{\bf chopping}} The secondary mirror is continuously moved on and off
source at approximately 7~Hz in order to remove the sky to zeroth order. This
is done in addition to standard jiggling.

\item{{\bf demodulation}} Removal of the chop signal by the transputers.
At this time, the raw data can not be accessed, only the demodulated data are
stored.

\item{{\bf exposure}} An exposure is the result from a complete set of
switches. For example, in a JIGGLE/MAP or PHOTOM observation where the
telescope is nodding the source between left and right beams, the data from
each nod position is a switch and the reduced result `left switch' - `right
switch' say, is an exposure. In a SCAN/MAP observation there is no beam
switching so, in this case, an exposure is the same as a switch.


\item{{\bf integration}} An integration means different things for different
observations.

For one of the mapping modes it means the data from one fully-sampled coverage
of the map area. In a JIGGLE/MAP, where full sampling is achieved by jiggling
the secondary mirror, an integration is generally the results from one pass
through the complete jiggle pattern. An integration is made up of one or more
exposures. 

Similarly, an integration for a SCAN/MAP observation is made up of data from
the raster scans that cover the map area once.

For PHOTOM observations an integration is usually the average of a 9 point
mini-jiggle.

For a SKYDIP observation, an integration is the data from a single revolution
of the sector chopper in front of the cryostat window.


\item{{\bf jiggle}} In order to sample an image fully the secondary mirror
is moved once a second (whilst chopping) to move the position of the array on
the sky; this is called `jiggling.'  There are a  complete set of jiggle
positions for each integration. A PHOTOM observation can also jiggle in order
to correct for seeing effects.

\item{{\bf measurement}} A measurement is a group of integrations. Most MAP or
PHOTOM observations will consist of only one measurement.

A FOCUS or ALIGN observation consists of five measurements (one for each
secondary mirror position). A SKYDIP observation consists of one measurement
at each elevation.

\item{{\bf nod}} In order to correct for atmospheric variation the telescope
is moved off-source in each exposure so that sky can be measured.

\item{{\bf ODF}} The observation definition file (ODF) is a file containing a
list of instructions for an observation with SCUBA.

\item{{\bf sub-instrument}} SCUBA contains bolometer arrays and photometric
pixels that can operate at several wavelengths simultaneously. Each of these
is called a sub-instrument. They are:

\begin{itemize}
\item SHORT - the short wave array containing 91 bolometers
\item LONG - the long wave array containing 37 bolometers
\item P1100 - the single bolometer optimised for 1100\micron.
\item P1350 - the single bolometer optimised for 1350\micron.
\item P2000 - the single bolometer optimised for 2000\micron.
\end{itemize}


\item{{\bf switch}} The switch is the fundamental unit of data-taking in an
observation. For example in a JIGGLE/MAP or PHOTOM observation each chunk of
jiggle positions measured with the object in the beam of a telescope is a
switch. Each scan across the source in a SCAN/MAP observation is also a switch.
 
\item{{\bf tau ($\tau$)}} Submillimetre extinction is measured using the
zenith optical depth, tau or $\tau$, this is a measure of the amount of water
vapour present in the atmosphere. For a tau, $\tau$ at a given airmass, $A$,
the attenuation due to the atmosphere is given as e$^{-A \tau}$. Note that tau
is wavelength dependent and that the value quoted by the 
\htmladdnormallink{Caltech Submillimetre
Observatory}{http://www.cco.caltech.edu/\~{}cso/} (CSO) is the $\tau$ at 225~GHz and will therefore be different at 
the other wavelengths used by SCUBA

\end{description}





\section{SURF coordinate frames\label{app:frames}}

\scusoft\ understands a number of different coordinate frames and uses
two letter codes\footnote{with the exception of PLANET frame} 
(as used by the JCMT) to describe them. The following
are available:

\begin{tabular}{ll}
RB & FK4 B1950\\
RJ & FK5 J2000\\
RD & Apparent RA/Dec, current epoch\\
GA & Galactic coordinates (J2000) \\
NA & Left hand Nasmyth coordinates \\
AZ & Az/El offsets \\
PLANET & JCMT input coordinate frame for moving source \\
PL & SURF output coordinate for offsets from tracking centre\\
\end{tabular}

These coordinate frames are used to specify the tracking centre
(\texttt{CENT\_CRD}), the coordinate frame of the chop (\texttt{CHOP\_CRD}),
the coordinate frame of any offsets (\texttt{LOCL\_CRD}) and the coordinate
frame of the sampling (\texttt{SAM\_CRDS}) (eg jiggle coordinate frame or scan
coordinate frame)

SURF currently assumes that the instrument is on the left hand Nasmyth
platform.

\section{Building SURF}

The \scusoft\ development tree is in CVS and has the layout shown in
Fig.\ \ref{surf:devel}. The development version can be built 
by using the default make target in the \texttt{surf} directory or
by running the individual makefiles in the \texttt{sculib}, \texttt{surflib},
\texttt{surf\_kap} and \texttt{src} directories. The \scusoft\ monolith is
built in the \texttt{src} directory along with compilation of the IFC files
(the monolith IFL file is created automatically from the task IFL files).

The current development version number (i.e.\ the version number written to
history components of the NDF files) is defined in \texttt{src/makefile}
via the \texttt{PKG\_VERS} variable. (see \texttt{surf\_set\_app\_name.F} for
more information on the setting of the version number in the monolith.

\begin{figure}
\begin{center}
\begin{myquote}
\begin{verbatim}
surf ---- src                          Top level surf_* and IFL files
      |
      |-- surflib                      surflib source code
      |
      |-- sculib                       sculib source code
      |
      |-- surf_kap                     Useful routines from Kappa
      |
      |-- scripts                      Scripts (eg surf.csh, remdbm.pl)
      |
      |-- misc                         Miscellaneous (eg ipfile.dat)
      |
      |-- dream                        DREAM mode import routine
      |
      |-- starlink                     Starlink makefiles
      |
      |-- docs --------- hlp           On-line help system
                     |
                     |-- sun216        Main SURF documentation
                     |
                     |-- sc10          Photometry cookbook
                     |
                     |-- sc11          Mapping cookbook
                     |
                     |-- sc11_mini     mini mapping cookbook
\end{verbatim}
\end{myquote}
\end{center}
\caption{Directory layout of SURF development system. CVS directories are not
shown.}
\label{surf:devel}
\end{figure}

In order to generate a Starlink distribution the \texttt{export} target must
be used from the root \texttt{makefile}. This generates tar files from each
directory and copies them to the \texttt{starlink} directory.  If required, the
version number must be updated manually in the Starlink \texttt{makefile}.
The platform dependencies are stored in the \texttt{mk} script that is used
to run \texttt{make} with the correct environment. The environment is selected
with the \texttt{SYSTEM} environment variable and can be one of 
\texttt{sun4\_Solaris}, \texttt{ix86\_Linux} or \texttt{alpha\_OSF1}.

Once the \texttt{SYSTEM} environment variable has been set the \texttt{mk}
script can be run to build and export \scusoft:

\begin{description}

\item[./mk export\_source] \mbox{} \\
   Exports the source into a compressed tar file ready for distribution.
\item[./mk build] \mbox{} \\
   Builds the monolith, extracting source files from tar files as required.
\item[./mk export] \mbox{} \\
   Exports the built system into a compressed tar file ready for distribution.
This will include the source code.
\item[./mk export\_run] \mbox{} \\
   Exports the run-only system into a compressed tar file ready for
distribution. This will not include the source code.
\item[./mk install] \mbox{} \\
  Installs the built system into the directory structure specified by the
\texttt{INSTALL} environment variable. Files are installed below this
directory into \texttt{bin/surf}, \texttt{docs}, \texttt{help/surf}
and \texttt{dates}. Soft links are made from the installed binary directory
to the current directory so that the installed files can be seen.
Usually, \texttt{INSTALL} is set to \texttt{/star}.
\item[./mk deinstall] \mbox{} \\
  Deinstall the package.
\end{description}

\section{SURF Constants\label{app:const}}

The SURF\_PAR include file defines the global constants that are used
in \scusoft\ for preallocating arrays and defining useful mathematical 
constants and strings. Only top level \scusoft\ routines (ie names of
\texttt{surf\_*} include this file; all other library routines 
obtain the value from argument.

\subsection{Constant strings}

\begin{quote}
\begin{description}
  \item[PACKAGE = 'SURF'] \mbox{} \\
    Name of the software package. Used in output messages. (CHAR*10)
\end{description}
\end{quote}

\subsection{Convenient numeric definitions}

\begin{quote}
\begin{description}
  \item[PI = 3.14159265359D0] \mbox{} \\
    $\pi$ [DOUBLE]
  \item[R2AS = 206264.806247D0] \mbox{} \\
    Used to convert radians to arcseconds:
    \begin{equation}
     \mathrm{arcsec} = 3600\frac{180}{\pi}\theta
    \end{equation}
  [DOUBLE]
\end{description}
\end{quote}

\subsection{Parameters controlling the observation limits}

\begin{quote}
\begin{description}
  \item[SCUBA\_\_MAX\_BEAM = 3] \mbox{} \\
     maximum number of bolometers in a single sub-instrument that actually
observer the source in a PHOTOM observation. [INTEGER]
  \item[SCUBA\_\_MAX\_MEAS = 200] \mbox{} \\
      maximum number of measurements in an observation. [INTEGER]
  \item[SCUBA\_\_MAX\_INT = 1000] \mbox{} \\
      maximum number of integrations in a measurement. [INTEGER]
  \item[SCUBA\_\_MAX\_EXP = 128] \mbox{} \\
      maximum number of exposures per integration. [INTEGER]
  \item[SCUBA\_\_MAX\_SWITCH = 3] \mbox{} \\
      maximum number of switches per exposure. [INTEGER]
  \item[SCUBA\_\_MAX\_JIGGLE = 512] \mbox{} \\
      maximum number of offsets in a jiggle pattern. [INTEGER]
\end{description} 
\end{quote}

\subsection{Parameters governing the instrument itself}

\begin{quote}
\begin{description}
  \item[SCUBA\_\_MAX\_SUB = 5] \mbox{} \\
    Maximum number of sub-instruments (i.e.\ for SCUBA there are LONG, SHORT,
P2000, P1350 and P1100 sub-instruments). [INTEGER]
  \item[SCUBA\_\_NUM\_ADC = 9] \mbox{} \\
    Number of A/D cards (called A, B, C, D, E, F, G, H and I for SCUBA).
  \item[SCUBA\_\_NUM\_CHAN = 16] \mbox{} \\
    Number of channels on each A/D card.
\end{description}
\end{quote}

\subsection{Parameters dealing with automatic output filename generation}

\begin{quote}
\begin{description}

  \item[SUFFIX\_ENV = 'SCUBA\_SUFFIX'] \mbox{} \\
    Name of the environment variable dealing with the format of the
automatically generated output filename. [CHAR*15]

  \item[SCUBA\_\_N\_SUFFIX = 3] \mbox{} \\
    Number of suffix options [INTEGER]

  \item[SUFFIX\_OPTIONS = 'LONG','SHORT','VERBOSE']\mbox{} \\
    The allowed suffix options i.e.\ these are the values that are valid for
the environment variable specified by SUFFIX\_ENV [CHAR*10 (SCUBA\_N\_SUFFIX)]

\end{description}
\end{quote}

\subsection{General limits}

\begin{quote}
\begin{description}

  \item[SCUBA\_\_MAX\_POINT = 20] \mbox{} \\
   Maximum number of pointing corrections.
  \item[SCUBA\_\_MAX\_FITS = 200] \mbox{} \\
   Maximum number of FITS items in header.
  \item[SCUBA\_\_MAX\_SECT = 10] \mbox{} \\
   Maximum number of SCUBA sections that can be specified.
  \item[SCUBA\_\_N\_TEMPS = 3] \mbox{} \\
    Number of temperatures per reading for SKYDIP observations: AMBIENT, SKY,
COLD. 

\end{description}
\end{quote}


\section{SCUBA FITS headers\label{app:fits}}

This section lists the FITS headers written by SCUBA. \scusoft\ uses a subset
of these keywords. Multiple headers are listed since the values are slightly
different for different observing modes.

\subsection{Skydip}

This is a typical SKYDIP FITS header:

\begin{myquote}
\begin{verbatim}
ACCEPT  = 'NO      '           / accept update; PROMPT, YES or NO
ALIGN_AX= 'not used'           / Alignment measurements in X or Y axis
ALIGN_SH=                   -1 / Distance between successive alignment offsets (
ALT-OBS =                 4092 / Height of observatory above sea level (metres)
AMEND   =             3.863703 / Airmass at end of observation
AMSTART =             1.015427 / Airmass at start of observation
APEND   =              627.851 / Air pressure at end of observation (mbar)
APSTART =              627.851 / Air pressure at start of observation (mbar)
ATEND   =             5.604397 / Air temp. at end of observation (C)
ATSTART =             5.677658 / Air temp. at start of observation (C)
BOLOMS  = 'SHORT_DC,LONG_DC'   / Names of bolometers measured
CALIBRTR=                    F / Internal calibrator is on or off
CAL_FRQ =                   -1 / Calibrator frequency (Hz)
CENT_CRD= 'AZ      '           / Centre coordinate system
CHOP_CRD= 'not used'           / Chopper coordinate system
CHOP_FRQ=                    2 / Chopper frequency (Hz)
CHOP_FUN= 'not used'           / Chopper waveform
CHOP_PA =                   -1 / Chopper P.A., 0 = in lat, 90 = in long
CHOP_THR=                   -1 / Chopper throw (arcsec)
DATA_DIR= '19990808'           / Sub-directory where datafile was stored
DATA_KPT= 'DEMOD   '           / The type of data stored to disk
END_AZD =              219.183 / Azimuth at end of observation (deg)
END_EL  =                   80 / Elevation of last SKYDIP point (deg)
END_ELD =              79.9896 / Elevation at end of observation
EQUINOX =                 2000 / Equinox of mean coordinate system
EXPOSED =                    5 / Exposure per pixel (seconds)
EXP_NO  =                    1 / Exposure number at end of observation
EXP_TIME=                  0.5 / Exposure time for each basic measurement (sec)
E_PER_I =                    1 / Number of exposures per integration
FILTER  = '450N:850N'          / Filters used
FOCUS_SH=                   -1 / Shift between focus measurements (mm)
GAIN    =                    1 / Programmable gain
HSTEND  = '20:05:49.00131'     / HST at end of observation
HSTSTART= '20:00:56.00281'     / HST at start of observation
HUMEND  =                   15 / Humidity (%) at end of observation
HUMSTART=                   15 / Humidity (%) at start of observation
INSTRUME= 'SCUBA   '           / Name of instrument used
INT_NO  =                    1 / Integration number at end of observation
JIGL_CNT=                   -1 / Number of offsets in jiggle pattern
JIGL_NAM= 'not used'           / File containing jiggle offsets
J_PER_S =                   -1 / Number of jiggles per switch position
J_REPEAT=                   -1 / No. jiggle pattern repeats in a switch
LAT     = 'not used'           / Object latitude
LAT-OBS =        19.8258323669 / Latitude of observatory (degrees)
LAT2    = 'not used'           / Object latitude at MJD2
LOCL_CRD= 'no      '           / Local offset coordinate system
LONG    = '+219:11:26.64'      / Object longitude
LONG-OBS=        204.520278931 / East longitude of observatory (degrees)
LONG2   = 'not used'           / Object Longitude at MJD2
MAP_HGHT=                   -1 / Height of rectangle to be mapped (arcsec)
MAP_PA  =                   -1 / P.A. of map vertical, +ve towards +ve long
MAP_WDTH=                   -1 / Width of rectangle to be mapped (arcsec)
MAP_X   =                   -1 / Map X offset from telescope centre (arcsec)
MAP_Y   =                   -1 / Map Y offset from telescope centre (arcsec)
MAX_EL  =                   80 / Max elevation of sky-dip (deg)
MEANDEC =            -17.93834 / -17:56:18.01437 = approx. mean Dec. (deg)
MEANRA  =             222.1444 / 222:08:39.76318 = approx. mean R.A. (deg)
MEAS_NO =                   10 / Measurement number at end of observation
MIN_EL  =                   15 / Min elevation of sky-dip (deg)
MJD1    =                   -1 / Modified Julian day planet at RA,DEC
MJD2    =                   -1 / Modified Julian day planet at RA2,DEC2
MODE    = 'SKYDIP  '           / The type of observation
N_INT   =                   10 / No. integrations in the observation
N_MEASUR=                   10 / No. measurements in the observation
OBJECT  = 'not used'           / Name of object
OBJ_TYPE= 'not used'           / Type of object
OBSDEF  = 'ss:odfsxskydip_850.obs_193300' / The observation definition file
OBSERVER= 'jmk     '           / The name of the observer
PROJ_ID = 'SCUBA   '           / The project identification
RUN     =                   26 / Run number of observation
SAM_CRDS= 'no      '           / Coordinatesystem of sampling mesh
SAM_DX  =                   -1 / Sample spacing along scan direction (arcsec)
SAM_DY  =                   -1 / Sample spacing perp. to scan (arcsec)
SAM_MODE= 'not used'           / Sampling method
SAM_PA  =                   -1 / Scan P.A. rel. to lat. line; 0=lat, 90=long
SCAN_REV=                    F / .TRUE. if alternate scans reverse direction
SPK_NSIG=                   -1 / N sigmas from fit of spike threshold
SPK_RMVL=                    F / Automatic spike removal
SPK_WDTH=                   -1 / Assumed width of spike
START_EL=                   15 / Elevation of first SKYDIP point (deg)
STATE   = 'Terminating         :   RESET_DATA_SIZE' / SCUCD state
STEND   = '16:49:09.149323'    / ST at end of observation
STRT_AZD=              219.174 / Azimuth at observation start (deg)
STRT_ELD=              14.9868 / Elevation at observation start (deg)
STSTART = '16:44:15.34744'     / ST at start of observation
SWTCH_MD= 'not     '           / Switch mode of observation
SWTCH_NO=                    1 / Switch number at end of observation
S_PER_E =                    1 / Number of switch positions per exposure
TELESCOP= 'JCMT    '           / Name of telescope
TEL_OPER= 'kkp     '           / Telescope operator
UTDATE  = '1999:8:8'           / UT date of observation
UTEND   = '6:05:48.9996'       / UT at end of observation
UTSTART = '6:00:55.99937'      / UT at start of observation
VERSION =                  1.1 / SCUCD version
ALIGN_DX=             0.027268 / SMU tables X axis alignment offset
ALIGN_DY=             0.482489 / SMU tables Y axis alignment offset
ALIGN_X =             -0.96279 / SMU tables X axis
ALIGN_Y =              3.79108 / SMU tables Y axis
AZ_ERR  =                    0 / Error in the telescope azimuth
CHOPPING=                    F / SMU Chopper chopping state
EL_ERR  =             0.154495 / Error in the telescope elevation
FOCUS_DZ=            -0.361142 / SMU tables Z axis focus offset
FOCUS_Z =             -17.2453 / SMU tables Z axis
SEEING  =                0.525 / SAO atmospheric seeing
SEE_DATE= '9908061715'         / Date and time of SAO seeing
TAU_225 =                0.066 / CSO tau
TAU_DATE= '9908080555'         / Date and time of CSO tau
TAU_RMS =              5.0E-03 / CSO tau rms
UAZ     =            -0.528088 / User azimuth pointing offset
UEL     =              4.83289 / User elevation pointing offset
UT_DATE = ' 8 Aug 1999'        / UT date at start of observation
BAD_LIM =                   32 / No. spikes before quality set bad
CALIB_LG=                   -1 / Lag of internal calibrator in samples
CALIB_PD=                   -1 / Period of internal calibrator in samples
CHOP_LG =                    0 / Chop lag in samples
CHOP_PD =             64.03689 / Chop period in samples
CNTR_DU3=                    0 / Nasmyth dU3 coord of instrument centre
CNTR_DU4=                    0 / Nasmyth dU4 coord of instrument centre
ETATEL_1=                 0.75 / Transmission of telescope
ETATEL_2=                  0.9 / Transmission of telescope
ETATEL_3=                   -1 / Transmission of telescope
ETATEL_4=                   -1 / Transmission of telescope
ETATEL_5=                   -1 / Transmission of telescope
FILT_1  = '450     '           / Filter name
FILT_2  = '850     '           / Filter name
FILT_3  = 'not used'           / Filter name
FILT_4  = 'not used'           / Filter name
FILT_5  = 'not used'           / Filter name
FLAT    = 'jcmtdata_dir:lwswphot.dat' / Name of flat-field file
JIG_DSCD=                   -1 / No. samples discarded after jiggle movement
L_GD_BOL= 'not used'           / Bol. to whose value LW guard ring is set
L_GUARD =                    F / Long wave guard ring on or off
MEAS_BOL= 'not used'           / Bolometers  actually measured in observation
N_BOLS  =                    2 / Number of bolometers selected
N_SUBS  =                    2 / Number of sub-instruments used
PHOT_BBF= 'not_used       LL,C14,NULL' / The bolometers on the source
PRE_DSCD=                    0 / No. of samples discarded before chop movement
PST_DSCD=                    0 / No. samples discarded after chop movement
REBIN   = 'LINEAR  '           / Rebinning method used by SCUIP
REF_ADC =                   -1 / A/D card of FLATFIELD reference bolometer
REF_CHAN=                   -1 / Channel of FLATFIELD reference bolometer
SAM_TIME=                  122 / A/D sample period in ticks (64musec)
SIMULATE=                    F / True if data is simulated
SKY     = 'jcmtdata_dir:skydip_startup.dat' / Name of sky opacities file
SUB_1   = 'SHORT   '           / SCUBA instrument being used
SUB_2   = 'LONG    '           / SCUBA instrument being used
SUB_3   = 'not used'           / SCUBA instrument being used
SUB_4   = 'not used'           / SCUBA instrument being used
SUB_5   = 'not used'           / SCUBA instrument being used
S_GD_BOL= 'not used'           / Bol. to whose value SW guard ring is set
S_GUARD =                    F / Short wave guard ring on or off
TAUZ_1  =                   -1 / Zenith sky optical depth
TAUZ_2  =                   -1 / Zenith sky optical depth
TAUZ_3  =                   -1 / Zenith sky optical depth
TAUZ_4  =                   -1 / Zenith sky optical depth
TAUZ_5  =                   -1 / Zenith sky optical depth
T_AMB   =             278.6532 / The ambient air temperature (K)
T_COLD_1=                  120 / Effective temperature of cold load (K)
T_COLD_2=                  100 / Effective temperature of cold load (K)
T_COLD_3=                   -1 / Effective temperature of cold load (K)
T_COLD_4=                   -1 / Effective temperature of cold load (K)
T_COLD_5=                   -1 / Effective temperature of cold load (K)
T_HOT   =                282.6 / The temperature of the hot load (K)
T_TEL   =             282.2908 / The temperature of the telescope
USE_CAL =                    F / .TRUE. if dividing chop by cal before rebin
WAVE_1  =                  442 / Wavelength of map (microns)
WAVE_2  =                  862 / Wavelength of map (microns)
WAVE_3  =                    0 / Wavelength of map (microns)
WAVE_4  =                    0 / Wavelength of map (microns)
WAVE_5  =                    0 / Wavelength of map (microns)
END
\end{verbatim}
\end{myquote}

\subsection{Jiggle map}

Jiggle map of a planet (note that for a moving source SCUBA supplies that
apparent RA and Dec of the source at two different times. These are specified
as \texttt{LON}G and \texttt{LAT} for \texttt{MJD1} and \texttt{LONG2} and
\texttt{LAT2} for \texttt{MJD2}):

\begin{myquote}
\begin{verbatim}
ACCEPT  = 'not used'           / accept update; PROMPT, YES or NO
ALIGN_AX= 'not used'           / Alignment measurements in X or Y axis
ALIGN_SH=                   -1 / Distance between successive alignment offsets (
ALT-OBS =                 4092 / Height of observatory above sea level (metres)
AMEND   =             1.296355 / Airmass at end of observation
AMSTART =             1.283144 / Airmass at start of observation
APEND   =             627.4847 / Air pressure at end of observation (mbar)
APSTART =             627.2405 / Air pressure at start of observation (mbar)
ATEND   =             7.655678 / Air temp. at end of observation (C)
ATSTART =             7.997562 / Air temp. at start of observation (C)
BOLOMS  = 'LONG,SHORT'         / Names of bolometers measured
CALIBRTR=                    F / Internal calibrator is on or off
CAL_FRQ =             2.929688 / Calibrator frequency (Hz)
CENT_CRD= 'PLANET  '           / Centre coordinate system
CHOP_CRD= 'LO      '           / Chopper coordinate system
CHOP_FRQ=               7.8125 / Chopper frequency (Hz)
CHOP_FUN= 'SQUARE  '           / Chopper waveform
CHOP_PA =                    0 / Chopper P.A., 0 = in lat, 90 = in long
CHOP_THR=                   68 / Chopper throw (arcsec)
DATA_DIR= '19990808'           / Sub-directory where datafile was stored
DATA_KPT= 'DEMOD   '           / The type of data stored to disk
END_AZD =              198.486 / Azimuth at end of observation (deg)
END_EL  =                   -1 / Elevation of last SKYDIP point (deg)
END_ELD =              50.3461 / Elevation at end of observation
EQUINOX =                 2000 / Equinox of mean coordinate system
EXPOSED =                0.128 / Exposure per pixel (seconds)
EXP_NO  =                   13 / Exposure number at end of observation
EXP_TIME=                0.128 / Exposure time for each basic measurement (sec)
E_PER_I =                   -1 / Number of exposures per integration
FILTER  = '450N:850N'          / Filters used
FOCUS_SH=                   -1 / Shift between focus measurements (mm)
GAIN    =                   10 / Programmable gain
HSTEND  = '18:52:56.00327'     / HST at end of observation
HSTSTART= '18:41:44.99954'     / HST at start of observation
HUMEND  =                   14 / Humidity (%) at end of observation
HUMSTART=                   14 / Humidity (%) at start of observation
INSTRUME= 'SCUBA   '           / Name of instrument used
INT_NO  =                    1 / Integration number at end of observation
JIGL_CNT=                   -1 / Number of offsets in jiggle pattern
JIGL_NAM= 'not used'           / File containing jiggle offsets
J_PER_S =                   -1 / Number of jiggles per switch position
J_REPEAT=                   -1 / No. jiggle pattern repeats in a switch
LAT     = '-017:55:40.89'      / Object latitude
LAT-OBS =        19.8258323669 / Latitude of observatory (degrees)
LAT2    = '-017:55:47.08'      / Object latitude at MJD2
LOCL_CRD= 'RD      '           / Local offset coordinate system
LONG    = '+014:48:26.75'      / Object longitude
LONG-OBS=        204.520278931 / East longitude of observatory (degrees)
LONG2   = '+014:48:28.04'      / Object Longitude at MJD2
MAP_HGHT=                  600 / Height of rectangle to be mapped (arcsec)
MAP_PA  =                    0 / P.A. of map vertical, +ve towards +ve long
MAP_WDTH=                  600 / Width of rectangle to be mapped (arcsec)
MAP_X   =                    0 / Map X offset from telescope centre (arcsec)
MAP_Y   =                    0 / Map Y offset from telescope centre (arcsec)
MAX_EL  =                   -1 / Max elevation of sky-dip (deg)
MEANDEC =            -17.92968 / -17:55:46.84753 = approx. mean Dec. (deg)
MEANRA  =             222.1171 / 222:07:01.490479 = approx. mean R.A. (deg)
MEAS_NO =                    1 / Measurement number at end of observation
MIN_EL  =                   -1 / Min elevation of sky-dip (deg)
MJD1    =    51398.19570601852 / Modified Julian day planet at RA,DEC
MJD2    =    51398.20612268519 / Modified Julian day planet at RA2,DEC2
MODE    = 'MAP     '           / The type of observation
N_INT   =                    1 / No. integrations in the observation
N_MEASUR=                    1 / No. measurements in the observation
OBJECT  = 'mars    '           / Name of object
OBJ_TYPE= 'UNKNOWN '           / Type of object
OBSDEF  = 'ss:scan_68dec.obs_mars_184023' / The observation definition file
OBSERVER= 'jmk     '           / The name of the observer
PROJ_ID = 'm99au34 '           / The project identification
RUN     =                   18 / Run number of observation
SAM_CRDS= 'NA      '           / Coordinatesystem of sampling mesh
SAM_DX  =                    3 / Sample spacing along scan direction (arcsec)
SAM_DY  =                   60 / Sample spacing perp. to scan (arcsec)
SAM_MODE= 'RASTER  '           / Sampling method
SAM_PA  =                 15.5 / Scan P.A. rel. to lat. line; 0=lat, 90=long
SCAN_REV=                    T / .TRUE. if alternate scans reverse direction
SPK_NSIG=                   -1 / N sigmas from fit of spike threshold
SPK_RMVL=                    F / Automatic spike removal
SPK_WDTH=                   -1 / Assumed width of spike
START_EL=                   -1 / Elevation of first SKYDIP point (deg)
STATE   = 'Terminating         :   STEP' / SCUCD state
STEND   = '15:36:04.176178'    / ST at end of observation
STRT_AZD=              194.923 / Azimuth at observation start (deg)
STRT_ELD=              51.0202 / Elevation at observation start (deg)
STSTART = '15:24:51.33911'     / ST at start of observation
SWTCH_MD= 'NOSW    '           / Switch mode of observation
SWTCH_NO=                    1 / Switch number at end of observation
S_PER_E =                    1 / Number of switch positions per exposure
TELESCOP= 'JCMT    '           / Name of telescope
TEL_OPER= 'kkp     '           / Telescope operator
UTDATE  = '1999:8:8'           / UT date of observation
UTEND   = '4:52:55.99983'      / UT at end of observation
UTSTART = '4:41:44.99954'      / UT at start of observation
VERSION =                  1.1 / SCUCD version
ALIGN_DX=             0.027268 / SMU tables X axis alignment offset
ALIGN_DY=             0.482489 / SMU tables Y axis alignment offset
ALIGN_X =             -2.32166 / SMU tables X axis
ALIGN_Y =              3.78719 / SMU tables Y axis
AZ_ERR  =                    0 / Error in the telescope azimuth
CHOPPING=                    T / SMU Chopper chopping state
EL_ERR  =                    0 / Error in the telescope elevation
FOCUS_DZ=            -0.361142 / SMU tables Z axis focus offset
FOCUS_Z =             -16.5076 / SMU tables Z axis
SEEING  =                0.525 / SAO atmospheric seeing
SEE_DATE= '9908061715'         / Date and time of SAO seeing
TAU_225 =                0.097 / CSO tau
TAU_DATE= '9908080435'         / Date and time of CSO tau
TAU_RMS =              8.0E-03 / CSO tau rms
UAZ     =            -0.528088 / User azimuth pointing offset
UEL     =              4.83289 / User elevation pointing offset
UT_DATE = ' 8 Aug 1999'        / UT date at start of observation
BAD_LIM =                   32 / No. spikes before quality set bad
CALIB_LG=                    6 / Lag of internal calibrator in samples
CALIB_PD=             42.66667 / Period of internal calibrator in samples
CHOP_LG =                    4 / Chop lag in samples
CHOP_PD =                   16 / Chop period in samples
CNTR_DU3=                    0 / Nasmyth dU3 coord of instrument centre
CNTR_DU4=                    0 / Nasmyth dU4 coord of instrument centre
ETATEL_1=                   -1 / Transmission of telescope
ETATEL_2=                   -1 / Transmission of telescope
ETATEL_3=                   -1 / Transmission of telescope
ETATEL_4=                   -1 / Transmission of telescope
ETATEL_5=                   -1 / Transmission of telescope
FILT_1  = '450     '           / Filter name
FILT_2  = '850     '           / Filter name
FILT_3  = 'not used'           / Filter name
FILT_4  = 'not used'           / Filter name
FILT_5  = 'not used'           / Filter name
FLAT    = 'jcmtdata_dir:lwswphot.dat' / Name of flat-field file
JIG_DSCD=                   -1 / No. samples discarded after jiggle movement
L_GD_BOL= 'not used'           / Bol. to whose value LW guard ring is set
L_GUARD =                    F / Long wave guard ring on or off
MEAS_BOL= 'not used'           / Bolometers  actually measured in observation
N_BOLS  =                  128 / Number of bolometers selected
N_SUBS  =                    2 / Number of sub-instruments used
PHOT_BBF= 'not_used       LL,C14,NULL' / The bolometers on the source
PRE_DSCD=                    0 / No. of samples discarded before chop movement
PST_DSCD=                    0 / No. samples discarded after chop movement
REBIN   = 'LINEAR  '           / Rebinning method used by SCUIP
REF_ADC =                   -1 / A/D card of FLATFIELD reference bolometer
REF_CHAN=                   -1 / Channel of FLATFIELD reference bolometer
SAM_TIME=                  125 / A/D sample period in ticks (64musec)
SIMULATE=                    F / True if data is simulated
SKY     = 'jcmtdata_dir:skydip_startup.dat' / Name of sky opacities file
SUB_1   = 'SHORT   '           / SCUBA instrument being used
SUB_2   = 'LONG    '           / SCUBA instrument being used
SUB_3   = 'not used'           / SCUBA instrument being used
SUB_4   = 'not used'           / SCUBA instrument being used
SUB_5   = 'not used'           / SCUBA instrument being used
S_GD_BOL= 'not used'           / Bol. to whose value SW guard ring is set
S_GUARD =                    F / Short wave guard ring on or off
TAUZ_1  =                    0 / Zenith sky optical depth
TAUZ_2  =                    0 / Zenith sky optical depth
TAUZ_3  =                    0 / Zenith sky optical depth
TAUZ_4  =                    0 / Zenith sky optical depth
TAUZ_5  =                    0 / Zenith sky optical depth
T_AMB   =                   -1 / The ambient air temperature (K)
T_COLD_1=                   -1 / Effective temperature of cold load (K)
T_COLD_2=                   -1 / Effective temperature of cold load (K)
T_COLD_3=                   -1 / Effective temperature of cold load (K)
T_COLD_4=                   -1 / Effective temperature of cold load (K)
T_COLD_5=                   -1 / Effective temperature of cold load (K)
T_HOT   =                   -1 / The temperature of the hot load (K)
T_TEL   =                   -1 / The temperature of the telescope
USE_CAL =                    F / .TRUE. if dividing chop by cal before rebin
WAVE_1  =                  442 / Wavelength of map (microns)
WAVE_2  =                  862 / Wavelength of map (microns)
WAVE_3  =                    0 / Wavelength of map (microns)
WAVE_4  =                    0 / Wavelength of map (microns)
WAVE_5  =                    0 / Wavelength of map (microns)
END
\end{verbatim}
\end{myquote}

\subsection{Photometry}

This is a photometry header:

\begin{myquote}
\begin{verbatim}
ACCEPT  = 'not used'           / accept update; PROMPT, YES or NO
ALIGN_AX= 'not used'           / Alignment measurements in X or Y axis
ALIGN_SH=                   -1 / Distance between successive alignment offsets (
ALT-OBS =                 4092 / Height of observatory above sea level (metres)
AMEND   =             1.279808 / Airmass at end of observation
AMSTART =             1.263818 / Airmass at start of observation
APEND   =             628.7057 / Air pressure at end of observation (mbar)
APSTART =             628.7057 / Air pressure at start of observation (mbar)
ATEND   =             6.678879 / Air temp. at end of observation (C)
ATSTART =             6.776558 / Air temp. at start of observation (C)
BOLOMS  = 'H7      '           / Names of bolometers measured
CALIBRTR=                    T / Internal calibrator is on or off
CAL_FRQ =             2.929688 / Calibrator frequency (Hz)
CENT_CRD= 'RB      '           / Centre coordinate system
CHOP_CRD= 'AZ      '           / Chopper coordinate system
CHOP_FRQ=               7.8125 / Chopper frequency (Hz)
CHOP_FUN= 'SCUBAWAVE'          / Chopper waveform
CHOP_PA =                   90 / Chopper P.A., 0 = in lat, 90 = in long
CHOP_THR=                   45 / Chopper throw (arcsec)
DATA_DIR= '19990808'           / Sub-directory where datafile was stored
DATA_KPT= 'DEMOD   '           / The type of data stored to disk
END_AZD =              274.745 / Azimuth at end of observation (deg)
END_EL  =                   -1 / Elevation of last SKYDIP point (deg)
END_ELD =              51.0018 / Elevation at end of observation
EQUINOX =                 2000 / Equinox of mean coordinate system
EXPOSED =                    0 / Exposure per pixel (seconds)
EXP_NO  =                    1 / Exposure number at end of observation
EXP_TIME=                1.024 / Exposure time for each basic measurement (sec)
E_PER_I =                    1 / Number of exposures per integration
FILTER  = '450N:850N'          / Filters used
FOCUS_SH=                   -1 / Shift between focus measurements (mm)
GAIN    =                   10 / Programmable gain
HSTEND  = '10:29:31.9986'      / HST at end of observation
HSTSTART= '10:25:34.0004'      / HST at start of observation
HUMEND  =                   36 / Humidity (%) at end of observation
HUMSTART=                   37 / Humidity (%) at start of observation
INSTRUME= 'SCUBA   '           / Name of instrument used
INT_NO  =                    8 / Integration number at end of observation
JIGL_CNT=                    9 / Number of offsets in jiggle pattern
JIGL_NAM= 'JCMTDATA_DIR:SQUEASY_9_2P0.JIG' / File containing jiggle offsets
J_PER_S =                    9 / Number of jiggles per switch position
J_REPEAT=                    1 / No. jiggle pattern repeats in a switch
LAT     = '+018:07:36.20'      / Object latitude
LAT-OBS =        19.8258323669 / Latitude of observatory (degrees)
LAT2    = 'not used'           / Object latitude at MJD2
LOCL_CRD= 'RB      '           / Local offset coordinate system
LONG    = '+004:28:44.42'      / Object longitude
LONG-OBS=        204.520278931 / East longitude of observatory (degrees)
LONG2   = 'not used'           / Object Longitude at MJD2
MAP_HGHT=                  180 / Height of rectangle to be mapped (arcsec)
MAP_PA  =                    0 / P.A. of map vertical, +ve towards +ve long
MAP_WDTH=                  180 / Width of rectangle to be mapped (arcsec)
MAP_X   =                    0 / Map X offset from telescope centre (arcsec)
MAP_Y   =                    0 / Map Y offset from telescope centre (arcsec)
MAX_EL  =                   -1 / Max elevation of sky-dip (deg)
MEANDEC =             18.23093 / 18:13:51.36292 = approx. mean Dec. (deg)
MEANRA  =             67.90475 / 67:54:17.08923 = approx. mean R.A. (deg)
MEAS_NO =                    1 / Measurement number at end of observation
MIN_EL  =                   -1 / Min elevation of sky-dip (deg)
MJD1    =                   -1 / Modified Julian day planet at RA,DEC
MJD2    =                   -1 / Modified Julian day planet at RA2,DEC2
MODE    = 'PHOTOM  '           / The type of observation
N_INT   =                    8 / No. integrations in the observation
N_MEASUR=                    1 / No. measurements in the observation
OBJECT  = 'hltau   '           / Name of object
OBJ_TYPE= 'UNKNOWN '           / Type of object
OBSDEF  = 'ss:odfsxphotom.t_hltau_100525' / The observation definition file
OBSERVER= 'jmk     '           / The name of the observer
PROJ_ID = 'scuba   '           / The project identification
RUN     =                  107 / Run number of observation
SAM_CRDS= 'AZ      '           / Coordinatesystem of sampling mesh
SAM_DX  =                   -1 / Sample spacing along scan direction (arcsec)
SAM_DY  =                   -1 / Sample spacing perp. to scan (arcsec)
SAM_MODE= 'JIGGLE  '           / Sampling method
SAM_PA  =                   -1 / Scan P.A. rel. to lat. line; 0=lat, 90=long
SCAN_REV=                    F / .TRUE. if alternate scans reverse direction
SPK_NSIG=                    0 / N sigmas from fit of spike threshold
SPK_RMVL=                    T / Automatic spike removal
SPK_WDTH=                    0 / Assumed width of spike
START_EL=                   -1 / Elevation of first SKYDIP point (deg)
STATE   = 'Terminating         :' / SCUCD state
STEND   = '7:15:14.03503'      / ST at end of observation
STRT_AZD=              274.523 / Azimuth at observation start (deg)
STRT_ELD=              51.8553 / Elevation at observation start (deg)
STSTART = '7:11:15.3828'       / ST at start of observation
SWTCH_MD= 'BMSW    '           / Switch mode of observation
SWTCH_NO=                    1 / Switch number at end of observation
S_PER_E =                    2 / Number of switch positions per exposure
TELESCOP= 'JCMT    '           / Name of telescope
TEL_OPER= 'kkp     '           / Telescope operator
UTDATE  = '1999:8:8'           / UT date of observation
UTEND   = '20:29:31.9986'      / UT at end of observation
UTSTART = '20:25:34.0004'      / UT at start of observation
VERSION =                  1.1 / SCUCD version
ALIGN_DX=             0.027268 / SMU tables X axis alignment offset
ALIGN_DY=             0.482489 / SMU tables Y axis alignment offset
ALIGN_X =             -2.98649 / SMU tables X axis
ALIGN_Y =              3.79108 / SMU tables Y axis
AZ_ERR  =                    0 / Error in the telescope azimuth
CHOPPING=                    T / SMU Chopper chopping state
EL_ERR  =                    0 / Error in the telescope elevation
FOCUS_DZ=            -0.504222 / SMU tables Z axis focus offset
FOCUS_Z =             -16.7523 / SMU tables Z axis
SEEING  =                0.525 / SAO atmospheric seeing
SEE_DATE= '9908061715'         / Date and time of SAO seeing
TAU_225 =                0.101 / CSO tau
TAU_DATE= '9908082019'         / Date and time of CSO tau
TAU_RMS =              3.0E-03 / CSO tau rms
UAZ     =             -5.09089 / User azimuth pointing offset
UEL     =              3.83056 / User elevation pointing offset
UT_DATE = ' 8 Aug 1999'        / UT date at start of observation
BAD_LIM =                   32 / No. spikes before quality set bad
CALIB_LG=                    6 / Lag of internal calibrator in samples
CALIB_PD=             42.66667 / Period of internal calibrator in samples
CHOP_LG =                    4 / Chop lag in samples
CHOP_PD =                   16 / Chop period in samples
CNTR_DU3=                    0 / Nasmyth dU3 coord of instrument centre
CNTR_DU4=                    0 / Nasmyth dU4 coord of instrument centre
ETATEL_1=                   -1 / Transmission of telescope
ETATEL_2=                   -1 / Transmission of telescope
ETATEL_3=                   -1 / Transmission of telescope
ETATEL_4=                   -1 / Transmission of telescope
ETATEL_5=                   -1 / Transmission of telescope
FILT_1  = '850     '           / Filter name
FILT_2  = '450     '           / Filter name
FILT_3  = 'not used'           / Filter name
FILT_4  = 'not used'           / Filter name
FILT_5  = 'not used'           / Filter name
FLAT    = 'jcmtdata_dir:lwswphot.dat' / Name of flat-field file
JIG_DSCD=                   -1 / No. samples discarded after jiggle movement
L_GD_BOL= 'not used'           / Bol. to whose value LW guard ring is set
L_GUARD =                    F / Long wave guard ring on or off
MEAS_BOL= 'LONG,SHORT,F12,F13' / Bolometers  actually measured in observation
N_BOLS  =                  130 / Number of bolometers selected
N_SUBS  =                    2 / Number of sub-instruments used
PHOT_BBF= 'NULL,H7,NULL,NULL,C14,NULL' / The bolometers on the source
PRE_DSCD=                    0 / No. of samples discarded before chop movement
PST_DSCD=                    0 / No. samples discarded after chop movement
REBIN   = 'LINEAR  '           / Rebinning method used by SCUIP
REF_ADC =                   -1 / A/D card of FLATFIELD reference bolometer
REF_CHAN=                   -1 / Channel of FLATFIELD reference bolometer
SAM_TIME=                  125 / A/D sample period in ticks (64musec)
SIMULATE=                    F / True if data is simulated
SKY     = 'jcmtdata_dir:skydip_startup.dat' / Name of sky opacities file
SUB_1   = 'LONG    '           / SCUBA instrument being used
SUB_2   = 'SHORT   '           / SCUBA instrument being used
SUB_3   = 'not used'           / SCUBA instrument being used
SUB_4   = 'not used'           / SCUBA instrument being used
SUB_5   = 'not used'           / SCUBA instrument being used
S_GD_BOL= 'not used'           / Bol. to whose value SW guard ring is set
S_GUARD =                    F / Short wave guard ring on or off
TAUZ_1  =                    0 / Zenith sky optical depth
TAUZ_2  =                    0 / Zenith sky optical depth
TAUZ_3  =                    0 / Zenith sky optical depth
TAUZ_4  =                    0 / Zenith sky optical depth
TAUZ_5  =                    0 / Zenith sky optical depth
T_AMB   =                   -1 / The ambient air temperature (K)
T_COLD_1=                   -1 / Effective temperature of cold load (K)
T_COLD_2=                   -1 / Effective temperature of cold load (K)
T_COLD_3=                   -1 / Effective temperature of cold load (K)
T_COLD_4=                   -1 / Effective temperature of cold load (K)
T_COLD_5=                   -1 / Effective temperature of cold load (K)
T_HOT   =                   -1 / The temperature of the hot load (K)
T_TEL   =                   -1 / The temperature of the telescope
USE_CAL =                    F / .TRUE. if dividing chop by cal before rebin
WAVE_1  =                  862 / Wavelength of map (microns)
WAVE_2  =                  442 / Wavelength of map (microns)
WAVE_3  =                   -1 / Wavelength of map (microns)
WAVE_4  =                   -1 / Wavelength of map (microns)
WAVE_5  =                   -1 / Wavelength of map (microns)
END
\end{verbatim}
\end{myquote}

\subsection{Scan Map}

This is a scan map FITS header. Note the presence of \texttt{MAP\_HGHT},
\texttt{MAP\_WDTH} and \texttt{MAP\_PA} to specify the map size (although
\scusoft\ does not use these since all map parameters are encoded into the
\texttt{RA1}, \texttt{RA2}, \texttt{DEC1} and \texttt{DEC2} arrays in the
\texttt{SCUCD} extension.

\begin{myquote}
\begin{verbatim}
ACCEPT  = 'not used'           / accept update; PROMPT, YES or NO
ALIGN_AX= 'not used'           / Alignment measurements in X or Y axis
ALIGN_SH=                   -1 / Distance between successive alignment offsets (
ALT-OBS =                 4092 / Height of observatory above sea level (metres)
AMEND   =             1.645293 / Airmass at end of observation
AMSTART =             1.602218 / Airmass at start of observation
APEND   =             628.3394 / Air pressure at end of observation (mbar)
APSTART =             628.3394 / Air pressure at start of observation (mbar)
ATEND   =             4.798538 / Air temp. at end of observation (C)
ATSTART =             4.847378 / Air temp. at start of observation (C)
BOLOMS  = 'LONG,SHORT'         / Names of bolometers measured
CALIBRTR=                    F / Internal calibrator is on or off
CAL_FRQ =             2.929688 / Calibrator frequency (Hz)
CENT_CRD= 'RB      '           / Centre coordinate system
CHOP_CRD= 'LO      '           / Chopper coordinate system
CHOP_FRQ=               7.8125 / Chopper frequency (Hz)
CHOP_FUN= 'SQUARE  '           / Chopper waveform
CHOP_PA =                   90 / Chopper P.A., 0 = in lat, 90 = in long
CHOP_THR=                   44 / Chopper throw (arcsec)
DATA_DIR= '19990808'           / Sub-directory where datafile was stored
DATA_KPT= 'DEMOD   '           / The type of data stored to disk
END_AZD =              214.199 / Azimuth at end of observation (deg)
END_EL  =                   -1 / Elevation of last SKYDIP point (deg)
END_ELD =               37.223 / Elevation at end of observation
EQUINOX =                 2000 / Equinox of mean coordinate system
EXPOSED =                0.128 / Exposure per pixel (seconds)
EXP_NO  =                   12 / Exposure number at end of observation
EXP_TIME=                0.128 / Exposure time for each basic measurement (sec)
E_PER_I =                   -1 / Number of exposures per integration
FILTER  = '450N:850N'          / Filters used
FOCUS_SH=                   -1 / Shift between focus measurements (mm)
GAIN    =                   10 / Programmable gain
HSTEND  = '21:46:43.99933'     / HST at end of observation
HSTSTART= '21:37:22.00058'     / HST at start of observation
HUMEND  =                   17 / Humidity (%) at end of observation
HUMSTART=                   16 / Humidity (%) at start of observation
INSTRUME= 'SCUBA   '           / Name of instrument used
INT_NO  =                    1 / Integration number at end of observation
JIGL_CNT=                   -1 / Number of offsets in jiggle pattern
JIGL_NAM= 'not used'           / File containing jiggle offsets
J_PER_S =                   -1 / Number of jiggles per switch position
J_REPEAT=                   -1 / No. jiggle pattern repeats in a switch
LAT     = '-024:22:00.00'      / Object latitude
LAT-OBS =        19.8258323669 / Latitude of observatory (degrees)
LAT2    = 'not used'           / Object latitude at MJD2
LOCL_CRD= 'RB      '           / Local offset coordinate system
LONG    = '+016:31:00.00'      / Object longitude
LONG-OBS=        204.520278931 / East longitude of observatory (degrees)
LONG2   = 'not used'           / Object Longitude at MJD2
MAP_HGHT=                  600 / Height of rectangle to be mapped (arcsec)
MAP_PA  =                    0 / P.A. of map vertical, +ve towards +ve long
MAP_WDTH=                  600 / Width of rectangle to be mapped (arcsec)
MAP_X   =                    0 / Map X offset from telescope centre (arcsec)
MAP_Y   =                    0 / Map Y offset from telescope centre (arcsec)
MAX_EL  =                   -1 / Max elevation of sky-dip (deg)
MEANDEC =            -24.46885 / -24:28:07.850189 = approx. mean Dec. (deg)
MEANRA  =             248.5074 / 248:30:26.53198 = approx. mean R.A. (deg)
MEAS_NO =                    1 / Measurement number at end of observation
MIN_EL  =                   -1 / Min elevation of sky-dip (deg)
MJD1    =                   -1 / Modified Julian day planet at RA,DEC
MJD2    =                   -1 / Modified Julian day planet at RA2,DEC2
MODE    = 'MAP     '           / The type of observation
N_INT   =                    1 / No. integrations in the observation
N_MEASUR=                    1 / No. measurements in the observation
OBJECT  = 'b31-22  '           / Name of object
OBJ_TYPE= 'UNKNOWN '           / Type of object
OBSDEF  = 'ss:scan_44ra.obs_b31x22_204433' / The observation definition file
OBSERVER= 'jmk     '           / The name of the observer
PROJ_ID = 'm99au34 '           / The project identification
RUN     =                   38 / Run number of observation
SAM_CRDS= 'NA      '           / Coordinatesystem of sampling mesh
SAM_DX  =                    3 / Sample spacing along scan direction (arcsec)
SAM_DY  =                   60 / Sample spacing perp. to scan (arcsec)
SAM_MODE= 'RASTER  '           / Sampling method
SAM_PA  =                 15.5 / Scan P.A. rel. to lat. line; 0=lat, 90=long
SCAN_REV=                    T / .TRUE. if alternate scans reverse direction
SPK_NSIG=                   -1 / N sigmas from fit of spike threshold
SPK_RMVL=                    F / Automatic spike removal
SPK_WDTH=                   -1 / Assumed width of spike
START_EL=                   -1 / Elevation of first SKYDIP point (deg)
STATE   = 'Terminating         :   STEP RECORD_DATA' / SCUCD state
STEND   = '18:30:20.72296'     / ST at end of observation
STRT_AZD=              212.051 / Azimuth at observation start (deg)
STRT_ELD=              38.3846 / Elevation at observation start (deg)
STSTART = '18:20:57.18613'     / ST at start of observation
SWTCH_MD= 'NOSW    '           / Switch mode of observation
SWTCH_NO=                    1 / Switch number at end of observation
S_PER_E =                    1 / Number of switch positions per exposure
TELESCOP= 'JCMT    '           / Name of telescope
TEL_OPER= 'kkp     '           / Telescope operator
UTDATE  = '1999:8:8'           / UT date of observation
UTEND   = '7:46:43.99933'      / UT at end of observation
UTSTART = '7:37:22.00058'      / UT at start of observation
VERSION =                  1.1 / SCUCD version
ALIGN_DX=             0.027268 / SMU tables X axis alignment offset
ALIGN_DY=             0.482489 / SMU tables Y axis alignment offset
ALIGN_X =             -2.02578 / SMU tables X axis
ALIGN_Y =              3.79108 / SMU tables Y axis
AZ_ERR  =                    0 / Error in the telescope azimuth
CHOPPING=                    T / SMU Chopper chopping state
EL_ERR  =                    0 / Error in the telescope elevation
FOCUS_DZ=            -0.361142 / SMU tables Z axis focus offset
FOCUS_Z =              -16.756 / SMU tables Z axis
SEEING  =                0.525 / SAO atmospheric seeing
SEE_DATE= '9908061715'         / Date and time of SAO seeing
TAU_225 =                0.063 / CSO tau
TAU_DATE= '9908080735'         / Date and time of CSO tau
TAU_RMS =              2.0E-03 / CSO tau rms
UAZ     =              1.42235 / User azimuth pointing offset
UEL     =              4.05503 / User elevation pointing offset
UT_DATE = ' 8 Aug 1999'        / UT date at start of observation
BAD_LIM =                   32 / No. spikes before quality set bad
CALIB_LG=                    6 / Lag of internal calibrator in samples
CALIB_PD=             42.66667 / Period of internal calibrator in samples
CHOP_LG =                    4 / Chop lag in samples
CHOP_PD =                   16 / Chop period in samples
CNTR_DU3=                    0 / Nasmyth dU3 coord of instrument centre
CNTR_DU4=                    0 / Nasmyth dU4 coord of instrument centre
ETATEL_1=                   -1 / Transmission of telescope
ETATEL_2=                   -1 / Transmission of telescope
ETATEL_3=                   -1 / Transmission of telescope
ETATEL_4=                   -1 / Transmission of telescope
ETATEL_5=                   -1 / Transmission of telescope
FILT_1  = '450     '           / Filter name
FILT_2  = '850     '           / Filter name
FILT_3  = 'not used'           / Filter name
FILT_4  = 'not used'           / Filter name
FILT_5  = 'not used'           / Filter name
FLAT    = 'jcmtdata_dir:lwswphot.dat' / Name of flat-field file
JIG_DSCD=                   -1 / No. samples discarded after jiggle movement
L_GD_BOL= 'not used'           / Bol. to whose value LW guard ring is set
L_GUARD =                    F / Long wave guard ring on or off
MEAS_BOL= 'not used'           / Bolometers  actually measured in observation
N_BOLS  =                  128 / Number of bolometers selected
N_SUBS  =                    2 / Number of sub-instruments used
PHOT_BBF= 'not_used       LL,C14,NULL' / The bolometers on the source
PRE_DSCD=                    0 / No. of samples discarded before chop movement
PST_DSCD=                    0 / No. samples discarded after chop movement
REBIN   = 'LINEAR  '           / Rebinning method used by SCUIP
REF_ADC =                   -1 / A/D card of FLATFIELD reference bolometer
REF_CHAN=                   -1 / Channel of FLATFIELD reference bolometer
SAM_TIME=                  125 / A/D sample period in ticks (64musec)
SIMULATE=                    F / True if data is simulated
SKY     = 'jcmtdata_dir:skydip_startup.dat' / Name of sky opacities file
SUB_1   = 'SHORT   '           / SCUBA instrument being used
SUB_2   = 'LONG    '           / SCUBA instrument being used
SUB_3   = 'not used'           / SCUBA instrument being used
SUB_4   = 'not used'           / SCUBA instrument being used
SUB_5   = 'not used'           / SCUBA instrument being used
S_GD_BOL= 'not used'           / Bol. to whose value SW guard ring is set
S_GUARD =                    F / Short wave guard ring on or off
TAUZ_1  =                    0 / Zenith sky optical depth
TAUZ_2  =                    0 / Zenith sky optical depth
TAUZ_3  =                    0 / Zenith sky optical depth
TAUZ_4  =                    0 / Zenith sky optical depth
TAUZ_5  =                    0 / Zenith sky optical depth
T_AMB   =                   -1 / The ambient air temperature (K)
T_COLD_1=                   -1 / Effective temperature of cold load (K)
T_COLD_2=                   -1 / Effective temperature of cold load (K)
T_COLD_3=                   -1 / Effective temperature of cold load (K)
T_COLD_4=                   -1 / Effective temperature of cold load (K)
T_COLD_5=                   -1 / Effective temperature of cold load (K)
T_HOT   =                   -1 / The temperature of the hot load (K)
T_TEL   =                   -1 / The temperature of the telescope
USE_CAL =                    F / .TRUE. if dividing chop by cal before rebin
WAVE_1  =                  442 / Wavelength of map (microns)
WAVE_2  =                  862 / Wavelength of map (microns)
WAVE_3  =                    0 / Wavelength of map (microns)
WAVE_4  =                    0 / Wavelength of map (microns)
WAVE_5  =                    0 / Wavelength of map (microns)
END
\end{verbatim}
\end{myquote}


\begin{landscape}
\section{Flatfield file format\label{app:ff}}

This is the flatfield file used for SCUBA at the time of writing:


\begin{myquote}
\begin{verbatim}
proc SET_BOLS
{ flat field data file written by SCUDR
{ Mon Jun 16 17:20:00 1997
{ observation run number 01
{      Name Type      dU3         dU4         Calib       Theta       A           B         Qual
 SETBOL G1  LONG      0.7830E+01 -0.7610E+02  0.1290E+01  0.4280E+01  0.7106E+01  0.7117E+01  0  66.0359490741  12 H8         
 SETBOL G2  LONG      0.2679E+02 -0.6138E+02  0.1086E+01  0.2370E+00  0.7387E+01  0.6795E+01  0  66.0482523148  12 H8         
 SETBOL G3  LONG      0.4785E+02 -0.4588E+02  0.1160E+01  0.1088E+00  0.7340E+01  0.6600E+01  0  66.0606250000  12 H8         
 SETBOL G4  LONG      0.7119E+02 -0.3377E+02  0.1080E+01  0.2125E-01  0.7499E+01  0.6642E+01  0  66.0729050926  12 H8         
 SETBOL G7  LONG     -0.1777E+02 -0.6524E+02  0.1089E+01  0.1843E+01  0.6980E+01  0.7279E+01  0  66.0852893519  12 H8         
 SETBOL G8  LONG      0.2193E+01 -0.4958E+02  0.1032E+01 -0.1701E+02 -0.4981E+01 -0.9338E+00  0  66.0979166667  12 H8         
 SETBOL G9  LONG      0.2291E+02 -0.3673E+02  0.1194E+01  0.9538E+01  0.7368E+01  0.6878E+01  0  66.1170254630  13 H8         
 SETBOL G10 LONG      0.4500E+02 -0.2298E+02  0.1187E+01  0.1407E+02  0.6766E+01  0.7412E+01  0  66.1295138889  13 H8         
 SETBOL G11 LONG      0.6850E+02 -0.1106E+02  0.1159E+01 -0.1926E+00  0.7540E+01  0.6623E+01  0  66.1417013889  13 H8         
 SETBOL G13 LONG     -0.3950E+02 -0.5281E+02  0.1039E+01 -0.3004E+01  0.7097E+01  0.6988E+01  0  66.1540625000  13 H8         
 SETBOL G14 LONG     -0.2157E+02 -0.3796E+02  0.1072E+01 -0.1470E+01  0.6906E+01  0.7204E+01  0  66.1663194444  13 H8         
 SETBOL G15 LONG      0.0510E+00 -0.2533E+02  0.1150E+01 -0.7553E-01  0.7276E+01  0.6851E+01  0  66.1786805556  13 H8         
 SETBOL G16 LONG      0.2174E+02 -0.1171E+02  0.9190E+00  0.6017E+03  0.7038E+01  0.7038E+01  0  66.1910532407  13 H8         
 SETBOL H1  LONG      0.4506E+02  0.1172E+01  0.1014E+01 -0.2208E+00  0.7407E+01  0.6631E+01  0  66.2103472222  14 H8         
 SETBOL H2  LONG      0.6772E+02  0.1328E+02  0.9490E+00 -0.1933E+00  0.7481E+01  0.6458E+01  0  66.2226736111  14 H8         
 SETBOL I12 LONG_DC   0.0000E+00  0.0000E+00  0.1000E+01  0.0000E+00  0.0000E+00  0.0000E+00  0   0.0000000000   0 NUL        
 SETBOL H4  LONG     -0.6336E+02 -0.3956E+02  0.1003E+01  0.5146E+02  0.7361E+01  0.6270E+01  0  66.2350000000  14 H8         
 SETBOL H5  LONG     -0.4256E+02 -0.2631E+02  0.9240E+00 -0.2044E+02  0.7017E+01  0.7017E+01  0  66.2473611111  14 H8         
 SETBOL H6  LONG     -0.2200E+02 -0.1264E+02  0.9870E+00 -0.3340E+01  0.7200E+01  0.6928E+01  0  66.2596296296  14 H8         
 SETBOL H7  LONG     -0.0000E+00  0.0000E+00  0.1000E+01  0.4839E+01  0.7151E+01  0.6991E+01  0  66.2720138889  14 H8         
 SETBOL H8  LONG      0.2262E+02  0.1231E+02  0.9270E+00 -0.3250E+00  0.7358E+01  0.6627E+01  0  66.5038773148  16 H8         
 SETBOL H9  LONG      0.4633E+02  0.2465E+02  0.1039E+01  0.1223E+02  0.7425E+01  0.6624E+01  0  66.2913541667  15 H8         
 SETBOL H10 LONG      0.7049E+02  0.3728E+02  0.1134E+01 -0.3192E+00  0.7535E+01  0.6472E+01  0  66.3037037037  15 H8         
 SETBOL H11 LONG     -0.6490E+02 -0.1333E+02  0.9580E+00  0.7245E+01  0.6991E+01  0.7178E+01  0  66.3161921296  15 H8         
 SETBOL H12 LONG     -0.4228E+02 -0.4700E+00  0.9160E+00  0.1142E+01  0.6943E+01  0.7179E+01  0  66.3286921296  15 H8         
 SETBOL H13 LONG     -0.2029E+02  0.1298E+02  0.1068E+01 -0.4099E+00  0.7245E+01  0.6759E+01  0  66.3409490741  15 H8         
 SETBOL H14 LONG      0.1250E+01  0.2529E+02  0.9650E+00  0.3552E+01  0.7087E+01  0.7087E+01  0  66.3532407407  15 H8         
 SETBOL H15 LONG      0.2519E+02  0.3783E+02  0.9380E+00 -0.2608E+05  0.7130E+01  0.7130E+01  0  66.3657407407  15 H8         
 SETBOL H16 LONG      0.4965E+02  0.4959E+02  0.9280E+00 -0.3867E+00  0.7513E+01  0.6436E+01  0  66.3780902778  15 H8         
 SETBOL I1  LONG     -0.6385E+02  0.1301E+02  0.9940E+00 -0.2134E+01  0.6749E+01  0.7152E+01  0  66.3975578704  16 H8         
 SETBOL I2  LONG     -0.4057E+02  0.2669E+02  0.1051E+01  0.4635E+03  0.7093E+01  0.7094E+01  0  66.4102546296  16 H8         
 SETBOL I3  LONG     -0.1803E+02  0.3894E+02  0.1135E+01  0.8583E+02  0.6732E+01  0.7409E+01  0  66.4227893519  16 H8         
 SETBOL I4  LONG      0.5930E+01  0.5154E+02  0.9620E+00 -0.5106E+01  0.6806E+01  0.7535E+01  0 192.7277662037   6 H7         
 SETBOL I5  LONG      0.2938E+02  0.6318E+02  0.9410E+00 -0.4519E+00  0.7451E+01  0.6512E+01  0  66.4477083333  16 H8         
 SETBOL I6  LONG     -0.6020E+02  0.4057E+02  0.9970E+00  0.3169E+01  0.7024E+01  0.7024E+01  0  66.4604976852  16 H8         
 SETBOL I7  LONG     -0.3669E+02  0.5404E+02  0.1054E+01 -0.6765E+02  0.7037E+01  0.7043E+01  0  66.4728125000  16 H8         
 SETBOL I8  LONG     -0.1578E+02  0.6389E+02  0.1065E+01 -0.2099E+02  0.6717E+01  0.7520E+01  0  66.4855787037  16 H8         
 SETBOL I9  LONG      0.9047E+01  0.7776E+02  0.1045E+01 -0.7307E-02  0.7092E+01  0.7032E+01  0  66.4977546296  16 H8         
												
{ short-wave array										
 SETBOL A1  SHORT     0.5708E+02 -0.2814E+02  0.1286E+01 -0.7433E-01  0.4129E+01  0.3690E+01  0  64.9820833333  43 C14        
 SETBOL A2  SHORT     0.4590E+02 -0.3591E+02  0.9311E+00 -0.3590E-01  0.4002E+01  0.3673E+01  0  64.9895138889  43 C14
 SETBOL A3  SHORT     0.3615E+02 -0.4146E+02  0.9499E+00  0.1720E+00  0.3987E+01  0.3677E+01  0  64.9969675926  43 C14
 SETBOL A4  SHORT     0.2464E+02 -0.4962E+02  0.1002E+01  0.2098E+00  0.3995E+01  0.3715E+01  0  65.0044212963  43 C14
 SETBOL A5  SHORT     0.1455E+02 -0.5628E+02  0.9585E+00 -0.4524E+00  0.3849E+01  0.3849E+01  0  65.0118055556  43 C14
 SETBOL A6  SHORT     0.3620E+01 -0.6315E+02  0.1115E+01 -0.1306E+01  0.3761E+01  0.3982E+01  0  65.0192824074  43 C14
 SETBOL A7  SHORT     0.5638E+02 -0.1661E+02  0.1000E+01 -0.2444E+00  0.4263E+01  0.3668E+01  0  65.0266898148  43 C14
 SETBOL A8  SHORT     0.4508E+02 -0.2315E+02  0.1126E+01 -0.1680E+00  0.3962E+01  0.3597E+01  0  65.0341319444  43 C14
 SETBOL A9  SHORT     0.3495E+02 -0.3012E+02  0.1089E+01 -0.2392E-01  0.3919E+01  0.3687E+01  0  65.5473379630   4 C14
 SETBOL A10 SHORT     0.2240E+02 -0.3762E+02  0.1107E+01 -0.8797E-02  0.3912E+01  0.3632E+01  0  65.5551041667   4 C14
 SETBOL A11 SHORT     0.1314E+02 -0.4396E+02  0.1203E+01 -0.2799E-02  0.3922E+01  0.3644E+01  0  65.5626851852   4 C14
 SETBOL A12 SHORT     0.2812E+01 -0.5070E+02  0.1067E+01 -0.6335E+00  0.3796E+01  0.3796E+01  0  65.5701967593   4 C14
 SETBOL A13 SHORT    -0.7590E+01 -0.5808E+02  0.9114E+00  0.3143E+02  0.3849E+01  0.3851E+01  0  65.5777777778   4 C14
 SETBOL A14 SHORT     0.5559E+02 -0.4683E+01  0.9639E+00 -0.1907E+00  0.4101E+01  0.3452E+01  0  65.5855555556   4 C14
 SETBOL A15 SHORT     0.4502E+02 -0.1081E+02  0.1068E+01 -0.2673E+00  0.3985E+01  0.3569E+01  0  65.5932754630   4 C14
 SETBOL A16 SHORT     0.3296E+02 -0.1807E+02  0.1170E+01 -0.1050E+00  0.3974E+01  0.3628E+01  0  65.6006944444   4 C14
 SETBOL B1  SHORT     0.2224E+02 -0.2456E+02  0.1102E+01 -0.9115E-01  0.3948E+01  0.3680E+01  0  65.6129629630   5 C14
 SETBOL B2  SHORT     0.1165E+02 -0.3120E+02  0.1157E+01 -0.6530E-02  0.3921E+01  0.3681E+01  0  65.6204513889   5 C14
 SETBOL B3  SHORT     0.2000E+01 -0.3855E+02  0.9660E+00  0.2064E-01  0.3904E+01  0.3688E+01  0  65.6279282407   5 C14
 SETBOL B4  SHORT    -0.9089E+01 -0.4490E+02  0.1076E+01  0.2914E-01  0.3910E+01  0.3702E+01  0  65.6354513889   5 C14
 SETBOL B5  SHORT    -0.1951E+02 -0.5160E+02  0.9298E+00  0.7764E+01  0.3756E+01  0.3901E+01  0  65.6433101852   5 C14
 SETBOL B6  SHORT     0.5585E+02  0.7780E+01  0.1088E+01 -0.2603E+00  0.4125E+01  0.3553E+01  0  72.3519444444  58 C14
 SETBOL B7  SHORT     0.4489E+02  0.9322E+00  0.1020E+01 -0.3361E+05  0.3717E+01  0.3717E+01  0  65.6584606481   5 C14
 SETBOL B8  SHORT     0.3283E+02 -0.5420E+01  0.1112E+01  0.5543E+00  0.3792E+01  0.3789E+01  0  65.6659953704   5 C14
 SETBOL B9  SHORT     0.2169E+02 -0.1181E+02  0.1019E+01  0.1420E+01  0.3657E+01  0.3946E+01  0  65.3514120370   1 C14
 SETBOL B10 SHORT     0.1166E+02 -0.1921E+02  0.1063E+01  0.1457E+01  0.3861E+01  0.3776E+01  0  65.3589467593   1 C14
 SETBOL B11 SHORT     0.7480E+00 -0.2610E+02  0.1017E+01  0.6716E+00  0.3820E+01  0.3829E+01  0  65.3664004630   1 C14
 SETBOL B12 SHORT    -0.1013E+02 -0.3228E+02  0.1085E+01 -0.1378E+00  0.3904E+01  0.3757E+01  0  65.3738888889   1 C14
 SETBOL B13 SHORT    -0.2043E+02 -0.3952E+02  0.1112E+01 -0.2830E+00  0.3875E+01  0.3700E+01  0  65.3814351852   1 C14
 SETBOL B14 SHORT    -0.2907E+02 -0.4586E+02  0.9829E+00 -0.2085E+00  0.3860E+01  0.3679E+01  0  65.3888657407   1 C14
 SETBOL B15 SHORT     0.5681E+02  0.2020E+02  0.9500E+00 -0.2601E+00  0.4107E+01  0.3470E+01  0  72.3643634259  58 C14
 SETBOL B16 SHORT     0.4512E+02  0.1323E+02  0.1090E+01 -0.3063E+00  0.3952E+01  0.3450E+01  0  72.3767245370  58 C14
 SETBOL C1  SHORT     0.3315E+02  0.7096E+01  0.1004E+01 -0.2430E+00  0.3927E+01  0.3510E+01  0  65.4163888889   2 C14
 SETBOL C2  SHORT     0.2194E+02  0.7400E+00  0.1066E+01 -0.2254E+00  0.3982E+01  0.3568E+01  0  65.4239004630   2 C14
 SETBOL C3  SHORT     0.1091E+02 -0.6110E+01  0.1044E+01  0.2300E+00  0.3858E+01  0.3727E+01  0  65.4313425926   2 C14
 SETBOL C4  SHORT     0.2756E+00 -0.1215E+02  0.1254E+01 -0.1419E+01  0.3763E+01  0.3830E+01  0  65.4390509259   2 C14
 SETBOL C5  SHORT    -0.1075E+02 -0.1977E+02  0.1111E+01  0.5111E+00  0.3827E+01  0.3827E+01  0  65.4468634259   2 C14
 SETBOL C6  SHORT    -0.2075E+02 -0.2666E+02  0.1075E+01 -0.2159E+00  0.3901E+01  0.3764E+01  0  65.4543634259   2 C14
 SETBOL C7  SHORT    -0.3146E+02 -0.3298E+02  0.1287E+01  0.7394E+01  0.3743E+01  0.3888E+01  0  65.4617824074   2 C14
 SETBOL C8  SHORT    -0.4124E+02 -0.4042E+02  0.9000E+00  0.4200E+01  0.3704E+01  0.3741E+01  0  65.4693287037   2 C14
 SETBOL C9  SHORT     0.5849E+02  0.3141E+02  0.1326E+01  0.3450E+00  0.3868E+01  0.3868E+01  0  72.3890740741  58 C14
 SETBOL C10 SHORT     0.4631E+02  0.2581E+02  0.1066E+01  0.1101E+01  0.3526E+01  0.4008E+01  0  65.4891782407   3 C14
 SETBOL C11 SHORT     0.3388E+02  0.2000E+02  0.1229E+01  0.4242E+00  0.3780E+01  0.3780E+01  0  65.4966782407   3 C14
 SETBOL C12 SHORT     0.2248E+02  0.1267E+02  0.1131E+01  0.1323E+01  0.3537E+01  0.3973E+01  0  65.5042824074   3 C14
 SETBOL C13 SHORT     0.1116E+02  0.6479E+01  0.1062E+01  0.5274E+00  0.3760E+01  0.3760E+01  0  65.5119328704   3 C14
 SETBOL C14 SHORT     0.0000E+00 -0.0000E+00  0.1000E+01 -0.2339E+00  0.3860E+01  0.3595E+01  0  72.4200810185  58 C14
 SETBOL C15 SHORT    -0.1116E+02 -0.6479E+01  0.1130E+01  0.1179E+01  0.3666E+01  0.3898E+01  0  65.5270949074   3 C14
 SETBOL C16 SHORT    -0.2215E+02 -0.1332E+02  0.1101E+01 -0.3823E+00  0.3928E+01  0.3760E+01  0  65.5345717593   3 C14
 SETBOL D1  SHORT    -0.3139E+02 -0.2045E+02  0.1157E+01  0.4295E+01  0.3729E+01  0.3857E+01  0  65.6783796296   6 C14
 SETBOL D2  SHORT    -0.4220E+02 -0.2734E+02  0.1175E+01  0.2061E+01  0.3803E+01  0.3755E+01  0  65.6859375000   6 C14
 SETBOL D3  SHORT    -0.5298E+02 -0.3271E+02  0.1142E+01 -0.2841E+00  0.3731E+01  0.3861E+01  0  65.6934606481   6 C14
 SETBOL D4  SHORT     0.4815E+02  0.3794E+02  0.9828E+00 -0.2879E+01  0.3816E+01  0.3816E+01  0  65.7009027778   6 C14
 SETBOL D5  SHORT     0.3558E+02  0.3207E+02  0.1027E+01 -0.2057E+01  0.3540E+01  0.3993E+01  0  65.7086342593   6 C14
 SETBOL D6  SHORT     0.2340E+02  0.2565E+02  0.1080E+01  0.2692E+01  0.3989E+01  0.3572E+01  0  65.7160416667   6 C14
 SETBOL D7  SHORT     0.1184E+02  0.1956E+02  0.1258E+01 -0.1224E+01  0.3779E+01  0.3779E+01  0  65.7237500000   6 C14
 SETBOL D8  SHORT     0.1100E+01  0.1301E+02  0.1052E+01  0.1119E+01  0.3643E+01  0.3878E+01  0  65.7313888889   6 C14
 SETBOL D9  SHORT    -0.1030E+02  0.6303E+01  0.1011E+01  0.1506E+02  0.3734E+01  0.3770E+01  0  65.7436805556   7 C14
 SETBOL D10 SHORT    -0.2141E+02 -0.4385E+00  0.1088E+01 -0.2669E+00  0.3951E+01  0.3611E+01  0  65.7512615741   7 C14
 SETBOL D11 SHORT    -0.3244E+02 -0.6844E+01  0.1173E+01  0.1050E+01  0.3711E+01  0.3913E+01  0  65.7587268519   7 C14
 SETBOL D12 SHORT    -0.4299E+02 -0.1380E+02  0.1172E+01 -0.8284E+00  0.3793E+01  0.3728E+01  0  65.7662037037   7 C14
 SETBOL D13 SHORT    -0.5295E+02 -0.2078E+02  0.9500E+00  0.5212E+01  0.3815E+01  0.3733E+01  0  65.7740393519   7 C14
 SETBOL D14 SHORT     0.3682E+02  0.4498E+02  0.9517E+00 -0.6182E+00  0.4028E+01  0.3525E+01  0  65.7814583333   7 C14
 SETBOL D15 SHORT     0.2543E+02  0.3744E+02  0.1036E+01 -0.6247E+01  0.3738E+01  0.3775E+01  0  65.7891203704   7 C14
 SETBOL D16 SHORT     0.1301E+02  0.3179E+02  0.1248E+01  0.2862E+01  0.3836E+01  0.3636E+01  0  65.7968518519   7 C14
 SETBOL E1  SHORT     0.1499E+01  0.2630E+02  0.1153E+01 -0.5981E+01  0.3772E+01  0.3773E+01  0  65.8090972222   8 C14
 SETBOL E2  SHORT    -0.9694E+01  0.1926E+02  0.1098E+01  0.6842E+01  0.3768E+01  0.3768E+01  0  72.4015509259  58 C14
 SETBOL E3  SHORT    -0.2096E+02  0.1302E+02  0.1142E+01 -0.3893E+00  0.3940E+01  0.3670E+01  0  65.8240972222   8 C14
 SETBOL E4  SHORT    -0.3153E+02  0.6590E+01  0.1126E+01  0.3139E+00  0.3796E+01  0.3796E+01  0  65.8315740741   8 C14
 SETBOL E5  SHORT    -0.4362E+02 -0.8783E+00  0.1026E+01  0.6693E+02  0.3609E+01  0.3792E+01  0  65.8392939815   8 C14
 SETBOL E6  SHORT    -0.5395E+02 -0.7067E+01  0.1020E+01  0.9356E+00  0.3635E+01  0.3792E+01  0  65.8469444444   8 C14
 SETBOL E7  SHORT     0.2731E+02  0.5184E+02  0.9742E+00  0.7678E+02  0.3756E+01  0.3897E+01  0  65.8546759259   8 C14
 SETBOL E8  SHORT     0.1511E+02  0.4443E+02  0.1204E+01 -0.5818E+00  0.4007E+01  0.3567E+01  0  65.8621875000   8 C14
 SETBOL E9  SHORT     0.3500E+01  0.3973E+02  0.1020E+01 -0.2097E+01  0.3605E+01  0.3976E+01  0  65.8745717593   9 C14
 SETBOL E10 SHORT    -0.8784E+01  0.3302E+02  0.1197E+01 -0.4411E+01  0.3800E+01  0.3767E+01  0  72.4138888889  58 C14
 SETBOL E11 SHORT    -0.1987E+02  0.2656E+02  0.1005E+01  0.4240E+01  0.3643E+01  0.3936E+01  0  65.8894791667   9 C14
 SETBOL E12 SHORT    -0.3045E+02  0.2024E+02  0.1067E+01  0.1210E+02  0.3940E+01  0.3664E+01  0  65.8971643519   9 C14
 SETBOL E13 SHORT    -0.4249E+02  0.1245E+02  0.9850E+00  0.8132E+00  0.3809E+01  0.3739E+01  0  65.9045254630   9 C14
 SETBOL E14 SHORT    -0.5245E+02  0.6274E+01  0.1177E+01 -0.1508E+01  0.3760E+01  0.3760E+01  0  65.9120717593   9 C14
 SETBOL E15 SHORT     0.1652E+02  0.5893E+02  0.9800E+00  0.3954E+01  0.3587E+01  0.4123E+01  0  65.9196643519   9 C14
 SETBOL E16 SHORT     0.3524E+01  0.5163E+02  0.1174E+01  0.9079E+00  0.3618E+01  0.4055E+01  0  65.9270949074   9 C14
 SETBOL F1  SHORT    -0.7665E+01  0.4592E+02  0.1168E+01  0.3404E+01  0.3775E+01  0.3842E+01  0  65.9394328704  10 C14
 SETBOL F2  SHORT    -0.1817E+02  0.3963E+02  0.1020E+01 -0.5761E+00  0.3991E+01  0.3608E+01  0  65.9469560185  10 C14
 SETBOL F3  SHORT    -0.3008E+02  0.3403E+02  0.1295E+01 -0.5724E+00  0.3977E+01  0.3659E+01  0  65.9544675926  10 C14
 SETBOL F4  SHORT    -0.4123E+02  0.2619E+02  0.1244E+01  0.8881E+00  0.3627E+01  0.3909E+01  0  65.9621064815  10 C14
 SETBOL F5  SHORT    -0.5155E+02  0.2029E+02  0.1069E+01  0.1093E+00  0.3760E+01  0.3760E+01  0  65.9695601852  10 C14
 SETBOL F6  SHORT     0.6895E+01  0.6560E+02  0.1245E+01  0.4535E-01  0.3981E+01  0.3981E+01  0  65.9770949074  10 C14
 SETBOL F7  SHORT    -0.4842E+01  0.6013E+02  0.9947E+00 -0.7757E+00  0.4069E+01  0.3691E+01  0  65.9846527778  10 C14
 SETBOL F8  SHORT    -0.1690E+02  0.5287E+02  0.1017E+01  0.8808E+00  0.3724E+01  0.4061E+01  0  65.9920138889  10 C14
 SETBOL F9  SHORT    -0.2931E+02  0.4752E+02  0.9485E+00 -0.5179E+01  0.3703E+01  0.3979E+01  0  66.0041087963  11 C14
 SETBOL F10 SHORT    -0.4009E+02  0.3977E+02  0.9474E+00  0.8256E+00  0.3667E+01  0.3968E+01  0  66.0116203704  11 C14
 SETBOL F11 SHORT    -0.5045E+02  0.3338E+02  0.1169E+01  0.7934E+00  0.3701E+01  0.3940E+01  0  66.0190625000  11 C14
 SETBOL I11 SHORT_DC  0.0000E+00  0.0000E+00  0.1000E+01  0.0000E+00  0.0000E+00  0.0000E+00  1   0.0000000000   0 NUL
{ P1100 bolometer										     		    
 SETBOL I10 P1100    -0.5465E+02  0.7116E+02  0.1000E+01 -0.1453E+03  0.1002E+02  0.9003E+01  1  69.6688541667  65 NUL
 SETBOL I15 P1100_DC  0.0000E+00  0.0000E+00  0.1000E+01  0.0000E+00  0.0000E+00  0.0000E+00  1   0.0000000000   0 NUL
{ P1350 bolometer										     		    
 SETBOL G12 P1350    -0.5827E+02 -0.7068E+02  0.1000E+01 -0.8953E+00  0.9681E+01  0.1101E+02  0  69.7100231481  66 NUL
 SETBOL I14 P1350_DC  0.0000E+00  0.0000E+00  0.1000E+01  0.0000E+00  0.0000E+00  0.0000E+00  0   0.0000000000   0 NUL
{ P2000 bolometer										     		    
 SETBOL G6  P2000     0.1000E+03  0.2375E+01  0.1000E+01 -0.1447E+00  0.1735E+02  0.1374E+02  0  69.7901157407  71 NUL
 SETBOL I13 P2000_DC  0.0000E+00  0.0000E+00  0.1000E+01  0.0000E+00  0.0000E+00  0.0000E+00  0   0.0000000000   0 NUL
end proc
\end{verbatim}
\end{myquote}
\end{landscape}

\section{The DREAM2SURF import task}

This section lists the complete source code for the \texttt{dream2surf} import 
task. This is an example of an import task for JIGGLE map type observations.

\begin{myquote}
\begin{verbatim}
      SUBROUTINE DREAM2SURF ( STATUS )
*+
*  Name:
*     DREAM2SURF

*  Purpose:
*     Convert DREAM output data into SURF data format

*  Language:
*     Starlink Fortran 77

*  Type of Module:
*     ADAM A-task

*  Invocation:
*     CALL DREAM2SURF ( STATUS )

*  Arguments:
*     STATUS = INTEGER (Given and Returned)
*        The global status
 
*  Description:
*     This routine reads in a DREAM format file and converts it
*     into a file format that can be processed by SURF.
*     This new file will have been 'REDUCE_SWITCH'ed and 
*     'FLATFIELD'ed.

*  Usage:
*     dream2surf in out

*  ADAM Parameters:
*     FLATFILE = CHAR (Read)
*        File containing the flatfield description
*     FSCYCLE = INTEGER (Read)
*        First cycle number to read from file
*     IN = FILE (Read)
*        The name of the DREAM input file to open.
*     MSG_FILTER = CHAR (Read)
*         Message filter level. Default is NORM.
*     NRCYCLE = INTEGER (Read)
*         Number of cycles to read from file
*     OUT = NDF (Write)
*        The NDF output file.

*  Authors:
*     TIMJ: T. Jenness (timj@jach.hawaii.edu)
*     GREVE: H.W. van Someren Greve (greve@nfra.nl)

*  Copyright:
*     Copyright (C) 1998 Particle Physics and Astronomy
*     Research Council. All Rights Reserved.


*  History:
*     Revision 1.4  1999/08/03 19:32:28  timj
*     Add copyright message to header.
*
*     Revision 1.3  1998/06/24 19:26:03  timj
*     Finally get jiggle pattern to work for new format sol files.
*
*     Revision 1.2  1998/06/19 19:28:26  timj
*     First released version
*
*     Revision 1.1  1998/05/14 20:41:03  timj
*     Initial revision
*

*  Bugs:
*     {note_any_bugs_here}
 
*-
 
*  Type Definitions:
      IMPLICIT NONE
 
*  Global constants:
      INCLUDE 'SAE_PAR'       ! Starlink status
      INCLUDE 'DAT_PAR'       ! DAT__ constants
      INCLUDE 'PRM_PAR'       ! For VAL__
      INCLUDE 'MSG_PAR'       ! For MSG_

      INCLUDE 'SCU_SOL'       ! Description of DREAM header file
      INCLUDE 'SURF_PAR'      ! Standard SCUBA include

*  COMMON data
      COMMON SOLPA            ! DREAM common block

*  Status:
      INTEGER STATUS

*  External references
      INTEGER CHR_LEN
      EXTERNAL CHR_LEN

*  Local Constants:
      INTEGER       DREAM__NBYTES
      PARAMETER     (DREAM__NBYTES = 4)  ! Number of bytes per DREAM record
      INTEGER       SRECSIZE                 
      PARAMETER     (SRECSIZE = 4096)          ! Small record size in bytes
      INTEGER       HEADER
      PARAMETER     (HEADER   = 5120)          ! Nr of I4 in Header records
      CHARACTER*(10) TSKNAME                   ! Task name
      PARAMETER (TSKNAME = 'DREAM2SURF')

      DOUBLE PRECISION LAT_OBS           ! Latitude of JCMT (degrees)
      PARAMETER (LAT_OBS = 19.8258323669)
      DOUBLE PRECISION LONG_OBS          ! Longitude of JCMT (degrees)
      PARAMETER (LONG_OBS = 204.520278931)
       


*  Local variables:
      INTEGER ADC             ! A to D number
      REAL             BOL_CALB (SCUBA__NUM_CHAN, SCUBA__NUM_ADC)
                                                 ! bolometer flatfield factors
      DOUBLE PRECISION BOL_DAY (SCUBA__NUM_CHAN, SCUBA__NUM_ADC)
                                                 ! time and day number on which
                                                 ! the bolometer flatfield was
                                                 ! measured
      REAL             BOL_DU3 (SCUBA__NUM_CHAN, SCUBA__NUM_ADC)
                                                 ! Nasmyth dU3 coords of 
                                                 ! bolometers
      REAL             BOL_DU4 (SCUBA__NUM_CHAN, SCUBA__NUM_ADC)
                                                 ! Nasmyth dU4 coords of 
                                                 ! bolometers
      INTEGER          BOL_QUAL (SCUBA__NUM_CHAN, SCUBA__NUM_ADC)
                                                 ! bolometer flatfield quality
      CHARACTER*20     BOL_REF (SCUBA__NUM_CHAN, SCUBA__NUM_ADC)
                                                 ! flatfield reference
                                                 ! bolometers
      REAL             BOL_RTEMP (SCUBA__NUM_CHAN, SCUBA__NUM_ADC)
                                                 ! scratch real bolometer data
      INTEGER          BOL_RUN (SCUBA__NUM_CHAN, SCUBA__NUM_ADC)
                                                 ! run number on which the
                                                 ! bolometer flatfield was
                                                 ! measured
      CHARACTER*20     BOL_TYPE (SCUBA__NUM_CHAN, SCUBA__NUM_ADC)
                                                 ! bolometer types
      INTEGER BOL_ADC ( SCUBA__NUM_ADC * SCUBA__NUM_CHAN ) ! ADC numbers
      INTEGER BOL_CHAN ( SCUBA__NUM_ADC * SCUBA__NUM_CHAN ) ! Channel numbers
      CHARACTER *10 CENT_CRD  ! Centre coordinate system
      INTEGER CHAN            ! Channel number
      CHARACTER * 80 CTEMP    ! Scratch string
      DOUBLE PRECISION DEC    ! Declination
      INTEGER DEMBNDS ( 4 )   ! Dimensions of DEM_PNTR
      INTEGER DEM_PNTR        ! Pointer to DEM_PNTR extension
      INTEGER DIM ( 4 )       ! Dimensions of an array
      DOUBLE PRECISION DLST   ! Time per cycle
      INTEGER DREAM_PTR       ! Pointer to DREAM input data
      DOUBLE PRECISION DTEMP  ! Scratch double
      CHARACTER * (80) FITS (SCUBA__MAX_FITS)  ! FITS extension
      INTEGER LBND ( 2 )      ! Lower bounds of output NDF
      INTEGER LUN             ! Logical unit number of input file
      INTEGER ERR             ! Error from DREAM system
      INTEGER FD              ! File descriptor of input
      CHARACTER*80 FLATFILE   ! Flatfield file name
      INTEGER FSCYCLE         ! First cycle to read from input file
      INTEGER I               ! Loop integer
      INTEGER IERR            ! For VEC_
      INTEGER IPOSN           ! Position in string
      INTEGER ITEMP           ! Scratch integer
      REAL    JIGL_X ( SCUBA__MAX_JIGGLE ) ! X jiggle positions
      REAL    JIGL_Y ( SCUBA__MAX_JIGGLE ) ! Y jiggle positions
      DOUBLE PRECISION LONGITUDE ! Longitude degrees west
      DOUBLE PRECISION LST    ! LST of observation
      DOUBLE PRECISION LST_STRT ! Start LST of observation (decimal hours)
      INTEGER MAXCYC          ! Maximum number of cycles that can be read
      DOUBLE PRECISION MJD    ! MJD of observation
      INTEGER NBYTES          ! Number of bytes per data record
      INTEGER NDIM            ! Number of dimensions in an array
      INTEGER NERR            ! For VEC_
      INTEGER NRCYCLE         ! Number of cycles to read
      INTEGER NREC            ! Number of records per cycle
      INTEGER NSVAL           ! Number of solved values per cycle
      INTEGER N_FITS          ! Number of FITS components
      INTEGER OFFSET          ! Y offset in output array
      INTEGER OUT_A_PTR       ! Pointer to axis
      INTEGER OUT_NDF         ! Output NDF identifier
      INTEGER OUT_PTR         ! Pointer to output data (NDF)
      INTEGER OUT_QUAL_PTR    ! Pointer to quality array
      INTEGER OUT_VAR_PTR     ! Pointer to variance array
      CHARACTER*(DAT__SZLOC) OUT_SCUCD_LOC ! Locator to SCUCD extension
      CHARACTER*(DAT__SZLOC) OUT_SCUBA_LOC ! Locator to SCUBA extension
      CHARACTER*(DAT__SZLOC) OUT_REDS_LOC ! Locator to REDS extension
      CHARACTER*(DAT__SZLOC) OUT_FIG_LOC ! Locator to FIGARO extension
      CHARACTER*(DAT__SZLOC) OUT_FITS_LOC ! Locator to FITS extension
      DOUBLE PRECISION RA     ! Right ascension of source
      INTEGER RECIN           ! Current input record
      INTEGER RECORD          ! Record number in input file
      INTEGER RECSS           ! Record size in words
      CHARACTER *20 RUNNO     ! Run number in 0 padded form (eg 0015) 
      CHARACTER * 80 STEMP    ! Scratch string
      INTEGER UBND ( 2 )      ! Upper bounds of output NDF

*-

      IF (STATUS .NE. SAI__OK) RETURN

*     Set the MSG output level (for use with MSG_OUTIF)
 
      CALL MSG_IFGET('MSG_FILTER', STATUS)


*     Read in the DREAM file (currently assume that 
*     user will supply full path or use DREAM_OUT env var.

      CALL RIO_ASSOC('IN', 'READ', 'UNFORMATTED', SRECSIZE, FD,
     :     STATUS)

*     Since we are using the DREAM I/O routines we need to get
*     the file unit number
      CALL FIO_UNIT(FD, LUN, STATUS)

*     Read header into common block

      IF (STATUS .EQ. SAI__OK) THEN
         RECORD = 1             ! Record number 1
         RECSS = SRECSIZE / DREAM__NBYTES ! Output record size in words

         CALL DISK_IO (2, LUN, RECSS, RECORD, SOLPA, HEADER, ERR)

         IF (ERR .NE. 0) THEN
            STATUS = SAI__ERROR
            CALL MSG_SETC('TSK', TSKNAME)
            CALL MSG_SETI('ERR', ERR)
            CALL MSG_SETI('REC', RECORD)
            CALL ERR_REP(' ','^TSK: Error ^ERR in reading record'//
     :           ' ^REC', STATUS)
         END IF
      END IF

*     Ask for the first cycle to select
      CALL PAR_GDR0I('FSCYCLE', 1, 1, r_Ncycle, .FALSE., FSCYCLE,
     :     STATUS)

*     Ask for the number of cycles to be selected
      MAXCYC = r_Ncycle - FSCYCLE + 1
      CALL PAR_GDR0I('NRCYCLE', MAXCYC, 1, MAXCYC, .FALSE., 
     :     NRCYCLE, STATUS)

*     Now we can start doing things.
*     Calculate total number of pixel values per cycle

      NSVAL = r_Nbol * MAX_PIX

      IF (NSVAL .LE. 0) THEN

         IF (STATUS .EQ. SAI__OK) THEN
            STATUS = SAI__ERROR
            CALL MSG_SETC('TSK', TSKNAME)
            CALL MSG_SETI('NPT', NSVAL)
            CALL ERR_REP(' ','^TSK: Number of observed points'//
     :           'is too small (^NPT)', STATUS)
         END IF
      END IF

*     Start NDF
      CALL NDF_BEGIN

*     Open an NDF file
*     Default file name constructed from the UT date

*     Year
      CALL CHR_ITOC(GR_YY, STEMP, ITEMP)
      IPOSN = CHR_LEN(STEMP)

*     Month (pad with leading zero)
      IF (GR_MN .LT. 10) CALL CHR_APPND('0', STEMP, IPOSN)
      CALL CHR_ITOC(GR_MN, CTEMP, ITEMP)
      CALL CHR_APPND(CTEMP, STEMP, IPOSN)

*     Day
      IF (GR_DD .LT. 10) CALL CHR_APPND('0', STEMP, IPOSN)
      CALL CHR_ITOC(GR_DD, CTEMP, ITEMP)
      CALL CHR_APPND(CTEMP,STEMP,IPOSN)

*     Find observation number from filename
*     Assumes string of form xxx_NNNN

      ITEMP = CHR_LEN ( SOLVE_DATA )
      RUNNO = SOLVE_DATA(ITEMP-3:)

*     Append an _
      CALL CHR_APPND('_', STEMP, IPOSN)

*     Append run number to out
      CALL CHR_APPND(RUNNO, STEMP, IPOSN)

*     Append DREAM signature
      CALL CHR_APPND('_drm', STEMP, IPOSN)


*     Set default
      CALL PAR_DEF0C('OUT', STEMP, STATUS)

*     Bounds of NDF
      LBND(1) = 1
      LBND(2) = 1
      UBND(1) = R_NBOL
      UBND(2) = NPIX * NRCYCLE

      CALL NDF_CREAT('OUT', '_REAL', 2, LBND, UBND, OUT_NDF, STATUS)

*     Map the output data array (plus the other arrays for SURF
*     compatibility)
      CALL NDF_MAP(OUT_NDF, 'QUALITY','_UBYTE', 'WRITE/ZERO',
     :     OUT_QUAL_PTR, ITEMP, STATUS)
      CALL NDF_MAP(OUT_NDF, 'DATA', '_REAL', 'WRITE', OUT_PTR,
     :     ITEMP, STATUS)
      CALL NDF_MAP(OUT_NDF, 'VARIANCE', '_REAL', 'WRITE/ZERO',
     :     OUT_VAR_PTR, ITEMP, STATUS)


*     Get some memory to store the data from each cycle
*     We are using 4 byte words but define in parameter

      NBYTES = MAX(DREAM__NBYTES * NSVAL, SRECSIZE)
      CALL PSX_MALLOC(NBYTES, DREAM_PTR, STATUS)

*     Loop over cycles

*     Calculate number of records per cycle
      NREC = (NSVAL + RECSS - 1 ) / RECSS

*     Loop
      DO I = FSCYCLE, NRCYCLE

         IF (STATUS .EQ. SAI__OK) THEN

*     Determine inpur record number
            RECIN = HEAD_S / RECSS + (I * NREC)

*     Read pixel intensities into memory
            CALL DISK_IO (2, LUN, RECSS, RECIN, %VAL(DREAM_PTR),
     :           NSVAL, ERR)

            IF (ERR .NE. 0) THEN
               print *,'in if ',ERR, STATUS
               STATUS = SAI__ERROR
               CALL MSG_SETC('TSK', TSKNAME)
               CALL MSG_SETI('ERR', ERR)
               CALL MSG_SETI('REC', RECIN)
               CALL ERR_REP(' ','^TSK: Error ^ERR in reading record'//
     :              ' ^REC', STATUS)
            END IF

*     Calculate current offset position in output array (time axis)
            OFFSET = (I - FSCYCLE) * NPIX + 1

*     Now we need to write this data to the NDF file
            CALL DREAM_DATA_TO_SURF_DATA(R_NBOL, MAX_PIX,
     :           NPIX * NRCYCLE,
     :           OFFSET, INT_POS, %VAL(DREAM_PTR), %VAL(OUT_PTR),
     :           STATUS)

         END IF

      END DO

*     Free the memory assoicated with each cycle
      CALL PSX_FREE(DREAM_PTR, STATUS)

*     Close the input file
      CALL RIO_CLOSE(FD, STATUS)

*     Unmap the data arrays
      CALL NDF_UNMAP(OUT_NDF, '*', STATUS)

*     Axis information is simply integration number.
*     Steal this from REDUCE_SWITCH
*
*     Axis 1: bolometers   2: Data

*     Deal with BOLOMETER axis
 
      CALL NDF_AMAP(OUT_NDF, 'CENTRE', 1, '_INTEGER', 'WRITE',
     :     OUT_A_PTR, ITEMP, STATUS)
      IF (STATUS .EQ. SAI__OK) THEN
         CALL SCULIB_NFILLI (R_NBOL, %val(OUT_A_PTR))
      END IF
      CALL NDF_ACPUT ('Bolometer', OUT_NDF, 'LABEL', 1, STATUS)
      CALL NDF_AUNMP (OUT_NDF, 'CENTRE', 1, STATUS)

*     Integrations
      CALL NDF_AMAP (OUT_NDF, 'CENTRE', 2, '_REAL', 'WRITE',
     :     OUT_A_PTR, ITEMP, STATUS)
 
      IF (STATUS .EQ. SAI__OK) THEN
         CALL SCULIB_NFILLR (ITEMP, %val(OUT_A_PTR))
         CALL SCULIB_MULCAR(ITEMP, %VAL(OUT_A_PTR), 1.0/REAL(NPIX),
     :        %VAL(OUT_A_PTR))
      END IF

      CALL NDF_ACPUT ('Integration', OUT_NDF, 'LABEL', 2, STATUS)
      CALL NDF_AUNMP (OUT_NDF, 'CENTRE', 2, STATUS)


*     Now we can start writing header information to the file

*     SCUCD extension
      CALL NDF_XNEW(OUT_NDF, 'SCUCD', 'SCUCD_ST', 0, 0,
     :     OUT_SCUCD_LOC, STATUS)

*     SCUBA extension
      CALL NDF_XNEW(OUT_NDF, 'SCUBA', 'SCUBA_ST', 0, 0,
     :     OUT_SCUBA_LOC, STATUS)

*     REDS extension
      CALL NDF_XNEW(OUT_NDF, 'REDS', 'SURF_EXTENSION', 0, 0,
     :     OUT_REDS_LOC, STATUS)

*     FIGARO extension (for completeness)
      CALL NDF_XNEW(OUT_NDF, 'FIGARO', 'FIGARO_EXT', 0, 0,
     :     OUT_FIG_LOC, STATUS)
      CALL DAT_ANNUL(OUT_FIG_LOC, STATUS)

*     First write DEM_PNTR array
*     Create the component
*     Note that DEM_PNTR is always 1 dimensional in this
*     since there are no switches, exposures or measurements.
*     Only has 3 dimensions since there are no switches

      DEMBNDS ( 1 ) = 1
      DEMBNDS ( 2 ) = NRCYCLE
      DEMBNDS ( 3 ) = 1
      CALL CMP_MOD(OUT_SCUBA_LOC, 'DEM_PNTR', '_INTEGER', 3,
     :     DEMBNDS, STATUS)

*     Map it (rather do this as otherwise I need to create
*     the array on the fly
      CALL CMP_MAPV(OUT_SCUBA_LOC, 'DEM_PNTR', '_INTEGER',
     :     'WRITE', DEM_PNTR, ITEMP, STATUS)

*     Loop over cycles again
      DO I = 1, NRCYCLE
         OFFSET = (I-1) * NPIX + 1
         CALL VEC_ITOI(.FALSE., 1, OFFSET,
     :        %VAL(DEM_PNTR + ((I-1) * VAL__NBI)), IERR, NERR,
     :        STATUS)
      END DO

*     Unmap DEM_PNTR
      CALL CMP_UNMAP(OUT_SCUBA_LOC, 'DEM_PNTR', STATUS)

*     Write LST information (Same dimensions as DEM_PNTR)
      DEMBNDS ( 1 ) = 1
      DEMBNDS ( 2 ) = 1
      DEMBNDS ( 3 ) = NRCYCLE
      DEMBNDS ( 4 ) = 1
      CALL CMP_MOD(OUT_SCUCD_LOC, 'LST_STRT', '_DOUBLE', 4,
     :     DEMBNDS, STATUS)

*     Map it (rather do this as otherwise I need to create
*     the array on the fly (re-use DEM_PNTR variable)
      CALL CMP_MAPV(OUT_SCUCD_LOC, 'LST_STRT', '_DOUBLE',
     :     'WRITE', DEM_PNTR, ITEMP, STATUS)

*     This is based on STIME
*     Returns LST in radians and the MJD
*     Dont know why John chose degrees east of meridian
*     convert back to degrees west.
      LONGITUDE =  LONG_OBS - 360.0D0

      CALL LST_FROM_UT(GR_YY, GR_MN, GR_DD, UT_HH, UT_MN, UT_SS,
     :     LONGITUDE, LST_STRT, MJD, STATUS)

*     increment in LST per cycle
*     I think this is stored in cycle_t and is in millisec in the header
*     Convert to radians
      DLST = (CYCLE_T / (1000.0D0 * 3600.0D0)) * 15.0D0 * PI / 180.0D0

*     Loop over cycles
      DO I = FSCYCLE, NRCYCLE
         LST = LST_STRT + (DBLE(I-1) * DLST)
         
         CALL VEC_DTOD(.FALSE., 1, LST,
     :        %VAL(DEM_PNTR + (I-FSCYCLE) * VAL__NBD), IERR, NERR,
     :        STATUS)

      END DO

*     Unmap LST_STRT
      CALL CMP_UNMAP(OUT_SCUCD_LOC, 'LST_STRT', STATUS)

*     Now deal with jiggle patterns
*     Create the JIGL_X and JIGL_Y components
      CALL CMP_MOD(OUT_SCUCD_LOC, 'JIGL_X', '_REAL', 1,
     :     NPIX, STATUS)
      CALL CMP_MOD(OUT_SCUCD_LOC, 'JIGL_Y', '_REAL', 1,
     :     NPIX, STATUS)

*     Now loop over jiggle positions and store the relevant
*     ones (> -1)

*     SURF requires that bolometer coordinates are derived by
*          Xpos = Bol_Xpos - Jigl_X
*          Ypos = Bol_Ypos - Jigl_Y
*
*     DREAM currently assumes
*          Xpos = -Bol_Xpos + Jigl_X
*          Ypos =  Bol_Ypos - Jigl_Y

*     I end up doing this
*          Invert x positions in flatfield (nasty)
*          Negate jiggle position for X and Y

      ITEMP = 0
      DO I = 0, MAX_PIX  - 1
         IF (INT_POS(I) .NE. - 1) THEN 
            ITEMP = ITEMP + 1
            JIGL_X(ITEMP) = JIG_X(I)
            JIGL_Y(ITEMP) = JIG_Y(I)
         END IF
      END DO


*     Write the jiggle pattern to the extensions
      CALL CMP_PUT1R(OUT_SCUCD_LOC, 'JIGL_X', NPIX, JIGL_X, STATUS)
      CALL CMP_PUT1R(OUT_SCUCD_LOC, 'JIGL_Y', NPIX, JIGL_Y, STATUS)


*     Now work out the order in which the bolometers are stored.

*     Create the BOL_CHAN and BOL_ADC extensions
      CALL CMP_MOD(OUT_SCUBA_LOC, 'BOL_CHAN', '_INTEGER', 1,
     :     R_NBOL, STATUS)
      CALL CMP_MOD(OUT_SCUBA_LOC, 'BOL_ADC', '_INTEGER', 1,
     :     R_NBOL, STATUS)


*     Need to loop over input bolometers and decode the name
*     DREAM stores the storage order in BOL_ORDER and the
*     name in BOL_NAME

      DO I = 0, R_NBOL - 1

*     Determine the ADC and Channel number from the name
         CALL SCULIB_BOLDECODE ( BOL_NAME(I), ADC, CHAN, STATUS)

*     Now put these values into the correct slot of BOL_CHAN and BOL_ADC
*     (Am I supposed to use BOL_ORDER?)
         BOL_CHAN ( I + 1) = CHAN
         BOL_ADC ( I + 1) = ADC
*         BOL_CHAN ( BOL_ORDER(I) + 1) = CHAN
*         BOL_ADC ( BOL_ORDER(I) + 1) = ADC

      END DO

*     Write the bolometer order to the extensions
      CALL CMP_PUT1I(OUT_SCUBA_LOC, 'BOL_ADC', R_NBOL, BOL_ADC, STATUS)
      CALL CMP_PUT1I(OUT_SCUBA_LOC, 'BOL_CHAN', R_NBOL, BOL_CHAN,STATUS)

*     Read flatfield from a text file
      CALL PAR_GET0C ('FLATFILE', FLATFILE, STATUS)

      CALL SCULIB_READBOLS (FLATFILE, SCUBA__NUM_CHAN,
     :     SCUBA__NUM_ADC, BOL_TYPE, BOL_DU3, BOL_DU4, BOL_CALB,
     :     BOL_RTEMP, BOL_RTEMP, BOL_RTEMP, BOL_QUAL, BOL_DAY,
     :     BOL_RUN, BOL_REF, STATUS)

*     Now write the flatfield information to the extensions
      NDIM = 2
      DIM (1) = SCUBA__NUM_CHAN
      DIM (2) = SCUBA__NUM_ADC

*     Create the extensions
      CALL CMP_MODC(OUT_SCUBA_LOC, 'BOL_TYPE', 20, NDIM, DIM, STATUS)
      CALL CMP_MOD(OUT_SCUBA_LOC, 'BOL_DU3', '_REAL', NDIM, DIM,
     :     STATUS)
      CALL CMP_MOD(OUT_SCUBA_LOC, 'BOL_DU4', '_REAL', NDIM, DIM,
     :     STATUS)

*     Only really interested in the BOL_TYPE and DU3, DU4 arrays
*     Write the data
      CALL CMP_PUTNC (OUT_SCUBA_LOC, 'BOL_TYPE', NDIM, DIM, BOL_TYPE,
     :     DIM, STATUS)
      CALL CMP_PUTNR (OUT_SCUBA_LOC, 'BOL_DU3', NDIM, DIM, BOL_DU3,
     :     DIM, STATUS)
      CALL CMP_PUTNR (OUT_SCUBA_LOC, 'BOL_DU4', NDIM, DIM, BOL_DU4,
     :     DIM, STATUS)



*     FITS EXTENSION -------------------
      N_FITS = 0  ! No FITS components to start with

*     Object name (default to something if nothing in header)
      CALL SCULIB_PUT_FITS_C(SCUBA__MAX_FITS, N_FITS, FITS,
     :     'OBJECT', 'DREAM', 'Name of object', STATUS)

*     RUN number
      CALL CHR_CTOI(RUNNO, ITEMP, STATUS)
      CALL SCULIB_PUT_FITS_I(SCUBA__MAX_FITS, N_FITS, FITS,
     :     'RUN', ITEMP, 'Run number of observation', STATUS)

*     Observation, Sample mode (Always MAP/JIGGLE) and sample coords
      CALL SCULIB_PUT_FITS_C(SCUBA__MAX_FITS, N_FITS, FITS,
     :     'MODE', 'MAP', 'The type of observation', STATUS)
      CALL SCULIB_PUT_FITS_C(SCUBA__MAX_FITS, N_FITS, FITS,
     :     'SAM_MODE', 'JIGGLE', 'Sampling method', STATUS)
      CALL SCULIB_PUT_FITS_C(SCUBA__MAX_FITS, N_FITS, FITS,
     :     'SAM_CRDS', 'NA', 'Coordinate system of sampling mesh', 
     :     STATUS)
      CALL SCULIB_PUT_FITS_I(SCUBA__MAX_FITS, N_FITS, FITS,
     :     'SAM_PA', -1, 'Scan P.A. rel. to lat. line; 0=lat, 90=long', 
     :     STATUS)

*     Coordinates of observation

*     Declination
*     Convert to string D:M:S
      IF (DECSN .GE. 0.0D0) THEN
         STEMP = '+'
      ELSE
         STEMP = '-'
      END IF
      IPOSN = 1

      DEC = DBLE(DECDD) + (DBLE(DECMN)/60.0D0) + (DBLE(DECSS)/3600.0D0)
      CALL CHR_RTOAN(REAL(DEC), 'DEGREES', STEMP, IPOSN)

*     Store in LAT
      CALL SCULIB_PUT_FITS_C(SCUBA__MAX_FITS, N_FITS, FITS,
     :     'LAT', STEMP, 'Object Latitude', STATUS)

*     Store in token
      CALL MSG_SETC('DEC', STEMP)

*     Determine Right ascension 
      RA = DBLE(RAHH) + (DBLE(RAMN)/60.0D0) + (DBLE(RASS)/3600.0D0)
      IPOSN = 0
      CALL CHR_RTOAN(REAL(RA), 'HOURS', STEMP, IPOSN)

*     Store in LONG
      CALL SCULIB_PUT_FITS_C(SCUBA__MAX_FITS, N_FITS, FITS,
     :     'LONG', STEMP, 'Object longitude', STATUS)

*     Store in message token
      CALL MSG_SETC('RA', STEMP)

*     Write coordinates to display
      CALL MSG_SETC('TSK',TSKNAME)
      CALL MSG_OUTIF(MSG__NORM, ' ', '^TSK: Coordinates: '//
     :     '^RA, ^DEC', STATUS)


*     Need to ask for coordinate frame of the tracking centre
      CALL PAR_CHOIC('COORDS', 'RB', 'AZ,RD,RB,GA,RJ', .TRUE.,
     :     CENT_CRD, STATUS)

      CALL SCULIB_PUT_FITS_C(SCUBA__MAX_FITS, N_FITS, FITS,
     :     'CENT_CRD',CENT_CRD,'Centre coordinate system', STATUS)



*     UT Date.
*     Convert Year
      CALL CHR_ITOC(GR_YY, STEMP, ITEMP)
      IPOSN = CHR_LEN(STEMP)
      CALL CHR_APPND(':',STEMP, IPOSN)

*     Convert month
      CALL CHR_ITOC(GR_MN, CTEMP, ITEMP)
      CALL CHR_APPND(CTEMP, STEMP, IPOSN)
      CALL CHR_APPND(':',STEMP, IPOSN)

*     Convert day
      CALL CHR_ITOC(GR_DD, CTEMP, ITEMP)
      CALL CHR_APPND(CTEMP, STEMP, IPOSN)

*     Write the MJD
      CALL SCULIB_PUT_FITS_D(SCUBA__MAX_FITS, N_FITS, FITS,
     :     'MJD-OBS', MJD, 'Modified Julian Date of obsstart',
     :     STATUS)

*     Write the date
*      STEMP = '1998:1:1'
      CALL SCULIB_PUT_FITS_C(SCUBA__MAX_FITS, N_FITS, FITS,
     :     'UTDATE', STEMP, 'UT Date of observation', STATUS)     



*     Write the UT time
*     Convert hour
      CALL CHR_ITOC(UT_HH, STEMP, ITEMP)
      IPOSN = CHR_LEN(STEMP)
      CALL CHR_APPND(':',STEMP, IPOSN)

*     Convert minute
      CALL CHR_ITOC(UT_MN, CTEMP, ITEMP)
      CALL CHR_APPND(CTEMP, STEMP, IPOSN)
      CALL CHR_APPND(':',STEMP, IPOSN)

*     Convert second
      CALL CHR_ITOC(UT_SS, CTEMP, ITEMP)
      CALL CHR_APPND(CTEMP, STEMP, IPOSN)

      CALL SCULIB_PUT_FITS_C(SCUBA__MAX_FITS, N_FITS, FITS,
     :     'UTSTART', STEMP, 'UT at start of observation', STATUS)

*     Convert LST (Hours minutes seconds) into a time
*     First convert LST_STRT to hours (should be radians to this point)
      LST_STRT = LST_STRT * 180.0D0 / (PI * 15.0D0)

      IPOSN = 0
      CALL CHR_DTOAN(LST_STRT, 'HOURS', STEMP, IPOSN)

*     and write it out
      CALL SCULIB_PUT_FITS_C(SCUBA__MAX_FITS, N_FITS, FITS,
     :     'STSTART', STEMP, 'ST at start of observation', STATUS)


*     Also store STEND
      IPOSN = 0

      DTEMP = LST_STRT + (DLST * DBLE(FSCYCLE + NRCYCLE - 1) 
     :     * 180.0D0 / (PI * 15.0D0))
      CALL CHR_DTOAN(DTEMP,
     :     'HOURS',STEMP, IPOSN)
      CALL SCULIB_PUT_FITS_C(SCUBA__MAX_FITS, N_FITS, FITS,
     :     'STEND', STEMP, 'ST at start of observation', STATUS)

*     exposure time per sample
      CALL SCULIB_PUT_FITS_D(SCUBA__MAX_FITS, N_FITS, FITS,
     :     'EXP_TIME', DBLE(CYCLE_T / (1000.0 * REAL(NPIX))),
     :     'Exposure time for each basic measurement (sec)',
     :     STATUS)


*     Set up MAP_X and MAP_Y offsets
      CALL SCULIB_PUT_FITS_I(SCUBA__MAX_FITS, N_FITS, FITS,
     :      'MAP_X', 0.0, 
     :      'Map X offset from telescope centre (arcsec)',
     :      STATUS)
      CALL SCULIB_PUT_FITS_I(SCUBA__MAX_FITS, N_FITS, FITS,
     :      'MAP_Y', 0.0, 
     :      'Map Y offset from telescope centre (arcsec)',
     :      STATUS)

*     Number of bolometers
      CALL SCULIB_PUT_FITS_I(SCUBA__MAX_FITS, N_FITS, FITS,
     :      'N_BOLS', R_NBOL, 
     :      'Number of bolometers selected', STATUS)

*     State of the observation (anything except ABORT)
      CALL SCULIB_PUT_FITS_C(SCUBA__MAX_FITS, N_FITS, FITS,
     :     'STATE','Terminating', 'SCUCD state', STATUS)

*     Version of SCUCD (set to 0 for now)
      CALL SCULIB_PUT_FITS_I(SCUBA__MAX_FITS, N_FITS, FITS,
     :     'VERSION', 0, 'SCUCD version (DREAM data)',STATUS)

*     Jiggle info
      CALL SCULIB_PUT_FITS_I(SCUBA__MAX_FITS, N_FITS, FITS,
     :     'JIGL_CNT', NPIX, 'Number of offsets in a jiggle pattern',
     :     STATUS)
      CALL SCULIB_PUT_FITS_I(SCUBA__MAX_FITS, N_FITS, FITS,
     :     'J_PER_S', NPIX, 'Number of jiggles per switch',
     :     STATUS)
      CALL SCULIB_PUT_FITS_I(SCUBA__MAX_FITS, N_FITS, FITS,
     :     'J_REPEAT', NPIX, 'No. of jiggle pattern repeats in switch',
     :     STATUS)

*     Center of array
      CALL SCULIB_PUT_FITS_D(SCUBA__MAX_FITS, N_FITS, FITS,
     :     'CNTR_DU3', 0.0D0, 'Nasmyth dU3 coord of instrument centre',
     :     STATUS)
      CALL SCULIB_PUT_FITS_D(SCUBA__MAX_FITS, N_FITS, FITS,
     :     'CNTR_DU4', 0.0D0, 'Nasmyth dU4 coord of instrument centre',
     :     STATUS)

*     Position of the telescope
      CALL SCULIB_PUT_FITS_D(SCUBA__MAX_FITS, N_FITS, FITS,
     :     'LAT-OBS', LAT_OBS, 'Latitude of observatory (degrees)',
     :     STATUS)
      CALL SCULIB_PUT_FITS_D(SCUBA__MAX_FITS, N_FITS, FITS,
     :     'LONG-OBS', LONG_OBS, 
     :     'East Longitude of observatory (degrees)',
     :     STATUS)

*     Need to supply some chop information (meaningless)
      CALL SCULIB_PUT_FITS_D(SCUBA__MAX_FITS, N_FITS, FITS,
     :     'CHOP_THR', 0.0D0, 'Chopper throw (arcsec)',
     :     STATUS)
      CALL SCULIB_PUT_FITS_D(SCUBA__MAX_FITS, N_FITS, FITS,
     :     'CHOP_PA', 0.0D0, 'Chopper P.A., 0 = in lat, 90 = in long',
     :     STATUS)
      CALL SCULIB_PUT_FITS_D(SCUBA__MAX_FITS, N_FITS, FITS,
     :     'CHOP_FRQ', SMU_F, 'Chopper frequency (Hz)',
     :     STATUS)
      CALL SCULIB_PUT_FITS_C(SCUBA__MAX_FITS, N_FITS, FITS,
     :     'CHOP_CRD', 'NA', 'Chopper coordinate system',
     :     STATUS)
      CALL SCULIB_PUT_FITS_C(SCUBA__MAX_FITS, N_FITS, FITS,
     :     'CHOP_FUN', 'DREAM', 'Chopper waveform',
     :     STATUS)



*     Sub instrument information
*     Assume that we have one sub-instrument and that it is LONG
      CALL SCULIB_PUT_FITS_I(SCUBA__MAX_FITS, N_FITS, FITS,
     :     'N_SUBS', 1, 'Number of sub-instruments used',
     :     STATUS)

      CALL SCULIB_PUT_FITS_C(SCUBA__MAX_FITS, N_FITS, FITS,
     :     'SUB_1','LONG', 'SCUBA sub-instrument being used', STATUS)
      CALL SCULIB_PUT_FITS_C(SCUBA__MAX_FITS, N_FITS, FITS,
     :     'SUB_2','not used', 'SCUBA sub-instrument being used',STATUS)
      CALL SCULIB_PUT_FITS_C(SCUBA__MAX_FITS, N_FITS, FITS,
     :     'SUB_3','not used', 'SCUBA sub-instrument being used',STATUS)
      CALL SCULIB_PUT_FITS_C(SCUBA__MAX_FITS, N_FITS, FITS,
     :     'SUB_4','not used', 'SCUBA sub-instrument being used',STATUS)
      CALL SCULIB_PUT_FITS_C(SCUBA__MAX_FITS, N_FITS, FITS,
     :     'SUB_5','not used', 'SCUBA sub-instrument being used',STATUS)

      CALL SCULIB_PUT_FITS_D(SCUBA__MAX_FITS, N_FITS, FITS,
     :     'TAUZ_1',0.0D0, 'Zenith sky optical depth',STATUS)

*     FILTERS and wavelength

      CALL SCULIB_PUT_FITS_C(SCUBA__MAX_FITS, N_FITS, FITS,
     :     'FILT_1','850', 'Filter name', STATUS)
      CALL SCULIB_PUT_FITS_D(SCUBA__MAX_FITS, N_FITS, FITS,
     :     'WAVE_1',862.0D0, 'Wavelength of map (microns)', STATUS)



*     Write FITS component
      CALL NDF_XNEW (OUT_NDF, 'FITS', '_CHAR*80', 1, N_FITS, 
     :     OUT_FITS_LOC, STATUS)
      CALL DAT_PUT1C (OUT_FITS_LOC, N_FITS, FITS, STATUS)

*     Annul extension locators
      CALL DAT_ANNUL(OUT_SCUCD_LOC, STATUS)
      CALL DAT_ANNUL(OUT_SCUBA_LOC, STATUS)
      CALL DAT_ANNUL(OUT_REDS_LOC, STATUS)
      CALL DAT_ANNUL(OUT_FITS_LOC, STATUS)

*     Write the HISTORY information. (DREAM info will be written
*     automatically when NDF is closed
      CALL NDF_HCRE(OUT_NDF, STATUS)

*     Close the NDF and shut down the NDF system (write DREAM history)
      CALL NDF_ANNUL(OUT_NDF, STATUS)
 
*     It seems that the only way to write multiple history
*     entries is to open and close the NDF multiple times!

*     REDUCE_SWITCH
*     Re-open the file
      CALL NDF_ASSOC('OUT', 'UPDATE', OUT_NDF, STATUS)

*     Need to write REDUCE_SWITCH and FLATFIELD tags to fool SURF
*     into thinking that the data have been processed by these tasks
      CALL NDF_HPUT(' ', 'REDUCE_SWITCH', .TRUE., 1, 
     :     'This is a dummy history component to fool SURF',
     :     .FALSE., .TRUE., .FALSE., OUT_NDF, STATUS)

*     Close the NDF
      CALL NDF_ANNUL(OUT_NDF, STATUS)

*     FLATFIELD
*     Re-open the file
      CALL NDF_ASSOC('OUT', 'UPDATE', OUT_NDF, STATUS)


      CALL NDF_HPUT(' ', 'FLATFIELD', .TRUE., 1, 
     :     'This is a dummy history component to fool SURF',
     :     .FALSE., .TRUE., .FALSE., OUT_NDF, STATUS)

*     Close the NDF
      CALL NDF_ANNUL(OUT_NDF, STATUS)


*     Shut down NDF system
      CALL NDF_END(STATUS)

      END
\end{verbatim}
\end{myquote}

\section{Library APIs}

This section lists the full API documentation for the \scusoft\ 
libraries.

\subsection{SURF}

These are the top-level \scusoft\ routines that are not simply task files.
These are routines for disk I/O, the top level A-task interface routine
and specialised code for \calcsky\ and \despike\ (since they are children
of the general purpose \rebin\ task interface).

\sstroutine{
   SURF\_GRID\_CALCSKY
}{
   Calculate sky contribution from median image
}{
   \sstdescription{
      This routine calculates the sky contribution by attempting
      to remove the source from the input data stream. The source
      signal can either be calculated by this routine or by reading
      in a model of the source from a file.

      When calculating the source structure internally a similar
      method to that used by DESPIKE is employed. The input data
      are placed into bins of size one quarter beamwidth. The median
      of each bin is calculated and this is treated as the source
      model (cf. REBIN\_METHOD=MEDIAN in REBIN).

      Once the source model is available, it is removed from
      all of the input data. The source-removed data are then analysed
      with the sky emission derived from the mean of the signal across
      the array for all the sample times.

      Since the sky signal is expected to vary on timesales of the
      order of one second, an option is included for smoothing the
      sky signal. This is especially useful for scan map data where
      samples are taken at 7.8~Hz.
   }
   \sstinvocation{
       CALL SURF\_GRID\_CALCSKY ( TSKNAME, N\_FILES, N\_PTS, N\_POS, N\_BOLS,
            N\_BOLS, WAVELENGTH, DIAMETER, IMNDF, N\_M\_FITS, MODEL\_FITS,
            CHOP\_THROW, CHOP\_PA, BOX\_SIZE, BOL\_RA\_PTR,  BOL\_DEC\_PTR,
            DATA\_PTR, QUALITY\_PTR, SKY\_PTR, SKY\_ERR, BADBIT, STATUS )
   }
   \sstarguments{
      \sstsubsection{
         TSKNAME = CHARACTER (Given)
      }{
         Name of task to be used in output messages
      }
      \sstsubsection{
         N\_FILES = INTEGER (Given)
      }{
         Number of data sets (ie files)
      }
      \sstsubsection{
         N\_PTS ( N\_FILES ) = INTEGER (Given)
      }{
         Total number of points in each map
      }
      \sstsubsection{
         N\_POS( N\_FILES ) = INTEGER (Given)
      }{
         Number of positions per set (Y positions)
      }
      \sstsubsection{
         N\_BOLS( N\_FILES ) = INTEGER (Given)
      }{
         Number of bolometers per set (X positions)
      }
      \sstsubsection{
         WAVELENGTH = REAL (Given)
      }{
         Wavelength of observation in microns
      }
      \sstsubsection{
         DIAMETER = REAL (Given)
      }{
         Diameter of telescope in metres
      }
      \sstsubsection{
         IMNDF = INTEGER (Given)
      }{
         NDF identifier of the supplied model. If this is equal to
         NDF\_\_NOID then no external model is used.
      }
      \sstsubsection{
         N\_M\_FITS = INTEGER (Given)
      }{
         Number of FITS keywords in the model
      }
      \sstsubsection{
         MODEL\_FITS = CHAR$*$(80) (Given)
      }{
         FITS keywords from the model
      }
      \sstsubsection{
         CHOP\_THROW = REAL (Given)
      }{
         Size of chop throw for the first input image (arcsec)
      }
      \sstsubsection{
         CHOP\_PA = REAL (Given)
      }{
         Chop position angle of first input image (degrees)
      }
      \sstsubsection{
         BOX\_SIZE ( N\_FILES ) = INTEGER (Given)
      }{
         Size of smoothing box - pixels
      }
      \sstsubsection{
         BOL\_RA\_PTR( N\_FILES ) = INTEGER (Given)
      }{
         Array of pointers to position information (X coords)
         Note that each data set has positions for N\_POS $*$ N\_BOLS
      }
      \sstsubsection{
         BOL\_RA\_PTR( N\_FILES ) = INTEGER (Given)
      }{
         Array of pointers to position information (Y coords)
      }
      \sstsubsection{
         DATA\_PTR( N\_FILES ) = INTEGER (Given, modified data)
      }{
         Pointers to actual data arrays. The arrays are modified
         and contain the source removed data on exit.
      }
      \sstsubsection{
         QUALITY\_PTR( N\_FILES ) = INTEGER (Given)
      }{
         Pointer to quality arrays
      }
      \sstsubsection{
         SKY\_PTR ( N\_FILES ) = INTEGER (Pointer Given, data returned)
      }{
         Storage space for the sky signal
      }
      \sstsubsection{
         SKY\_ERR ( N\_FILES ) = INTEGER (Pointer Given, data returned)
      }{
         Storage space for the error on sky signal
      }
      \sstsubsection{
         BADBIT ( N\_FILES ) = BYTE (Given)
      }{
         Bad bit mask for identifying bad pixels from quality
      }
      \sstsubsection{
         STATUS = INTEGER (Given \& Returned)
      }{
         Global Status
      }
   }
   \sstexamples{
      \sstexamplesubsection{
         calcsky test\_rlb model=! $\backslash$$\backslash$
      }{
         Calculate sky for test\_rlb.sdf. Only read in one file and
         don{\tt '}t use an external source model.
      }
      \sstexamplesubsection{
         calcsky list.inp model=m82 noloop$\backslash$$\backslash$
      }{
         Read in the files specified in list.inp and use m82.sdf
         as a model of the source.
      }
      \sstexamplesubsection{
         calcsky file nosrc=nosrc boxsz=10.0 $\backslash$$\backslash$
      }{
         Calculate sky for file.sdf. Store the source subtracted image
         in nosrc.sdf. Use a smoothing size of 10 seconds.
      }
   }
   \sstnotes{
      \sstitemlist{

         \sstitem
         The model itself is only an approximation
           to the data (since the data points can fall anywhere within
           a given cell) so some source signal will remain after source
           subtraction.

         \sstitem
         If a model is supplied externally (via MODEL parameter) the
           cell size and the map centre of the model
           is used for the source subtraction.

         \sstitem
         The sky signal is stored in an NDF extension (.MORE.REDS.SKY).
           The file must be processed by REMSKY to actually remove the
           sky contribution.
      }
   }
   \sstdiytopic{
      Related Applications
   }{
      SURF: REMSKY
   }
   \sstdiytopic{
      Copyright
   }{
      Copyright \copyright 1995,1996,1997,1998,1999 Particle Physics and Astronomy
      Research Council. All Rights Reserved.
   }
}
\sstroutine{
   SURF\_GRID\_DESPIKE
}{
   Despike data by sorting into a given grid position
}{
   \sstdescription{
      For each data point this routine places it into a bin in the
      output grid depending on the position of the data point on the sky.
      The position in the input data array is stored.
      This is done in two stages:

      \sstitemlist{

         \sstitem
         Find the size of the output grid from the maximum extent of
            the input data.

         \sstitem
         Loop through data. Find I,J coordinate of each point in the
            output grid.

         \sstitem
         Find out maximum number of points for an I,J position.

         \sstitem
         Put data onto grid in array (I,J,N) [REALS].
            We also need to store positions of these data.
            We can either do it by storing the file number, bolometer and
            position (time) index OR we can just store some index in a merged
            data array that goes from 1..TOT\_PTS.
             First method is easy but memory hungry. Second method is
             more efficient but does need some reconstruction to work
             out where the point was in the original data.
           Use the second method.

      }
      Once the data is gridded, it is first displayed and then
      despiked. Currently despiking is done on a simple sigma clipping
      basis for each bin.
   }
   \sstinvocation{
       CALL SURF\_GRID\_DESPIKE ( N\_FILES, N\_PTS, N\_POS, N\_BOLS, WAVELENGTH,
            DIAMETER, BOL\_RA\_PTR, BOL\_DEC\_PTR, DATA\_PTR, QUALITY\_PTR,
            NX, NY, ICEN, JCEN, NSPIKES,
            BADBIT, STATUS )
   }
   \sstarguments{
      \sstsubsection{
         N\_FILES = INTEGER (Given)
      }{
         Number of data sets (ie files)
      }
      \sstsubsection{
         N\_PTS ( N\_FILES ) = INTEGER (Given)
      }{
         Total number of points in each map
      }
      \sstsubsection{
         N\_POS( N\_FILES ) = INTEGER (Given)
      }{
         Number of positions per set (Y positions)
      }
      \sstsubsection{
         N\_BOLS( N\_FILES ) = INTEGER (Given)
      }{
         Number of bolometers per set (X positions)
      }
      \sstsubsection{
         BITNUM = INTEGER (Given)
      }{
         Bit number to be affected by this routine
      }
      \sstsubsection{
         WAVELENGTH = REAL (Given)
      }{
         Wavelength of data (microns)
      }
      \sstsubsection{
         DIAMETER = REAL (Given)
      }{
         Diameter of telescope (metres)
      }
      \sstsubsection{
         BOL\_RA\_PTR( N\_FILES ) = INTEGER (Given)
      }{
         Array of pointers to position information (X coords)
         Note that each data set has positions for N\_POS $*$ N\_BOLS
      }
      \sstsubsection{
         BOL\_RA\_PTR( N\_FILES ) = INTEGER (Given)
      }{
         Array of pointers to position information (Y coords)
      }
      \sstsubsection{
         DATA\_PTR( N\_FILES ) = INTEGER (Given)
      }{
         Pointers to actual data arrays
      }
      \sstsubsection{
         QUALITY\_PTR( N\_FILES ) = INTEGER (Given)
      }{
         Pointer to quality arrays
      }
      \sstsubsection{
         NX = INTEGER (Returned)
      }{
         Number of points in grid (X)
      }
      \sstsubsection{
         NY = INTEGER (Returned)
      }{
         Number of points in grid (Y)
      }
      \sstsubsection{
         ICEN = INTEGER (Returned)
      }{
         Reference pixel (X)
      }
      \sstsubsection{
         JCEN = INTEGER (Returned)
      }{
         Reference pixel (Y)
      }
      \sstsubsection{
         NSPIKES ( N\_FILES ) = INTEGER (Returned)
      }{
         Number of spikes detected (and removed) in each file
      }
      \sstsubsection{
         BADBIT ( N\_FILES ) = BYTE (Given)
      }{
         Bad bit mask for identifying bad pixels from quality
      }
      \sstsubsection{
         STATUS = INTEGER (Given \& Returned)
      }{
         Global Status
      }
   }
   \sstnotes{
      For SMODE=NONE, DMODE is only requested if a plot is required.
   }
   \sstdiytopic{
      Copyright
   }{
      Copyright \copyright 1995,1996,1997,1998,1999 Particle Physics and Astronomy
      Research Council. All Rights Reserved.
   }
}
\sstroutine{
   SURF\_MON
}{
   main routine for SCUBA offline data reduction package
}{
   \sstdescription{
      This is the main routine for the SCUBA reduction A-task.
   }
   \sstinvocation{
      CALL SURF\_MON( STATUS )
   }
   \sstarguments{
      \sstsubsection{
         STATUS = INTEGER (Given and Returned)
      }{
         The global status
      }
   }
   \sstnotes{
      This routine is not seen by the user
   }
   \sstdiytopic{
      Copyright
   }{
      Copyright \copyright 1995,1996,1997,1998,1999 Particle Physics and Astronomy
      Research Council. All Rights Reserved.
   }
}
\sstroutine{
   SURF\_READ\_REBIN\_NDF
}{
   Read an NDF into memory prior to regridding
}{
   \sstdescription{
      This routines reads all necessary information from an NDF (via
      the given NDF identifier) for rebinning. Optionally, returns
      the quality array (needed for despiking).
   }
   \sstinvocation{
       CALL SURF\_READ\_REBIN\_NDF( IN\_NDF, MAX\_FILE, NSPEC,
            DATA\_SPEC, OUT\_COORDS, N\_FILE, USE\_SECTION,
            N\_BOL, N\_POS, N\_INTS, N\_BEAMS, MJD\_STANDARD, IN\_UT1,
            OUT\_RA\_CEN, OUT\_DEC\_CEN, FITS, N\_FITS, WAVELENGTH,
            SUB\_INSTRUMENT, SOBJECT, SUTDATE, SUTSTART,
            BOL\_ADC, BOL\_CHAN,
            BOL\_RA\_PTR, BOL\_RA\_END, BOL\_DEC\_PTR,
            BOL\_DEC\_END, DATA\_PTR, DATA\_END, VARIANCE\_PTR, VARIANCE\_END,
            QMF, QUALITY\_PTR, QUALITY\_END, BADBITS,
            USE\_LST, LST\_PTR, ANG\_INT, ANG\_MEAS, INT\_LIST, MEAS\_LIST, BOLWT,
            STATUS)
   }
   \sstarguments{
      \sstsubsection{
         IN\_NDF = INTEGER (Given)
      }{
         NDF identifier of input NDF
      }
      \sstsubsection{
         MAX\_FILE = INTEGER (Given)
      }{
         Max number of files allowed [used for INT\_LIST/MEAS only]
      }
      \sstsubsection{
         NSPEC = INTEGER (Given)
      }{
         Number of SCUBA sections in DATA\_SPEC
      }
      \sstsubsection{
         DATA\_SPEC( NSPEC ) = CHAR (Given)
      }{
         SCUBA sections
      }
      \sstsubsection{
         OUT\_COORDS = CHAR (Given)
      }{
         Output coordinates system. (Passed into SURF\_READ\_REBIN\_NDF)
      }
      \sstsubsection{
         N\_FILE = INTEGER (Given)
      }{
         Current file number (less than MAX\_FILE and greater than 0).
      }
      \sstsubsection{
         USE\_SECTION = LOGICAL (Given)
      }{
         Determines whether we are using the section or the invers
      }
      \sstsubsection{
         N\_BOL = INTEGER (Returned)
      }{
         Number of bolometers associated with this observation
      }
      \sstsubsection{
         N\_POS = INTEGER (Returned)
      }{
         Number of samples in observation
      }
      \sstsubsection{
         N\_INTS = INTEGER (Returned)
      }{
         Total Number of integrations in observation (INT$*$MEAS)
      }
      \sstsubsection{
         N\_MEAS = INTEGER (Returned)
      }{
         Number of measurements in observation
      }
      \sstsubsection{
         N\_BEAMS = INTEGER (Given \& Returned)
      }{
         Number of beams requested in the positions arrays.
         N\_BEAMS = 1  will return the middle-beam (standard position)
         N\_BEAMS = 2 will return L and R beams for RASTER and JIGGLE
         N\_BEAMS = 3 will return L, M and R beams (as M,L,R) for JIGGLE
                    data that has been reduce\_switched and for SCAN
                    data. N\_BEAM is set to 2 if a single switch of a
                    JIGGLE map is requested.
      }
      \sstsubsection{
         MJD\_STANDARD = DOUBLE (Given \& Returned)
      }{
         Modified Julian data of observation. Returned if N\_FILE=1.
      }
      \sstsubsection{
         IN\_UT1 = DOUBLE (Returned)
      }{
         Modified Julian data of observation UT1 at start of an input
         observation, expressed as modified Julian day
      }
      \sstsubsection{
         IN\_RA\_CEN = DOUBLE (Returned)
      }{
         RA of centre
      }
      \sstsubsection{
         IN\_DEC\_CEN = DOUBLE (Returned)
      }{
         Dec of centre
      }
      \sstsubsection{
         FITS ( SCUBA\_MAX\_FITS ) = CHARACTER $*$80 (Returned)
      }{
         FITS header entries
      }
      \sstsubsection{
         N\_FITS = INTEGER (Returned)
      }{
         Number of FITS entries
      }
      \sstsubsection{
         WAVELENGTH = REAL (Given \& Returned)
      }{
         Wavelength of map
      }
      \sstsubsection{
         SUB\_INSTRUMENT = CHAR (Given \& Returned)
      }{
         Sub instrument of map
      }
      \sstsubsection{
         SOBJECT = CHAR (Returned)
      }{
         Name of object
      }
      \sstsubsection{
         SUTDATE = CHAR (Returned)
      }{
         UT date of the observation
      }
      \sstsubsection{
         SUTSTART = CHAR (Returned)
      }{
         UT time of the observation
      }
      \sstsubsection{
         BOL\_ADC = INTEGER (Returned)
      }{
         ADC information for bolometers - only used by BOLREBIN
      }
      \sstsubsection{
         BOL\_CHAN = INTEGER (Returned)
      }{
         Channel information for bolometers - only used by BOLREBIN
      }
      \sstsubsection{
         BOL\_RA\_PTR = INTEGER (Returned)
      }{
         Pointer to RA bolometer positions [Pointer is given, data is
         returned]
      }
      \sstsubsection{
         BOL\_RA\_END = INTEGER (Returned)
      }{
         Pointer to end of RA bolometer positions
      }
      \sstsubsection{
         BOL\_DEC\_PTR = INTEGER (Returned)
      }{
         Pointer to DEC bolometer positions
      }
      \sstsubsection{
         BOL\_DEC\_END = INTEGER (Returned)
      }{
         Pointer to end of DEC bol positions
      }
      \sstsubsection{
         DATA\_PTR = INTEGER (Returned)
      }{
         Pointer to data values
      }
      \sstsubsection{
         DATA\_END = INTEGER (Returned)
      }{
         Pointer to end of data values
      }
      \sstsubsection{
         VARIANCE\_PTR = INTEGER (Returned)
      }{
         Pointer to variance values
      }
      \sstsubsection{
         VARIANCE\_END = INTEGER (Returned)
      }{
         Pointer to end of variance values
      }
      \sstsubsection{
         QMF = LOGICAL (Given)
      }{
         Flag to decide whether quality is being stored (.FALSE.) or
         being folded into the data array (.true.). See NDF\_SQMF
      }
      \sstsubsection{
         QUALITY\_PTR = INTEGER (Returned)
      }{
         Pointer to quality array
      }
      \sstsubsection{
         QUALITY\_END = INTEGER (Returned)
      }{
         Pointer to end of quality array
      }
      \sstsubsection{
         BADBITS = BYTE (Returned)
      }{
         Bad bits mask for quality array
      }
      \sstsubsection{
         USE\_LST = LOGICAL (Given)
      }{
         Governs whether we want an LST array returned
      }
      \sstsubsection{
         LST\_PTR( 2 ) = INTEGER (Returned)
      }{
         Array of pointers (begin and end)  to array of LSTs
      }
      \sstsubsection{
         ANG\_INT( MAX\_FILE, SCUBA\_\_MAX\_INT,2)  = REAL (Returned)
      }{
         Array containing the polarimetry angles for each integration
         The 2 dimensions are for WPLATE and ANGROT
      }
      \sstsubsection{
         ANG\_MEAS( MAX\_FILE, SCUBA\_\_MAX\_MEAS,2) = REAL (Returned)
      }{
         Array containing the pol angles for each measurement
         The 2 dimensions are for WPLATE and ANGROT
      }
      \sstsubsection{
         INT\_LIST( MAX\_FILE, SCUBA\_\_MAX\_INT$+$1) = INTEGER (Returned)
      }{
         Position of integrations in each data file
      }
      \sstsubsection{
         MEAS\_LIST(MAX\_FILE, SCUBA\_\_MAX\_MEAS$+$1) = INTEGER (Returned)
      }{
         Position of measurements in each data file
      }
      \sstsubsection{
         BOLWT ( N\_BOL ) = REAL (Returned)
      }{
         Relative weights of each bolometer
      }
      \sstsubsection{
         STATUS = INTEGER (Given and Returned)
      }{
         The global status
      }
   }
   \sstdiytopic{
      Copyright
   }{
      Copyright \copyright 1995,1996,1997,1998,1999 Particle Physics and Astronomy
      Research Council. All Rights Reserved.
   }
}
\sstroutine{
   SURF\_RECURSE\_READ
}{
   Allow the recursive read of REBIN text files and NDFs
}{
   \sstdescription{
      This routine takes an input name and reads in all NDFs resulting
      from this name. SCUBA sections are parsed. Text files are expanded
      into NDF$+$parameters. Nested text files are allowed.
   }
   \sstinvocation{
       CALL SURF\_RECURSE\_READ( RLEV, NAME, MAX\_FILE,
            OUT\_COORDS, N\_FILE, N\_BOL, N\_POS, N\_INTS,
            IN\_UT1, IN\_RA\_CEN, IN\_DEC\_CEN, FITS, N\_FITS, WAVELENGTH,
            SUB\_INSTRUMENT, OBJECT, UTDATE, UTSTART, FILENAME,
            BOL\_ADC, BOL\_CHAN,
            BOL\_RA\_PTR, BOL\_RA\_END, BOL\_DEC\_PTR,
            BOL\_DEC\_END, DATA\_PTR, DATA\_END, VARIANCE\_PTR, VARIANCE\_END,
            QMF, QUALITY\_PTR, QUALITY\_END, QBITS, ANG\_INT, ANG\_MEAS,
            INT\_LIST, MEAS\_LIST, BOLWT, WEIGHT, SHIFT\_DX, SHIFT\_DY,
            NPARS, PARS, STATUS)
   }
   \sstarguments{
      \sstsubsection{
         RLEV = INTEGER (Given)
      }{
         Recursion level. This is the number of times that the routine
         has called itself. A value of 1 should be passed in by the
         external routine.
      }
      \sstsubsection{
         NAME = CHAR (Given)
      }{
         Name of input file. This may also contain scuba section
         specifications. The input file can be a NDF or a text file
         containing a list of NDFs.
      }
      \sstsubsection{
         MAX\_FILE = INTEGER (Given)
      }{
         Maximum number of NDFs that can be read by the system.
      }
      \sstsubsection{
         OUT\_COORDS = CHAR (Given)
      }{
         Output coordinates system. (Passed into SURF\_READ\_REBIN\_NDFS)
      }
      \sstsubsection{
         N\_FILE = INTEGER (Given \& Returned)
      }{
         Current file number (less than MAX\_FILE and greater than 0).
         This counter is incremented when an NDF has been read successfully.
      }
      \sstsubsection{
         N\_BOL( MAX\_FILE ) = INTEGER (Returned)
      }{
         Number of bolometers associated with each file
      }
      \sstsubsection{
         N\_POS( MAX\_FILE ) = INTEGER (Returned)
      }{
         Number of samples associated with each file
      }
      \sstsubsection{
         N\_INTS( MAX\_FILE ) = INTEGER (Returned)
      }{
         Total Number of integrations associated with each file (INT$*$MEAS)
      }
      \sstsubsection{
         N\_MEAS( MAX\_FILE ) = INTEGER (Returned)
      }{
         Number of measurements associated with each file
      }
      \sstsubsection{
         IN\_UT1( MAX\_FILE ) = DOUBLE (Returned)
      }{
         Modified Julian data of observation for each file
      }
      \sstsubsection{
         IN\_RA\_CEN( MAX\_FILE ) = DOUBLE (Returned)
      }{
         RA of centre for each file
      }
      \sstsubsection{
         IN\_DEC\_CEN( MAX\_FILE ) = DOUBLE (Returned)
      }{
         Dec of centre for each file
      }
      \sstsubsection{
         FITS ( N\_FITS, MAX\_FILE) = CHARACTER$*$(80) (Returned)
      }{
         FITS entries for each file
      }
      \sstsubsection{
         N\_FITS = INTEGER (Given)
      }{
         Size of FITS array for each file
      }
      \sstsubsection{
         WAVELENGTH = REAL (Given \& Returned)
      }{
         Wavelength of map
      }
      \sstsubsection{
         SUB\_INSTRUMENT = CHAR (Given \& Returned)
      }{
         Sub instrument of map
      }
      \sstsubsection{
         OBJECT( MAX\_FILE ) = CHAR (Returned)
      }{
         Name of object in each file
      }
      \sstsubsection{
         UTDATE( MAX\_FILE ) = CHAR (Returned)
      }{
         UT date of each observation
      }
      \sstsubsection{
         UTSTART( MAX\_FILE ) = CHAR (Returned)
      }{
         UT time of each observation
      }
      \sstsubsection{
         FILENAME( MAX\_FILE ) = CHAR (Returned)
      }{
         Actual filename of each file read.
      }
      \sstsubsection{
         BOL\_ADC ( ) = INTEGER (Returned)
      }{
         ADC information for bolometers - only used by BOLREBIN
      }
      \sstsubsection{
         BOL\_CHAN ( ) = INTEGER (Returned)
      }{
         Channel information for bolometers - only used by BOLREBIN
      }
      \sstsubsection{
         BOL\_RA\_PTR( MAX\_FILE ) = INTEGER (Returned)
      }{
         Array of pointers to RA bolometer positions read from each file
      }
      \sstsubsection{
         BOL\_RA\_END( MAX\_FILE ) = INTEGER (Returned)
      }{
         Array of pointers to end of RA bolometer positions read from each file
      }
      \sstsubsection{
         BOL\_DEC\_PTR( MAX\_FILE ) = INTEGER (Returned)
      }{
         Array of pointers to DEC bolometer positions read from each file
      }
      \sstsubsection{
         BOL\_DEC\_END( MAX\_FILE ) = INTEGER (Returned)
      }{
         Array of pointers to end of DEC bol positions read from each file
      }
      \sstsubsection{
         DATA\_PTR( MAX\_FILE ) = INTEGER (Returned)
      }{
         Array of pointers to data values read from each file
      }
      \sstsubsection{
         DATA\_END( MAX\_FILE ) = INTEGER (Returned)
      }{
         Array of pointers to end of data values
      }
      \sstsubsection{
         VARIANCE\_PTR( MAX\_FILE ) = INTEGER (Returned)
      }{
         Array of pointers to variance values read from each file
      }
      \sstsubsection{
         VARIANCE\_END( MAX\_FILE ) = INTEGER (Returned)
      }{
         Array of pointers to end of variance values
      }
      \sstsubsection{
         QMF = LOGICAL (Given)
      }{
         Flag to decide whether quality is being stored (.FALSE.) or
         being folded into the data array (.true.). See NDF\_SQMF
      }
      \sstsubsection{
         QUALITY\_PTR(MAX\_FILE) = INTEGER (Returned)
      }{
         Pointer to quality array
      }
      \sstsubsection{
         QUALITY\_END(MAX\_FILE) = INTEGER (Returned)
      }{
         Pointer to end of quality array
      }
      \sstsubsection{
         QBITS(MAX\_FILE) = BYTE (Returned)
      }{
         Bad bits mask for each file
      }
      \sstsubsection{
         ANG\_INT( MAX\_FILE, SCUBA\_\_MAX\_INT,2)  = REAL (Returned)
      }{
         Array containing the polarimetry angles for each integration
         The 2 dimensions are for WPLATE and ANGROT
      }
      \sstsubsection{
         ANG\_MEAS( MAX\_FILE, SCUBA\_\_MAX\_MEAS,2) = REAL (Returned)
      }{
         Array containing the pol angles for each measurement
         The 2 dimensions are for WPLATE and ANGROT
      }
      \sstsubsection{
         INT\_LIST( MAX\_FILE, SCUBA\_\_MAX\_INT$+$1) = INTEGER (Returned)
      }{
         Position of integrations in each data file
      }
      \sstsubsection{
         MEAS\_LIST(MAX\_FILE, SCUBA\_\_MAX\_MEAS$+$1) = INTEGER (Returned)
      }{
         Position of measurements in each data file
      }
      \sstsubsection{
         BOLWT (max num of bols, MAX\_FILE) = REAL (Returned)
      }{
         Relative Weights of each bolometer for each file
      }
      \sstsubsection{
         WEIGHT( MAX\_FILE ) = REAL (Returned)
      }{
         Weight of each input file
      }
      \sstsubsection{
         SHIFT\_DX( MAX\_FILE ) = REAL (Returned)
      }{
         X Shift of each input file
      }
      \sstsubsection{
         SHIFT\_DY( MAX\_FILE ) = REAL (Returned)
      }{
         Y Shift of each input file
      }
      \sstsubsection{
         NPARS = INTEGER (Given)
      }{
         Number of parameters in PARS array.
      }
      \sstsubsection{
         PARS( NPARS ) = REAL (Given)
      }{
         Values of input parameters. 1: WEIGHT, 2: SHIFT\_DX, 3: SHIFT\_DY
      }
      \sstsubsection{
         STATUS = INTEGER (Given and Returned)
      }{
         The global status
      }
   }
   \sstnotes{
      \sstitemlist{

         \sstitem
         This subroutine can be called recursively. A limit of 5 recursion
           levels is imposed. Note that recursion is not approved of in
           Fortran 77 (mainly because local variables remember their state
           on entry!).

         \sstitem
         Text files with extension .txt are converted to NDFs by the
           NDF\_OPEN command. This is bad - do not use .txt files as batch
           files.
      }
   }
   \sstdiytopic{
      Copyright
   }{
      Copyright \copyright 1995,1996,1997,1998,1999 Particle Physics and Astronomy
      Research Council. All Rights Reserved.
   }
}
\sstroutine{
   SURF\_REQUEST\_OUTPUT\_COORDS
}{
   Request longitude and latitude of output map from user
}{
   \sstdescription{
      Prompts user for output coordinates of map in the specified
      coordinate frame. The default values are derived from the
      supplied apparent RA/Dec centre. The selected centre
      is returned (in apparent ra/dec).
   }
   \sstinvocation{
       CALL SURF\_REQUEST\_OUTPUT\_COORDS( TASK, PARLONG, PARLAT,
            OUT\_COORDS, LAT\_OBS, DEF\_RA\_CEN, DEF\_DEC\_CEN, MJD, HOURS,
            OUT\_RA\_CEN, OUT\_DEC\_CEN, OUT\_ROTATION, OUT\_LONG, OUT\_LAT,
            STATUS)
   }
   \sstarguments{
      \sstsubsection{
         TASK = CHARACTER (Given)
      }{
         Description of task to be used in error messages
      }
      \sstsubsection{
         PARLONG = CHARACTER (Given)
      }{
         Name of the parameter used to request the LONGITUDE
      }
      \sstsubsection{
         PARLAT = CHARACTER (Given)
      }{
         Name of the parameter used to request the LATITUDE
      }
      \sstsubsection{
         OUT\_COORDS = CHARACTER (Given)
      }{
         Coordinate frame of output coordinates.
      }
      \sstsubsection{
         LAT\_OBS = DOUBLE (Given)
      }{
         Latitude of observatory (Radians)
      }
      \sstsubsection{
         DEF\_RA\_CEN = DOUBLE (Given)
      }{
         Apparent RA of the default map centre (radians)
      }
      \sstsubsection{
         DEF\_DEC\_CEN = DOUBLE (Given)
      }{
         Apparent Dec of the default map centre (radians)
      }
      \sstsubsection{
         MJD = DOUBLE (Given)
      }{
         Modified Julian date to be used as reference for apparent
         RA/Dec coordinates
      }
      \sstsubsection{
         HOURS = LOGICAL (Given)
      }{
         Flag to decide whether longitude is expressed as hours or degrees
      }
      \sstsubsection{
         OUT\_RA\_CEN = DOUBLE (Returned)
      }{
         Apparent RA of output map centre (radians)
      }
      \sstsubsection{
         OUT\_DEC\_CEN = DOUBLE (Returned)
      }{
         Apparent Dec of output map centre (radians)
      }
      \sstsubsection{
         OUT\_ROTATION = DOUBLE (Returned)
      }{
         angle between apparent N and N of output coord system (radians)
      }
      \sstsubsection{
         OUT\_LONG = DOUBLE (Returned)
      }{
         longitude of output map centre in the selected coordinate frame (rad)
      }
      \sstsubsection{
         OUT\_LAT = DOUBLE (Returned)
      }{
         latitude of output map centre in the selected coordinate frame (rad)
      }
      \sstsubsection{
         STATUS = INTEGER (Given \& Returned)
      }{
         Global status
      }
   }
   \sstnotes{
      \sstitemlist{

         \sstitem
         The output coordinates are always slightly different from the
           when defaults are accepted since there is a loss of precision
           converting the ra/dec to and from a string form.
      }
   }
   \sstdiytopic{
      Copyright
   }{
      Copyright \copyright 1995,1996,1997,1998,1999 Particle Physics and Astronomy
      Research Council. All Rights Reserved.
   }
}
\sstroutine{
   SURF\_SET\_APP\_NAME
}{
   Set application name for NDF history
}{
   \sstdescription{
      This routine sets the application name to be of the form

          TASKNAME   ( PACKAGE VN.n-n )

      That is, the taskname (which would be the default history
      if we did not change it) with the package name and the
      version number. The version number is specified at compile time.
   }
   \sstinvocation{
      CALL SURF\_SET\_APP\_NAME( TSKNAME, STATUS )
   }
   \sstarguments{
      \sstsubsection{
         TASKNAME = \_CHAR (Given)
      }{
         Name of the task.
      }
      \sstsubsection{
         STATUS = \_INTEGERS (Given \& Returned)
      }{
         Global status
      }
   }
   \sstnotes{
      \sstitemlist{

         \sstitem
         This routine should be compiled using the C pre-processor
           with a define for the PKG\_VERS cpp variable. This variable
           should include the required quotes (since CPP will not change
           quoted strings). An example could be:
                     f77 -DPKG\_VERS={\tt "}{\tt '}5.2-4{\tt '}{\tt "} surf\_set\_app\_name.F
           The V is prepended automatically.
      }
   }
   \sstdiytopic{
      Copyright
   }{
      Copyright \copyright 1995,1996,1997,1998,1999 Particle Physics and Astronomy
      Research Council. All Rights Reserved.
   }
}
\sstroutine{
   SURF\_WRITE\_DATA
}{
   Write bolometer positions and values to text file
}{
   \sstdescription{
      This routine writes the value, variance and position of each
      data point to a ASCII file. It is called as part of the EXTRACT\_DATA
      task.
      The interface is the same as that used in the REBIN task.
      The data and variance are in volts. The positions are in radians.
      The data are written out as columns (RA DEC DATA VAR) and subsets
      can be extracted by using SCUBA sections.
   }
   \sstinvocation{
       CALL SURF\_WRITE\_DATA( FD, NPTS, IN\_DATA, IN\_VARIANCE,
        BOL\_RA, BOL\_DEC, STATUS )
   }
   \sstarguments{
      \sstsubsection{
         FD = INTEGER (Given)
      }{
         Output file descriptor
      }
      \sstsubsection{
         NPTS = INTEGER (Given)
      }{
         Number of points in input data
      }
      \sstsubsection{
         IN\_DATA( NPTS ) = REAL (Given)
      }{
         Data value
      }
      \sstsubsection{
         IN\_VARIANCE( NPTS ) = REAL (Given)
      }{
         Variance on data
      }
      \sstsubsection{
         BOL\_RA( NPTS ) = DOUBLE (Given)
      }{
         Bolometer position (radians offset from map centre)
      }
      \sstsubsection{
         BOL\_DEC( NPTS ) = DOUBLE (Given)
      }{
         Bolometer position (radians offset from map centre)
      }
      \sstsubsection{
         STATUS = INTEGER (Given and Returned)
      }{
         The global status
      }
   }
   \sstnotes{
      For each file name that is entered, values for the parameters
      SELECT\_INTS, WEIGHT, SHIFT\_DX and SHIFT\_DY are requested.
      \sstitemlist{

         \sstitem
         The application can read in up to 100 separate input datasets.

         \sstitem
         No data is returned if the DATA or positions are bad.
         Data is still returned if Variance is bad.
      }
   }
   \sstdiytopic{
      ASCII input files
   }{
      The REF and IN parameters accept ASCII text files as input. These
      text files may contain comments (signified by a \#), NDF names,
      values for the parameters WEIGHT, SHIFT\_DX and SHIFT\_DY,
      and names of other ASCII files. There is one data file per line.
      An example file is:

\texttt{
\begin{tabular}{lcccl}
        file1\{b5\}  & 1.0& 0.5&   0.0&  \# Read bolometer 5 from file1.sdf\\
        file2        &    &    &      &\# Read file 2 but you will still be\\
                     &    &    &      &  \# prompted for WEIGHT, and shifts.\\
        file3\{i3\}- & 1.0& 0.0&   0.0&  \# Use everything except int 3\\
        test.bat     &    &    &      &\# Read in another text file\\
\end{tabular}
}

      Note that the parameters are position dependent and are not necessary.
      Missing parameters are requested. This means it is not possible to
      specify SHIFT\_DX (position 3) without specifying the WEIGHT.
      Also note that SCUBA sections can be specified with any input NDF.
   }
   \sstdiytopic{
      Related Applications
   }{
      SURF: REBIN, BOLREBIN, INTREBIN, CHANGE\_QUALITY
   }
   \sstdiytopic{
      Copyright
   }{
      Copyright \copyright 1995,1996,1997,1998,1999 Particle Physics and Astronomy
      Research Council. All Rights Reserved.
   }
}
\sstroutine{
   SURF\_WRITE\_MAP\_INFO
}{
   Calculate FITS and WCS information for rebinned image
}{
   \sstdescription{
      This routine writes Axis, FITS and WCS information to a rebinned
      map file. Note that this routine does not write WCS information
      to the FITS header - that is stored in the WCS information.
      Keywords CRVAL, CRPIX and CTYPE are used here to generate the
      WCS only.
   }
   \sstinvocation{
       CALL SURF\_WRITE\_MAP\_INFO( OUT\_NDF, OUT\_COORDS, OUT\_TITLE,
            MJD\_STANDARD, NFILE, FILENAME, OUT\_LONG, OUT\_LAT,
            OUT\_PIXEL, I\_CENTRE, J\_CENTRE, NX\_OUT, NY\_OUT, WAVELENGTH
            FILTER, WRITE\_CHOP, CHOP\_PA, CHOP\_THROW, WPLATE, ANGROT,
            STATUS )
   }
   \sstarguments{
      \sstsubsection{
         OUT\_NDF = INTEGER (Given)
      }{
         NDF identifier of output file
      }
      \sstsubsection{
         OUT\_COORDS = CHAR (Given)
      }{
         Output coordinate system
      }
      \sstsubsection{
         OUT\_TITLE = CHAR (Given)
      }{
         Title of map
      }
      \sstsubsection{
         OUT\_UNITS = CHAR (Given)
      }{
         Output units
      }
      \sstsubsection{
         MJD\_STANDARD = DOUBLE (Given)
      }{
         Modified Julian date of map
      }
      \sstsubsection{
         NFILE = INTEGER (Given)
      }{
         Number of files in input map
      }
      \sstsubsection{
         FILENAME = CHAR(NFILE) (Given)
      }{
         Names of input files
      }
      \sstsubsection{
         OUT\_PIXEL = REAL (Given)
      }{
         Size of pixel in output map (radians)
      }
      \sstsubsection{
         OUT\_LONG  = DOUBLE (Given)
      }{
         Longitude of output map (radians)
      }
      \sstsubsection{
         OUT\_LAT  = DOUBLE (Given)
      }{
         Latitude of output map (radians)
      }
      \sstsubsection{
         I\_CENTRE = INTEGER (Given)
      }{
         Position of X reference pixel in output map
      }
      \sstsubsection{
         J\_CENTRE = INTEGER (Given)
      }{
         Position of Y reference pixel in output map
      }
      \sstsubsection{
         NX\_OUT = INTEGER (Given)
      }{
         Size of output map in X direction
      }
      \sstsubsection{
         NY\_OUT = INTEGER (Given)
      }{
         Size of output map in Y direction
      }
      \sstsubsection{
         WAVELENGTH = REAL (Given)
      }{
         Wavelength of the map (microns)
      }
      \sstsubsection{
         FILTER = CHARACTER (Given)
      }{
         Name of filter
      }
      \sstsubsection{
         WRITE\_CHOP = LOGICAL (Given)
      }{
         Write the CHOP information to the header
      }
      \sstsubsection{
         CHOP\_CRD = CHARACTER (Given)
      }{
         Chop coordinate frame
      }
      \sstsubsection{
         CHOP\_PA = REAL (Given)
      }{
         Chop position angle in degrees east of north
      }
      \sstsubsection{
         CHOP\_THROW = REAL (Given)
      }{
         Chop throw in arcsec
      }
      \sstsubsection{
         WPLATE = REAL (Given)
      }{
         Wave plate angle (ignored if bad value).
      }
      \sstsubsection{
         ANGROT = REAL (Given)
      }{
         POLPACK rotation angle between waveplate and X pixel axis
         (ignored if bad value)
      }
      \sstsubsection{
         TELESCOPE = CHARACTER (Given)
      }{
         Name of telescope to write to FITS header
      }
      \sstsubsection{
         INSTRUMENT = CHARACTER (Given)
      }{
         Name of instrument to write to FITS header
      }
      \sstsubsection{
         STATUS = INTEGER (Given and Returned)
      }{
         The global status
      }
   }
   \sstdiytopic{
      Copyright
   }{
      Copyright \copyright 1995,1996,1997,1998,1999 Particle Physics and Astronomy
      Research Council. All Rights Reserved.
   }
}
\sstroutine{
   SURF\_WRITE\_PHOTOM
}{
   Routine to output ASCII results of PHOTOM reduction
}{
   \sstdescription{
      This routine writes out the results for 1 sub-instrument of a PHOTOM
      observation.
      For each bolometer that measured the source:-

\texttt{
\begin{tabular}{cccccc}
\multicolumn{2}{l}{Bolometer} &:&\multicolumn{3}{l}{$<$bolometer name$>$}\\
\multicolumn{2}{l}{Weight}    &:&\multicolumn{3}{l}{$<$weight to be given to its results$>$}\\
&&&&&\\
Integration &Peak  &   Peak\_sig&      Peak\_x &     Peak\_y  &    Quality\\
$<$integration$>$& $<$peak$>$&   $<$Error$>$ &  $<$x of peak$>$& $<$y of
peak$>$ &  $<$quality$>$\\
\multicolumn{6}{c}{for all the integrations taken in the observation}\\
&&&&&\\
\multicolumn{6}{l}{Measurement results   :}\\
\multicolumn{2}{c}{Fit to coadded jiggle:}&:&&\\
&$<$peak$>$ &  $<$variance$>$ &  $<$x of peak$>$&$<$y of peak$>$&$<$quality$>$\\
\multicolumn{2}{c}{Coadded fit results :}&:&&\\
&$<$peak$>$ &  $<$variance$>$&  &           &                  $<$quality$>$\\
\end{tabular}   
}
}
   \sstinvocation{
       CALL SURF\_WRITE\_PHOTOM (FD, MAX\_BEAM,
         N\_BOLS, BOL\_CHAN, BOL\_ADC, PHOT\_BB, MAX\_INT, N\_MEASUREMENTS,
         N\_INTEGRATIONS,PEAK\_D, PEAK\_V, PEAK\_X, PEAK\_Y, PEAK\_Q, BEAM\_WEIGHT,
         MEAS\_1\_D, MEAS\_1\_V, MEAS\_1\_X, MEAS\_1\_Y, MEAS\_1\_Q,
         MEAS\_2\_D, MEAS\_2\_V, MEAS\_2\_Q, STATUS)
   }
   \sstarguments{
      \sstsubsection{
         FD                     = INTEGER (Given)
      }{
         ASCII file descriptor
      }
      \sstsubsection{
         MAX\_BEAM               = INTEGER (Given)
      }{
         the maximum number of bolometers that can observe the source
         in a single observation
      }
      \sstsubsection{
         PARABOLA               = LOGICAL (Given)
      }{
         True if we fitted the coaaded jiggle with a parabola
      }
      \sstsubsection{
         N\_BOLS                 = INTEGER (Given)
      }{
         the number of bolometers used in the sub-instrument
      }
      \sstsubsection{
         BOL\_CHAN (N\_BOLS)      = INTEGER (Given)
      }{
         the channel numbers of the bolometers observing the object
      }
      \sstsubsection{
         BOL\_ADC (N\_BOLS)       = INTEGER (Given)
      }{
         the A/D numbers of the bolometers observing the object
      }
      \sstsubsection{
         PHOT\_BB (MAX\_BEAM)     = INTEGER (Given)
      }{
         the indices of the bolometers used to observe the source in
         each beam in the BOL\_CHAN and BOL\_ADC arrays
      }
      \sstsubsection{
         MAX\_INT                = INTEGER (Given)
      }{
         the maximum number of integrations in an observation
      }
      \sstsubsection{
         N\_MEASUREMENTS         = INTEGER (Given)
      }{
         the number of measurements in the observation
      }
      \sstsubsection{
         N\_INTEGRATIONS         = INTEGER (Given)
      }{
         the number of integrations in the observation
      }
      \sstsubsection{
         PEAK\_D (MAX\_INT, MAX\_BEAM) = REAL (Given)
      }{
         the fitted peak value for each integration with each
         bolometer
      }
      \sstsubsection{
         PEAK\_V (MAX\_INT, MAX\_BEAM) = REAL (Given)
      }{
         the variance on PEAK
      }
      \sstsubsection{
         PEAK\_X (MAX\_INT, MAX\_BEAM) = REAL (Given)
      }{
         the x offset of the fitted peak for each integration with
         each bolometer
      }
      \sstsubsection{
         PEAK\_Y (MAX\_INT, MAX\_BEAM) = REAL (Given)
      }{
         the y offset of the fitted peak
      }
      \sstsubsection{
         PEAK\_Q (MAX\_INT, MAX\_BEAM) = BYTE (Given)
      }{
         the quality of each fitted peak (0 is good)
      }
      \sstsubsection{
         BEAM\_WEIGHT (MAX\_BEAM) = REAL (Given)
      }{
         the weights assigned to the measurements with each bolometer
      }
      \sstsubsection{
         MEAS\_1\_D (MAX\_BEAM)    = REAL (given)
      }{
         the fitted peak to the coadded integrations
      }
      \sstsubsection{
         MEAS\_1\_V (MAX\_BEAM)    = REAL (Given)
      }{
         the variance on MEAS\_1\_D
      }
      \sstsubsection{
         MEAS\_1\_X (MAX\_BEAM)    = REAL (Given)
      }{
         the x offset of the peak fitted to the coadd
      }
      \sstsubsection{
         MEAS\_1\_Y (MAX\_BEAM)    = REAL (Given)
      }{
         the y offset
      }
      \sstsubsection{
         MEAS\_1\_Q (MAX\_BEAM)    = BYTE (Given)
      }{
         the quality on MEAS\_1\_D
      }
      \sstsubsection{
         MEAS\_2\_D (MAX\_BEAM)    = REAL (Given)
      }{
         the coadd of the peaks fitted to the individual integrations
         for each bolometer
      }
      \sstsubsection{
         MEAS\_2\_V (MAX\_BEAM)    = REAL (Given)
      }{
         the variance on MEAS\_2\_D
      }
      \sstsubsection{
         MEAS\_2\_Q (MAX\_BEAM)    = BYTE (Given)
      }{
         the quality on MEAS\_2\_D
      }
      \sstsubsection{
         STATUS                 = INTEGER (Given and returned)
      }{
         global status
      }
   }
   \sstdiytopic{
      Copyright
   }{
      Copyright \copyright 1995,1996,1997,1998,1999 Particle Physics and Astronomy
      Research Council. All Rights Reserved.
   }
}
\sstroutine{
   SURF\_WRITE\_PHOTOM\_HEADER
}{
   Write photom information to text file
}{
   \sstdescription{
      This routine writes out the header for 1 sub-instrument of a PHOTOM
      observation.
         If status is good on entry the routine will call FIO\_ASSOC to open
      the ASCII file to hold the results; the filename will be read from
      parameter {\tt '}FILE{\tt '}. FIO\_WRITE will then be called to write out the results
      in the following format:-

\texttt{
\begin{tabular}{lcl}
\multicolumn{3}{l}{Output from SURF reduction of a PHOTOM observation}\\
       Reduction date        &:& $<$date$>$\\
       Observation definition&:& $<$ODF name$>$\\
       Date of observation   &:& $<$date$>$\\
       Time of observation   &:& $<$time$>$\\
       Run number            &:& $<$number$>$\\
       Object                &:& $<$object$>$\\
       Sub-instrument        &:& $<$name$>$\\
       Filter                &:& $<$name$>$\\
       Centre coords         &:& $<$coord system of centre$>$\\
       Latitude              &:& $<$latitude$>$\\
       Longitude             &:& $<$longitude$>$\\
\end{tabular}
}

If the centre coordinate system is PLANET then:-

\texttt{
\begin{tabular}{lcl}
       2nd latitude          &:& $<$latitude on MJD2$>$\\
       2nd longitude         &:& $<$longitude on MJD2$>$\\
       date of 1st position  &:& $<$modified Julian day when source at lat,long$>$\\
       date of 2nd position  &:& $<$modified Julian day when source at lat2,long2$>$\\
\end{tabular}
}

      end if

\texttt{
\begin{tabular}{lcl}
       Offset coords         &:& $<$coord system of source offset$>$\\
       x offset              &:& $<$x offset of source$>$\\
       y offset              &:& $<$y offset of source$>$\\
       Sampe coords          &:& $<$coord system of jiggle offsets$>$\\
       Sample position angle &:& $<$angle that x axis of jiggle offsets is rotated\\
    &&                           anticlockwise from the x axis of the sample\\
    &&                           coord system$>$\\
       Sky error removal     &:& $<$TRUE if the SURF REMSKY application has \\
   &&                            been run on the data$>$\\
       Analysis mode         &:& AVERAGE or PARABOLA\\
\end{tabular}
}

   }
   \sstinvocation{
       CALL SURF\_WRITE\_PHOTOM\_HEADER (ODF, OBS\_DATE, OBS\_TIME,
            ANALYSIS, RUN\_NUMBER, OBJECT, SUB\_INSTRUMENT, FILTER,
            CENTRE\_COORDS, LAT, LONG, LAT2, LONG2, MJD1, MJD2,
            OFFSET\_COORDS, MAP\_X, MAP\_Y, SAMPLE\_COORDS, SAMPLE\_PA,
            SKY\_SUBTRACTION, FD, STATUS)
   }
   \sstarguments{
      \sstsubsection{
         ODF                    = CHARACTER$*$($*$) (Given)
      }{
         the name of the observation definition file
      }
      \sstsubsection{
         OBS\_DATE               = CHARACTER$*$($*$) (Given)
      }{
         the date of the observation
      }
      \sstsubsection{
         OBS\_TIME               = CHARACTER$*$($*$) (Given)
      }{
         the UT of the observation
      }
      \sstsubsection{
         RUN\_NUMBER             = INTEGER (Given)
      }{
         the observation number
      }
      \sstsubsection{
         OBJECT                 = CHARACTER$*$($*$) (Given)
      }{
         the name of the object observed
      }
      \sstsubsection{
         SUB\_INSTRUMENT         = CHARACTER$*$($*$) (Given)
      }{
         the name of the sub-instrument used
      }
      \sstsubsection{
         FILTER                 = CHARACTER$*$($*$) (Given)
      }{
         the name of the filter used
      }
      \sstsubsection{
         CENTRE\_COORDS          = CHARACTER$*$($*$) (Given)
      }{
         the coordinate system of the telescope centre coords
      }
      \sstsubsection{
         LAT                    = CHARACTER$*$($*$) (Given)
      }{
         the latitude of the telescope centre
      }
      \sstsubsection{
         LONG                   = CHARACTER$*$($*$) (Given)
      }{
         the longitude of the telescope centre
      }
      \sstsubsection{
         LAT2                   = CHARACTER$*$($*$) (Given)
      }{
         the second source latitude if CENTRE\_COORDS is PLANET
      }
      \sstsubsection{
         LONG2                  = CHARACTER$*$($*$) (Given)
      }{
         the second source longitude if CENTRE\_COORDS is PLANET
      }
      \sstsubsection{
         MJD1                   = DOUBLE PRECISION (Given)
      }{
         the modified Julian date of LAT,LONG if CENTRE\_COORDS is PLANET
      }
      \sstsubsection{
         MJD2                   = DOUBLE PRECISION (Given)
      }{
         the modified Julian date of LAT2,LONG2 if CENTRE\_COORDS is PLANET
      }
      \sstsubsection{
         OFFSET\_COORDS          = CHARACTER$*$($*$) (Given)
      }{
         the coordinate system of the map centre offsets
      }
      \sstsubsection{
         MAP\_X                  = REAL (Given)
      }{
         the x offset of the map centre
      }
      \sstsubsection{
         MAP\_Y                  = REAL (Given)
      }{
         the y offset of the map centre
      }
      \sstsubsection{
         SAMPLE\_COORDS          = CHARACTER$*$($*$) (Given)
      }{
         the coordinate system of the jiggle offsets
      }
      \sstsubsection{
         SAMPLE\_PA              = REAL (Given)
      }{
         the position angle by which the x-axis of the jiggle offsets is
         rotated (anti-clockwise) from the x-axis of the SAMPLE\_COORDS
         system
      }
      \sstsubsection{
         SKY\_SUBTRACTION        = LOGICAL (Given)
      }{
         .TRUE. if the sky error has been removed
      }
      \sstsubsection{
         MAX\_BEAM               = INTEGER (Given)
      }{
         the maximum number of bolometers that can observe the source
         in a single observation
      }
      \sstsubsection{
         PHOT\_BB (MAX\_BEAM)     = INTEGER (Given)
      }{
         the indices of the bolometers used to observe the source in
         each beam in the BOL\_CHAN and BOL\_ADC arrays
      }
      \sstsubsection{
         FD                     = INTEGER (Returned)
      }{
         File descriptor of output file
      }
      \sstsubsection{
         STATUS                 = INTEGER (Given and returned)
      }{
         global status
      }
   }
   \sstdiytopic{
      Copyright
   }{
      Copyright \copyright 1995,1996,1997,1998,1999 Particle Physics and Astronomy
      Research Council. All Rights Reserved.
   }
}


\subsection{SURFLIB}

This section describes the subroutines that are available in the
\texttt{surflib} distribution (\texttt{libsurflib.a}).

\sstroutine{
   SURFLIB\_2DFT\_CHOP
}{
   Calculate the FT of the 2D chop throw
}{
   \sstdescription{
      Given the chop parameters calculate the FT of the chop.
      The dual beam chop is treated as two delta functions of
      opposite sign separated by the chop throw and centred at the
      middle of the data array. The Fourier transform of this function
      is a sine wave in the direction of the chop (zero at the middle
      beam) with wavelength related to the spatial frequency of the
      array (pixel spacing).
   }
   \sstinvocation{
       CALL SURFLIB\_2DFT\_CHOP (CHOP\_THROW, CHOP\_PA, PIXSIZE,
            NX, NY, FT\_DATA, FT\_VARIANCE, WT\_DATA, WT\_VARIANCE, STATUS)
   }
   \sstarguments{
      \sstsubsection{
         CHOP\_THROW                  = REAL (Given)
      }{
         chop throw in arcsec
      }
      \sstsubsection{
         CHOP\_PA                     = REAL (Given)
      }{
         position angle of chop. Positive is anti-clockwise
         from North. This angle is in degrees.
      }
      \sstsubsection{
         PIXSIZE                     = REAL (Given)
      }{
         pixel size in arcsec
      }
      \sstsubsection{
         NX                          = INTEGER (Given)
      }{
         Number of pixels in X
      }
      \sstsubsection{
         NY                          = INTEGER (Given)
      }{
         Number of pixels in Y
      }
      \sstsubsection{
         FT\_DATA (IDIMS(1), IDIMS(2)) = REAL (Returned)
      }{
         data array containing F.T. of chop
      }
      \sstsubsection{
         FT\_VARIANCE (IDIMS(1), IDIMS(2))       = REAL (Returned)
      }{
         variance on DATA; set to 0
      }
      \sstsubsection{
         WT\_DATA (IDIMS(1), IDIMS(2)) = REAL (Returned)
      }{
         data array containing weight of chop
      }
      \sstsubsection{
         WT\_VARIANCE (IDIMS(1), IDIMS(2))       = REAL (Returned)
      }{
         variance on DATA; set to 0
      }
      \sstsubsection{
         STATUS                      = INTEGER (Given and returned)
      }{
         global status
      }
   }
   \sstdiytopic{
      Copyright
   }{
      Copyright \copyright 1995,1996,1997,1998,1999 Particle Physics and Astronomy
      Research Council. All Rights Reserved.
   }
}
\sstroutine{
   SURFLIB\_CALC\_DUAL\_BEAM
}{
   Calculates a chopped map from a single beam map
}{
   \sstdescription{
      This routine can be used to added a chopped beam response
      to a single beam image. When NBEAMS=2 this routine
      calculates a middle-beam dual-beam map by
      calculating the difference between the left and right
      beams for each pixel in the map.

      When NBEAMS=3 a triple beam map is constructed where the central
      beam is the difference between the central beam and the average of
      the positions a chop distance away on either side.

      All pixels off the edge of the image are assumed to
      have flux 0 and variance 0.

      Bad pixels are assumed to be equivalent to flux of 0
      and variance 0, although for the triple beam image a bad pixel
      is retained if it coincide with the middle beam position.
   }
   \sstinvocation{
       CALL SURFLIB\_CALC\_CHOPPED\_IMAGE( NBEAMS, CHOP\_THR, CHOP\_PA,
            DIM1, DIM2, IN\_DATA, OUT\_DATA, USEVAR, IN\_VAR, OUT\_VAR,
            STATUS)
   }
   \sstarguments{
      \sstsubsection{
         NBEAMS = INTEGER (Given)
      }{
         Number of beams to add. (2 or 3)
      }
      \sstsubsection{
         CHOP\_THR = REAL (Given)
      }{
         Chop throw in pixels
      }
      \sstsubsection{
         CHOP\_PA  = REAL (Given)
      }{
         Chop position angle (east of north) in degrees
      }
      \sstsubsection{
         DIM1 = INTEGER (Given)
      }{
         First dimension of image
      }
      \sstsubsection{
         DIM2 = INTEGER (Given)
      }{
         Second dimension of image [can not use DIMS as an array
         since the linux fortran compiler doesnt like it]
      }
      \sstsubsection{
         IN\_DATA = REAL (Given)
      }{
         Input single-beam image
      }
      \sstsubsection{
         OUT\_DATA = REAL (Returned)
      }{
         Output dual-beam image
      }
      \sstsubsection{
         USEVAR = LOGICAL (Given)
      }{
         Are we processing a variance array?
      }
      \sstsubsection{
         IN\_VAR = REAL (Given)
      }{
         Input variance image (if USEVAR=TRUE)
      }
      \sstsubsection{
         OUT\_VAR = REAL (Returned)
      }{
         Output variance image (if USEVAR=TRUE)
      }
      \sstsubsection{
         STATUS = INTEGER (Given and Returned)
      }{
         Global status
      }
   }
   \sstdiytopic{
      Copyright
   }{
      Copyright \copyright 1995,1996,1997,1998,1999 Particle Physics and Astronomy
      Research Council. All Rights Reserved.
   }
}
\sstroutine{
   SURFLIB\_CALC\_GRIDIJ
}{
   Calculate the position of a pixel number in an IJ grid
}{
   \sstdescription{
      Calculates the I,J grid positions related to pixel numbers
      1 to NX$*$NY.
      Must calculate all positions in one go since it is very inefficient
      to calculate the spiral positions on demand
      Three modes types are currently available.
   }
   \sstinvocation{
      CALL SURFLIB\_CALC\_GRIDIJ( TYPE, NX, NY, ICEN, JCEN, I, J, STATUS)
   }
   \sstarguments{
      \sstsubsection{
         TYPE = CHAR (Given)
      }{
         Type of unwrapping. Can be one of:
         XLINEAR - unfold each X strip in turn for each Y
         YLINEAR - unfold each Y strip in turn for each X
         SPIRAL - unfold in a spiral from the reference pixels (ICEN,JCEN)
         DIAG1  - Diagonal starting at 1,1
         DIAG2  - Diagonal starting at NX,1
      }
      \sstsubsection{
         NX = INTEGER (Given)
      }{
         Number of pixels in X
      }
      \sstsubsection{
         NY = INTEGER (Given)
      }{
         Number of pixels in Y
      }
      \sstsubsection{
         ICEN = INTEGER (Given)
      }{
         Location of reference pixel in X
      }
      \sstsubsection{
         JCEN = INTEGER (Given)
      }{
         Location of reference pixel in Y
      }
      \sstsubsection{
         I ( NX $*$ NY ) = INTEGER (Returned)
      }{
         X Positions in grid, indexed by pixel number
      }
      \sstsubsection{
         J ( NX $*$ NY ) = INTEGER (Returned)
      }{
         Y Positions in grid, indexed by pixel number
      }
      \sstsubsection{
         STATUS = INTEGER (Given \& Returned)
      }{
         Global Status
      }
   }
   \sstdiytopic{
      Copyright
   }{
      Copyright \copyright 1995,1996,1997,1998,1999 Particle Physics and Astronomy
      Research Council. All Rights Reserved.
   }
}
\sstroutine{
   SURFLIB\_CALC\_IJPOS
}{
   Calculate the position of data on an XY grid
}{
   \sstdescription{
      Given an XY position, calculates the IJ position in the grid.
      The I,J{\tt '}s are then placed into the output array (of dimension
      2, N\_PTS)
   }
   \sstinvocation{
       CALL SURFLIB\_CALC\_IJPOS (N\_PTS, PIXEL\_SZ,
            I\_CENTRE, J\_CENTRE, X\_POS, Y\_POS, IJ,
            STATUS )
   }
   \sstarguments{
      \sstsubsection{
         N\_PTS = INTEGER (Given)
      }{
         Number of X,Y pairs supplied
      }
      \sstsubsection{
         PIXEL\_SZ = DOUBLE (Given)
      }{
         Size of each pixel (radians)
      }
      \sstsubsection{
         I\_CENTRE = INTEGER (Given)
      }{
         Location of reference pixel in X
      }
      \sstsubsection{
         J\_CENTRE = INTEGER (Given)
      }{
         Location of reference pixel in Y
      }
      \sstsubsection{
         X\_POS = DOUBLE PRECISION (Given)
      }{
         X positions
      }
      \sstsubsection{
         Y\_POS = DOUBLE PRECISION (Given)
      }{
         Y positions
      }
      \sstsubsection{
         IJ ( 2, N\_PTS ) = INTEGER (Given)
      }{
         Positions of each point (I,J) in output grid
      }
      \sstsubsection{
         STATUS = INTEGER (Given \& Returned)
      }{
         Global Status
      }
   }
   \sstdiytopic{
      Copyright
   }{
      Copyright \copyright 1995,1996,1997,1998,1999 Particle Physics and Astronomy
      Research Council. All Rights Reserved.
   }
}
\sstroutine{
   SURFLIB\_CALC\_OUTPUT\_GRID
}{
   Calculate the grid required to contain data given the X,Y positions
}{
   \sstdescription{
      Finds the extent of the data (given X and Y) and then
      reports back the size of the grid that is necessary to contain
      the data. Also returns the reference pixel position.
   }
   \sstinvocation{
       CALL SURFLIB\_CALC\_OUTPUT\_GRID (N\_FILES, N\_PTS, PIXEL\_SZ,
            X\_PTR, Y\_PTR, NX, NY, I\_CENTRE, J\_CENTRE,
            STATUS )
   }
   \sstarguments{
      \sstsubsection{
         N\_FILES = INTEGER (Given)
      }{
         Number of data sets (ie files)
      }
      \sstsubsection{
         N\_PTS = INTEGER (Given)
      }{
         Number of X,Y pairs supplied
      }
      \sstsubsection{
         PIXEL\_SZ = REAL (Given)
      }{
         Size of each pixel (radians)
      }
      \sstsubsection{
         X\_PTR( N\_FILES ) = INTEGER (Given)
      }{
         Pointers to arrays containing DOUBLE PRECISION X positions
      }
      \sstsubsection{
         Y\_PTR( N\_FILES ) = INTEGER (Given)
      }{
         Pointers to arrays containing DOUBLE PRECISION Y positions
      }
      \sstsubsection{
         NX = INTEGER (Given)
      }{
         Required number of points in X in output grid
      }
      \sstsubsection{
         NY = INTEGER (Given)
      }{
         Required number of points in Y in output grid
      }
      \sstsubsection{
         I\_CENTRE = INTEGER (Given)
      }{
         Location of reference pixel in X
      }
      \sstsubsection{
         J\_CENTRE = INTEGER (Given)
      }{
         Location of reference pixel in Y
      }
      \sstsubsection{
         STATUS = INTEGER (Given \& Returned)
      }{
         Global Status
      }
   }
   \sstnotes{
      \sstitemlist{

         \sstitem
         This routine requires double precision arguments for the X and
           Y coordinates.
      }
   }
   \sstdiytopic{
      Copyright
   }{
      Copyright \copyright 1995,1996,1997,1998,1999 Particle Physics and Astronomy
      Research Council. All Rights Reserved.
   }
}
\sstroutine{
   SURFLIB\_CALC\_POLPACK\_ANGROT
}{
   Average rotations angles over integration
}{
   \sstdescription{
      This routine calculates the mean rotation angle for each
      integration. The input data contains one angle per sample
      and this is averaged to calculate the rotation angle per
      integration.
      The input angle is the rotation angle (anti-clockwise)
      between the RA frame and the nasmyth (waveplate) frame.
      The output angle is the angle between the waveplate 0 and
      the X PIXEL axis (not the RA axis)
      Since the supplied angle is the angle from the Y-axis
      (ie all angles are from Nasmyth Y axis) we need to add
      90 degrees
   }
   \sstinvocation{
      CALL SURFLIB\_CALC\_POLPACK\_ANGROT(N\_POS, N\_EXPOSURES,
           N\_INTEGRATIONS, N\_MEASUREMENTS, DEM\_PNTR, ANGROT\_IN,
           ANGROT\_OUT, STATUS)
   }
   \sstarguments{
      \sstsubsection{
         N\_POS = INTEGER (Given)
      }{
         Size of ANGROT\_IN array
      }
      \sstsubsection{
         N\_EXPOSURES = INTEGER (Given)
      }{
         Number of exposures in DEM\_PNTR
      }
      \sstsubsection{
         N\_INTEGRATIONS = INTEGER (Given)
      }{
         Number of integrations in DEM\_PNTR
      }
      \sstsubsection{
         N\_MEASUREMENTS = INTEGER (Given)
      }{
         Number of measurements in DEM\_PNTR
      }
      \sstsubsection{
         DEM\_PNTR(1,N\_EXPOSURES,N\_INTEGRATIONS,N\_MEASUREMENTS) = INTEGER (Given)
      }{
         Start position for each EXP, INT, MEAS
      }
      \sstsubsection{
         ANGROT\_IN (N\_POS) = REAL (Given)
      }{
         Rotation for each sample (N\_POS)
      }
      \sstsubsection{
         ANGROT\_OUT (N\_INTEGRATIONS, N\_MEASUREMENTS) = REAL (Returned)
      }{
         Rotation angle averaged over each integration
      }
      \sstsubsection{
         ANGROT\_VAR (N\_INTEGRATIONS, N\_MEASUREMENTS) = REAL (Returned)
      }{
         Variance on averaged angle
      }
      \sstsubsection{
         STATUS = INTEGER (Given \& Returned)
      }{
         Global status
      }
   }
   \sstdiytopic{
      Copyright
   }{
      Copyright \copyright 1995,1996,1997,1998,1999 Particle Physics and Astronomy
      Research Council. All Rights Reserved.
   }
}
\sstroutine{
   SURFLIB\_CLIP\_GRID
}{
   Remove spikes from grid
}{
   \sstdescription{
      Given the binned data calculated by SURFLIB\_FILL\_GRID
      and the limits calculated by SURFLIB\_STATS\_GRID,
      clip all points that lie outside the statistical limits provided
      by the stats array.
   }
   \sstinvocation{
       CALL SURFLIB\_CLIP\_GRID(N\_FILES, N\_PTS, NX, NY, NMAX, NSIGMA,
            N\_BOLS, BITNUM, QUALITY\_PTR, BINS, BIN\_POS, STATS, IPOS,
            JPOS, NSPIKES, STATUS)
   }
   \sstarguments{
      \sstsubsection{
         N\_FILES = INTEGER (Given)
      }{
         Number of input data sets
      }
      \sstsubsection{
         N\_PTS(N\_FILES) = INTEGER (Given)
      }{
         Number of points in each data set
      }
      \sstsubsection{
         NX = INTEGER (Given)
      }{
         Size of X dimension
      }
      \sstsubsection{
         NY = INTEGER (Given)
      }{
         Size of Y dimension
      }
      \sstsubsection{
         NMAX = INTEGER (Given)
      }{
         Maximum value allowed for third dimension of BINS
      }
      \sstsubsection{
         N\_BOLS ( N\_FILES ) = INTEGER (Given)
      }{
         Number of bolometers per file (used for constructing message for users)
      }
      \sstsubsection{
         BITNUM = INTEGER (Given)
      }{
         Bit number that is modified when set bad
      }
      \sstsubsection{
         DATA\_PTR( N\_FILES ) = INTEGER (Given \& Returned)
      }{
         The input data with each spike replaced with a median value
      }
      \sstsubsection{
         QUALITY\_PTR(N\_FILES) = INTEGER (Given \& Returned)
      }{
         Pointers to quality arrays. Note that quality arrays are modified.
         (The whole point of this routine) even though the pointers are not.
      }
      \sstsubsection{
         BINS(NX, NY, NMAX) = REAL (Given)
      }{
         The data stored in relation to its position
      }
      \sstsubsection{
         BIN\_POS (NX, NY, NMAX) = INTEGER (Given)
      }{
         Position of each data point in quality arrays
      }
      \sstsubsection{
         STATS (NX, NY, 3) = REAL (Given)
      }{
         Statistics of each bin. 1=Median, 2=High, 3=Low
      }
      \sstsubsection{
         PNTS(NMAX) = REAL (Given)
      }{
         Scratch space for copying in the data from each I,J
      }
      \sstsubsection{
         IPOS(NX$*$NY) = INTEGER (Given)
      }{
         I position in array
      }
      \sstsubsection{
         JPOS(NX$*$NY) = INTEGER (Given)
      }{
         J position in array
      }
      \sstsubsection{
         NSPIKES = INTEGER (Returned)
      }{
         Number of spikes removed for a given file
      }
      \sstsubsection{
         STATUS = INTEGER (Given \& Returned)
      }{
         Global Status
      }
   }
   \sstdiytopic{
      Copyright
   }{
      Copyright \copyright 1995,1996,1997,1998,1999 Particle Physics and Astronomy
      Research Council. All Rights Reserved.
   }
}
\sstroutine{
   SURFLIB\_DECODE\_REMSKY\_STRING
}{
   Calculate bolometer list from remsky input
}{
   \sstdescription{
      Given an array of bolometer descriptions, return an array
      of bolometer numbers.
      Input array can contain a variety of forms for identifying bolometers
      or lists of bolometers. This string is generated by, for example, the
      REMSKY task.
   }
   \sstinvocation{
       CALL SURFLIB\_DECODE\_REMSKY\_STRING( SUB\_INSTRUMENT,
            N\_ELEMENTS, BOL\_DESC, N\_BOLS, BOL\_ADC, BOL\_CHAN,
            BOL\_LIST, N\_BOLS\_OUT,
            STATUS )
   }
   \sstarguments{
      \sstsubsection{
         SUB\_INSTRUMENT = CHARACTER (Given)
      }{
         Sub instrument name. This governs the bolometers that
         are allowed to be returned.
      }
      \sstsubsection{
         N\_ELEMENTS = INTEGER (Given)
      }{
         Number of elements in the input array
      }
      \sstsubsection{
         BOL\_DESC( N\_ELEMENTS ) = CHARACTER$*$($*$) (Given)
      }{
         Array of bolometer descriptions
         Can take the form of:

\begin{tabular}{ll}
     \texttt{n}     &   bolometer number\\
     \texttt{id}    &   bolometer id (eg H7)\\
     \texttt{all}   &   all bolometers\\
     \texttt{rn}    &   Nth ring of bolometers\\
\end{tabular}

         Can be prefixed with a minus sign to indicate that a bolometer
         should be removed from the final list. The calculation is sequential.
      }
      \sstsubsection{
         N\_BOLS = INTEGER (Given)
      }{
         Maximum allowed size for return array containing bolometer numbers
         An error is given if this number is exceeded. In effect this
         is equivalent to the actual number of bolometers on the array.
      }
      \sstsubsection{
         BOL\_ADC( N\_BOLS ) = INTEGER (Given)
      }{
         ADC numbers for identifying bolometers
      }
      \sstsubsection{
         BOL\_CHAN( N\_BOLS ) = INTEGER (Given)
      }{
         Channel number for identifying bolometers
      }
      \sstsubsection{
         BOL\_LIST( N\_BOLS ) = INTEGER (Returned)
      }{
         List of bolometer numbers corresponding to the supplied strings.
      }
      \sstsubsection{
         N\_BOLS\_OUT = INTEGER (Returned)
      }{
         Actual number of bolometers returned in BOL\_LIST
      }
      \sstsubsection{
         STATUS = INTEGER (Given \& Returned)
      }{
         Global status
      }
   }
   \sstdiytopic{
      Syntax
   }{
\begin{tabular}{ll}
\texttt{[all]}              & Whole array\\
\texttt{[r0]}               & Ring zero (central pixel)\\
\texttt{[h7,r1]}            & inner ring and h7\\
\texttt{[r1,-h8]}           & inner ring without h8\\
\texttt{[r1,-18]}           & inner ring without bolometer 18\\
\texttt{[all,-r1,-h7]}      & all pixels except the inner ring/h7\\
\texttt{[all,-r3,g1]}       & all pixels except ring 3 but with
                             g1 (g1 is in r3)\\
\texttt{[all,-r1,-r2,-r3,-r4,-r5]} & Selects the central pixel!!\\
\end{tabular}

       Note that the sequence is evaluated in turn so that \texttt{[all,-all,h7]} will
       still end up with pixel \texttt{h7}.
   }
   \sstdiytopic{
      Copyright
   }{
      Copyright \copyright 1995,1996,1997,1998,1999 Particle Physics and Astronomy
      Research Council. All Rights Reserved.
   }
}
\sstroutine{
   SURFLIB\_DIFF\_DESPIKE
}{
   Despike data scan by detecting points
   more than nsigma from running 2-point mean
}{
   \sstdescription{
      This routine removes spikes from SCAN/MAP observations.
      The scan map differential despiking algorithm uses 2 criteria
      to decide which points are spikes.

      First, for each bolometer used a pass is made through each
      scan calculating for each point:-

\begin{equation}
        diff(i) = point(i) - \frac{(point(i-1) + point(i+1))}{2.0}
\end{equation}

      Values of {\tt '}diff{\tt '} for the first and last points in the scan are
      calculated in a similar way but subtracting the mean of points
      2 and 3 and points n-1 and n-2 respectively.

      The mean and standard deviation of {\tt '}diff{\tt '} are calculated by
      coadding the 10 points at each end of the scan where,
      hopefully, there is no source emission. Spikes in these
      regions are handled by removing points from the coadd that lie
      further than 3 sigma from the mean, then redoing the
      calculation recursively until no further points need be
      removed.

      The first criterion for a spike is that it{\tt '}s {\tt '}diff{\tt '} value
      should be further from the mean of {\tt '}diff{\tt '} by NSIGMA times the
      sigma derived from the endpoints.

      The problem with this simple approach is that bright sources
      in the scan themselves lead to excursions in {\tt '}diff{\tt '} that can
      be wrongly identified as spikes. To prevent this happening a
      second criterion is used. In this the scan values are
      convolved with a 3 sample wide box so that each {\tt '}box{\tt '} point is
      the average of the point itself and the points on either side of
      it. {\tt '}Box{\tt '} is expected to increase faster for real sources than
      for spikes because in them the increase will be spread over
      all 3 averaged points rather than just 1.

      The second criterion for a spike is met, therefore, if a
      point{\tt '}s {\tt '}diff{\tt '} is further from the {\tt '}diff{\tt '} mean than the value
      of {\tt '}box{\tt '} at that point.

      Fixed-up values for points that have identified as spikes are
      calculated by interpolating between the closest healthy points
      on either side.

      The second spike criterion also means unfortunately that the
      technique is less sensitive to spikes on bright sources than
      elsewhere. In addition, it is still possible to clip bright
      sources if too low a value for NSIGMA is used. It is
      recommended to run despike several times with different values
      of NSIGMA. Begin with NSIGMA=5, look at the result to see how
      effective despiking has been, then repeat the process with
      NSIGMA=4.5, 4.0 etc. until you start to clip source
      information.
   }
   \sstinvocation{
       CALL SURFLIB\_DIFF\_DESPIKE (N\_EXPOSURES, N\_INTEGRATIONS,
         N\_MEASUREMENTS, DEMOD\_POINTER, N\_BOL, N\_POS, IN\_DATA,
         IN\_VARIANCE, IN\_QUALITY, BADBIT, NSIGMA, OUT\_DATA,
         OUT\_VARIANCE, OUT\_QUALITY, NUM\_SPIKES, STATUS)
   }
   \sstarguments{
      \sstsubsection{
         N\_EXPOSURES                 = INTEGER (Given)
      }{
         maximum number of exposures per integration
      }
      \sstsubsection{
         N\_INTEGRATIONS              = INTEGER (Given)
      }{
         number of integrations in the observation
      }
      \sstsubsection{
         N\_MEASUREMENTS              = INTEGER (Given)
      }{
         number of measurements in the observation
      }
      \sstsubsection{
         DEMOD\_POINTER (N\_EXPOSURES, N\_INTEGRATIONS, N\_MEASUREMENTS)
      }{
         array pointing to start and finish of scans in IN\_DATA
      }
      \sstsubsection{
         N\_BOL                       = INTEGER (Given)
      }{
         the number of bolometers for which data was taken
      }
      \sstsubsection{
         N\_POS                       = INTEGER (Given)
      }{
         the number of positions measured in the scan
      }
      \sstsubsection{
         IN\_DATA (N\_BOL, N\_POS)      = REAL (Given)
      }{
         the measured data
      }
      \sstsubsection{
         IN\_VARIANCE (N\_BOL, N\_POS)  = REAL (Given)
      }{
         the variance on IN\_DATA
      }
      \sstsubsection{
         IN\_QUALITY (N\_BOL, N\_POS)   = BYTE (Given)
      }{
         the quality on IN\_DATA
      }
      \sstsubsection{
         BADBIT                      = BYTE (Given)
      }{
         bad bit mask
      }
      \sstsubsection{
         NSIGMA                      = REAL (Given)
      }{
         cut-off sigma for spikes
      }
      \sstsubsection{
         OUT\_DATA (N\_BOL, N\_POS)     = REAL (Returned)
      }{
         the deconvolved data
      }
      \sstsubsection{
         OUT\_VARIANCE (N\_BOL, N\_POS) = REAL (Returned)
      }{
         the variance on OUT\_DATA
      }
      \sstsubsection{
         OUT\_QUALITY (N\_BOL, N\_POS)  = BYTE (Returned)
      }{
         the quality on OUT\_DATA
      }
      \sstsubsection{
         NUM\_SPIKES                  = INTEGER (Returned)
      }{
         Total number of spikes detected
      }
      \sstsubsection{
         STATUS                      = INTEGER (Given and returned)
      }{
         global status
      }
   }
}
\sstroutine{
   SURFLIB\_FILL\_GRID
}{
   Populate grid with data points corresponding to position
}{
   \sstdescription{
      Given an array of I,J coordinates associated with data values,
      this routine places each data point onto the rectangular grid
      and the position in the data array into a corresponding grid.
      Grid is of size (NX,NY,NMAX) and each data point is placed into
      a different layer (up to NMAX) for each I,J. We already know the
      value of NMAX since the output array has been pre-allocated.
   }
   \sstinvocation{
       CALL SURFLIB\_FILL\_GRID (N\_PTS, NX, NY, NMAX, OFFSET,
            IN\_DATA, IN\_QUALITY, BADBIT, IJ, GRID, BINS, BIN\_POS,
            STATUS )
   }
   \sstarguments{
      \sstsubsection{
         N\_PTS = INTEGER (Given)
      }{
         Number of I,J pairs supplied
      }
      \sstsubsection{
         NX = INTEGER (Given)
      }{
         Size of X dimension
      }
      \sstsubsection{
         NY = INTEGER (Given)
      }{
         Size of Y dimension
      }
      \sstsubsection{
         NMAX = INTEGER (Given)
      }{
         Maximum value allowed for third dimension of BINS
      }
      \sstsubsection{
         OFFSET = INTEGER (Given)
      }{
         This is the offset in the IJ data array. Not used for IJ itself
         (Since this is added to the pointer of the input array if needed)
         in this array but is added onto each of the values placed into
         the BIN\_POS array so that we can keep track of the input data
         without using a further two pieces of information.
      }
      \sstsubsection{
         IN\_DATA(N\_PTS) = REAL (Given)
      }{
         Input data
      }
      \sstsubsection{
         IN\_QUALITY(N\_PTS) = BYTE (Given)
      }{
         Input quality
      }
      \sstsubsection{
         BADBIT = BYTE (Given)
      }{
         Bad bits mask for quality array
      }
      \sstsubsection{
         IJ ( 2, N\_PTS ) = INTEGER (Given)
      }{
         Positions of each point (I,J) in output grid for each N\_PTS
      }
      \sstsubsection{
         GRID (NX, NY) = INTEGER (Given \& Returned)
      }{
         Scratch space for keeping track of the number of points added
         into each I,J
      }
      \sstsubsection{
         BINS(NX, NY, NMAX) = REAL (Given \& Returned)
      }{
         The data stored in relation to its position
      }
      \sstsubsection{
         BIN\_POS (NX, NY, NMAX) = INTEGER (Given \& Returned)
      }{
         The position in the input data arrays associated with each data point
         in BINS
      }
      \sstsubsection{
         STATUS = INTEGER (Given \& Returned)
      }{
         Global Status
      }
   }
   \sstnotes{
      \sstitemlist{

         \sstitem
         GRID is not initialised by this routine since it is incremented
           each time a point is placed in BINS

         \sstitem
         BINS and BIN\_POS are not initialised here but should be filled
           with bad before this routine is first called.

         \sstitem
         a status of SAI\_\_WARN is returned if the indices lie outside
           the model area. It is up to the calling routine to decide
           whether to halt.
      }
   }
   \sstdiytopic{
      Copyright
   }{
      Copyright \copyright 1995,1996,1997,1998,1999 Particle Physics and Astronomy
      Research Council. All Rights Reserved.
   }
}
\sstroutine{
   SURFLIB\_FILL\_POLPACK\_ANGLES
}{
   Average
}{
   \sstdescription{
      Copy the waveplate and rotation angles from WPLATE and ANGROT
      (Which have dimensions (N\_INT,N\_MEAS) to a 2-dim array containing
      the angles (waveplate and rotation) for each integration and
      (in a separate array) the angles for each measurement.
      Note that since the rotation angles are stored in the file
      as one per integration, they need to be averaged over the
      integration in order to calculate the correct value for the
      measurement. ANG\_MEAS and ANG\_INT should be set to VAL\_\_BADR
      before entry to this routine.
   }
   \sstinvocation{
       CALL SURFLIB\_FILL\_POLPACK\_ANGLES( MAX\_FILE, MAX\_INT,
            MAX\_MEAS, N\_FILE, N\_INT, N\_MEAS,
            WPLATE, ANGROT, FAST\_AXIS, ANG\_INT, ANG\_MEAS, STATUS)
   }
   \sstarguments{
      \sstsubsection{
         MAX\_FILE = INTEGER (Given)
      }{
         First Dimension of ANG\_INT,ANG\_MEAS
      }
      \sstsubsection{
         MAX\_INT = INTEGER (given)
      }{
         Second dimension of ANG\_INT
      }
      \sstsubsection{
         MAX\_MEAS = INTEGER (given)
      }{
         Second dimension of ANG\_MEAS
      }
      \sstsubsection{
         N\_FILE = INTEGER (Given)
      }{
         Current slice in ANG\_INT, ANG\_MEAS
      }
      \sstsubsection{
         N\_INT = INTEGER (Given)
      }{
         Number of integrations per measurement
      }
      \sstsubsection{
         N\_MEAS = INTEGER (Given)
      }{
         Number of measurements
      }
      \sstsubsection{
         WPLATE(N\_INT, N\_MEAS) = REAL (Given)
      }{
         Waveplate angles for each integration
      }
      \sstsubsection{
         ANGROT(N\_INT, N\_MEAS) = REAL (Given)
      }{
         Waveplate angles for each integration
      }
      \sstsubsection{
         FAST\_AXIS = REAL (Given)
      }{
         Angle of the fast axis of the waveplate and the zero position
         of the waveplate. Added to each waveplate angle to calculate
         the true angle of the waveplate.
      }
      \sstsubsection{
         ANG\_INT(MAX\_FILE, MAX\_INT, 2) = REAL (Returned)
      }{
         Waveplate and rotation angles for each integration
         1=Waveplate, 2=Rotation. The waveplate angle is corrected for
         fast axis.
      }
      \sstsubsection{
         ANG\_MEAS(MAX\_FILE, MAX\_MEAS, 2) = REAL (Returned)
      }{
         Waveplate and rotation angles for each measurement.
         1=Waveplate, 2=Rotation. The waveplate angle is corrected for
         fast axis. The rotation angle is average over each integration
         during the measurement. The waveplate angle is set to bad
         if it is not constant for the given measurement.
      }
      \sstsubsection{
         STATUS = INTEGER (Given and Returned)
      }{
         Global status
      }
   }
   \sstdiytopic{
      Copyright
   }{
      Copyright \copyright 1995,1996,1997,1998,1999 Particle Physics and Astronomy
      Research Council. All Rights Reserved.
   }
}
\sstroutine{
   SURFLIB\_FILL\_WPLATE
}{
   Populate the WavePlate array
}{
   \sstdescription{
      Populate an array containing the waveplate angle for each
      sample. This can be calculated in two ways. If the SCUCD\_WPLATE
      array is present (N\_WPLATE$>$0), then this array contains the waveplate
      position for each measurement. If this array is not present
      the Waveplate angle is assumed to start at 0 degrees and increment
      by 22.5 degrees for each integration. If the number of
      measurements exceeds the number of positions in the SCUCD\_WPLATE
      array the sequence wraps round to the first value.
      The final output array will contain a single number (waveplate
      position) per sample (up to N\_POS). The output value will be
      in degrees.

      Since integrations and measurements are stored sequentially
      in a data file (in DEM\_PNTR) there is no need to use
      SCULIB\_FIND\_SWITCH to determine the start and end points
      in the data array - use SCULIB\_FIND\_INT.

      The waveplate angles are also stored in a small
      array containing a single wave plate angle for each integration
      and each measurement. This is created to make it easier to
      extract the individual waveplate positions later on when
      storing them in output file
   }
   \sstinvocation{
       CALL SURFLIB\_FILL\_WPLATE(USE\_WP, N\_WPLATE, SCUCD\_WPLATE, N\_POS, N\_EXP,
            N\_INT, N\_MEAS, DEM\_PNTR, WPLATE\_OUT, WPLATE\_ANG, STATUS)
   }
   \sstarguments{
      \sstsubsection{
         N\_WPLATE = INTEGER (Given)
      }{
         Size of SCUCD\_WPLATE array. If the value is zero then
         SCUCD\_WPLATE will not be used.
      }
      \sstsubsection{
         SCUCD\_WPLATE(N\_WPLATE) = REAL (Given)
      }{
         Waveplate positions for each measurement (degrees)
      }
      \sstsubsection{
         N\_POS = INTEGER (Given)
      }{
         Number of samples in WPLATE\_OUT
      }
      \sstsubsection{
         N\_EXP = INTEGER (Given)
      }{
         Number of exposures in DEM\_PNTR array
      }
      \sstsubsection{
         N\_INT = INTEGER (Given)
      }{
         Number of integrations in DEM\_PNTR array
      }
      \sstsubsection{
         N\_MEAS= INTEGER (Given)
      }{
         Number of measurements in DEM\_PNTR array
      }
      \sstsubsection{
         DEM\_PNTR ( 1, N\_EXP, N\_INT, N\_MEAS ) = INTEGER (Given)
      }{
         Start position for each EXP, INT, MEAS
      }
      \sstsubsection{
         WPLATE\_OUT ( N\_POS ) = REAL (Returned)
      }{
         Waveplate position for each sample (degrees)
      }
      \sstsubsection{
         WPLATE\_ANG( N\_INT, N\_MEAS ) = REAL (Given)
      }{
         Waveplate position for each integration/measurement
      }
      \sstsubsection{
         STATUS = INTEGER (Given \& Returned)
      }{
         Global status
      }
   }
   \sstdiytopic{
      Copyright
   }{
      Copyright \copyright 1995,1996,1997,1998,1999 Particle Physics and Astronomy
      Research Council. All Rights Reserved.
   }
}
\sstroutine{
   SURFLIB\_HISTOGRAM\_GRID
}{
   Calculate a 2-D histogram of I,J coordinates on a grid
}{
   \sstdescription{
      Given an array of I,J coordinates calculate the histogram
      of all points onto this grid of size NX $*$ NY. Just tells you
      how many data points there are per cell. Returns the histogram
      and the highest value in the histogram
   }
   \sstinvocation{
       CALL SURFLIB\_HISTOGRAM\_GRID (N\_PTS, NX, NY, USEDATA,
            IN\_DATA, IN\_QUALITY, BADBIT, IJ, GRID, IMAX, JMAX, NMAX,
            STATUS )
   }
   \sstarguments{
      \sstsubsection{
         N\_PTS = INTEGER (Given)
      }{
         Number of I,J pairs supplied
      }
      \sstsubsection{
         NX = INTEGER (Given)
      }{
         Size of X dimension
      }
      \sstsubsection{
         NY = INTEGER (Given)
      }{
         Size of Y dimension
      }
      \sstsubsection{
         USEDATA = LOGICAL (Given)
      }{
         If TRUE the histogram takes bad data points into account by not
         including them in the histogram. If false the straight histogram
         of the I,J{\tt '}s is constructed.
      }
      \sstsubsection{
         IN\_DATA(N\_PTS) = REAL (Given)
      }{
         Input data
      }
      \sstsubsection{
         IN\_QUALITY(N\_PTS) = BYTE (Given)
      }{
         Input quality
      }
      \sstsubsection{
         BADBIT = BYTE (Given)
      }{
         Bad bits mask for quality array
      }
      \sstsubsection{
         IJ ( 2, N\_PTS ) = INTEGER (Given)
      }{
         Positions of each point (I,J) in output grid
      }
      \sstsubsection{
         GRID (NX, NY) = INTEGER (Given \& Returned)
      }{
         Histogram. Note that array is not cleared by this routine leaving
         the possibility that the routine can be called multiple times
         to include more than one data set.
      }
      \sstsubsection{
         IMAX = INTEGER (Returned)
      }{
         I position of maximum [first one encountered]
      }
      \sstsubsection{
         JMAX = INTEGER (Returned)
      }{
         J position of maximum [first one encountered]
      }
      \sstsubsection{
         NMAX = INTEGER (Returned)
      }{
         Maximum value in histogram
      }
      \sstsubsection{
         STATUS = INTEGER (Given \& Returned)
      }{
         Global Status
      }
   }
   \sstnotes{
      \sstitemlist{

         \sstitem
         If USEDATA is TRUE the data and quality array are used so that bad
           pixels can be identified and not included into the histogram of
           positions.

         \sstitem
         GRID is not initialised by this routine

         \sstitem
         A warning error is raised if the bounds of the histogram
           are exceeded. It is up to the caller to decide whether this
           is a fatal error.
      }
   }
   \sstdiytopic{
      Copyright
   }{
      Copyright \copyright 1995,1996,1997,1998,1999 Particle Physics and Astronomy
      Research Council. All Rights Reserved.
   }
}
\sstroutine{
   SURFLIB\_MEDIAN\_REGRID
}{
   Generate image with each pixel the median of all pixels in bin
}{
   \sstdescription{
      This is done in two stages:

      1) Find the size of the output grid from the maximum extent of
         the input data.
      2) Loop through data. Find I,J coordinate of each point in the
         output grid.
      3) Find out maximum number of points for an I,J position.
      4) Put data onto grid in array (I,J,N) [REALS].
         We also need to store positions of these data.
         We can either do it by storing the file number, bolometer and
         position (time) index OR we can just store some index in a merged
         data array that goes from 1..TOT\_PTS.
          First method is easy but memory hungry. Second method is
          more efficient but does need some reconstruction to work
          out where the point was in the original data.
        Use the second method.

      Once the data is gridded, it is first displayed and then
      despiked. Currently despiking is done on a simple sigma clipping
      basis for each bin.
   }
   \sstinvocation{
       CALL SURFLIB\_MEDIAN\_REGRID( N\_FILES, N\_PTS, DIAMETER,
            WAVELENGTH, OUT\_PIXEL, NX, NY, ICEN, JCEN, BOL\_RA\_PTR,
            BOL\_DEC\_PTR, DATA\_PTR,  OUT\_DATA, OUT\_VAR, OUT\_QUAL,
            STATUS )
   }
   \sstarguments{
      \sstsubsection{
         N\_FILES = INTEGER (Given)
      }{
         Number of data sets (ie files)
      }
      \sstsubsection{
         N\_PTS ( N\_FILES ) = INTEGER (Given)
      }{
         Total number of points in each map
      }
      \sstsubsection{
         N\_POS( N\_FILES ) = INTEGER (Given)
      }{
         Number of positions per set (Y positions)
      }
      \sstsubsection{
         N\_BOLS( N\_FILES ) = INTEGER (Given)
      }{
         Number of bolometers per set (X positions)
      }
      \sstsubsection{
         BOL\_RA\_PTR( N\_FILES ) = INTEGER (Given)
      }{
         Array of pointers to position information (X coords)
         Note that each data set has positions for N\_POS $*$ N\_BOLS
      }
      \sstsubsection{
         BOL\_RA\_PTR( N\_FILES ) = INTEGER (Given)
      }{
         Array of pointers to position information (Y coords)
      }
      \sstsubsection{
         DATA\_PTR( N\_FILES ) = INTEGER (Given)
      }{
         Pointers to actual data arrays
      }
      \sstsubsection{
         QUALITY\_PTR( N\_FILES ) = INTEGER (Given)
      }{
         Pointer to quality arrays
      }
      \sstsubsection{
         NX = INTEGER (Returned)
      }{
         Number of points in grid (X)
      }
      \sstsubsection{
         NY = INTEGER (Returned)
      }{
         Number of points in grid (Y)
      }
      \sstsubsection{
         ICEN = INTEGER (Returned)
      }{
         Reference pixel (X)
      }
      \sstsubsection{
         JCEN = INTEGER (Returned)
      }{
         Reference pixel (Y)
      }
      \sstsubsection{
         NSPIKES ( N\_FILES ) = INTEGER (Returned)
      }{
         Number of spikes detected (and removed) in each file
      }
      \sstsubsection{
         BADBIT ( N\_FILES ) = BYTE (Given)
      }{
         Bad bit mask for identifying bad pixels from quality
      }
      \sstsubsection{
         STATUS = INTEGER (Given \& Returned)
      }{
         Global Status
      }
   }
   \sstnotes{
      For SMODE=NONE, DMODE is only requested if a plot is required.
   }
   \sstdiytopic{
      Copyright
   }{
      Copyright \copyright 1995,1996,1997,1998,1999 Particle Physics and Astronomy
      Research Council. All Rights Reserved.
   }
}
\sstroutine{
   SURFLIB\_PLOT\_GRID
}{
   Plot unwrapped grid
}{
   \sstdescription{
      Given the binned data calculated by SURFLIB\_FILL\_GRID
      Plot the unwrapped data as {\tt '}data{\tt '} against {\tt '}bin{\tt '}
      Use this as a first attempt at spike detection
   }
   \sstinvocation{
       CALL SURFLIB\_PLOT\_GRID(UNIT, NX, NY, NMAX, NSIGMA, IPOS, JPOS
           BINS, STATS, PNTS, POS, STATUS)
   }
   \sstarguments{
      \sstsubsection{
         UNIT = INTEGER (Given)
      }{
         Display number
      }
      \sstsubsection{
         NX = INTEGER (Given)
      }{
         Size of X dimension
      }
      \sstsubsection{
         NY = INTEGER (Given)
      }{
         Size of Y dimension
      }
      \sstsubsection{
         NMAX = INTEGER (Given)
      }{
         Maximum value allowed for third dimension of BINS
      }
      \sstsubsection{
         NSIGMA = REAL (Given)
      }{
         How many sigma away we plot the guide lines. If NSIGMA is negative
         we dont plot the sigma ranges on the plot.
      }
      \sstsubsection{
         IPOS(NX $*$ NY) = INTEGER (Given)
      }{
         I coordinate for each pixel
      }
      \sstsubsection{
         JPOS(NX $*$ NY) = INTEGER (Given)
      }{
         J coordinate for each pixe
      }
      \sstsubsection{
         STATS(NX, NY, 3) = REAL (Given)
      }{
         Statistics for each bin. 1=Median, 2=high, 3=low
      }
      \sstsubsection{
         BINS(NX, NY, NMAX) = REAL (Given)
      }{
         The data stored in relation to its position
      }
      \sstsubsection{
         PNTS(NMAX) = REAL (Given)
      }{
         Scratch space for copying in the data from each I,J
      }
      \sstsubsection{
         POS(NMAX) = REAL (Given)
      }{
         Scratch space for storing the X positions for each marker on the plot
      }
      \sstsubsection{
         STATUS = INTEGER (Given \& Returned)
      }{
         Global Status
      }
   }
   \sstdiytopic{
      Copyright
   }{
      Copyright \copyright 1995,1996,1997,1998,1999 Particle Physics and Astronomy
      Research Council. All Rights Reserved.
   }
}
\sstroutine{
   SURFLIB\_PROCESS\_BOLS
}{
   Calculate apparent RA/Dec of bolometers and some extra processing
}{
   \sstdescription{
      This routine calculates the apparent RA and Dec of each bolometer
      for all measurements, integrations and exposures. Both JIGGLE and
      SCAN data are supported.
      In addition, if this is called from EXTINCTION (ie EXTINCTION=.TRUE.)
      the data is corrected for extinction using the calculated elevation
      of the bolometers. If EXTINCTION is .FALSE. all the bolometers positions
      are converted to apparent RA-Dec at the reference modified Julian
      date (given by MJD\_STANDARD).

      There is some complication for moving sources (use PL output coords):

       For JIGGLE:
      \sstitemlist{

         \sstitem
           [1] Calculate map RA/Dec at new MJD

         \sstitem
           [2] Convert to tangent plane offsets for this map centre

      }
       For SCAN:
      \sstitemlist{

         \sstitem
           [1] Calculate array centre for reference MJD

         \sstitem
           [2] Calculate tangent offsets of array centre relative to map centre

         \sstitem
           [3] Calculate map centre at the interpolated MJD

         \sstitem
           [4] Work out array centre given new map centre

      }
       This is all because the file stores RA start and end relative
       to the fixed centre at the start of the observation.
   }
   \sstinvocation{
       CALL SURFLIB\_PROCESS\_BOLS(TSKNAME, N\_BEAMS,
            N\_BOL, N\_POS, N\_POS\_BEAMS, N\_SWITCHES, N\_EXPOSURES, N\_INTEGRATIONS,
            N\_MEASUREMENTS,START\_EXP, END\_EXP, START\_INT, END\_INT,
            START\_MEAS, END\_MEAS,N\_MAP, N\_FITS, FITS, DEM\_PNTR, LST\_STRT,
            IN\_ROTATION, SAMPLE\_MODE, SAMPLE\_COORDS, OUT\_COORDS,
            JIGGLE\_REPEAT, JIGGLE\_COUNT, JIGGLE\_X, JIGGLE\_Y,
            JIGGLE\_P\_SWITCH, RA\_CEN, DEC\_CEN,
            RA1, RA2, DEC1, DEC2, MJD\_STANDARD, IN\_UT1,
            MJD1, LONG1, LAT1, MJD2, LONG2, LAT2,
            LOCAL\_COORDS, MAP\_X, MAP\_Y,
            N\_POINT, POINT\_LST, POINT\_DAZ, POINT\_DEL,
            NUM\_CHAN, NUM\_ADC, BOL\_ADC, BOL\_CHAN, BOL\_DU3, BOL\_DU4,
            SCAN\_REVERSAL, FIRST\_LST, SECOND\_LST, FIRST\_TAU, SECOND\_TAU,
            BOL\_DEC, BOL\_RA, NDATA, NVARIANCE, USE\_LST, LST\_DATA,
            WAVEPLATE\_ANG, BOL\_IP\_DATA,
            STATUS)
   }
   \sstarguments{
      \sstsubsection{
         TSKNAME    = \_CHAR (Given)
      }{
         Taskname that is used to determine whether certain parts
         of the code are executed. Special values are:
         SCUOVER - when using SCUOVER this prevents the EXP\_START
                   loop from running properly. Only a single set
                   of bolometer positions are returned.
         EXTINCTION - corrects the input data for extinction
         REMIP   - Corrects the input data for instrumental polarisation
         All other values for this parameter assume that the bolometer
         positions are to be calculated and returned (converted to
         a standard MJD)
      }
      \sstsubsection{
         N\_BEAMS = \_INTEGER (Given)
      }{
         Number of beams in the input data. Only used if EXTINCTION=.TRUE.
      }
      \sstsubsection{
         N\_BOL = \_INTEGER (Given)
      }{
         Number of bolometers in the input data
      }
      \sstsubsection{
         N\_POS = \_INTEGER (Given)
      }{
         Number of `samples{\tt '} taken
      }
      \sstsubsection{
         N\_POS\_BEAMS = INTEGER (Given \& Returned)
      }{
         Number of beam positions to be returned per point. Can only be
         reduced. If equals 3 then order is M,L,R
      }
      \sstsubsection{
         N\_SWITCHES = \_INTEGER (Given)
      }{
         Number of switches
      }
      \sstsubsection{
         N\_EXPOSURES = \_INTEGER (Given)
      }{
         Number of exposures
      }
      \sstsubsection{
         N\_INTEGRATIONS = \_INTEGER (Given)
      }{
         Number of integrations
      }
      \sstsubsection{
         N\_MEASUREMENTS = \_INTEGER (Given)
      }{
         Number of measurements
      }
      \sstsubsection{
         START\_EXP  = INTEGER (Given)
      }{
         First exposure to process
      }
      \sstsubsection{
         END\_EXP    = INTEGER (Given)
      }{
         Last exposure to process
      }
      \sstsubsection{
         START\_INT  = INTEGER (Given)
      }{
         First integration to process
      }
      \sstsubsection{
         END\_INT    = INTEGER (Given)
      }{
         End integration to process
      }
      \sstsubsection{
         START\_MEAS = INTEGER (Given)
      }{
         First measurement to process
      }
      \sstsubsection{
         END\_MEAS   = INTEGER (Given)
      }{
         End measurement to process
      }
      \sstsubsection{
         N\_MAP = \_INTEGER (Given)
      }{
         Map number of input data (not used for EXTINCTION or REMIP)
      }
      \sstsubsection{
         N\_FITS = \_INTEGER (Given)
      }{
         Number of FITS items
      }
      \sstsubsection{
         FITS() = \_CHAR$*$80 (Given)
      }{
         FITS array
      }
      \sstsubsection{
         DEM\_PNTR() = \_INTEGER (Given)
      }{
         DEM\_PNTR array - position in file of each exposure
      }
      \sstsubsection{
         LST\_STRT() = \_DOUBLE (Given)
      }{
         LST of each exposure
      }
      \sstsubsection{
         IN\_ROTATION = \_DOUBLE (Given)
      }{
         Angle between apparent N and N of input coord system (radians)
      }
      \sstsubsection{
         SAMPLE\_MODE = \_CHAR (Given)
      }{
         Sample mode of input file
      }
      \sstsubsection{
         SAMPLE\_COORDS = \_CHAR (Given)
      }{
         Coordinate system of sample offsets
      }
      \sstsubsection{
         OUT\_COORDS = \_CHAR (Given)
      }{
         Output coordinate system
      }
      \sstsubsection{
         JIGGLE\_REPEAT = \_INTEGER (Given)
      }{
         Number of times jiggle pattern is repeated in a switch
      }
      \sstsubsection{
         JIGGLE\_COUNT = \_INTEGER (Given)
      }{
         Number of jiggle in pattern
      }
      \sstsubsection{
         JIGGLE\_X(JIGGLE\_COUNT) = \_REAL (Given)
      }{
         X jiggle offsets (arcsec)
      }
      \sstsubsection{
         JIGGLE\_Y(JIGGLE\_COUNT) = \_REAL (Given)
      }{
         Y jiggle offsets (arcsec)
      }
      \sstsubsection{
         JIGGLE\_P\_SWITCH = \_INTEGER
      }{
         Number of jiggles per switch
      }
      \sstsubsection{
         RA\_CEN = \_DOUBLE (Given)
      }{
         apparent RA of output map centre (radians). Used mainly for JIGGLE
         but also for scan/map data pre Dec 1997
      }
      \sstsubsection{
         DEC\_CEN = \_DOUBLE (Given)
      }{
         apparent Dec of output map centre (radians) Used mainly for JIGGLE
         but also for scan/map data pre Dec 1997
      }
      \sstsubsection{
         RA1 = \_REAL (Given)
      }{
         RA (in RD, RJ or AZ) at start of scan for each exposure (SCAN only)
      }
      \sstsubsection{
         RA2 = \_REAL (Given)
      }{
         RA (in RD, RJ or AZ) at end of scan for each exposure (SCAN only)
      }
      \sstsubsection{
         DEC1 = \_REAL (Given)
      }{
         DEC (in RJ, RD or AZ) at start of scan for each exposure (SCAN only)
      }
      \sstsubsection{
         DEC2 = \_REAL (Given)
      }{
         DEC (in RD, RJ, AZ) at end of scan for each exposure (SCAN only)
      }
      \sstsubsection{
         MJD\_STANDARD = \_DOUBLE (Given)
      }{
         Standard MJD to which each input map is referenced (EXTINCTION=FALSE)
      }
      \sstsubsection{
         IN\_UT1 = \_DOUBLE (Given)
      }{
         MJD of input data.
      }
      \sstsubsection{
         MJD1 = DOUBLE (Given)
      }{
         MJD of first planet position
      }
      \sstsubsection{
         LONG1 = DOUBLE (Given)
      }{
         Longitude (apparent RA) at MJD1
      }
      \sstsubsection{
         LAT1 = DOUBLE (Given)
      }{
         Latitude (apparent Dec) at MJD1
      }
      \sstsubsection{
         MJD2 = DOUBLE (Given)
      }{
         MJD of second planet position
      }
      \sstsubsection{
         LONG2 = DOUBLE (Given)
      }{
         Longitude (apparent RA) at MJD2
      }
      \sstsubsection{
         LAT2 = DOUBLE (Given)
      }{
         Latitude (apparent Dec) at MJD2
      }
      \sstsubsection{
         LOCAL\_COORDS = CHARACTER (Given)
      }{
         Coordinate system of offsets
      }
      \sstsubsection{
         MAP\_X = DOUBLE (Given)
      }{
         X offset in LOCAL\_COORDS (radians)
      }
      \sstsubsection{
         MAP\_Y = DOUBLE (Given)
      }{
         Y offset in LOCAL\_COORDS (radians)
      }
      \sstsubsection{
         N\_POINT = \_INTEGER (Given)
      }{
         Number of pointing corrections (should be zero if EXTINCTION)
      }
      \sstsubsection{
         POINT\_DEL = \_REAL (Given)
      }{
         Elevation pointing corrections (radians) [only if EXTINCTION=FALSE]
      }
      \sstsubsection{
         POINT\_DAZ = \_REAL (Given)
      }{
         Azimuth pointing corrections (radians) [only if EXTINCTION=FALSE]
      }
      \sstsubsection{
         POINT\_LST = \_DOUBLE (Given)
      }{
         LST of pointing corrections (radians) [only if EXTINCTION=FALSE]
      }
      \sstsubsection{
         NUM\_CHAN = \_INTEGER (Given)
      }{
         Number of channels in DAQ
      }
      \sstsubsection{
         NUM\_ADC = \_INTEGER (Given)
      }{
         Number of AtoD cards.
      }
      \sstsubsection{
         BOL\_ADC = \_INTEGER (Given)
      }{
         A/D numbers of bolometers measured in input file
      }
      \sstsubsection{
         BOL\_CHAN = \_INTEGER (Given)
      }{
         channel numbers of bolometers measured in input file
      }
      \sstsubsection{
         BOL\_DU3 = \_REAL (Given)
      }{
         dU3 Nasmyth coordinates of bolometers
      }
      \sstsubsection{
         BOL\_DU4 = \_REAL (Given)
      }{
         dU4 Nasmyth coordinates of bolometers
      }
      \sstsubsection{
         SCAN\_REVERSAL = LOGICAL (Given)
      }{
         Multiply alternate exposures by -1 if SCANning
      }
      \sstsubsection{
         FIRST\_LST = \_DOUBLE (Given)
      }{
         LST of first tau value (EXTINCTION only)
      }
      \sstsubsection{
         SECOND\_LST = \_DOUBLE (Given)
      }{
         LST of second tau value (EXTINCTION only)
      }
      \sstsubsection{
         FIRST\_TAU = \_REAL (Given)
      }{
         First tau value (EXTINCTION only)
      }
      \sstsubsection{
         SECOND\_TAU = \_REAL (Given)
      }{
         Second tau value (EXTINCTION only)
      }
      \sstsubsection{
         BOL\_DEC(N\_BOL, N\_POS, N\_POS\_BEAMS) = \_DOUBLE (Returned)
      }{
         Apparent DEC of bolometers for each measurement for MJD\_STANDARD
         Depending on N\_POS\_BEAMS can contain M, L and R beam positions.
      }
      \sstsubsection{
         BOL\_RA(N\_BOL, N\_POS, N\_POS\_BEAMS) = \_DOUBLE (Returned)
      }{
         Apparent RA of bolometers for each measurement for MJD\_STANDARD
         Depending on N\_POS\_BEAMS can contain M, L and R beam positions.
      }
      \sstsubsection{
         NDATA(N\_BOL, N\_POS, N\_BEAMS) = \_REAL (Given \& Returned)
      }{
         corrected data (EXTINCTION and REMIP only)
      }
      \sstsubsection{
         NVARIANCE(N\_BOL, N\_POS, N\_BEAMS) = \_REAL (Given \& Returned)
      }{
         Extinction corrected variance (EXTINCTION only)
      }
      \sstsubsection{
         USE\_LST = LOGICAL (Given)
      }{
         Do we return the LST positions? (TRUE  then we do)
      }
      \sstsubsection{
         LST\_DATA(N\_POS) = DOUBLE (Returned)
      }{
         LST for each position.
      }
      \sstsubsection{
         WAVEPLATE\_ANG = REAL (Given)
      }{
         waveplate position angle (Nasmyth degrees) (Given \& Returned)
         The supplied value is the nasmyth angle in degrees and the
         return value the sky rotation angle for each position.
         (ie parallactic angle - elevation) in degrees (REMIP only)
      }
      \sstsubsection{
         BOL\_IP\_DATA(8, NUM\_CHAN, NUM\_ADC) = REAL (Given)
      }{
         The IP data. This array contains the ip data for each
         bolometer (specified by a combination of CHAN and ADC).
         The 4 slices represent: P0, Pslope, Theta0 and ThetaSlope (REMIP only)
         The next 4 slices are for the related variance (same order)
      }
      \sstsubsection{
         STATUS = \_INTEGER (Given \& Returned)
      }{
         Global status
      }
   }
   \sstnotes{
      MAP\_X and MAP\_Y are only used for JIGGLE modes.
      The SCAN/MAP offsets are wrapped into the RA1,DEC1,RA2,DEC2
      numbers by the on-line system.
      The offsets are added to the map centre every time round the
      loop. This is because it is possible to have AZ offsets for
      RA,Dec centres.

      This routine tries to deal with  the different versions of
      SCUCD (only affects SCAN/MAP data).
      For version 0:

        RA/Decs of the scan were incorrectly stored as RJ. They
        are converted to RD before further processing.
        The coordinate frame of RA1, RA2, DEC1 and DEC2 fro SCAN/MAP
        depends on the CENTRE\_COORDS of the observation. For CENTRE=RD
        the scan positions are in RD; for CENTRE=RB,RJ,GA the scans
        are in RJ and for centre=AZ the scans are in AZ.
        (SCULIB\_SCAN\_2\_RD)

      For version 1.0:

        A bug was introduced concerning the calculation of the
        ends of the scans. The bug is recreated and inverted in order
        to compensate. (SCULIB\_FIX\_SCAN\_V10)

      Also, LO chopping for jiggle maps was broken until 19980730
      such that the initial chop pa was calculated in Az but never
      updated as the sky rotated. This fix is only important for SCUBA2MEM
      where the positions of the off-beams are returned.
   }
   \sstbugs{
      Currently the IN\_UT1 is assumed to be the MJD when the data taking
      begins. As of 24-NOV-1997 IN\_UT1 is actually the MJD of the start
      of the observation (ie when the telescope begins to slew). This is
      a problem for data using a moving centre.
   }
}
\sstroutine{
   SURFLIB\_READ\_IPFILE
}{
   Reads IP data from text file
}{
   \sstdescription{
      This routine reads IP data from a text file and stores
      it in a multi-dimensional array.
      The format of the file is:
      Any line starting with the string {\tt '}FAST{\tt '} is assumed
      to be specifying the angle between 0 degrees on the waveplate
      and the angle to the fast axis. This is wavelength dependent.
      The line should be of the format

      \texttt{FAST filter angle}

      eg   ``FAST 450 -3.0''

      The list of valid filter names is passed in to the routine.
      An error is raised if a fast axis angle has not been specified
      for each filter.

      For each bolometer the following information is required:


\texttt{BolName P0 P0\_Err Pslope\_Err Theta0 Theta0\_Err ThetaSlope ThetaSlope\_Err}

\begin{tabular}{ll}
BolName   & Name of bolometer (eg H7,G12..)\\
P0        & IP percentage polarization at the horizon\\
P0Err     & Error in P0\\
PSlope    & Gradient in P (ie P = P0 $+$ Pslope $*$ Elevation)\\
PSlopeErr & Error in PSlope\\
Theta0    & polarization angle of IP (degrees)\\
Theta0Err & Error in polarisation angle\\
ThetaSlope& Gradient in Theta (ie Theta = Theta0 $+$ ThetaSlope $*$ El)\\
ThSlErr   & Error in ThetaSlope\\
\end{tabular}

      For example:
      \texttt{H7  2.0 0.01 0.02 0.01 167 2 1.0 0.01}

      \{ is the comment character. Blank lines are ignored.
      The columns must be space separated.

      All angles are converted to radians before further processing
      (including Pslope which is in degrees$^{-1}$)
   }
   \sstinvocation{
       CALL SURFLIB\_READ\_IPFILE( FD, NUM\_CHAN, NUM\_ADC,
            N\_FILT, FILTERS, FAST\_ANG, BOL\_IP\_DATA, STATUS)
   }
   \sstarguments{
      \sstsubsection{
         FD = FILE (Given)
      }{
         Identifier to the open file (opened for read access)
      }
      \sstsubsection{
         NUM\_CHAN = INTEGER (Given)
      }{
         Number of available bolometer channels
      }
      \sstsubsection{
         NUM\_ADC  = INTEGER (Given)
      }{
         Number of available AtoD cards
      }
      \sstsubsection{
         N\_FILT   = INTEGER (Given)
      }{
         Number of filters requested
      }
      \sstsubsection{
         FILTERS( N\_FILT ) = CHARACTER (Given)
      }{
         Array of requested filter names
      }
      \sstsubsection{
         FAST\_ANG( N\_FILT ) = REAL (Returned)
      }{
         Fast axis angle for each filter. Will contain bad values
         if no angle was found
      }
      \sstsubsection{
         BOL\_IP\_DATA(8, NUM\_CHAN, NUM\_ADC) = REAL (Returned)
      }{
         The IP data. This array contains the ip data for each
         bolometer (specified by a combination of CHAN and ADC).
         The 4 slices represent: P0, Pslope, Theta0 and ThetaSlope
         The last 4 slices contain the corresponding variance
         All angles are converted to radians (Pslope has to be converted
         from /degree to /radian).
      }
   }
   \sstdiytopic{
      Copyright
   }{
      Copyright \copyright 1995,1996,1997,1998,1999 Particle Physics and Astronomy
      Research Council. All Rights Reserved.
   }
}
\sstroutine{
   SURFLIB\_REM\_GRID
}{
   Remove an image grid of data from demodulated data
}{
   \sstdescription{
      Removes a gridded image from demodulated data using a
      lookup table containing the (I,J) coordinates (in the grid)
      of each input data point.
      Can be used to remove the source from an input data set.
   }
   \sstinvocation{
       CALL SURFLIB\_REM\_GRID( N\_POS, N\_BOLS, NX, NY,
            IJ, GRID, IN\_DATA, STATUS)
   }
   \sstarguments{
      \sstsubsection{
         N\_PTS = INTEGER (Given)
      }{
         Number of data points in input data
      }
      \sstsubsection{
         NX = INTEGER (Given)
      }{
         X dimension of input grid
      }
      \sstsubsection{
         NY = INTEGER (Given)
      }{
         Y dimension of input grid
      }
      \sstsubsection{
         IJ (2, N\_PTS)  INTEGER (Given)
      }{
         Positions of each point (I,J) in input grid
      }
      \sstsubsection{
         GRID ( NX, NY) = REAL (Given)
      }{
         Image grid to be subtracted from data
      }
      \sstsubsection{
         IN\_DATA( N\_PTS ) = REAL (Given \& Returned)
      }{
         Input data to be modified
      }
      \sstsubsection{
         STATUS = INTEGER (Given \& Returned)
      }{
         Global Status
      }
   }
   \sstdiytopic{
      Copyright
   }{
      Copyright \copyright 1995,1996,1997,1998,1999 Particle Physics and Astronomy
      Research Council. All Rights Reserved.
   }
}
\sstroutine{
   SURFLIB\_REM\_TIMESERIES
}{
   Remove a time series  from a 2D array
}{
   \sstdescription{
      Remove a time series (Y data) from a 2D array where the
      second dimension corresponds to time. It is assumed that
      both the time series and the 2D array have the same
      number of points in time (N\_POS)
   }
   \sstinvocation{
       CALL SURFLIB\_REM\_TIMESERIES ( N\_BOLS, N\_POS,
            TIME\_SER, TIME\_VAR, IN\_DATA, IN\_VAR, STATUS)
   }
   \sstarguments{
      \sstsubsection{
         N\_BOLS = INTEGER (Given)
      }{
         Number of bolometers in data array (x dimension)
      }
      \sstsubsection{
         N\_POS = INTEGER (Given)
      }{
         Number of time data values (y dim)
      }
      \sstsubsection{
         TIME\_SER ( N\_POS ) = REAL (Given)
      }{
         Time series to be removed
      }
      \sstsubsection{
         TIME\_VAR ( N\_POS ) = REAL (Given)
      }{
         Variance associated with each time value
      }
      \sstsubsection{
         IN\_DATA (N\_BOLS, N\_POS ) = REAL (Given \& Returned)
      }{
         Input data to be modified
      }
      \sstsubsection{
         IN\_VAR (N\_BOLS, N\_POS ) = REAL (Given \& Returned)
      }{
         Input variance
      }
      \sstsubsection{
         STATUS = INTEGER (Given \& Returned)
      }{
         Global Status
      }
   }
   \sstdiytopic{
      Copyright
   }{
      Copyright \copyright 1995,1996,1997,1998,1999 Particle Physics and Astronomy
      Research Council. All Rights Reserved.
   }
}
\sstroutine{
   SURFLIB\_REMOVE\_DC\_FROM\_EXP
}{
   Remove linear baseline from each exposure
}{
   \sstdescription{
      This routine takes a data array. It then removes a
      DC offset (either the mean or the median derived from the scan)
      from each scan.
   }
   \sstinvocation{
       CALL SCULIB\_REMOVE\_DC\_FROM\_EXP(DORLB, N\_EXPOSURES,
            N\_INTEGRATIONS, N\_MEASUREMENTS, METHOD, DEM\_PNTR, N\_BOL, N\_POS,
            IN\_DATA, IN\_QUALITY, SAMPLE\_DX, CHOP\_THROW, OUT\_DATA,
            OUT\_QUALITY, BADBIT, STATUS)
   }
   \sstarguments{
      \sstsubsection{
         DORLB                       = LOGICAL (Given)
      }{
         control whether we are subtracting the baseline (TRUE)
         or storing the basline (FALSE)
      }
      \sstsubsection{
         N\_EXPOSURES                 = INTEGER (Given)
      }{
         maximum number of exposures per integration
      }
      \sstsubsection{
         N\_INTEGRATIONS              = INTEGER (Given)
      }{
         number of integrations in the observation
      }
      \sstsubsection{
         N\_MEASUREMENTS              = INTEGER (Given)
      }{
         number of measurements in the observation
      }
      \sstsubsection{
         METHOD                      = CHAR (Given)
      }{
         Removal method. MEDIAN or MEAN supported.
      }
      \sstsubsection{
         DEMOD\_POINTER (N\_EXPOSURES, N\_INTEGRATIONS, N\_MEASUREMENTS)
      }{
                               = INTEGER (Given)
         array pointing to start and finish of scans in IN\_DATA
      }
      \sstsubsection{
         N\_BOL                       = INTEGER (Given)
      }{
         the number of bolometers for which data was taken
      }
      \sstsubsection{
         N\_POS                       = INTEGER (Given)
      }{
         the number of positions measured in the scan
      }
      \sstsubsection{
         IN\_DATA (N\_BOL, N\_POS)      = REAL (Given)
      }{
         the measured data
      }
      \sstsubsection{
         IN\_VARIANCE (N\_BOL, N\_POS)      = REAL (Given)
      }{
         the measured variance
      }
      \sstsubsection{
         IN\_QUALITY (N\_BOL, N\_POS)   = BYTE (Given)
      }{
         the quality on IN\_DATA
      }
      \sstsubsection{
         OUT\_DATA (N\_BOL, N\_POS)     = REAL (Returned)
      }{
         the data with baseline removed
      }
      \sstsubsection{
         OUT\_VARIANCE (N\_BOL, N\_POS)      = REAL (Given)
      }{
         the output variance
      }
      \sstsubsection{
         OUT\_QUALITY (N\_BOL, N\_POS)  = BYTE (Returned)
      }{
         the quality on OUT\_DATA
      }
      \sstsubsection{
         BADBIT                      = BYTE (Given)
      }{
         bad bit mask
      }
      \sstsubsection{
         STATUS = INTEGER (Given and Returned)
      }{
         Global Status value
      }
   }
   \sstdiytopic{
      Copyright
   }{
      Copyright \copyright 1995,1996,1997,1998,1999 Particle Physics and Astronomy
      Research Council. All Rights Reserved.
   }
}
\sstroutine{
   SURFLIB\_REMOVE\_DC\_VIA\_SECT
}{
   Remove median baseline from each exposure
}{
   \sstdescription{
      This routine takes a data array and removes a
      DC offset from each scan. The DC level is the median of the
      data specified in the SCUBA section for that particular integration.
      The same level is removed for each scan of an integration
      with a different value for each bolometer.
   }
   \sstinvocation{
       CALL SCULIB\_REMOVE\_DC\_FROM\_EXP(DORLB,N\_SPEC,SECTION,USE\_SECT,N\_EXPOSURES,
            N\_INTEGRATIONS, N\_MEASUREMENTS, METHOD, DEM\_PNTR, N\_BOL, N\_POS,
            IN\_DATA, IN\_QUALITY, SAMPLE\_DX, CHOP\_THROW, OUT\_DATA,
            OUT\_QUALITY, BADBIT, STATUS)
   }
   \sstarguments{
      \sstsubsection{
         DORLB                       = LOGICAL (Given)
      }{
         control whether we are subtracting the baseline (TRUE)
         or storing the basline (FALSE)
      }
      \sstsubsection{
         N\_SPEC                      = INTEGER (Given)
      }{
         Number of sections specified
      }
      \sstsubsection{
         SECTION ( N\_SPEC )          = CHAR (Given)
      }{
         Array of section specifications
      }
      \sstsubsection{
         USE\_SECT                    = LOGICAL (Given)
      }{
         Is this an inverse section
      }
      \sstsubsection{
         N\_EXPOSURES                 = INTEGER (Given)
      }{
         maximum number of exposures per integration
      }
      \sstsubsection{
         N\_INTEGRATIONS              = INTEGER (Given)
      }{
         number of integrations in the observation
      }
      \sstsubsection{
         N\_MEASUREMENTS              = INTEGER (Given)
      }{
         number of measurements in the observation
      }
      \sstsubsection{
         DEMOD\_POINTER (N\_EXPOSURES, N\_INTEGRATIONS, N\_MEASUREMENTS)
      }{
                               = INTEGER (Given)
         array pointing to start and finish of scans in IN\_DATA
      }
      \sstsubsection{
         N\_BOL                       = INTEGER (Given)
      }{
         the number of bolometers for which data was taken
      }
      \sstsubsection{
         N\_POS                       = INTEGER (Given)
      }{
         the number of positions measured in the scan
      }
      \sstsubsection{
         IN\_DATA (N\_BOL, N\_POS)      = REAL (Given)
      }{
         the measured data
      }
      \sstsubsection{
         IN\_VARIANCE (N\_BOL, N\_POS)      = REAL (Given)
      }{
         the measured variance
      }
      \sstsubsection{
         IN\_QUALITY (N\_BOL, N\_POS)   = BYTE (Given)
      }{
         the quality on IN\_DATA
      }
      \sstsubsection{
         OUT\_DATA (N\_BOL, N\_POS)     = REAL (Returned)
      }{
         the data with baseline removed
      }
      \sstsubsection{
         OUT\_VARIANCE (N\_BOL, N\_POS)      = REAL (Given)
      }{
         the output variance
      }
      \sstsubsection{
         OUT\_QUALITY (N\_BOL, N\_POS)  = BYTE (Returned)
      }{
         the quality on OUT\_DATA
      }
      \sstsubsection{
         BADBIT                      = BYTE (Given)
      }{
         bad bit mask
      }
      \sstsubsection{
         STATUS = INTEGER (Given and Returned)
      }{
         Global Status value
      }
   }
   \sstdiytopic{
      Copyright
   }{
      Copyright \copyright 1995,1996,1997,1998,1999 Particle Physics and Astronomy
      Research Council. All Rights Reserved.
   }
}
\sstroutine{
   SURFLIB\_REMOVE\_IP
}{
   Remove instrumental polarisation
}{
   \sstdescription{
      Remove the instrumental polarisation signal from the data
      The following formula is used:

      New flux = Measured flux / (1 $+$ \%age IP )

      IP = \%age polarisation $*$ cos (4$*$waveplate - 2$*$ IP angle)

      where \%age polarisation and IP angle vary linearly with
      elevation.

      This is an approximation of

        Actual flux = measured flux - mean flux $*$ \%age IP

      where IP = P $*$ (1 $+$ cos (4 WP - 2 THETA ) )
      because the mean flux can not be calculated trivially since
      the bolometers are jiggling on and off the source

      This routine corrects the data point for each bolometer
      by looking up the IP data in the supplied array. Only one
      set of bolometers can be processed at any one time (since the
      elevation of the array changes during the observation). The
      routine should be called repeatedly for each time/elevation.
   }
   \sstinvocation{
       CALL SURFLIB\_REMOVE\_IP( ELEVATION, NUM\_CHAN, NUM\_ADC,
            N\_BOLS, WPLATE\_ANG, BOL\_CHAN, BOL\_ADC,
            BOL\_IP\_DATA, BOL\_DATA, STATUS)
   }
   \sstarguments{
      \sstsubsection{
         ELEVATION = DOUBLE (Given)
      }{
         Elevation of the source (radians)
      }
      \sstsubsection{
         NUM\_CHAN = INTEGER (Given)
      }{
         Number of channels per A/D card
      }
      \sstsubsection{
         NUM\_ADC = INTEGER (Given)
      }{
         the number of A/D cards
      }
      \sstsubsection{
         N\_BOLS = INTEGER (Given)
      }{
         the actual number of bolometers
      }
      \sstsubsection{
         WPLATE\_ANG = REAL (Given)
      }{
         nasmyth angle of waveplate for these data (degrees)
      }
      \sstsubsection{
         BOL\_CHAN (N\_BOLS) = INTEGER (Given)
      }{
         channel numbers of bolometers
      }
      \sstsubsection{
         BOL\_ADC (N\_BOLS) = INTEGER (Given)
      }{
         ADC numbers of bolometers
      }
      \sstsubsection{
         BOL\_IP\_DATA( 8, NUM\_CHAN, NUM\_ADC ) = REAL (Given)
      }{
         The IP data. This array contains the ip data for each
         bolometer (specified by a combination of CHAN and ADC).
         The 4 slices represent: P0, Pslope, Theta0 and ThetaSlope
         The last 4 slices represent variance of each data point
         All angles are in radians.
      }
      \sstsubsection{
         BOL\_DATA (N\_BOLS) = REAL (Given and returned)
      }{
         bolometer data
      }
      \sstsubsection{
         BOL\_VAR ( N\_BOLS ) = REAL (Given and returned)
      }{
         variance of bolometer data
      }
      \sstsubsection{
         STATUS = INTEGER (Given and returned)
      }{
         global status
      }
   }
   \sstnotes{
      \sstitemlist{

         \sstitem
         The elevation of each bolometer is assumed to be the elevation
           of the central bolometer.

         \sstitem
         Output variance is calculated

         \sstitem
         degrees are used throughout since this simplifies the calculations
           when the gradients and ThetaZero are specified in degrees
           (or inverse degrees for the slope) from the input IP file.
      }
   }
   \sstdiytopic{
      Copyright
   }{
      Copyright \copyright 1995,1996,1997,1998,1999 Particle Physics and Astronomy
      Research Council. All Rights Reserved.
   }
}
\sstroutine{
   SURFLIB\_STATS\_GRID
}{
   Calculate statistics of binned data
}{
   \sstdescription{
      Calculate the statistics of the binned data and return the
      median and upper and lower limits of acceptable data.
   }
   \sstinvocation{
       CALL SURFLIB\_STATS\_GRID(SMODE, NX, NY, NMAX, NSIGMA,
            IPOS, JPOS, BINS, PNTS, STATS, STATUS)
   }
   \sstarguments{
      \sstsubsection{
         SMODE = CHARACTER (Given)
      }{
         Smoothing mode. The statistics can be smoothed by adjacent pixels
         in order to remove spikes from the clipping envelope.
         Options are:
            NONE - No smoothing (ie the statistics of each bin)
            HANN - Hanning smoothing. Triangular envolope across 3 points.
         Note that smoothing depends on the way the grid was unwrapped.
         since that determines which pixels are adjacent (except for
         SMODE=NONE.
      }
      \sstsubsection{
         NX = INTEGER (Given)
      }{
         Size of X dimension
      }
      \sstsubsection{
         NY = INTEGER (Given)
      }{
         Size of Y dimension
      }
      \sstsubsection{
         NMAX = INTEGER (Given)
      }{
         Maximum value allowed for third dimension of BINS
      }
      \sstsubsection{
         NSIGMA = REAL (Given)
      }{
         Standard deviation limit.
      }
      \sstsubsection{
         IPOS(NX $*$ NY) = INTEGER (Given)
      }{
         I coordinate for each pixel
      }
      \sstsubsection{
         JPOS(NX $*$ NY) = INTEGER (Given)
      }{
         J coordinate for each pixel
      }
      \sstsubsection{
         BINS(NX, NY, NMAX) = REAL (Given)
      }{
         The data stored in relation to its position
      }
      \sstsubsection{
         PNTS(NMAX) = REAL (Given)
      }{
         Scratch space for copying in the data from each I,J
      }
      \sstsubsection{
         STATS(NX, NY, 3) = REAL (Returned)
      }{
         Statistics of each bin. 3 components are: Median, Upper limit
         lower limit.
      }
      \sstsubsection{
         STATUS = INTEGER (Given \& Returned)
      }{
         Global status.
      }
   }
   \sstdiytopic{
      Copyright
   }{
      Copyright \copyright 1995,1996,1997,1998,1999 Particle Physics and Astronomy
      Research Council. All Rights Reserved.
   }
}
\sstroutine{
   SURFLIB\_TRIM\_IMAGE
}{
   Sets quality bit for specified distance from exisiting bad data
}{
   \sstdescription{
      Will take the quality array of a 2-D image and set the
      quality bit of all pixels within the specified trimming
      distance of a pixel that has bit 0 set. The output pixel
      bit is only set explicitly if bit 0 is not already set.
      This is so that the effect of the trim can be removed
      simply by changing the bad bits mask.
   }
   \sstinvocation{
      CALL SURFLIB\_TRIM\_IMAGE( TRIM, NX, NY, BITNUM,
      IN\_QUAL, OUT\_QUAL, STATUS)
   }
   \sstarguments{
      \sstsubsection{
         TRIM = REAL (Given)
      }{
         The distance in pixels inside which quality bits will be
         set.
      }
      \sstsubsection{
         NX = INTEGER (Given)
      }{
         Number of X pixels in image
      }
      \sstsubsection{
         NY = INTEGER (Given)
      }{
         Number of Y pixels in image
      }
      \sstsubsection{
         BITNUM = INTEGER (Given)
      }{
         Bit number to set to bad if within TRIM distance of a
         pixel that has bit 0 set. Start counting at 0.
      }
      \sstsubsection{
         IN\_QUAL(NX,NY) = UBYTE (Given)
      }{
         Input image.
      }
      \sstsubsection{
         OUT\_QUAL(NX,NY) = UBYTE (Returned)
      }{
         Output image quality.
      }
      \sstsubsection{
         STATUS = INTEGER (Given and Returned)
      }{
         Global status
      }
   }
   \sstdiytopic{
      Copyright
   }{
      Copyright \copyright 1995,1996,1997,1998,1999 Particle Physics and Astronomy
      Research Council. All Rights Reserved.
   }
}



\subsection{SCULIB}

This section describes the subroutines that are available in the
\texttt{sculib} distribution (\texttt{libsculib.a}).

\sstroutine{
   SCULIB\_1D2\_JIGGLE
}{
   routine to unpack a 1-d jiggle dataset into a 2-d
   image
}{
   \sstdescription{
      This routine unpacks the data for a specified bolometer from a
      dataset stored according to jiggle number into a 2-d dataset
      with data stored according to jiggle offset.
        The input dataset can contain data either from all or part of
      ONE run through the jiggle pattern or for several REPEATS of
      the entire pattern. In the first case, the routine copies the
      input data into the appropriate part of the 2-d image. In the
      second, the 2-d datum for a given pixel will be the average of the
      input values for it, and the variance will calculated from the
      spread of input points around the mean unless only one input point
      contributed, in which case the input variance will be used.
   }
   \sstinvocation{
       CALL SCULIB\_1D2\_JIGGLE (BOL, N\_BOLS, J\_START, N\_JIG,
         J\_REPEAT, J\_COUNT, J\_DATA, J\_VARIANCE, J\_QUALITY, IDIM,
         JDIM, I\_JIGGLE, J\_JIGGLE, DATA\_2D, VARIANCE\_2D, QUALITY\_2D,
         STATUS)
   }
   \sstarguments{
      \sstsubsection{
         BOL                          = INTEGER (Given)
      }{
         the number of the bolometer whose data is to be unpacked
      }
      \sstsubsection{
         N\_BOLS                       = INTEGER (Given)
      }{
         the number of bolometers measured
      }
      \sstsubsection{
         J\_START                      = INTEGER (Given)
      }{
         the index within the pattern of the first jiggle position
         measured
      }
      \sstsubsection{
         N\_JIG                        = INTEGER (Given)
      }{
         the total number of jiggle positions measured
      }
      \sstsubsection{
         J\_REPEAT                     = INTEGER (Given)
      }{
         the number of times the jiggle pattern was repeated in the
         dataset
      }
      \sstsubsection{
         J\_COUNT                      = INTEGER (Given)
      }{
         the number of jiggles in the pattern
      }
      \sstsubsection{
         J\_DATA (N\_BOLS, N\_JIG)
      }{
                                = REAL (Given)
         the measured data at each jiggle position
      }
      \sstsubsection{
         J\_VARIANCE (N\_BOLS, N\_JIG)
      }{
                                = REAL (Given)
         the measured variance
      }
      \sstsubsection{
         J\_QUALITY (N\_BOLS, N\_JIG)
      }{
                                = INTEGER (Given)
         the quality on the measured data
      }
      \sstsubsection{
         IDIM                         = INTEGER (Given)
      }{
         the x dimension of the 2-d map
      }
      \sstsubsection{
         JDIM                         = INTEGER (Given)
      }{
         the y dimension of the 2-d map
      }
      \sstsubsection{
         I\_JIGGLE (J\_COUNT)           = INTEGER (Given)
      }{
         the `i{\tt '} index on the 2-d map of each jiggle position
      }
      \sstsubsection{
         J\_JIGGLE (J\_COUNT)           = INTEGER (Given)
      }{
         the `j{\tt '} index on the 2-d map of each jiggle position
      }
      \sstsubsection{
         DATA\_2D (IDIM, JDIM)         = REAL (Returned)
      }{
         the data plane of the 2-d map
      }
      \sstsubsection{
         VARIANCE\_2D (IDIM, JDIM)     = REAL (Returned)
      }{
         the variance plane
      }
      \sstsubsection{
         QUALITY\_2D (IDIM, JDIM)      = INTEGER (Returned)
      }{
         the quality plane
      }
      \sstsubsection{
         STATUS                       = INTEGER (Given and returned)
      }{
         global status
      }
   }
   \sstdiytopic{
      Copyright
   }{
      Copyright \copyright 1995,1996,1997,1998,1999 Particle Physics and Astronomy
      Research Council. All Rights Reserved.
   }
}
\sstroutine{
   SCULIB\_2POS\_CONFN
}{
   Generate a convolution function to remove the 2-position
   chop function from raster scans
}{
   \sstdescription{
       This routine will return the asymmetric convolution function
       required for deconvolving a dual beam map into a single beam map.
       It is the convolution function that is derived in Emerson Klein
       and Haslam (1979) Astron. Astrophys. 76 p92 paper.

      For the case where UNBAL = 1, the two beams are of equal strength:-

       The ideal convolution function would be the one whose FT was the
       inverse of the FT of the chop function. Unfortunately, the chop
       function FT has zeroes in it, at which the inverse FT will tend to
       infinity. The ideal convolution function, therefore, does not exist.
       This problem is avoided by generating a function whose FT is the
       same as the ideal function except at the problem points, where it is
       set to zero. This function consists of a series of delta functions
       separated by the chop spacing, the central 2 points being half the
       chop spacing on either side of the centre of the function.

       The function is normalised such that convolving this function with
       an original chop function of 2 delta functions of unit height will
       give a delta function of height 1.

      Don{\tt '}t understand what happens when UNBAL $<$$>$ 1, but the code has left
      doing the same as NOD2.

       The convolution must cover the map even when the centre of the function
       is at the left or right hand extremity of the map. Hence the convolution
       function must be twice the length of the raw map.

       Since the raw data is not sampled such that the chop spacing is an
       integer number of samples, the actual convolution function must be
       rebinned onto the sample mesh by sinc interpolation.

       This routine is essentially a rewritten version of CONF22 from the
       RESTOR program in the NOD2 package by Haslam (1974) Astron.
       Astrophys. Suppl. 15 p333
   }
   \sstinvocation{
       CALL SCULIB\_2POS\_CONFN (BSEP, PIXSEP, NPIX, UNBAL, NCFN, CONF,
         STATUS)
   }
   \sstarguments{
      \sstsubsection{
         BSEP = REAL (Given)
      }{
         The beam separation in arcseconds
      }
      \sstsubsection{
         PIXSEP = REAL (Given)
      }{
         The pixel separation in arcseconds
      }
      \sstsubsection{
         NPIX = INTEGER (Given)
      }{
         The number of pixels in the x direction
      }
      \sstsubsection{
         UNBAL = REAL (Given)
      }{
         The relative amplitudes of the right and left hand beams
         UNBAL=amp(lhb)/abs(amp(rhb))
      }
      \sstsubsection{
         NCFN = INTEGER (Returned)
      }{
         The length of the convolution array
      }
      \sstsubsection{
         CONF( $*$ ) = REAL (Returned)
      }{
         The convolution function
      }
      \sstsubsection{
         STATUS = INTEGER (Given and returned)
      }{
         Global status
      }
   }
   \sstdiytopic{
      Copyright
   }{
      Copyright \copyright 1995,1996,1997,1998,1999 Particle Physics and Astronomy
      Research Council. All Rights Reserved.
   }
}
\sstroutine{
   SCULIB\_2POS\_DECONV
}{
   deconvolve square chop from scan
}{
   \sstdescription{
      This routine deconvolves the chopped beam response from the
      individual scans of a SCUBA raster map.
         To achieve this it cycles through the scans making up the
      observation, calling SCULIB\_FIND\_SWITCH to locate the start and
      finish indices of each scan in the demodulated data array. If
      SCULIB\_FIND\_SWITCH indicates that there is no data for this scan,
      which might happen if the observation was aborted, then no further
      action is taken. If the scan is too long to be handled by the
      routine an error message will be output and the routine will return
      with bad status.
         Otherwise, SCULIB\_GENSYCONFN and SCULIB\_2POS\_CONFN will be
      called to generate the convolution functions needed to deconvolve
      the chop. The routine then cycles through the bolometers, calling
      SCULIB\_CONVOLVE to do the required convolutions with the scan data
      for each.
   }
   \sstinvocation{
       CALL SCULIB\_2POS\_DECONV (N\_EXPOSURES, N\_INTEGRATIONS,
         N\_MEASUREMENTS, DEMOD\_POINTER, N\_BOL, N\_POS, IN\_DATA,
         IN\_VARIANCE, IN\_QUALITY, SAMPLE\_DX, BEAM\_SEP,
         OUT\_DATA, OUT\_VARIANCE, OUT\_QUALITY, STATUS)
   }
   \sstarguments{
      \sstsubsection{
         N\_EXPOSURES                 = INTEGER (Given)
      }{
         maximum number of exposures per integration
      }
      \sstsubsection{
         N\_INTEGRATIONS              = INTEGER (Given)
      }{
         number of integrations in the observation
      }
      \sstsubsection{
         N\_MEASUREMENTS              = INTEGER (Given)
      }{
         number of measurements in the observation
      }
      \sstsubsection{
         DEMOD\_POINTER (N\_EXPOSURES, N\_INTEGRATIONS, N\_MEASUREMENTS)
      }{
         array pointing to start and finish of scans in IN\_DATA
      }
      \sstsubsection{
         N\_BOL                       = INTEGER (Given)
      }{
         the number of bolometers for which data was taken
      }
      \sstsubsection{
         N\_POS                       = INTEGER (Given)
      }{
         the number of positions measured in the scan
      }
      \sstsubsection{
         IN\_DATA (N\_BOL, N\_POS)      = REAL (Given)
      }{
         the measured data
      }
      \sstsubsection{
         IN\_VARIANCE (N\_BOL, N\_POS)  = REAL (Given)
      }{
         the variance on IN\_DATA
      }
      \sstsubsection{
         IN\_QUALITY (N\_BOL, N\_POS)   = BYTE (Given)
      }{
         the quality on IN\_DATA
      }
      \sstsubsection{
         SAMPLE\_DX                   = REAL (Given)
      }{
         the measurement spacing along the scan (arcseconds)
      }
      \sstsubsection{
         BEAM\_SEP                    = REAL (Given)
      }{
         the beam separation (arcseconds)
      }
      \sstsubsection{
         OUT\_DATA (N\_BOL, N\_POS)     = REAL (Returned)
      }{
         the deconvolved data
      }
      \sstsubsection{
         OUT\_VARIANCE (N\_BOL, N\_POS) = REAL (Returned)
      }{
         the variance on OUT\_DATA
      }
      \sstsubsection{
         OUT\_QUALITY (N\_BOL, N\_POS)  = BYTE (Returned)
      }{
         the quality on OUT\_DATA
      }
      \sstsubsection{
         BADBIT                      = BYTE (Given)
      }{
         bad bit mask
      }
      \sstsubsection{
         STATUS                      = INTEGER (Given and returned)
      }{
         global status
      }
   }
   \sstdiytopic{
      Copyright
   }{
      Copyright \copyright 1995,1996,1997,1998,1999 Particle Physics and Astronomy
      Research Council. All Rights Reserved.
   }
}
\sstroutine{
   SCULIB\_3POS\_CONFN
}{
   Generate a convolution function to remove the 3-position
   chop function from raster scans
}{
   \sstdescription{
      This routine will return the convolution function
      required for deconvolving the 3-position chop from a scan.

      This function consists of a series of delta functions at a spacing
      equal to that between a -ve spike and the central spike of the
      3-position chop. The height of the delta dunctions peaks at the
      centre of the convolution function and falls off by 1 at a time for
      the delta functions on either side, so that the envelope of the
      function is an isosceles triangle.

      The convolution function must cover the scan even when the centre of
      the function is at one end of it. Hence, the convolution function
      must be twice the length of the scan.

      Since the raw data are not sampled such that the chop spacing is an
      integer number of samples, the actual convolution function must be
      rebinned onto the sample mesh by sinc interpolation.
   }
   \sstinvocation{
       CALL SCULIB\_3POS\_CONFN (BSEP, PIXSEP, NPIX, NCFN, CONF,
         STATUS)
   }
   \sstarguments{
      \sstsubsection{
         BSEP = REAL (Given)
      }{
         The beam separation in arcseconds
      }
      \sstsubsection{
         PIXSEP = REAL (Given)
      }{
         The pixel separation in arcseconds
      }
      \sstsubsection{
         NPIX = INTEGER (Given)
      }{
         The number of pixels in the x direction
      }
      \sstsubsection{
         NCFN = INTEGER (Returned)
      }{
         The length of the convolution array
      }
      \sstsubsection{
         CONF( $*$ ) = REAL (Returned)
      }{
         The convolution function
      }
      \sstsubsection{
         STATUS = INTEGER (Given and returned)
      }{
         Global status
      }
   }
   \sstdiytopic{
      Copyright
   }{
      Copyright \copyright 1995,1996,1997,1998,1999 Particle Physics and Astronomy
      Research Council. All Rights Reserved.
   }
}
\sstroutine{
   SCULIB\_3POS\_DECONV
}{
   deconvolve 3-position chop from scan
}{
   \sstdescription{
      This routine deconvolves the 3 position chopped beam response from the
      individual scans of a SCUBA raster map.
         To achieve this it cycles through the scans making up the
      observation, calling SCULIB\_FIND\_SWITCH to locate the start and
      finish indices of each scan in the demodulated data array. If
      SCULIB\_FIND\_SWITCH indicates that there is no data for this scan,
      which may happen if the observation was aborted, then no further
      action is taken. If a scan is too long to be handled by the routine
      an error message will be output and the routine will return with
      bad status.
         Otherwise, SCULIB\_GENSYCONFN and SCULIB\_3POS\_CONFN will be
      called to generate the convolution functions needed to deconvolve
      the chop. The routine then cycles through the bolometers, calling
      SCULIB\_CONVOLVE to do the required convolutions with the scan data
      for each.
   }
   \sstinvocation{
       CALL SCULIB\_3POS\_DECONV (N\_EXPOSURES, N\_INTEGRATIONS,
         N\_MEASUREMENTS, DEMOD\_POINTER, N\_BOL, N\_POS, IN\_DATA,
         IN\_VARIANCE, IN\_QUALITY, SAMPLE\_DX, BEAM\_SEP,
         OUT\_DATA, OUT\_VARIANCE, OUT\_QUALITY, STATUS)
   }
   \sstarguments{
      \sstsubsection{
         N\_EXPOSURES                 = INTEGER (Given)
      }{
         maximum number of exposures per integration
      }
      \sstsubsection{
         N\_INTEGRATIONS              = INTEGER (Given)
      }{
         number of integrations in the observation
      }
      \sstsubsection{
         N\_MEASUREMENTS              = INTEGER (Given)
      }{
         number of measurements in the observation
      }
      \sstsubsection{
         DEMOD\_POINTER (N\_EXPOSURES, N\_INTEGRATIONS, N\_MEASUREMENTS)
      }{
         array pointing to start and finish of scans in IN\_DATA
      }
      \sstsubsection{
         N\_BOL                       = INTEGER (Given)
      }{
         the number of bolometers for which data was taken
      }
      \sstsubsection{
         N\_POS                       = INTEGER (Given)
      }{
         the number of positions measured in the scan
      }
      \sstsubsection{
         IN\_DATA (N\_BOL, N\_POS)      = REAL (Given)
      }{
         the measured data
      }
      \sstsubsection{
         IN\_VARIANCE (N\_BOL, N\_POS)  = REAL (Given)
      }{
         the variance on IN\_DATA
      }
      \sstsubsection{
         IN\_QUALITY (N\_BOL, N\_POS)   = BYTE (Given)
      }{
         the quality on IN\_DATA
      }
      \sstsubsection{
         SAMPLE\_DX                   = REAL (Given)
      }{
         the measurement spacing along the scan (arcseconds)
      }
      \sstsubsection{
         BEAM\_SEP                    = REAL (Given)
      }{
         the beam separation (arcseconds)
      }
      \sstsubsection{
         OUT\_DATA (N\_BOL, N\_POS)     = REAL (Returned)
      }{
         the deconvolved data
      }
      \sstsubsection{
         OUT\_VARIANCE (N\_BOL, N\_POS) = REAL (Returned)
      }{
         the variance on OUT\_DATA
      }
      \sstsubsection{
         OUT\_QUALITY (N\_BOL, N\_POS)  = BYTE (Returned)
      }{
         the quality on OUT\_DATA
      }
      \sstsubsection{
         STATUS                      = INTEGER (Given and returned)
      }{
         global status
      }
   }
   \sstdiytopic{
      Copyright
   }{
      Copyright \copyright 1995,1996,1997,1998,1999 Particle Physics and Astronomy
      Research Council. All Rights Reserved.
   }
}
\sstroutine{
   SCULIB\_ADD\_CHOP
}{
   Add chop throw to apparent RA/Dec
}{
   \sstdescription{
      This routine takes a chop throw and adds it on the
      supplied apparent RA/Dec centre position returning a new
      apparent RA/Dec. Works with scan map and jiggle map
      but does not actually know the difference itself (unless
      using SC chop). For EKH (SC) chopping we have to convert
      the scan ends to tangent plane offsets from the reference
      centre before calculating the angle.
   }
   \sstinvocation{
       CALL SCULIB\_ADD\_CHOP(BEAM, RA\_REF\_CENTRE, DEC\_REF\_CENTRE, RA\_CENTRE,
            DEC\_CENTRE, CHOP\_CRD, CHOP\_PA, CHOP\_FUN, CHOP\_THROW,
            LST, MJD, LAT\_OBS, RA\_START, RA\_END, DEC\_START,
            DEC\_END, OUT\_RA\_CEN, OUT\_DEC\_CEN,
            STATUS)
   }
   \sstarguments{
      \sstsubsection{
         BEAM        = CHARACTER $*$ ($*$) (Given)
      }{
         Beam name ({\tt '}M{\tt '},{\tt '}L{\tt '} or {\tt '}R{\tt '})
      }
      \sstsubsection{
         RA\_REF\_CENTRE = DOUBLE PRECISION (Given)
      }{
         Apparent RA of the map centre. This is the centre that
         scan map tangent plane offsets are calculated from.
      }
      \sstsubsection{
         DEC\_REF\_CENTRE = DOUBLE PRECISION (Given)
      }{
         Apparent dec of the map centre. This is the centre that
         scan map tangent plane offsets are calculated from.
      }
      \sstsubsection{
         RA\_CENTRE              = DOUBLE PRECISION (Given)
      }{
         the apparent RA of the `centre{\tt '} of the array (radians)
      }
      \sstsubsection{
         DEC\_CENTRE             = DOUBLE PRECISION (Given)
      }{
         the apparent dec of the `centre{\tt '} of the array (radians)
      }
      \sstsubsection{
         CHOP\_CRD               = CHARACTER $*$($*$) (Given)
      }{
         Coordinate system of CHOP throw (SC,AZ,RB,RJ,GA)
      }
      \sstsubsection{
         CHOP\_PA                = REAL (Given)
      }{
         Position angle of chop (in radians)
      }
      \sstsubsection{
         CHOP\_FUN               = CHARACTER $*$ ($*$) (Given)
      }{
         Chop function (CENTER, SQUARE, TRIPOS)
      }
      \sstsubsection{
         CHOP\_THROW             = REAL (Given)
      }{
         Chop throw in radians
      }
      \sstsubsection{
         LST                    = DOUBLE PRECISION (Given)
      }{
         the local sidereal time (radians)
      }
      \sstsubsection{
         MJD                    = DOUBLE PRECISION (Given)
      }{
         MJD of observation
      }
      \sstsubsection{
         LAT\_OBS                = DOUBLE PRECISION (Given)
      }{
         the latitude of the observatory (radians)
      }
      \sstsubsection{
         RA\_START               = REAL (Given)
      }{
         RA of start of scan (if SAMPLE\_MODE=RASTER)
      }
      \sstsubsection{
         RA\_END                 = REAL (Given)
      }{
         RA of end of scan (if SAMPLE\_MODE=RASTER)
      }
      \sstsubsection{
         DEC\_START              = REAL (Given)
      }{
         DEC of start of scan (if SAMPLE\_MODE=RASTER)
      }
      \sstsubsection{
         DEC\_END                = REAL (Given)
      }{
         DEC of end of scan (if SAMPLE\_MODE=RASTER)
      }
      \sstsubsection{
         OUT\_RA\_CEN             = DOUBLE (Returned)
      }{
         New position with CHOP
      }
      \sstsubsection{
         OUT\_DEC\_CEN            = DOUBLE (Returned)
      }{
         New DEC with CHOP
      }
      \sstsubsection{
         STATUS                 = INTEGER (Given and returned)
      }{
         The global status
      }
   }
   \sstnotes{
      It is assumed that the chop throw is divided by two
      prior to calling this routine for scan map mode since the
      chop is effectively half a chop either side of the middle.
      For jiggle modes the chop is effectively a full chop throw
      from the centre since the 3-beam chopping is done by a combination
      of chopping and nodding.
      \sstitemlist{

         \sstitem
         This routine does not yet support TRIPOS chopping (I think)

         \sstitem
         AZ, LO and SC chopping are supported

         \sstitem
         The negative beam is always the {\tt '}left{\tt '} beam, the $+$ve beam is the
           {\tt '}right{\tt '}.

         \sstitem
         The chop tracking {\tt '}droopy{\tt '} beam problem should be dealt with
           in an earlier subroutine such that the chop is converted from
           LO to AZ before calling this routine
      }
   }
   \sstdiytopic{
      Copyright
   }{
      Copyright \copyright 1995,1996,1997,1998,1999 Particle Physics and Astronomy
      Research Council. All Rights Reserved.
   }
}
\sstroutine{
   SCULIB\_ADD\_DEMOD\_EXPOSURE
}{
   add demodulated data for an exposure into the
   integration result
}{
   \sstdescription{
      This routine adds the reduced demodulated data for an exposure in a
      jiggle-sampled observation into the integration result.
        The exposure can contain data either for all or part of ONE run
      through the jiggle pattern or for several REPEATS of the entire jiggle
      pattern. In the first case, the routine copies the exposure data into
      the appropriate part of the integration arrays. In the second, the
      integration data will be the average of the input exposure values and the
      integration variance will be set equal to the exposure value if there
      was just one valid measurement, or calculated from the dispersion of
      the exposure data about the mean otherwise.
   }
   \sstinvocation{
       CALL SCULIB\_ADD\_DEMOD\_EXPOSURE (N\_BOLS, J\_START, N\_JIG,
         J\_REPEAT, J\_COUNT, EXP\_DATA, EXP\_VARIANCE, EXP\_QUALITY,
         INT\_DATA, INT\_VARIANCE, INT\_QUALITY, STATUS)
   }
   \sstarguments{
      \sstsubsection{
         N\_BOLS                           = INTEGER (Given)
      }{
         the number of bolometers measured
      }
      \sstsubsection{
         J\_START                          = INTEGER (Given)
      }{
         the index within the pattern of the first jiggle position measured
      }
      \sstsubsection{
         N\_JIG                            = INTEGER (Given)
      }{
         the number of jiggle positions measured
      }
      \sstsubsection{
         J\_REPEAT                         = INTEGER (Given)
      }{
         the number of times the measured pattern was repeated
      }
      \sstsubsection{
         J\_COUNT                          = INTEGER (Given)
      }{
         the number of jiggle positions in the entire pattern
      }
      \sstsubsection{
         EXP\_DATA (N\_BOLS, N\_JIG)         = REAL (Given)
      }{
         the data for the exposure
      }
      \sstsubsection{
         EXP\_VARIANCE (N\_BOLS, N\_JIG)     = REAL (Given)
      }{
         the exposure variance
      }
      \sstsubsection{
         EXP\_QUALITY (N\_BOLS, N\_JIG)      = INTEGER (Given)
      }{
         the exposure quality
      }
      \sstsubsection{
         INT\_DATA (N\_BOLS, J\_COUNT)       = REAL (Returned)
      }{
         the integration data
      }
      \sstsubsection{
         INT\_VARIANCE (N\_BOLS, J\_COUNT)   = REAL (Returned)
      }{
         the integration variance
      }
      \sstsubsection{
         INT\_QUALITY (N\_BOLS, J\_COUNT)    = INTEGER (Returned)
      }{
         the integration quality
      }
      \sstsubsection{
         STATUS                           = INTEGER (Given and returned)
      }{
         global status
      }
   }
   \sstdiytopic{
      Copyright
   }{
      Copyright \copyright 1995,1996,1997,1998,1999 Particle Physics and Astronomy
      Research Council. All Rights Reserved.
   }
}
\sstroutine{
   SCULIB\_ADDARE
}{
   add one real array to another into a third
}{
   \sstdescription{
      Adds two real arrays. Note that any of the arrays may be the same.
   }
   \sstinvocation{
       CALL SCULIB\_ADDARE (N, ARRAY1, ARRAY2, ARRAY3, Q1DATA,
          Q2DATA, Q3DATA, V1DATA, V2DATA, V3DATA, QUALITY, FLAGGED,
          VARIANCE)
   }
   \sstarguments{
      \sstsubsection{
         N            = INTEGER (Given)
      }{
         Number of elements in each array
      }
      \sstsubsection{
         ARRAY1 (N)   = REAL (Given)
      }{
         Input array
      }
      \sstsubsection{
         ARRAY2 (N)   = REAL (Given)
      }{
         Second input array
      }
      \sstsubsection{
         ARRAY3 (N)   = REAL (Returned)
      }{
         Result array.  ARRAY3 = ARRAY1 $+$ ARRAY2
      }
      \sstsubsection{
         Q1DATA (N)   = INTEGER (Given)
      }{
         Quality array for first input array
      }
      \sstsubsection{
         Q2DATA (N)   = INTEGER (Given)
      }{
         Quality array for second input array
      }
      \sstsubsection{
         Q3DATA (N)   = INTEGER (Returned)
      }{
         Quality array for output array
      }
      \sstsubsection{
         V1DATA (N)   = REAL (Given)
      }{
         Variance array for first input array
      }
      \sstsubsection{
         V2DATA (N)   = REAL (Given)
      }{
         Variance array for second input array
      }
      \sstsubsection{
         V3DATA (N)   = REAL (Returned)
      }{
         Variance array for output array
      }
      \sstsubsection{
         QUALITY      = LOGICAL (Given)
      }{
         True if input has quality information
      }
      \sstsubsection{
         FLAGGED      = LOGICAL (Given)
      }{
         True if input has flagged data values
      }
      \sstsubsection{
         VARIANCE     = LOGICAL (Given)
      }{
         True if both input arrays have variance arrays
      }
   }
   \sstnotes{
      \sstitemlist{

         \sstitem
         Does not use Quality correctly.
           Uses INTEGER quality rather than UBYTE
      }
   }
   \sstdiytopic{
      Copyright
   }{
      Copyright \copyright 1995,1996,1997,1998,1999 Particle Physics and Astronomy
      Research Council. All Rights Reserved.
   }
   \sstimplementationstatus{
      \sstitemlist{

         \sstitem
         Propagates variance

         \sstitem
         Checks for bad values
      }
   }
}
\sstroutine{
   SCULIB\_ADDCAD
}{
   add a constant to a double array
}{
   \sstdescription{
      adds a real constant to a double array
   }
   \sstinvocation{
      CALL SCULIB\_ADDCAD (N, IN, RVAL, OUT)
   }
   \sstarguments{
      \sstsubsection{
         N              = INTEGER (Given)
      }{
         number of elements in arrays
      }
      \sstsubsection{
         IN (N)         = DOUBLE PRECISION (Given)
      }{
         input array
      }
      \sstsubsection{
         RVAL           = DOUBLE PRECISION (Given)
      }{
         real constant to be added to array
      }
      \sstsubsection{
         OUT (N)        = DOUBLE PRECISION (Returned)
      }{
         output array (may be same as input)
      }
   }
   \sstnotes{
      \sstitemlist{

         \sstitem
         No range checks are performed

         \sstitem
         No status checking

         \sstitem
         No checks for bad values
      }
   }
   \sstdiytopic{
      Copyright
   }{
      Copyright \copyright 1995,1996,1997,1998,1999 Particle Physics and Astronomy
      Research Council. All Rights Reserved.
   }
}
\sstroutine{
   SCULIB\_ADDCAI
}{
   add a constant to an integer array
}{
   \sstdescription{
      adds an integer constant to an integer array
   }
   \sstinvocation{
      CALL SCULIB\_ADDCAI (N, IN, IVAL, OUT)
   }
   \sstarguments{
      \sstsubsection{
         N              = INTEGER (Given)
      }{
         number of elements in arrays
      }
      \sstsubsection{
         IN (N)         = INTEGER (Given)
      }{
         input array
      }
      \sstsubsection{
         IVAL           = INTEGER (Given)
      }{
         constant to be added to array
      }
      \sstsubsection{
         OUT (N)        = INTEGER (Returned)
      }{
         output array (may be same as input)
      }
   }
   \sstnotes{
      \sstitemlist{
         \sstitem
         No range checks are performed

         \sstitem
         No status checking

         \sstitem
         No checks for bad values
      }
   }
   \sstdiytopic{
      Copyright
   }{
      Copyright \copyright 1995,1996,1997,1998,1999 Particle Physics and Astronomy
      Research Council. All Rights Reserved.
   }
}
\sstroutine{
   SCULIB\_ADDCAR
}{
   add a constant to a real array
}{
   \sstdescription{
      adds a real constant to a real array
   }
   \sstinvocation{
      CALL SCULIB\_ADDCAR (N, IN, RVAL, OUT)
   }
   \sstarguments{
      \sstsubsection{
         N              = INTEGER (Given)
      }{
         number of elements in arrays
      }
      \sstsubsection{
         IN (N)         = REAL (Given)
      }{
         input array
      }
      \sstsubsection{
         RVAL           = REAL (Given)
      }{
         real constant to be added to array
      }
      \sstsubsection{
         OUT (N)        = REAL (Returned)
      }{
         output array (may be same as input)
      }
   }
   \sstnotes{
      \sstitemlist{

         \sstitem
         No range checks are performed

         \sstitem
         No status checking

         \sstitem
         No checks for bad values
      }
   }
   \sstdiytopic{
      Copyright
   }{
      Copyright \copyright 1995,1996,1997,1998,1999 Particle Physics and Astronomy
      Research Council. All Rights Reserved.
   }
}
\sstroutine{
   SCULIB\_AIRMASS
}{
   calculate the airmass for a given zenith distance
}{
   \sstdescription{
      This routine calculates the airmass corresponding to the input
      zenith distance. The airmass is just sec(Z) until it reaches 2,
      after which a more complex algorithm developed by Ian Coulson
      is used (it is described by comments in the code). There may be
      a discontinuity in returned airmasses around the value 2, I
      haven{\tt '}t checked.
   }
   \sstinvocation{
      CALL SCULIB\_AIRMASS (Z, AIRMASS, STATUS)
   }
   \sstarguments{
      \sstsubsection{
         Z               = REAL (Given)
      }{
         zenith distance (radians)
      }
      \sstsubsection{
         AIRMASS         = REAL (Returned)
      }{
         airmass
      }
      \sstsubsection{
         STATUS          = INTEGER (Given and returned)
      }{
         global status
      }
   }
   \sstdiytopic{
      Copyright
   }{
      Copyright \copyright 1993-1999 Particle Physics and Astronomy
      Research Council. All Rights Reserved.
   }
}
\sstroutine{
   SCULIB\_ANALYSE\_PHOTOM\_JIGGLE
}{
   analyse the jiggle map made by a
   bolometer during a PHOTOM observation
}{
   \sstdescription{
      This routine analyses the jiggle map made by a bolometer during a
      PHOTOM observation.
         If status is good on entry the routine checks that the data array
      holds data for the specified bolometer. If it does not an error
      message will be output and the bad status returned.
         All being well the routine will analyse the data by the method
       specified in PHOTOM\_ANALYSIS:-

      AVERAGE -    the `result{\tt '} will be the mean of the jiggle points.
                   The variance is the sum of the variances of the averaged
                   data divided by the number of them squared. The result
                   will be flat-fielded and reported to the user. If no
                   valid data were available for the calculation of the
                   mean then a warning message will be output but good
                   status returned.

      PARABOLA -   the routine will call SCULIB\_FIT\_2D\_PARABOLA to fit a
                   2-d parabola to the data and the `result{\tt '} will be the
                   peak of the fitted curve. If the fit is good the results
                   will be flat-fielded and reported to the user. Otherwise
                   the result quality is set bad but good status will be
                   returned.
   }
   \sstinvocation{
       CALL SCULIB\_ANALYSE\_PHOTOM\_JIGGLE (PHOTOM\_ANALYSIS, BOL, N\_BOLS,
         J\_COUNT, JDATA, VARIANCE, QUALITY, RESULT\_D, RESULT\_V, RESULT\_Q,
         BADBIT, STATUS)
   }
   \sstarguments{
      \sstsubsection{
         PHOTOM\_ANALYSIS              = CHARACTER$*$($*$) (Given)
      }{
         the method of analysis to be applied to the data
      }
      \sstsubsection{
         BOL                          = INTEGER (Given)
      }{
         the index of the bolometer whose data is to be analysed
      }
      \sstsubsection{
         N\_BOLS                       = INTEGER (Given)
      }{
         the number of bolometers being measured
      }
      \sstsubsection{
         J\_COUNT                      = INTEGER (Given)
      }{
         the number of jiggles in the pattern
      }
      \sstsubsection{
         JIGGLE\_X                     = REAL (Given)
      }{
         x jiggle offsets
      }
      \sstsubsection{
         JIGGLE\_Y                     = REAL (Given)
      }{
         y jiggle offsets
      }
      \sstsubsection{
         JDATA (N\_BOLS, J\_COUNT)       = REAL (Given)
      }{
         the measured data
      }
      \sstsubsection{
         VARIANCE (N\_BOLS, J\_COUNT)   = REAL (Given)
      }{
         the variance
      }
      \sstsubsection{
         QUALITY (N\_BOLS, J\_COUNT)    = BYTE (Given)
      }{
         the quality
      }
      \sstsubsection{
         RESULT\_D                     = REAL (Returned)
      }{
         the sum of the input data
      }
      \sstsubsection{
         RESULT\_V                     = REAL (Returned)
      }{
         the sum of the input variances
      }
      \sstsubsection{
         RESULT\_Q                     = BYTE (Returned)
      }{
         the quality on the sum
      }
      \sstsubsection{
         A0                           = REAL (Returned)
      }{
         fitted parabola parameter
      }
      \sstsubsection{
         A1                           = REAL (Returned)
      }{
         fitted parabola parameter
      }
      \sstsubsection{
         X0                           = REAL (Returned)
      }{
         x offset of peak of fitted parabola
      }
      \sstsubsection{
         Y0                           = REAL (Returned)
      }{
         y offset of peak of fitted parabola
      }
      \sstsubsection{
         BADBIT                       = BYTE (Given)
      }{
         bad bit mask
      }
      \sstsubsection{
         STATUS                       = INTEGER (Given and returned)
      }{
         global status
      }
   }
   \sstdiytopic{
      Copyright
   }{
      Copyright \copyright 1993-1999 Particle Physics and Astronomy
      Research Council. All Rights Reserved.
   }
}
\sstroutine{
   SCULIB\_APPARENT\_2\_MP
}{
   Calculate mean place from a given apparent RA,Dec
}{
   \sstdescription{
      This routine takes the apparent RA,Dec as input and converts
      them to the position in a specified output coordinate system.
      Allowed OUTPUT\_COORDS are AZ, RJ, RB, GA, EQ.
      This is the reverse of SCULIB\_CALC\_APPARENT
   }
   \sstinvocation{
       CALL SCULIB\_APPARENT\_2\_MP(RA\_APP, DEC\_APP, OUT\_COORDS,
            LST, MJD, LAT\_OBS, LONG, LAT,
            STATUS )
   }
   \sstarguments{
      \sstsubsection{
         RA\_APP                 = DOUBLE PRECISION (Given)
      }{
         Apparent RA of point at date (radians)
      }
      \sstsubsection{
         DEC\_APP                = DOUBLE PRECISION (Given)
      }{
         Apparent Dec
      }
      \sstsubsection{
         OUT\_COORDS             = CHARACTER $*$ ($*$) (Given)
      }{
      }
      \sstsubsection{
         LST                    = DOUBLE PRECISION (Given)
      }{
         LST for requested coordinates (for AZ and HA)
      }
      \sstsubsection{
         MJD                    = DOUBLE PRECISION (Given)
      }{
         Modified Julian date of observation
      }
      \sstsubsection{
         LAT\_OBS                = DOUBLE PRECISION (Given)
      }{
         Latitude of observatory in radians.
         For JCMT this value is 3.46026051751D-1
      }
      \sstsubsection{
         LONG                   = DOUBLE PRECISION (Returned)
      }{
         longitude of centre in input coord system (radians)
      }
      \sstsubsection{
         LAT                    = DOUBLE PRECISION (Returned)
      }{
         latitude of centre in input coord system (radians)
      }
      \sstsubsection{
         STATUS                 = INTEGER (Given and returned)
      }{
         Global status
      }
   }
   \sstdiytopic{
      Copyright
   }{
      Copyright \copyright 1995,1996,1997,1998,1999 Particle Physics and Astronomy
      Research Council. All Rights Reserved.
   }
}
\sstroutine{
   SCULIB\_APPARENT\_2\_TP
}{
   calculate tangent plane coordinates from apparent
   RA, Decs
}{
   \sstdescription{
      This routine converts a list of apparent RA,Decs to tangent plane
      offsets from a tangent point whose position is also given in apparent
      RA,Dec. In addition, a rotation is applied to the tangent plane so that
      it can be aligned with a coordinate system other than apparent RA,Dec,
      and a shift is added to allow the map to be moved about in that output
      frame.
   }
   \sstinvocation{
       CALL SCULIB\_APPARENT\_2\_TP (N\_POS, BOL\_XPOS, BOL\_YPOS, RA\_CEN,
         DEC\_CEN, ROTATION, SHIFT\_DX, SHIFT\_DY, STATUS)
   }
   \sstarguments{
      \sstsubsection{
         N\_POS                  = INTEGER (Given)
      }{
         the number of positions to be converted
      }
      \sstsubsection{
         BOL\_XPOS (N\_POS)       = DOUBLE PRECISION (Given and returned)
      }{
         apparent RA on input (radians), x tangent plane offset on
         output (radians)
      }
      \sstsubsection{
         BOL\_YPOS (N\_POS)       = DOUBLE PRECISION (Given and returned)
      }{
         apparent Dec on input (radians), y tangent plane offset on
         output (radians)
      }
      \sstsubsection{
         RA\_CEN                 = DOUBLE PRECISION (Given)
      }{
         apparent RA of tangent point (radians)
      }
      \sstsubsection{
         DEC\_CEN                = DOUBLE PRECISION (Given)
      }{
         apparent Dec of tangent point (radians)
      }
      \sstsubsection{
         ROTATION               = DOUBLE PRECISION (Given)
      }{
         angle between output North and apparent North (radians,
         measured anticlockwise from output North)
      }
      \sstsubsection{
         SHIFT\_DX               = DOUBLE PRECISION (Given)
      }{
         value to be added to x offsets (radians)
      }
      \sstsubsection{
         SHIFT\_DY               = DOUBLE PRECISION (Given)
      }{
         value to be added to y offsets (radians)
      }
      \sstsubsection{
         STATUS                 = INTEGER (Given and returned)
      }{
      }
   }
   \sstdiytopic{
      Copyright
   }{
      Copyright \copyright 1995,1996,1997,1998,1999 Particle Physics and Astronomy
      Research Council. All Rights Reserved.
   }
}
\sstroutine{
   SCULIB\_BESSEL\_WTINIT
}{
   Generate a weighting function for rebinning
}{
   \sstdescription{
      This is a FORTRAN version of the C code written for the transputer
      rebinning. Here is the C description:
      Initialise the BESSEL weighting function for Bessel interpolation.
      The 1-D function 2.0$*$J1(x)/x is initialised in the declaration, and the
      weighting function lookup table is calculated from this. The
      tabulation is at intervals of pi/RES for the range 0.0 to RADIUS$*$pi,
      resulting in (RES$*$RADIUS)$*$$*$2 values. The weight function is also
      multiplied by a cosine over the oute third to reduce edge effects.

      The look up table is initialised such that one can access the
      weight directly given the square of the distance between the input
      and output pixel. We use a fixed resolution out to 10PI.
   }
   \sstinvocation{
      CALL SCULIB\_BESSEL\_WTINIT( WTFN, RADIUES, RES, STATUS)
   }
   \sstarguments{
      \sstsubsection{
         WTFN ( RADIUS $*$ RADIUS $*$ RES $*$ RES $+$ 1) = REAL (Returned)
      }{
         The weighting function generated by this routine.
         The index corresponds to the square of the distance from the
         centre in scale units.
      }
      \sstsubsection{
         RADIUS = INTEGER (Given)
      }{
         Size of the weighting function in scale units.
      }
      \sstsubsection{
         RES = INTEGER (Given)
      }{
         Number of points per scale lenghth.
      }
      \sstsubsection{
         STATUS = INTEGER (Given \& Returned)
      }{
         Global status
      }
   }
   \sstdiytopic{
      Copyright
   }{
      Copyright (C) 1995,1996,1997,1998,1999 Particle Physics and Astronomy
      Research Council. All Rights Reserved.
   }
}
\sstroutine{
   SCULIB\_BESSJ1
}{
   calculates Bessel function J1(x)
}{
   \sstdescription{
      If status is good on entry this function returns the Bessel function
      J1(x) for any real x.
   }
   \sstinvocation{
      Y = SCULIB\_BESSJ1 (X, STATUS)
   }
   \sstarguments{
      \sstsubsection{
         X                      = REAL (Given)
      }{
         the argument of the J1 Bessel function
      }
      \sstsubsection{
         STATUS                 = INTEGER (Given and returned)
      }{
         global status
      }
   }
   \sstreturnedvalue{
      \sstsubsection{
         SCULIB\_BESSJ1 = REAL
      }{
         Bessel function J1(X)
      }
   }
   \sstnotes{
      Probably should be changed to use PDA\_DBESJ1 (SUN/194). That routine
      is DOUBLE PRECISION.
   }
   \sstdiytopic{
      Method
   }{
      Uses an algorithm from Numerical Recipes in Fortran.
   }
   \sstdiytopic{
      Copyright
   }{
      Copyright \copyright 1995,1996,1997,1998,1999 Particle Physics and Astronomy
      Research Council. All Rights Reserved.
   }
}
\sstroutine{
   SCULIB\_BITON
}{
   turn on a bit
}{
   \sstdescription{
      Turns on a bit in a byte.
   }
   \sstinvocation{
      NEWBYTE = SCULIB\_BITON( VAL, BIT )
   }
   \sstarguments{
      \sstsubsection{
         VAL              = BYTE (Given)
      }{
         The byte to be changed
      }
      \sstsubsection{
         BIT              = INT (Given)
      }{
         The bit to be turned on
      }
   }
   \sstreturnedvalue{
      \sstsubsection{
         SCULIB\_BITON = BYTE
      }{
         Modified byte.
      }
   }
   \sstdiytopic{
      Copyright
   }{
      Copyright \copyright 1995,1996,1997,1998,1999 Particle Physics and Astronomy
      Research Council. All Rights Reserved.
   }
}
\sstroutine{
   SCULIB\_BITOFF
}{
   Turn off a bit
}{
   \sstdescription{
      Turns off a bit in a byte
   }
   \sstinvocation{
      NEWBYTE = SCULIB\_BITOFF ( VAL, BIT )
   }
   \sstarguments{
      \sstsubsection{
         VAL              = BYTE (Given)
      }{
         The byte to be changed
      }
      \sstsubsection{
         BIT              = INT (Given)
      }{
         The bit to be turned off
      }
   }
   \sstreturnedvalue{
      \sstsubsection{
         SCULIB\_BITOFF = BYTE
      }{
         Modified byte.
      }
   }
   \sstdiytopic{
      Copyright
   }{
      Copyright \copyright 1995,1996,1997,1998,1999 Particle Physics and Astronomy
      Research Council. All Rights Reserved.
   }
}
\sstroutine{
   SCULIB\_BITOR
}{
   Returns the bitwise OR of two bytes
}{
   \sstdescription{
      Returns the bitwise OR
   }
   \sstinvocation{
      LOGOR = SCULIB\_BITOR(VAL1, VAL2)
   }
   \sstarguments{
      \sstsubsection{
         VAL1              = BYTE (Given)
      }{
         First byte
      }
      \sstsubsection{
         VAL2              = BYTE (Given)
      }{
         Second byte
      }
   }
   \sstreturnedvalue{
      \sstsubsection{
         SCULIB\_BITOR = BYTE
      }{
         OR of bytes
      }
   }
   \sstdiytopic{
      Copyright
   }{
      Copyright \copyright 1995,1996,1997,1998,1999 Particle Physics and Astronomy
      Research Council. All Rights Reserved.
   }
}
\sstroutine{
   SCULIB\_BITAND
}{
   Calculate the bitwise AND of two bytes
}{
   \sstdescription{
      Returns the bitwise AND
   }
   \sstinvocation{
      LOGAND = SCULIB\_BITAND( VAL1, VAL2 )
   }
   \sstarguments{
      \sstsubsection{
         VAL1              = BYTE (Given)
      }{
         First byte
      }
      \sstsubsection{
         VAL2              = BYTE (Given)
      }{
         Second byte
      }
   }
   \sstreturnedvalue{
      \sstsubsection{
         SCULIB\_BITAND = BYTE
      }{
         AND of bytes.
      }
   }
   \sstdiytopic{
      Copyright
   }{
      Copyright \copyright 1995,1996,1997,1998,1999 Particle Physics and Astronomy
      Research Council. All Rights Reserved.
   }
}
\sstroutine{
   SCULIB\_BITTEST
}{
   Test whether a bit is on
}{
   \sstdescription{
      Tests whether specified bit is turned on or not.
      Returns TRUE if it is on, false otherwise.
   }
   \sstinvocation{
      ISON = SCULIB\_BITTEST( VAL, BIT )
   }
   \sstarguments{
      \sstsubsection{
         VAL              = BYTE (Given)
      }{
         The byte to be tested
      }
      \sstsubsection{
         BIT              = INT (Given)
      }{
         The bit to be tested
      }
   }
   \sstreturnedvalue{
      \sstsubsection{
         SCULIB\_BITTEST = LOGICAL
      }{
         True if specified bit is turned on.
      }
   }
   \sstdiytopic{
      Copyright
   }{
      Copyright \copyright 1995,1996,1997,1998,1999 Particle Physics and Astronomy
      Research Council. All Rights Reserved.
   }
}
\sstroutine{
   SCULIB\_BOLDECODE
}{
   decode a bolometer ID into ADC and channel number
}{
   \sstdescription{
      Given a character string containing a bolometer ID, put ADC and
      CHANNEL number into output variables. The syntax of the bolometer
      ID is a3, A3 or 3. Allowed ADCs run from a-i. Channels
      must be in range 1-16. Bolometer IDs outside these ranges will cause
      an error to be returned.
   }
   \sstinvocation{
      CALL SCULIB\_BOLDECODE (BOLCODE, ADC, CHANNEL, STATUS)
   }
   \sstarguments{
      \sstsubsection{
         BOLCODE                     = CHARACTER$*$($*$) (given)
      }{
         bolometer ID
      }
      \sstsubsection{
         ADC                         = INTEGER (Returned)
      }{
         ADC number
      }
      \sstsubsection{
         CHANNEL                     = INTEGER (Returned)
      }{
         channel number
      }
      \sstsubsection{
         STATUS                      = INTEGER (Given and returned)
      }{
         global status
      }
   }
   \sstdiytopic{
      Copyright
   }{
      Copyright \copyright 1995,1996,1997,1998,1999 Particle Physics and Astronomy
      Research Council. All Rights Reserved.
   }
}
\sstroutine{
   SCULIB\_BOLNAME
}{
   generate bolometer name from ADC and channel number
}{
   \sstdescription{
      Given bolometer ADC and channel, generate the bolometer name. E.g.
      ADC=2, CHAN=5 gives B5. If the ADC number is outside the range 1-9, or
      the channel number outside 1-16, and error will be reported and bad status
      returned.
   }
   \sstinvocation{
      CALL SCULIB\_BOLNAME (ADC, CHANNEL, BOLCODE, STATUS)
   }
   \sstarguments{
      \sstsubsection{
         ADC                         = INTEGER (Given)
      }{
         ADC number
      }
      \sstsubsection{
         CHANNEL                     = INTEGER (Given)
      }{
         channel number
      }
      \sstsubsection{
         BOLCODE                     = CHARACTER$*$($*$) (Returned)
      }{
         bolometer ID
      }
      \sstsubsection{
         STATUS                      = INTEGER (Given and returned)
      }{
         global status
      }
   }
   \sstdiytopic{
      Copyright
   }{
      Copyright \copyright 1995,1996,1997,1998,1999 Particle Physics and Astronomy
      Research Council. All Rights Reserved.
   }
}
\sstroutine{
   SCULIB\_BOLSELECT
}{
   interpret a list of selected bolometers
}{
   \sstdescription{
      Interpret a character string containing a list of bolometers,
      put number of bolometers to be used in N\_BOL\_SELECT, channel
      and ADC of selected bolometers in BOL\_SELECT\_CHAN and BOL\_SELECT\_ADC
      respectively. BOL\_ENABLED entries will also be set .TRUE. for bolometers
      that have been selected. The number of sub-instruments involved
      is returned in N\_SUB, and the names of the sub-instruments are given
      in SUB\_INSTRMNT. The Nasmyth coords of the point on the focal plane
      to be treated as the axis of the instrument will be returned in
      CENTRE\_DU3, CENTRE\_DU4.

      The given string can either be ALL, in which case all bolometer channels
      will be selected, or be a list of words specifying parts of the 144
      channel array. Each word in such a list must be separated from the others
      by {\tt '},{\tt '}, and the words can select bolometers by type - SHORT, LONG,
      P1100, P1300, P2000, SHORT\_DC, LONG\_DC, P1100\_DC, P1300\_DC, P2000\_DC,
      or by channel number, either individually or as a range, specified by
      ADC letter and channel number e.g. a3. Allowed ADCs are A-I and channels
      1-16.

      Repeated words in the list will be ignored, as will repeated selections
      of the same bolometer by different routes.

      The instrument sections returned in SUB\_INSTRMNT will be arranged in the
      order SHORT, LONG, P1100, P1300, P2000.

      CENTRE\_DU3, CENTRE\_DU4 will be set equal to ARRAY\_CENTRE\_DU3, ARRAY\_
      CENTRE\_DU4 if the first word in the selection is SHORT, LONG,
      SHORT\_DC or LONG\_DC. They will be set to the offsets of the bolometer
      concerned if the first word in the selection is P1100, P1300, P2000
      or a bolometer name like A7. If either of the coordinates ends up
      being set to the {\tt '}bad{\tt '} value, they will be reset to zero and a
      warning message output.

      Errors will be reported and bad status returned if -

      \sstitemlist{

         \sstitem
          there are too many words in BOLOMETERS

         \sstitem
          BOLOMETERS is empty

         \sstitem
          if when bolometers are selected by type
              no detectors of the desired type are found in the database

         \sstitem
          if the type selected refers to a single bolometer (P1100, P1300,
           P2000, SHORT\_DC, LONG\_DC, P1100\_DC, P1300\_DC, P2000\_DC)
              more than bolometer of the required type is found in the database

         \sstitem
          if bolometers are selected by ID, e.g. a15,
              the ID or range of IDs is bad

      }
      Warnings will be reported but good status returned if -

      \sstitemlist{

         \sstitem
          any selected bolometer has undefined attributes (type, calib, dU3, or
           dU4)

         \sstitem
          CENTRE\_DU3, CENTRE\_DU4 are set to {\tt '}bad{\tt '} values (they will be reset
           to 0 too)
      }
   }
   \sstinvocation{
       CALL SCULIB\_BOLSELECT (BOLOMETERS, BOL\_TYPE, BOL\_CALIB,
         BOL\_DU3, BOL\_DU4, BOL\_QUAL, BOL\_ENABLED, NUM\_CHAN, NUM\_ADC,
         ARRAY\_CENTRE\_DU3, ARRAY\_CENTRE\_DU4, BOL\_SELECT\_CHAN,
         BOL\_SELECT\_ADC, N\_BOL\_SELECT, MAX\_SUB, SUB\_INSTRMNT, N\_SUB,
         CENTRE\_DU3, CENTRE\_DU4, STATUS)
   }
   \sstarguments{
      \sstsubsection{
         BOLOMETERS                  = CHARACTER$*$($*$) (Given)
      }{
         list of bolometer selections
      }
      \sstsubsection{
         BOL\_TYPE (NUM\_CHAN, NUM\_ADC)
      }{
                               = CHARACTER$*$($*$) (Given)
         type of bolometer
      }
      \sstsubsection{
         BOL\_CALIB (NUM\_CHAN, NUM\_ADC)
      }{
                               = REAL (Given)
         target calibrator values for bolometers
      }
      \sstsubsection{
         BOL\_DU3 (NUM\_CHAN, NUM\_ADC) = REAL (Given)
      }{
         Nasmyth dU3 offset of bolometer from field centre
      }
      \sstsubsection{
         BOL\_DU4 (NUM\_CHAN, NUM\_ADC) = REAL (Given)
      }{
         Nasmyth dU4 offset of bolometer from field centre
      }
      \sstsubsection{
         BOL\_QUAL (NUM\_CHAN, NUM\_ADC)
      }{
                               = INTEGER (Given)
         quality of bolometers
      }
      \sstsubsection{
         BOL\_ENABLED (NUM\_CHAN, NUM\_ADC)
      }{
                               = LOGICAL (Returned)
         .TRUE. if bolometer was selected
      }
      \sstsubsection{
         NUM\_CHAN                    = INTEGER (Given)
      }{
         number of channels per A/D
      }
      \sstsubsection{
         NUM\_ADC                     = INTEGER (Given)
      }{
         number of A/D cards
      }
      \sstsubsection{
         ARRAY\_CENTRE\_DU3            = REAL (Given)
      }{
         the DU3 offset of the centre to be used for either array
      }
      \sstsubsection{
         ARRAY\_CENTRE\_DU4            = REAL (Given)
      }{
         the DU4 offset of the centre to be used for either array
      }
      \sstsubsection{
         BOL\_SELECT\_CHAN (NUM\_CHAN $*$ NUM\_ADC)
      }{
                               = INTEGER (Returned)
         channel numbers of selected bolometers
      }
      \sstsubsection{
         BOL\_SELECT\_ADC (NUM\_CHAN $*$ NUM\_ADC)
      }{
                               = INTEGER (Returned)
         A/D card numbers of selected bolometers
      }
      \sstsubsection{
         N\_BOL\_SELECT                = INTEGER (Returned)
      }{
         total number of bolometers selected
      }
      \sstsubsection{
         MAX\_SUB                     = INTEGER (Given)
      }{
         maximum number of sub-instruments
      }
      \sstsubsection{
         SUB\_INSTRMNT (MAX\_SUB)      = CHARACTER$*$($*$) (Returned)
      }{
         names of instrument sections being used
      }
      \sstsubsection{
         N\_SUB                       = INTEGER (Returned)
      }{
         the number of sub-instruments being used
      }
      \sstsubsection{
         CENTRE\_DU3                  = REAL (Returned)
      }{
         the DU3 of the instrument {\tt '}centre{\tt '}
      }
      \sstsubsection{
         CENTRE\_DU4                  = REAL (Returned)
      }{
         the DU4 of the instrument {\tt '}centre{\tt '}
      }
      \sstsubsection{
         STATUS                      = INTEGER (Given and returned)
      }{
         global status
      }
   }
   \sstnotes{
      The routine will not work properly if the input BOLOMETERS string
      is more than 400 characters long.
   }
   \sstdiytopic{
      Copyright
   }{
      Copyright \copyright 1995,1996,1997,1998,1999 Particle Physics and Astronomy
      Research Council. All Rights Reserved.
   }
}
\sstroutine{
   SCULIB\_CALC\_APPARENT
}{
   calculate apparent RA, Dec of plate centre and angle
   of input coord system N relative to apparent N
}{
   \sstdescription{
      This routine takes the input coordinates and coordinate system of the
      map centre and converts them to the apparent coords at the time of the
      observation. In addition, the angle between the north direction in
      the input coordinate frame and that in the apparent frame is calculated
      (measured anti-clockwise from input north, in radians). See SCU/3.0/JFL/
      0393.

      \sstitemlist{

         \sstitem
         [AZ] coords:
            Calculate the apparent RA and DEC at the LST when the routine is
            called by:-

\begin{equation}
\sin(dec\_app) = \sin(lat\_obs) * \sin(el) + \cos(lat\_obs) * \cos(el) * \cos(az)
\end{equation}

\begin{equation}
\sin(H\_A) = - \frac{\sin(az) * \cos(el)}{\cos(dec\_app)}
\end{equation}

\begin{equation}
\cos(H\_A) = \frac{\sin(el) - \sin(dec\_app) * \sin(lat\_obs)}{\cos(dec\_app) * \cos(lat\_obs)}
\end{equation}

$\cos(dec\_app)$ is present in the denominator of both the
$\sin(H\_A)$ and $\cos(H\_A)$ terms and it could cause both to blow up -- so
it{\tt '}s left out as only the ratio is important

\begin{equation}
RA\_app = LST - H\_A
\end{equation}

\begin{equation}
\sin(rotation) =   \frac{\sin(az) * \cos(lat\_obs)}{\cos(dec\_app)}
\end{equation}

\begin{equation}
\cos(rotation) = \frac{\sin(lat\_obs) - \sin(dec\_app) * \sin(el)}{\cos(dec\_app) * \cos(el)}
\end{equation}

  Again, $\cos(dec\_app)$ appears in the denominator of both $\sin$ and $\cos$
         expressions and is left out of the calculation because on ly the ratio
         is important.

         \sstitem
         [RB] coords:
            Use SLA\_FK54Z to convert to RJ.
            Use SLA\_MAP to convert to apparent, giving ra\_app, dec\_app.
            Use same method to calculate apparent position of RB N pole, giving
            ra\_N\_app, dec\_N\_app.
            Then calculate rotation from:-

\begin{equation}
\mathrm{d}RA =  ra\_app - ra\_N\_app
\end{equation}

\begin{equation}
\sin (rotation) =   \frac{\sin(\mathrm{d}RA) * \cos(dec\_N\_app)}{\cos(lat)}
\end{equation}

\begin{equation}
\cos (rotation) = \frac{\sin(dec\_N\_app) - \sin(dec\_app) * \sin(lat)}{\cos(dec\_app) * \cos(lat)}
\end{equation}

     Since $\cos(lat)$ is in the denominator for both $\sin$ and $\cos$ terms,
         is always $+$ve except for $\pm \pi/2$ where it goes to zero and blows
         up the equations, leave it out in the calculations. The ratio
         of sin and cos will be unaffected.

         If dec\_app = $\pi/2$ (i.e. at N pole of apparent system) then

\begin{equation}
rotation = \pi - (ra\_N\_app - ra\_app)
\end{equation}

         \sstitem
         [RJ] coords:
            Use SLA\_MAP to convert to apparent.
            Use same method to calculate apparent position of RB N pole.
            Derive rotation angle in the same way as for RB.

         \sstitem
         [GA] coords:
            Use SLA\_GALEQ to convert to RJ.
            Use SLA\_MAP to convert to apparent, giving ra\_app, dec\_app.
            Use same method to calculate apparent position of GA N pole, giving
            ra\_N\_app, dec\_N\_app.
            Derive rotation angle in the same way as for RB.

         \sstitem
         [EQ] coords:
            Use SLA\_ECLEQ to convert to RJ.
            Use SLA\_MAP to convert to apparent, giving ra\_app, dec\_app.
            Use same method to calculate apparent position of EQ N pole, giving
            ra\_N\_app, dec\_N\_app.
            Derive rotation angle in the same way as for RB.

         \sstitem
         [HA] coords:
            Apparent RA = LST - LONG

            Apparent Dec = LAT

            Rotation = 0.0D0

         \sstitem
         [RD] coords:
            Apparent RA, Dec set to input values. Rotation = 0.0.

         \sstitem
         [PLANET] coords:
            If MJD1 = MJD2 then apparent RA, Dec set to input LONG, LAT.
               Rotation = 0.0.
            Otherwise apparent RA, Dec interpolated (or extrapolated) between
               LONG, LAT, MJD1 and LONG2, LAT2, MJD2 according to MJD.
               Rotation = 0.
      }
   }
   \sstinvocation{
       CALL SCULIB\_CALC\_APPARENT (LAT\_OBS, LONG, LAT, LONG2, LAT2, MAP\_X, MAP\_Y,
         COORD\_TYPE, LST, MJD, MJD1, MJD2, RA\_APP, DEC\_APP, ROTATION, STATUS)
   }
   \sstarguments{
      \sstsubsection{
         LAT\_OBS                = DOUBLE PRECISION (Given)
      }{
         latitude of observatory (radians)
      }
      \sstsubsection{
         LONG                   = DOUBLE PRECISION (Given)
      }{
         longitude of centre in input coord system (radians)
      }
      \sstsubsection{
         LAT                    = DOUBLE PRECISION (Given)
      }{
         latitude of centre in input coord system (radians)
      }
      \sstsubsection{
         LONG2                  = DOUBLE PRECISION (Given)
      }{
         longitude of second centre in PLANET coord system (radians)
      }
      \sstsubsection{
         LAT2                   = DOUBLE PRECISION (Given)
      }{
         latitude of second centre in PLANET coord system (radians)
      }
      \sstsubsection{
         MAP\_X                  = DOUBLE PRECISION (Given)
      }{
         x tangent plane offset of point from centre (radians)
         The offset must be in the same coordinate as COORD\_TYPE
      }
      \sstsubsection{
         MAP\_Y                  = DOUBLE PRECISION (Given)
      }{
         y tangent plane offset of point from centre (radians)
         The offset must be in the same coordinate as COORD\_TYPE
      }
      \sstsubsection{
         COORD\_TYPE             = CHARACTER$*$($*$) (Given)
      }{
         Coord system of input centre, RD, RB, RJ, GA, EQ, PLANET
      }
      \sstsubsection{
         LST                    = DOUBLE PRECISION (Given)
      }{
         LST for requested coordinates (for AZ and HA)
      }
      \sstsubsection{
         MJD                    = DOUBLE PRECISION (Given)
      }{
         Modified Julian date of observation
      }
      \sstsubsection{
         MJD1                   = DOUBLE PRECISION (Given)
      }{
         Modified Julian date of first centre in PLANET coord system
      }
      \sstsubsection{
         MJD2                   = DOUBLE PRECISION (Given)
      }{
         Modified Julian date of second centre in PLANET coord system
      }
      \sstsubsection{
         RA\_APP                 = DOUBLE PRECISION (Returned)
      }{
         Apparent RA of point at date (radians)
      }
      \sstsubsection{
         DEC\_APP                = DOUBLE PRECISION (Returned)
      }{
         Apparent Dec
      }
      \sstsubsection{
         ROTATION               = DOUBLE PRECISION (Returned)
      }{
         Angle between apparent north and north of input coord system
         (radians, measured clockwise from input north)
      }
      \sstsubsection{
         STATUS                 = INTEGER (Given and returned)
      }{
         Global status
      }
   }
   \sstnotes{
      Does not handle LOCAL\_COORDS for MAP\_X and MAP\_Y
      (see SCULIB\_APARRENT\_2\_MP for information on how to do this)
   }
   \sstdiytopic{
      Copyright
   }{
      Copyright \copyright 1995,1996,1997,1998,1999 Particle Physics and Astronomy
      Research Council. All Rights Reserved.
   }
}
\sstroutine{
   SCULIB\_CALC\_BOL\_COORDS
}{
   Calculate the bolometer offsets in (apparent RA,DEC), AzEl or NA
}{
   \sstdescription{
      This routine calculates the apparent RA and dec of a specified set
      of bolometers. It does this by:-

      \sstitemlist{

         \sstitem
            calculating the offsets that must be added to the bolometer
              positions to cater for the fact that the origin of the bolometer
              coordinate system may be offset from the `centre{\tt '} specified by
              RA\_CENTRE , DEC\_CENTRE.

         \sstitem
            Three types of offsets may be added; Nasmyth offsets which are
              simply subtracted from the bolometer coordinates, azimuth offsets
              which are subtracted from the bolometer positions in az and el,
              or offsets in a coordinate system fixed relative to the sky,
              rotated relative to apparent RA,Dec by the angle ROTATION. The
              latter are added to the bolometer coordinates when they are in
              the form of tangent plane coords in apparent RA,Dec.

         \sstitem
            working out the elevation and parallactic angle for the sidereal
              time and apparent RA, dec of the `centre{\tt '}.

         \sstitem
            calculating the Nasmyth offset of each bolometer, then rotating
              them into tangent plane offsets in azimuth and elevation.

         \sstitem
            adding pointing corrections in azimuth and elevation, then
              rotating the coords into the apparent RA,Dec tangent plane.

         \sstitem
            calling SLA\_DTP2S to work out the apparent RA, dec of the offset
              positions.
      }
   }
   \sstinvocation{
       CALL SCULIB\_CALC\_BOL\_COORS (OUT\_COORDS,RA\_CENTRE, DEC\_CENTRE, LST,
         LAT\_OBS, OFFSET\_COORDS, OFFSET\_X, OFFSET\_Y, ROTATION,
         N\_POINT, MAX\_POINT, POINT\_LST, POINT\_DAZ, POINT\_DEL,
         NUM\_CHAN, NUM\_ADC, N\_BOL, BOL\_CHAN, BOL\_ADC, U3, U4, U3\_CENTRE,
         U4\_CENTRE, X\_BOL, Y\_BOL, ELEVATION, PAR\_ANGLE, STATUS)
   }
   \sstarguments{
      \sstsubsection{
         OUT\_COORDS = CHARACTER $*$ ($*$) (Given)
      }{
         Output coordinate system (NA, AZ, RA)
      }
      \sstsubsection{
         RA\_CENTRE              = DOUBLE PRECISION (Given)
      }{
         the apparent RA of the `centre{\tt '} (radians)
      }
      \sstsubsection{
         DEC\_CENTRE             = DOUBLE PRECISION (Given)
      }{
         the apparent dec of the `centre{\tt '} (radians)
      }
      \sstsubsection{
         LST                    = DOUBLE PRECISION (Given)
      }{
         the local sidereal time (radians)
      }
      \sstsubsection{
         LAT\_OBS                = DOUBLE PRECISION (Given)
      }{
         the latitude of the observatory (radians)
      }
      \sstsubsection{
         OFFSET\_COORDS          = CHARACTER$*$($*$) (Given)
      }{
         the coordinate system of the offset of the array origin from
         the `centre{\tt '}; NA or RD or AZ
      }
      \sstsubsection{
         OFFSET\_X               = REAL (Given)
      }{
         the x offset of the array origin from the `centre{\tt '}
         (arcseconds)
      }
      \sstsubsection{
         OFFSET\_Y               = REAL (Given)
      }{
         the y offset of the array origin from the `centre{\tt '}
         (arcseconds)
      }
      \sstsubsection{
         ROTATION               = DOUBLE PRECISION (Given)
      }{
          for OFFSET\_COORDS other than NA, this gives the angle
         from N in the offset coordinate system to N in apparent
         RA,Dec (radians, increasing clockwise)
      }
      \sstsubsection{
         N\_POINT                = INTEGER (Given)
      }{
         number of elements used in pointing correction arrays
      }
      \sstsubsection{
         MAX\_POINT              = INTEGER (Given)
      }{
         dimension of pointing correction arrays
      }
      \sstsubsection{
         POINT\_LST (MAX\_POINT)  = DOUBLE PRECISION (Given)
      }{
         LST of measured corrections (radians)
      }
      \sstsubsection{
         POINT\_DAZ (MAX\_POINT)  = REAL (Given)
      }{
         correction to be added in azimuth (arcsec)
      }
      \sstsubsection{
         POINT\_DEL (MAX\_POINT)  = REAL (Given)
      }{
         correction to be added in elevation (arcsec)
      }
      \sstsubsection{
         NUM\_CHAN               = INTEGER (Given)
      }{
         the number of channels per A/D card
      }
      \sstsubsection{
         NUM\_ADC                = INTEGER (Given)
      }{
         the number of A/D cards
      }
      \sstsubsection{
         N\_BOL                  = INTEGER (Given)
      }{
         the actual number of bolometers
      }
      \sstsubsection{
         BOL\_CHAN (N\_BOL)       = INTEGER (Given)
      }{
         channel numbers of bolometers
      }
      \sstsubsection{
         BOL\_ADC (N\_BOL)        = INTEGER (Given)
      }{
         ADC numbers of bolometers
      }
      \sstsubsection{
         U3 (NUM\_ADC,NUM\_CHAN)  = REAL (Given)
      }{
         the U3 offsets of the bolometers (arcsec)
      }
      \sstsubsection{
         U4 (NUM\_ADC,NUM\_CHAN)  = REAL (Given)
      }{
         the U4 offsets of the bolometers (arcsec)
      }
      \sstsubsection{
         U3\_CENTRE              = REAL (Given)
      }{
         the U3 offset of the tracking `centre{\tt '} on the array (arcsec)
      }
      \sstsubsection{
         U4\_CENTRE              = REAL (Given)
      }{
         the U4 offset of the tracking `centre{\tt '} on the array (arcsec)
      }
      \sstsubsection{
         X\_BOL (N\_BOL)         = DOUBLE PRECISION (Returned)
      }{
         the X offset (apparent RA or X) of the bolometer (radians)
      }
      \sstsubsection{
         Y\_BOL (N\_BOL)        = DOUBLE PRECISION (Returned)
      }{
         the Y offset (apparent dec or Y) of the bolometer (radians)
      }
      \sstsubsection{
         ELEVATION       = DOUBLE (Returned)
      }{
         elevation of each position
      }
      \sstsubsection{
         PAR\_ANGLE       = DOUBLE (Returned)
      }{
         parallactic angle of each position
      }
      \sstsubsection{
         STATUS                 = INTEGER (Given and returned)
      }{
         The global status
      }
   }
   \sstdiytopic{
      Copyright
   }{
      Copyright \copyright 1995,1996,1997,1998,1999 Particle Physics and Astronomy
      Research Council. All Rights Reserved.
   }
}
\sstroutine{
   SCULIB\_CALC\_CLOCKERR
}{
   Calculate start up time and error in system clock
}{
   \sstdescription{
      Checks the self-consistency of the FITS headers by calculating
      the expected LST from the Azimuth and Elevation values stored
      in the header and comparing this with the supplied reference
      LST. Returns the difference between the two values and the
      LST calculated from the headers.
   }
   \sstinvocation{
       CALL SCULIB\_CALC\_CLOCKERR( N\_FITS, FITS, RA\_CEN, LST\_REF,
          CLOCK\_ERR, LST\_AZEL, STATUS )
   }
   \sstarguments{
      \sstsubsection{
         N\_FITS = INTEGER (Given)
      }{
         Number of items in the FITS array
      }
      \sstsubsection{
         FITS() = CHAR$*$80 (Given)
      }{
         FITS array
      }
      \sstsubsection{
         RA\_CEN = DOUBLE (Given)
      }{
         Apparent RA corresponding to the Az/El defined in the FITS header.
         In radians.
      }
      \sstsubsection{
         LST\_REF = DOUBLE (Given)
      }{
         Local sidereal time at the start of data acquisition. This
         usually comes from the LST\_STRT array. In radians.
      }
      \sstsubsection{
         CLOCK\_ERR = DOUBLE (Returned)
      }{
         The time difference (in radians) between the times stored
         in the data file (header items and LST\_STRT) and the actual
         time of the observation derived from the azimuth and elevation
         of the observed source. The value is the number that must be
         added to the supplied LST\_REF in order to correct it.
      }
      \sstsubsection{
         LST\_AZEL = DOUBLE (Returned)
      }{
         Local Sidereal time calculated from the azimuth, elevation
         and apparent RA/Dec (radians)
      }
      \sstsubsection{
         STATUS = INTEGER (Given \& Returned)
      }{
         Global status. An error will be returned if the observation
         was not done in a tracking frame (eg NA or AZ) since those
         frames can not be used to determine the time from the telescope
         az and el.
      }
   }
   \sstnotes{
      Requires the following FITS headers to be defined:
      \sstitemlist{

         \sstitem
         LAT-OBS

         \sstitem
         STRT\_ELD

         \sstitem
         STRT\_AZD

         \sstitem
         CENT\_CRD

      }
      The clock error returned must be added to all times read from
      the data file in order to make those times correct.

      The centre coordinates of the observation must corresponde to
      an astronomical frame. AZ can not be used (for obvious reasons).
   }
   \sstdiytopic{
      Copyright
   }{
      Copyright (C) 2000 Particle Physics and Astronomy
      Research Council. All Rights Reserved.
   }
}
\sstroutine{
   SCULIB\_CALC\_FLATFIELD
}{
   calculate flat-field results for a bolometer
}{
   \sstdescription{
      This routine calculates the image parameters from data taken as a
      flat-field measurement of a SCUBA bolometer. It takes as input data a
      map made of a point source with pixels on a square grid.

         If status is good on entry, the routine will set default return values
      to VAL\_\_BADR except QUALITY set to 0. As progress is made through the
      routine the various image quantities will be set to their derived values,
      and QUALITY will be set to 1 if any problem is encountered.

         An attempt will then be made to estimate the 0 level of the image by
      averaging valid data points in the corners of the map area. If there are
      any such points then the 0 level will be subtracted from the map.

         Next the routine will calculate the 0th and 1st order moments of the
      image on the map. Map pixels with bad quality are ignored.

  The 0th order is calculated as
\begin{equation}
U_0 = \sum_{ij} f_{ij},
\end{equation}
where i,j are the pixel indices. The 1st order moment in X is calculated as
\begin{equation}
U_x = \sum_{ij} x f_{ij},
\end{equation}
where x is the x position of pixel i,j; and the 1st order moment in Y is
calculated as
\begin{equation}
U_y = \sum_{ij} y f_{ij},
\end{equation}
where y is the y position of pixel i,j.

      If the 0th order moment of the image was 0, i.e. there is no image on
      the map, then a warning message will be output and the routine will
      return with good status and QUALITY set to 1. Otherwise, the x,y centre
      of the image is calculated from:-

\begin{equation}
X_\mathrm{cen} = \frac{U_x}{U_0}
\end{equation}
\begin{equation}
Y_\mathrm{cen} = \frac{U_y}{U_0}
\end{equation}

      With this information the image can be analysed further by one of
      2 methods. The method used depends on the value of parameter
      FLAT\_ANALYSIS. If there is an error reading this parameter, or
      its value is not either MOMENTS or FIT, then a warning message will
      be issued and a value of MOMENTS will be assumed.

      For a MOMENTS analysis the 2nd order moments can be calculated about the
      centre of the image:-

\begin{equation}
          U_{xx} = \sum_{ij} (x - X_{cen})^2 f_{ij}
\end{equation}
\begin{equation}
          U_{xy} = \sum_{ij} (x - X_{cen})(y - Y\_{cen}) f_{ij}
\end{equation}
\begin{equation}
          U_{yy} = \sum_{ij} (y - Y_{cen})^2  f_{ij}
\end{equation}

      The parameters of the weighted ellipse describing the data are then
      (following `Analysis of Astronomical Images using Moments{\tt '}, R.Stobie,
      J.B.I.S., 33, 323):-

\begin{equation}
\tan(2\theta) = \frac{2 U_{xy}}{U_{xx} - U_{yy}}
\end{equation}
\begin{equation}
A^2 = 2 ( U_{xx} + U_{yy} ) + 2 \sqrt{(U_{xx} - U_{yy})^2 + 4U_{xy}^2}
\end{equation}
\begin{equation}
B^2 = 2 ( U_{xx} + U_{yy} ) - 2 \sqrt{(U_{xx} - U_{yy})^2 + 4U_{xy}^2}
\end{equation}

      For a FIT analysis, a 2-d Gaussian is fitted by a least-squares
      routine to the data. SCULIB\_FIT\_ROUTINE is the routine that performs
      the fit while SCULIB\_GAUSSIAN\_XISQ calculates the chi-squared of the
      fit to the data. An error will be reported and QUALITY set to 1
      if there are any problems with the fit process. The fitted function
      is :-

\begin{equation}
  P \exp( \frac{-(x-X_{cen})^2 - (y-Y_{cen})^2}{\sigma^2})
\end{equation}

      where SIGMA ($\sigma$) is an ellipse with semi-axis lengths A and B, with
THETA ($\theta$)
      the angle between the x axis and A (THETA in radians, measured
      anti-clockwise).

      If all is well still, the volume under the image and the variance on it
      are calculated. The volume will be the sum of all map pixels under a
      circle of 12 arcsec radius for short-wave array bolometers, or 25
      arcsec for other types. If this area laps over the edges of the map, or
      any pixels inside it have bad quality then a warning message will be
      output but QUALITY will stay good.
   }
   \sstinvocation{
       CALL SCULIB\_CALC\_FLATFIELD (BOLNAME, BOL\_TYPE, IDIM, JDIM, MAP\_DATA,
         MAP\_VARIANCE, MAP\_QUALITY, X, Y, VOLUME, VOLUME\_VAR, X\_CENTRE,
         Y\_CENTRE, THETA, A, B, QUALITY, STATUS)
   }
   \sstarguments{
      \sstsubsection{
         BOLNAME                           = CHARACTER$*$($*$) (Given)
      }{
         the name of the bolometer being measured
      }
      \sstsubsection{
         BOL\_TYPE                          = CHARACTER$*$($*$) (Given)
      }{
         the type of the bolometer being measured
      }
      \sstsubsection{
         IDIM                              = INTEGER (Given)
      }{
         `I{\tt '} dimension of map
      }
      \sstsubsection{
         JDIM                              = INTEGER (Given)
      }{
         `J{\tt '} dimension of map
      }
      \sstsubsection{
         MAP\_DATA (IDIM, JDIM)             = REAL (Given)
      }{
         map data values
      }
      \sstsubsection{
         MAP\_VARIANCE (IDIM, JDIM)         = REAL (Given)
      }{
         map variance
      }
      \sstsubsection{
         MAP\_QUALITY (IDIM, JDIM)          = REAL (Given)
      }{
         map quality
      }
      \sstsubsection{
         X (IDIM)                          = REAL (Given)
      }{
         x axis of map
      }
      \sstsubsection{
         Y (JDIM)                          = REAL (Given)
      }{
         y axis of map
      }
      \sstsubsection{
         VOLUME                            = REAL (Returned)
      }{
         volume under image within MULT times the fitted ellipse
      }
      \sstsubsection{
         VOLUME\_VAR                        = REAL (Returned)
      }{
         variance on VOLUME\_DATA
      }
      \sstsubsection{
         X\_CENTRE                          = REAL (Returned)
      }{
         x of image centre
      }
      \sstsubsection{
         Y\_CENTRE                          = REAL (Returned)
      }{
         y of image centre
      }
      \sstsubsection{
         THETA                             = REAL (Returned)
      }{
         angle of tilt of ellipse (radians)
      }
      \sstsubsection{
         A                                 = REAL (Returned)
      }{
         semi-major axis of ellipse
      }
      \sstsubsection{
         B                                 = REAL (Returned)
      }{
         semi-minor axis of ellipse
      }
      \sstsubsection{
         STATUS                            = INTEGER (Given and returned)
      }{
         global status
      }
   }
   \sstdiytopic{
      Copyright
   }{
      Copyright \copyright 1995,1996,1997,1998,1999 Particle Physics and Astronomy
      Research Council. All Rights Reserved.
   }
}
\sstroutine{
   SCULIB\_CALC\_GRID
}{
   calculate the minimum rectangular grid that would
   contain the input jiggle pattern
}{
   \sstdescription{
      This routine takes the x,y coords of the mapped points and attempts to
      fit them onto a rectangular grid. The grid contains the minimum number
      of points required to hold the measured positions.
   }
   \sstinvocation{
       CALL SCULIB\_CALC\_GRID (N, X, Y, XMIN, XMAX, XSPACE, NX,
         YMIN, YMAX, YSPACE, NY, IPOS, JPOS, STATUS)
   }
   \sstarguments{
      \sstsubsection{
         N                                     = INTEGER (Given)
      }{
         the number of positions in the jiggle pattern
      }
      \sstsubsection{
         X (N)                                 = REAL (Given)
      }{
         the x offsets of the jiggle
      }
      \sstsubsection{
         Y (N)                                 = REAL (Given)
      }{
         the y offsets
      }
      \sstsubsection{
         XMIN                                  = REAL (Returned)
      }{
         the minimum of the output map{\tt '}s x-axis
      }
      \sstsubsection{
         XMAX                                  = REAL (Returned)
      }{
         the maximum of the x-axis
      }
      \sstsubsection{
         XSPACE                                = REAL (Returned)
      }{
         the pixel spacing in x
      }
      \sstsubsection{
         NX                                    = INTEGER (Returned)
      }{
         the map dimension in x
      }
      \sstsubsection{
         YMIN                                  = REAL (Returned)
      }{
         the minimum of the output map{\tt '}s y-axis
      }
      \sstsubsection{
         YMAX                                  = REAL (Returned)
      }{
         the maximum of the y-axis
      }
      \sstsubsection{
         YSPACE                                = REAL (Returned)
      }{
         the pixel spacing in y
      }
      \sstsubsection{
         NY                                    = INTEGER (Returned)
      }{
         the map dimension in y
      }
      \sstsubsection{
         IPOS (N)                              = INTEGER (Returned)
      }{
         the I index of each jiggle position in the map
      }
      \sstsubsection{
         JPOS (N)                              = INTEGER (Returned)
      }{
         the J index of each jiggle position in the map
      }
      \sstsubsection{
         STATUS                                = INTEGER (Given and returned)
      }{
         global status
      }
   }
   \sstdiytopic{
      Copyright
   }{
      Copyright \copyright 1995,1996,1997,1998,1999 Particle Physics and Astronomy
      Research Council. All Rights Reserved.
   }
}
\sstroutine{
   SCULIB\_CALC\_OUTPUT\_COORDS
}{
   calculate output coords of map centre and angle
   of output coord system N relative to apparent N
}{
   \sstdescription{
      This routine takes the apparent centre coords at the time of the
      observation and converts them to the output coordinate system. In
      addition, the angle between the north direction in the output
      coordinate frame and that in the apparent frame is calculated
      (measured anticlockwise from output north, in radians).
   }
   \sstinvocation{
       CALL SCULIB\_CALC\_OUTPUT\_COORDS (RA\_APP, DEC\_APP, MJD,
         OUTPUT\_COORDS, LONG, LAT, STATUS)
   }
   \sstarguments{
      \sstsubsection{
         RA\_APP                 = DOUBLE PRECISION (Given)
      }{
         apparent RA of map centre on MJD (radians)
      }
      \sstsubsection{
         DEC\_APP                = DOUBLE PRECISION (Given)
      }{
         apparent declination of map centre on MJD (radians)
      }
      \sstsubsection{
         MJD                    = DOUBLE PRECISION (Given)
      }{
         U1 of observation expressed as modified Julian day
      }
      \sstsubsection{
         OUTPUT\_COORDS          = CHARACTER$*$($*$) (Given)
      }{
         output coordinated system; RB, RJ, RD, GA or EQ
      }
      \sstsubsection{
         LONG                   = DOUBLE PRECISION (Returned)
      }{
         longitude of map centre in output coord system (radians)
      }
      \sstsubsection{
         LAT                    = DOUBLE PRECISION (Returned)
      }{
         latitude of map centre in output coord system (radians)
      }
      \sstsubsection{
         STATUS                 = INTEGER (Given and returned)
      }{
         global status
      }
   }
   \sstnotes{
      SLA routines are used to perform the coordinate conversions.
   }
   \sstdiytopic{
      Copyright
   }{
      Copyright \copyright 1995,1996,1997,1998,1999 Particle Physics and Astronomy
      Research Council. All Rights Reserved.
   }
}
\sstroutine{
   SCULIB\_CALC\_SKYDIP\_TEMPS
}{
   Calculate all the skydip temps from raw data
}{
   \sstdescription{
      This routine accepts the SKYDIP (3 dimensional) raw data and
      works out the temperature associated with each integration.
      The 2 dimensional array (bolometers, integrations) is then returned.
   }
   \sstinvocation{
       CALL SCULIB\_CALC\_SKYDIP\_TEMPS( N\_TEMPS, N\_FITS, FITS, PARAM,
            N\_INTEGRATIONS, N\_MEASUREMENTS, N\_POS, N\_BOLS, DEM\_PNTR,
            IN\_DATA, OUT\_DATA, OUT\_VAR, OUT\_QUAL, OUT\_DEM\_PNTR,
            OUT\_DATA\_AV, OUT\_VAR\_AV, OUT\_QUAL\_AV, SCRATCH,
            STATUS)
   }
   \sstarguments{
      \sstsubsection{
         N\_TEMPS = INTEGER (Given)
      }{
         Number of SKYDIP temperatures.
      }
      \sstsubsection{
         SUB\_REQUIRED = INTEGER (Given)
      }{
         The number of the sub-instrument that is requested.
         If this number is less than 1 or greater than N\_BOLS then
         all subinstrumentes are returned. This should be 0 when called
         from REDUCE\_SWITCH and the required number if called from
         SKYDIP. Exists purely to prevent multiple requests for T\_COLD
         if we are only interested in one wavelength
      }
      \sstsubsection{
         N\_FITS = INTEGER (Given)
      }{
         Number of FITS items
      }
      \sstsubsection{
         FITS = CHARACTER$*$80 (Given)
      }{
         The FITS items
      }
      \sstsubsection{
         PARAM = CHARACTER$*$(PAR\_\_SZNAM) (Given)
      }{
         Name of parameter from which T\_COLD can be read
      }
      \sstsubsection{
         N\_INTEGRATIONS = INTEGER (Given)
      }{
         Number of integrations
      }
      \sstsubsection{
         N\_MEASUREMENTS = INTEGER (Given)
      }{
         Number of measurements
      }
      \sstsubsection{
         N\_POS = INTEGER (Given)
      }{
         Number of data points
      }
      \sstsubsection{
         N\_BOLS = INTEGER (Given)
      }{
         Number of bolometers (sub instruments)
      }
      \sstsubsection{
         NUM\_CHAN                   = INTEGER (Given)
      }{
         number of A/D channels per A/D card
      }
      \sstsubsection{
         NUM\_ADC                    = INTEGER (Given)
      }{
         number of A/D cards
      }
      \sstsubsection{
         BOL\_TYPE (NUM\_CHAN, NUM\_ADC)
      }{
                              = CHARACTER$*$($*$) (Given)
         the types of the bolometers
      }
      \sstsubsection{
         BOL\_CHAN (N\_BOLS)          = INTEGER (Given)
      }{
         channel numbers of selected bolometers
      }
      \sstsubsection{
         BOL\_ADC (N\_BOLS)           = INTEGER (Given)
      }{
         ADC numbers of selected bolometers
      }
      \sstsubsection{
         DEM\_PNTR(1, 1, N\_INTEGRATIONS, N\_MEASUREMENTS) = INTEGER (Given)
      }{
         Pointer to positions in the input data
      }
      \sstsubsection{
         IN\_DATA(N\_TEMPS, N\_BOLS, N\_POS) = REAL (Given)
      }{
         Input data [NTEMPS, NBOLS, NPOS]
      }
      \sstsubsection{
         OUT\_DATA(N\_BOLS, N\_POS) = REAL (Returned)
      }{
         Processed sky brightness temperatures.
      }
      \sstsubsection{
         OUT\_VAR(N\_BOLS, N\_POS) = REAL (Returned)
      }{
         Variance on each point in OUT\_DATA
      }
      \sstsubsection{
         OUT\_QUAL(N\_BOLS, N\_POS) = BYTE (Returned)
      }{
         Quality of OUT\_DATA
      }
      \sstsubsection{
         OUT\_DEM\_PNTR(1, 1, N\_INTEGRATIONS, N\_MEASUREMENTS) = INTEGER (Returned)
      }{
         Modified DEM\_PNTR to reflect changes in data array
      }
      \sstsubsection{
         OUT\_DATA\_AV(N\_MEASUREMENTS, N\_BOLS) = REAL (Returned)
      }{
         Average sky temp for each measurement
      }
      \sstsubsection{
         OUT\_VAR\_AV(N\_MEASUREMENTS, N\_BOLS) = REAL (Returned)
      }{
         Variance of the Average sky temp for each measurement
      }
      \sstsubsection{
         OUT\_QUAL\_AV(N\_MEASUREMENTS, N\_BOLS) = BYTE (Returned)
      }{
         Quality of the Average sky temp for each measurement
      }
      \sstsubsection{
         SCRATCH(N\_TEMPS, N\_BOLS, N\_INTEGRATIONS) = REAL (Given)
      }{
         Work space to store the data for each measurement (since I
         cant simply pass a pointer in as the data is non-contiguous)
      }
      \sstsubsection{
         STATUS = INTEGER (Given)
      }{
         Global Status
      }
   }
   \sstdiytopic{
      Copyright
   }{
      Copyright \copyright 1995,1996,1997,1998,1999 Particle Physics and Astronomy
      Research Council. All Rights Reserved.
   }
}
\sstroutine{
   SCULIB\_CALC\_SUB\_BOLS
}{
   find the positions of bolometers belonging to a
   specified sub-instrument in a demodulated data array
}{
   \sstdescription{
      This subroutine finds the location of data from bolometers belonging to
      a specified sub-instrument in a demodulated data array that may contain
      data for several sub-instruments. The number of bolometers found, their
      indices in the data array, and their ADC,channel numbers are returned
      by the routine.
   }
   \sstinvocation{
       CALL  SCULIB\_CALC\_SUB\_BOLS (N\_BOL\_IN, IN\_BOL\_ADC, IN\_BOL\_CHAN,
         NUM\_CHAN, NUM\_ADC, BOL\_TYPE, SUB\_INSTRUMENT, N\_BOL\_OUT,
         OUT\_BOL\_ADC, OUT\_BOL\_CHAN, IN\_POINTER, STATUS)
   }
   \sstarguments{
      \sstsubsection{
         N\_BOL\_IN                     = INTEGER (Given)
      }{
         number of bolometers in data array
      }
      \sstsubsection{
         IN\_BOL\_ADC (N\_BOL\_IN)        = INTEGER (Given)
      }{
         ADC numbers of bolometers in data array
      }
      \sstsubsection{
         IN\_BOL\_CHAN (N\_BOL\_IN)       = INTEGER (Given)
      }{
         channel numbers of bolometers in data array
      }
      \sstsubsection{
         NUM\_CHAN                     = INTEGER (Given)
      }{
         number of channels per A/D card
      }
      \sstsubsection{
         NUM\_ADC                      = INTEGER (Given)
      }{
         number of A/D cards
      }
      \sstsubsection{
         BOL\_TYPE (NUM\_CHAN,NUM\_ADC)  = CHARACTER$*$($*$) (Given)
      }{
         types of bolometers
      }
      \sstsubsection{
         SUB\_INSTRUMENT               = CHARACTER$*$($*$) (Given)
      }{
         the name of the sub-instrument for which data are to be extracted
      }
      \sstsubsection{
         N\_BOL\_OUT                    = INTEGER (Returned)
      }{
         the number of bolometers in the data array that belong to the
         specified sub-instrument
      }
      \sstsubsection{
         OUT\_BOL\_ADC (N\_BOL\_IN)       = INTEGER (Returned)
      }{
         ADC numbers of bolometers in specified sub-instrument
      }
      \sstsubsection{
         OUT\_BOL\_CHAN (N\_BOL\_IN)      = INTEGER (Returned)
      }{
         channel numbers of bolometers in specified sub-instrument
      }
      \sstsubsection{
         IN\_POINTER (N\_BOL\_IN)        = INTEGER (Returned)
      }{
         array pointing to indices in data array that contain
         data for bolometers in the specified sub-instrument
      }
      \sstsubsection{
         STATUS                       = INTEGER (Given and returned)
      }{
         global status
      }
   }
   \sstdiytopic{
      Copyright
   }{
      Copyright \copyright 1995,1996,1997,1998,1999 Particle Physics and Astronomy
      Research Council. All Rights Reserved.
   }
}
\sstroutine{
   SCULIB\_CFILLB
}{
   fill a byte array with a constant
}{
   \sstdescription{
      fills an byte array with a constant
   }
   \sstinvocation{
      CALL SCULIB\_CFILLB (N, IVAL, ARRAY)
   }
   \sstarguments{
      \sstsubsection{
         N                  = INTEGER (Given)
      }{
         number of array elements
      }
      \sstsubsection{
         IVAL               = BYTE (Given)
      }{
         constant to which array is to be set
      }
      \sstsubsection{
         ARRAY (N)          = BYTE (Returned)
      }{
         array to be set
      }
   }
   \sstdiytopic{
      Copyright
   }{
      Copyright \copyright 1995,1996,1997,1998,1999 Particle Physics and Astronomy
      Research Council. All Rights Reserved.
   }
   \sstimplementationstatus{
      \sstitemlist{

         \sstitem
         No status checking
      }
   }
}
\sstroutine{
   SCULIB\_CFILLD
}{
   fill a double precision array with a constant
}{
   \sstdescription{
      fills a double precision array with a constant
   }
   \sstinvocation{
      CALL SCULIB\_CFILLD (N, DVAL, ARRAY)
   }
   \sstarguments{
      \sstsubsection{
         N                  = INTEGER (Given)
      }{
         number of array elements
      }
      \sstsubsection{
         DVAL               = DOUBLE PRECISION (Given)
      }{
         constant to which array is to be set
      }
      \sstsubsection{
         ARRAY (N)          = DOUBLE PRECISION (Returned)
      }{
         array to be set
      }
   }
   \sstdiytopic{
      Copyright
   }{
      Copyright \copyright 1995,1996,1997,1998,1999 Particle Physics and Astronomy
      Research Council. All Rights Reserved.
   }
   \sstimplementationstatus{
      \sstitemlist{

         \sstitem
         No status checking
      }
   }
}
\sstroutine{
   SCULIB\_CFILLI
}{
   fill an integer array with a constant
}{
   \sstdescription{
      fills an integer array with a constant
   }
   \sstinvocation{
      CALL SCULIB\_CFILLI (N, IVAL, ARRAY)
   }
   \sstarguments{
      \sstsubsection{
         N                  = INTEGER (Given)
      }{
         number of array elements
      }
      \sstsubsection{
         IVAL               = INTEGER (Given)
      }{
         constant to which array is to be set
      }
      \sstsubsection{
         ARRAY (N)          = INTEGER (Returned)
      }{
         array to be set
      }
   }
   \sstdiytopic{
      Copyright
   }{
      Copyright \copyright 1995,1996,1997,1998,1999 Particle Physics and Astronomy
      Research Council. All Rights Reserved.
   }
   \sstimplementationstatus{
      \sstitemlist{

         \sstitem
         No status checking
      }
   }
}
\sstroutine{
   SCULIB\_CFILLR
}{
   fill a real array with a constant
}{
   \sstdescription{
      fills a real array with a constant
   }
   \sstinvocation{
      CALL SCULIB\_CFILLR (N, RVAL, ARRAY)
   }
   \sstarguments{
      \sstsubsection{
         N                  = INTEGER (Given)
      }{
         number of array elements
      }
      \sstsubsection{
         RVAL               = REAL (Given)
      }{
         constant to which array is to be set
      }
      \sstsubsection{
         ARRAY (N)          = REAL (Returned)
      }{
         array to be set
      }
   }
   \sstdiytopic{
      Copyright
   }{
      Copyright \copyright 1995,1996,1997,1998,1999 Particle Physics and Astronomy
      Research Council. All Rights Reserved.
   }
   \sstimplementationstatus{
      \sstitemlist{

         \sstitem
         No status checking
      }
   }
}
\sstroutine{
   SCULIB\_CLIP\_BOL
}{
   To clip a bolometer at $+$/- nsigma
}{
   \sstdescription{
      This routine despikes a single bolometer at the $+$/- n sigma
      level. A 1-dimensional data set is assumed.
   }
   \sstinvocation{
       CALL SCULIB\_CLIP\_BOL(N\_POS, SCUDATA, SCUQUAL,
            N\_SIGMA, BADBIT, NSPIKES, STATUS)
   }
   \sstarguments{
      \sstsubsection{
         N\_POS = INTEGER (Given)
      }{
         Number of jiggle positions in data set
      }
      \sstsubsection{
         SCUDATA(N\_POS) = REAL (Given)
      }{
         The data
      }
      \sstsubsection{
         SCUQUAL(N\_POS) = BYTE (Given \& Returned)
      }{
         The data quality
      }
      \sstsubsection{
         N\_SIGMA = DOUBLE (Given)
      }{
         Number of sigma to clip at. If this number is positive
         then a clipped mean and standard deviation are calculated.
         If it is negative the statistics are calculated without
         iterative clipping.
      }
      \sstsubsection{
         BADBIT = BYTE (Given \& Returned)
      }{
         Bad bit mask
      }
      \sstsubsection{
         NSPIKES = INTEGER (Returned)
      }{
         Number of spikes that were detected.
      }
      \sstsubsection{
         STATUS = INTEGER (Given and Returned)
      }{
         Global Status value
      }
   }
   \sstdiytopic{
      Copyright
   }{
      Copyright \copyright 1995,1996,1997,1998,1999 Particle Physics and Astronomy
      Research Council. All Rights Reserved.
   }
}
\sstroutine{
   SCULIB\_COADD
}{
   coadd exposure into coadded result
}{
   \sstdescription{
      This routine coadds the current exposure to the input coadd arrays
      and puts the result in the output coadd arrays. The input and output
      arrays can be the same. The coadd result is the average of exposures
      that have been input. The coadd variance is calculated from
      the spread of the input exposures about the mean if more than
      one exposure has been added in. If only one exposure is
      available then the coadd variance is set equal to variance on the
      input data if present, otherwise it is set equal to zero.
      Input pixels with bad quality are ignored.
   }
   \sstinvocation{
       CALL SCULIB\_COADD (N, IN\_DATA, IN\_VARIANCE, IN\_QUALITY,
          INCOADD\_DATA, INCOADD\_VAR, INCOADD\_QUAL, INCOADD\_NUMBER,
          OUTCOADD\_DATA, OUTCOADD\_VAR, OUTCOADD\_QUAL, OUTCOADD\_NUMBER,
          BADBIT, VARIANCE, STATUS)
   }
   \sstarguments{
      \sstsubsection{
         N                       = INTEGER (Given)
      }{
         Number of elements in arrays.
      }
      \sstsubsection{
         IN\_DATA (N)             = REAL (Given)
      }{
         data to be added to coadd
      }
      \sstsubsection{
         IN\_VARIANCE (N)         = REAL (Given)
      }{
         variance on input data
      }
      \sstsubsection{
         IN\_QUALITY (N)          = BYTE (Given)
      }{
         quality on input data
      }
      \sstsubsection{
         INCOADD\_DATA (N)        = REAL (Given)
      }{
         Input coadd.
      }
      \sstsubsection{
         INCOADD\_VAR (N)         = REAL (Given)
      }{
         Input coadd variance.
      }
      \sstsubsection{
         INCOADD\_QUAL (N)       = BYTE (returned)
      }{
         Input coadd quality.
      }
      \sstsubsection{
         INCOADD\_NUMBER (N)      = INTEGER (Given)
      }{
         Number of exposures coadded in input coadd.
      }
      \sstsubsection{
         OUTCOADD\_DATA (N)       = REAL (Returned)
      }{
         Output coadd.
      }
      \sstsubsection{
         OUTCOADD\_VAR (N)        = REAL (Returned)
      }{
         Output coadd variance.
      }
      \sstsubsection{
         OUTCOADD\_QUAL (N)       = BYTE (returned)
      }{
         Output coadd quality.
      }
      \sstsubsection{
         OUTCOADD\_NUMBER (N)     = INTEGER (Returned)
      }{
         Number of exposures coadded in output coadd.
      }
      \sstsubsection{
         BADBIT                  = BYTE (Given)
      }{
         Bad bit mask
      }
      \sstsubsection{
         VARIANCE                = LOGICAL (Given)
      }{
         T if input data has variance associated with it
      }
      \sstsubsection{
         STATUS                  = INTEGER (Given \& Returned)
      }{
         Global status
      }
   }
   \sstdiytopic{
      Copyright
   }{
      Copyright \copyright 1995,1996,1997,1998,1999 Particle Physics and Astronomy
      Research Council. All Rights Reserved.
   }
}
\sstroutine{
   SCULIB\_COADD\_MAPS
}{
   Average together integration and output map
}{
   \sstdescription{
      This routine coadds the current integration image to the output image.
      The coadd variance is calculated from the spread of the integrations
      about the mean if more than one integration has been added in. The
      averaging is done over integrations, so that each point in an image
      is separate. An output point is good if either of the inputs is good.
      If only one dataset is coadded th variance is set zero.
      Input pixels with bad quality are ignored.
   }
   \sstinvocation{
       CALL SCULIB\_COADD\_MAPS(NPTS, DATA\_IN, QUALITY\_IN, WEIGHT,
            NCOADD, DATA\_OUT, VARIANCE\_OUT, QUALITY\_OUT, WEIGHT\_OUT,
            STATUS)
   }
   \sstarguments{
      \sstsubsection{
         NPTS = INTEGER (Given)
      }{
         Number of input data points
      }
      \sstsubsection{
         DATA\_IN( NPTS ) = REAL (Given)
      }{
         Input data
      }
      \sstsubsection{
         QUALITY\_IN( NPTS ) = BYTE (Given)
      }{
         Input data quality
      }
      \sstsubsection{
         WEIGHT = REAL (Given)
      }{
         Relative weight of input data
      }
      \sstsubsection{
         NCOADD( NPTS ) = INTEGER (Given \& Returned)
      }{
         Number of points coadded into each data point
      }
      \sstsubsection{
         DATA\_OUT( NPTS ) = REAL (Given \& Returned)
      }{
         Coadded output data
      }
      \sstsubsection{
         VARIANCE\_OUT( NPTS ) = REAL (Given \& Returned)
      }{
         Variance of coadded output data.
      }
      \sstsubsection{
         QUALITY\_OUT( NPTS ) = BYTE (Given \& Returned)
      }{
         Quality of coadded output data
      }
      \sstsubsection{
         WEIGHT\_OUT( NPTS ) = REAL (Given \& Returned)
      }{
         Total weight of output data
      }
      \sstsubsection{
         STATUS = INTEGER (Given \& Returned)
      }{
         Global status
      }
   }
   \sstdiytopic{
      Copyright
   }{
      Copyright \copyright 1995,1996,1997,1998,1999 Particle Physics and Astronomy
      Research Council. All Rights Reserved.
   }
}
\sstroutine{
   SCULIB\_COADD\_REMOVE
}{
   remove exposure from coadded result
}{
   \sstdescription{
      This routine removes an exposure from the input coadd and puts the
      result in the output coadd arrays. The input and output
      arrays can be the same. The coadd result is the average of exposures
      remaining in the coadd. The coadd variance is calculated from
      the spread of the exposures about the mean if more than one exposure
      remains, 0 otherwise. Exposure pixels with bad quality are ignored.
   }
   \sstinvocation{
       CALL SCULIB\_COADD\_REMOVE (N, IN\_DATA, IN\_VARIANCE, IN\_QUALITY,
         INCOADD\_DATA, INCOADD\_VAR, INCOADD\_QUAL, INCOADD\_NUMBER,
         OUTCOADD\_DATA, OUTCOADD\_VAR, OUTCOADD\_QUAL, OUTCOADD\_NUMBER,
         VARIANCE)
   }
   \sstarguments{
      \sstsubsection{
         N                       = INTEGER (Given)
      }{
         Number of elements in arrays.
      }
      \sstsubsection{
         IN\_DATA (N)             = REAL (Given)
      }{
         data to be added to coadd
      }
      \sstsubsection{
         IN\_VARIANCE (N)         = REAL (Given)
      }{
         variance on input data
      }
      \sstsubsection{
         IN\_QUALITY (N)          = INTEGER (Given)
      }{
         quality on input data
      }
      \sstsubsection{
         INCOADD\_DATA (N)        = REAL (Given)
      }{
         Input coadd.
      }
      \sstsubsection{
         INCOADD\_VAR (N)         = REAL (Given)
      }{
         Input coadd variance.
      }
      \sstsubsection{
         INCOADD\_QUAL (N)        = INTEGER (Given)
      }{
         Input coadd quality
      }
      \sstsubsection{
         INCOADD\_NUMBER (N)      = INTEGER (Given)
      }{
         Number of exposures coadded in input coadd.
      }
      \sstsubsection{
         OUTCOADD\_DATA (N)       = REAL (Returned)
      }{
         Output coadd.
      }
      \sstsubsection{
         OUTCOADD\_VAR (N)        = REAL (Returned)
      }{
         Output coadd variance.
      }
      \sstsubsection{
         OUTCOADD\_QUAL (N)       = INTEGER (returned)
      }{
         Output coadd quality.
      }
      \sstsubsection{
         OUTCOADD\_NUMBER (N)     = INTEGER (Returned)
      }{
         Number of exposures coadded in output coadd.
      }
      \sstsubsection{
         VARIANCE                = LOGICAL (Given)
      }{
         T if input data has variance associated with it
      }
   }
   \sstdiytopic{
      Copyright
   }{
      Copyright \copyright 1995,1996,1997,1998,1999 Particle Physics and Astronomy
      Research Council. All Rights Reserved.
   }
}
\sstroutine{
   SCULIB\_COMPRESS\_DEMOD
}{
   get demodulated data for a bolometer and coadd
   jiggles in each integration
}{
   \sstdescription{
      This routine selects data for a specified bolometer from a demodulated
      data array, and coadds the results for each jiggle in each integration
      to give an average for that integration. Data with bad quality are
      ignored. The variance on the average will also be derived; set equal
      to the variance on the input data if only one jiggle contributes to the
      average, otherwise calculated from the spread of the input points about
      the mean.
   }
   \sstinvocation{
       CALL SCULIB\_COMPRESS\_DEMOD (N\_BOLS, N\_JIGS, N\_INTS,
         IN\_DATA, ADC\_INDEX, CHAN\_INDEX, INT\_INDEX, ADC, CHAN,
         OUT\_DATA, OUT\_VARIANCE, OUT\_QUALITY, STATUS)
   }
   \sstarguments{
      \sstsubsection{
         N\_BOLS                 = INTEGER (Given)
      }{
         number of bolometers measured
      }
      \sstsubsection{
         N\_JIGS                 = INTEGER (Given)
      }{
         number of jiggles in pattern
      }
      \sstsubsection{
         N\_INTS                 = INTEGER (Given)
      }{
         number of integrations taken
      }
      \sstsubsection{
         IN\_DATA (4, N\_BOLS, N\_JIGS $*$ N\_INTS)
      }{
                          = REAL (Given)
         the demodulated data; 1=data, 2=variance, 3=calibrator, 4=quality
      }
      \sstsubsection{
         ADC\_INDEX (N\_BOLS)     = INTEGER (Given)
      }{
         the ADC numbers of the measured bolometers
      }
      \sstsubsection{
         CHAN\_INDEX (N\_BOLS)    = INTEGER (Given)
      }{
         the channel numbers of the measured bolometers
      }
      \sstsubsection{
         INT\_INDEX (N\_INTS)     = INTEGER (Given)
      }{
         the k index in IN\_DATA of the first jiggle of each integration
      }
      \sstsubsection{
         ADC                    = INTEGER (Given)
      }{
         the ADC number of the requested bolometer
      }
      \sstsubsection{
         CHAN                   = INTEGER (Given)
      }{
         the channel number of the requested bolometer
      }
      \sstsubsection{
         OUT\_DATA (N\_INTS)      = REAL (Returned)
      }{
         the values for each integration averaged over the jiggles
      }
      \sstsubsection{
         OUT\_VARIANCE (N\_INTS)  = REAL (Returned)
      }{
         the variance on OUT\_DATA
      }
      \sstsubsection{
         OUT\_QUALITY (N\_INTS)   = INTEGER (Returned)
      }{
         the quality on OUT\_DATA
      }
      \sstsubsection{
         STATUS                 = INTEGER (Given and returned)
      }{
         global status
      }
   }
   \sstnotes{
      Only works for data where SWITCH\_PER\_EXP=1, EXP\_PER\_INT=1 and
      N\_MEASUREMENTS=1.
   }
   \sstdiytopic{
      Copyright
   }{
      Copyright \copyright 1995,1996,1997,1998,1999 Particle Physics and Astronomy
      Research Council. All Rights Reserved.
   }
}
\sstroutine{
   SCULIB\_CONSTRUCT\_OUT
}{
   Append a identification string to a string
}{
   \sstdescription{
      Appends an ID string to an input string given the contents
      of the environment variable specified by ENV\_NAME
      Allowed values of ENV\_NAME are stored in OPTIONS and the
      corresponding string stored in OPTION\_STRINGS.
      Default is to go for the first entry in the IDSTRING array.
      if the value of ENV\_NAME is either not set or unrecognized.
   }
   \sstinvocation{
       CALL SCULIB\_CONSTRUCT\_OUT(IN\_STRING, ENV\_NAME,
         N\_OPTIONS, OPTIONS, OPTION\_STRINGS, OUT\_STRING, STATUS)
   }
   \sstarguments{
      \sstsubsection{
         IN\_STRING = CHAR (Given)
      }{
         Input string
      }
      \sstsubsection{
         ENV\_NAME = CHAR (Given)
      }{
         Name of environment variable to examine for suffix preference
      }
      \sstsubsection{
         N\_OPTIONS = INTEGER (Given)
      }{
         Number of suffix options
      }
      \sstsubsection{
         OPTIONS(N\_OPTIONS) = CHAR (Given)
      }{
         Names of different options
      }
      \sstsubsection{
         OPTION\_STRINGS(N\_OPTIONS) = CHAR (Given)
      }{
         Suffixes for each option. If the string starts with a {\tt '}!{\tt '}
         the input string is chopped from the last {\tt '}\_{\tt '} before adding
         the new string.
      }
      \sstsubsection{
         OUT\_STRING = CHAR (Returned)
      }{
         Input string with suffix
      }
      \sstsubsection{
         STATUS = INTEGER (Given \& Returned)
      }{
         Global status
      }
   }
   \sstdiytopic{
      Copyright
   }{
      Copyright \copyright 1995,1996,1997,1998,1999 Particle Physics and Astronomy
      Research Council. All Rights Reserved.
   }
}
\sstroutine{
   SCULIB\_CONVOLVE
}{
   convolve array A with array B to give result in R
}{
   \sstdescription{
      This routine convolves array A with array B using multiplication
      only (i.e. no FFTs are used). The routine assumes that the A arrays
      are N\_A long and that the convolution array B is N\_B long, centred
      at N\_MIDDLE. The final result is normalised by dividing by NORM.
      Elements of the A array that have bad quality are ignored. The input
      variances are propagated as:-

\begin{equation}
R_{Variance} = \frac{ \sum B^2A_{Variance}}{norm^2}
\end{equation}

      After this process the errors can no longer be considered independent.
   }
   \sstinvocation{
       CALL SCULIB\_CONVOLVE (A\_DATA, A\_VARIANCE, A\_QUALITY, B, N\_A, N\_B,
         N\_MIDDLE, NORM, R\_DATA, R\_VARIANCE, R\_QUALITY, STATUS)
   }
   \sstarguments{
      \sstsubsection{
         A\_DATA (N\_A)                  = REAL (Given)
      }{
         input data array
      }
      \sstsubsection{
         A\_VARIANCE (N\_A)              = REAL (Given)
      }{
         variance on A\_DATA
      }
      \sstsubsection{
         A\_QUALITY (N\_A)               = BYTE (Given)
      }{
         quality on A\_DATA
      }
      \sstsubsection{
         B (N\_B)                       = REAL (Given)
      }{
         convolution function
      }
      \sstsubsection{
         N\_A                           = INTEGER (Given)
      }{
         the size of A\_DATA
      }
      \sstsubsection{
         N\_B                           = INTEGER (Given)
      }{
         length of convolution function
      }
      \sstsubsection{
         N\_MIDDLE                      = INTEGER (Given)
      }{
         index of zero point of convolution function
      }
      \sstsubsection{
         NORM                          = REAL (Given)
      }{
         normalisation factor for convolution function
      }
      \sstsubsection{
         R\_DATA (N\_A)                  = REAL (Returned)
      }{
         result of the convolution
      }
      \sstsubsection{
         R\_VARIANCE (N\_A)              = REAL (Returned)
      }{
         variance on R\_DATA
      }
      \sstsubsection{
         R\_QUALITY (N\_A)               = QUALITY (Returned)
      }{
         quality on R\_DATA
      }
      \sstsubsection{
         BADBIT                        = BYTE (Given)
      }{
         badbit mask
      }
      \sstsubsection{
         STATUS                        = INTEGER (Given and returned)
      }{
         global status
      }
   }
   \sstdiytopic{
      Copyright
   }{
      Copyright \copyright 1995,1996,1997,1998,1999 Particle Physics and Astronomy
      Research Council. All Rights Reserved.
   }
}
\sstroutine{
   SCULIB\_COPY\_DEMOD\_SWITCH
}{
   copy a switch of demodulated data into
   arrays holding the switch chop data, chop
   variance, calibrator, cal variance and quality
   separately
}{
   \sstdescription{
      This routine copies the demodulated data for the jiggle positions
      measured in the last switch into separate arrays for the switch
      chop data, chop variance, calibrator, cal variance and quality.
      Quality is marked bad if BIT 1 is set and LEVEL is greater than requested
   }
   \sstinvocation{
       CALL SCULIB\_COPY\_DEMOD\_SWITCH (NBOLS, NDATA, NJIG,
         SWITCH, SWITCH\_DATA, SWITCH\_VARIANCE, SWITCH\_CALIBRATOR,
         SWITCH\_CAL\_VARIANCE, SWITCH\_QUALITY, STATUS)
   }
   \sstarguments{
      \sstsubsection{
         NBOLS                           = INTEGER (Given)
      }{
         number of bolometers being measured
      }
      \sstsubsection{
         NDATA                           = INTEGER (Given)
      }{
         number of data per measured position; 5 for jiggle data, 4 for
         raster
      }
      \sstsubsection{
         NJIG                            = INTEGER (Given)
      }{
         number of jiggle positions measured
      }
      \sstsubsection{
         LEVEL                           = INTEGER (Given)
      }{
         level at which spikes are marked as bad quality
      }
      \sstsubsection{
         SWITCH (5, NBOLS, NJIG)         = REAL (Given)
      }{
         switch datablock
      }
      \sstsubsection{
         SWITCH\_DATA (NBOLS, NJIG)       = REAL (Returned)
      }{
         chop data for this switch of the exposure
      }
      \sstsubsection{
         SWITCH\_VARIANCE (NBOLS, NJIG)   = REAL (Returned)
      }{
         variance on the chop signal
      }
      \sstsubsection{
         SWITCH\_CALIBRATOR (NBOLS, NJIG) = REAL (Returned)
      }{
         calibrator for this switch of the exposure
      }
      \sstsubsection{
         SWITCH\_CAL\_VARIANCE (NBOLS, NJIG)= REAL (Returned)
      }{
         variance on the calibrator signal
      }
      \sstsubsection{
         SWITCH\_QUALITY (NBOLS, NJIG)    = BYTE (Returned)
      }{
         quality for this switch of the exposure
      }
      \sstsubsection{
         STATUS                          = INTEGER (Given and returned)
      }{
         global status
      }
   }
   \sstdiytopic{
      Copyright
   }{
      Copyright \copyright 1995,1996,1997,1998,1999 Particle Physics and Astronomy
      Research Council. All Rights Reserved.
   }
}
\sstroutine{
   SCULIB\_COPY\_GOOD
}{
   Go through data set and throw away bad values
}{
   \sstdescription{
      This routine copies good data and variance from an input array to
      to an output array, leaving behind BAD values.
   }
   \sstinvocation{
       CALL SCULIB\_COPY\_GOOD(NPTS, INDATA, INVAR, INARR1, INARR2,
            NGOOD, OUTDATA, OUTVAR, OUTARR1, OUTARR2, STATUS)
   }
   \sstarguments{
      \sstsubsection{
         NPTS = INTEGER (Given)
      }{
         The number of points in the input array
      }
      \sstsubsection{
         INDATA ( NPTS ) = REAL (Given)
      }{
         Input data
      }
      \sstsubsection{
         INVAR ( NPTS ) = REAL (Given)
      }{
         Input variance
      }
      \sstsubsection{
         INARR1 ( NPTS ) = DOUBLE (Given)
      }{
         Auxilliary double precision array
      }
      \sstsubsection{
         INARR2 ( NPTS ) = DOUBLE (Given)
      }{
         Auxilliary double precision array
      }
      \sstsubsection{
         NGOOD = INTEGER (Returned)
      }{
         Number of good points in input data
      }
      \sstsubsection{
         OUTDATA ( NPTS ) = REAL (Returned)
      }{
         Output data
      }
      \sstsubsection{
         OUTVAR ( NPTS ) = REAL (Returned)
      }{
         Output variance
      }
      \sstsubsection{
         OUTARR1 ( NPTS ) = REAL (Returned)
      }{
         Good auxiliary data converted to real
      }
      \sstsubsection{
         OUTARR2 ( NPTS ) = REAL (Returned)
      }{
         Good auxiliary data converted to real
      }
      \sstsubsection{
         STATUS = INTEGER (Given \& Returned)
      }{
         Global status value
      }
   }
   \sstnotes{
      \sstitemlist{

         \sstitem
         Data are still good even if VARIANCE or AUX are bad
      }
   }
   \sstdiytopic{
      Copyright
   }{
      Copyright \copyright 1995,1996,1997,1998,1999 Particle Physics and Astronomy
      Research Council. All Rights Reserved.
   }
   \sstimplementationstatus{
      \sstitemlist{

         \sstitem
         Only a routine for REAL numbers exists
      }
   }
}
\sstroutine{
   SCULIB\_COPYB
}{
   copy one byte array to another
}{
   \sstdescription{
      copies one integer array into another
   }
   \sstinvocation{
      CALL SCULIB\_COPYB (N, FROM, TO)
   }
   \sstarguments{
      \sstsubsection{
         N          = INTEGER (Given)
      }{
         number of integers in arrays
      }
      \sstsubsection{
         FROM (N)   = BYTE (Given)
      }{
         array copied from
      }
      \sstsubsection{
         TO (N)     = BYTE (Returned)
      }{
         array copied to
      }
   }
   \sstnotes{
      This routine is deprecated. Use VEC\_BTOB (SUN/39) instead.
   }
   \sstdiytopic{
      Copyright
   }{
      Copyright \copyright 1995,1996,1997,1998,1999 Particle Physics and Astronomy
      Research Council. All Rights Reserved.
   }
}
\sstroutine{
   SCULIB\_COPYD
}{
   copy one double precision array to another
}{
   \sstdescription{
      copies one double precision array into another
   }
   \sstinvocation{
      CALL SCULIB\_COPYD (N, FROM, TO)
   }
   \sstarguments{
      \sstsubsection{
         N          = INTEGER (Given)
      }{
         number of elements in arrays
      }
      \sstsubsection{
         FROM (N)   = DOUBLE PRECISION (Given)
      }{
         array copied from
      }
      \sstsubsection{
         TO (N)     = DOUBLE PRECISION (Returned)
      }{
         array copied to
      }
   }
   \sstnotes{
      This routine is deprecated. Use VEC\_DTOD (SUN/39) instead.
   }
   \sstdiytopic{
      Copyright
   }{
      Copyright \copyright 1995,1996,1997,1998,1999 Particle Physics and Astronomy
      Research Council. All Rights Reserved.
   }
}
\sstroutine{
   SCULIB\_COPYI
}{
   copy one integer array to another
}{
   \sstdescription{
      copies one integer array into another
   }
   \sstinvocation{
      CALL SCULIB\_COPYI (N, FROM, TO)
   }
   \sstarguments{
      \sstsubsection{
         N          = INTEGER (Given)
      }{
         number of integers in arrays
      }
      \sstsubsection{
         FROM (N)   = INTEGER (Given)
      }{
         array copied from
      }
      \sstsubsection{
         TO (N)     = INTEGER (Returned)
      }{
         array copied to
      }
   }
   \sstnotes{
      This routine is deprecated. Use VEC\_ITOI (SUN/39) instead.
   }
   \sstdiytopic{
      Copyright
   }{
      Copyright \copyright 1995,1996,1997,1998,1999 Particle Physics and Astronomy
      Research Council. All Rights Reserved.
   }
}
\sstroutine{
   SCULIB\_COPYR
}{
   copy one real array to another
}{
   \sstdescription{
      copies one real array into another
   }
   \sstinvocation{
      CALL SCULIB\_COPYR (N, FROM, TO)
   }
   \sstarguments{
      \sstsubsection{
         N          = INTEGER (Given)
      }{
         number of reals in arrays
      }
      \sstsubsection{
         FROM (N)   = REAL (Given)
      }{
         array copied from
      }
      \sstsubsection{
         TO (N)     = REAL (Returned)
      }{
         array copied to
      }
   }
   \sstnotes{
      This routine is deprecated. Use VEC\_RTOR (SUN/39) instead.
   }
   \sstdiytopic{
      Copyright
   }{
      Copyright \copyright 1995,1996,1997,1998,1999 Particle Physics and Astronomy
      Research Council. All Rights Reserved.
   }
}
\sstroutine{
   SCULIB\_CORRECT\_EXTINCTION
}{
   correct bolometers for sky opacity
}{
   \sstdescription{
      This routine corrects bolometer data for the effect of sky opacity.
      It does this by calculating the airmass of the point that each
      bolometer was looking at, then multiplying the data by
      exp (airmass $*$ TAUZ). Bolometers with bad data quality will be
      ignored.
   }
   \sstinvocation{
       CALL SCULIB\_CORRECT\_EXTINCTION (SIZE\_BOL, N\_BOL, BOL\_DATA,
         BOL\_VARIANCE, BOL\_RA, BOL\_DEC, LST, LAT\_OBS,
         TAUZ, STATUS)
   }
   \sstarguments{
      \sstsubsection{
         SIZE\_BOL                       = INTEGER (Given)
      }{
         dimension of arrays
      }
      \sstsubsection{
         N\_BOL                          = INTEGER (Given)
      }{
         used size of arrays
      }
      \sstsubsection{
         BOL\_DATA (SIZE\_BOL)            = REAL (Given and returned)
      }{
         bolometer data
      }
      \sstsubsection{
         BOL\_VARIANCE (SIZE\_BOL)        = REAL (Given and returned)
      }{
         variance on BOL\_DATA
      }
      \sstsubsection{
         BOL\_RA (SIZE\_BOL)              = DOUBLE PRECISION (Given)
      }{
         apparent RA of bolometer (radians)
      }
      \sstsubsection{
         BOL\_DEC (SIZE\_BOL)             = DOUBLE PRECISION (Given)
      }{
         apparent dec of bolometer (radians)
      }
      \sstsubsection{
         LST                            = DOUBLE PRECISION (Given)
      }{
         sidereal time (radians)
      }
      \sstsubsection{
         LAT\_OBS                        = DOUBLE PRECISION (Given)
      }{
         latitude of observatory (radians)
      }
      \sstsubsection{
         TAUZ                           = REAL (Given)
      }{
         the zenith sky opacity
      }
      \sstsubsection{
         STATUS                         = INTEGER (Given and returned)
      }{
         global status
      }
   }
   \sstdiytopic{
      Copyright
   }{
      Copyright \copyright 1995,1996,1997,1998,1999 Particle Physics and Astronomy
      Research Council. All Rights Reserved.
   }
}
\sstroutine{
   SCULIB\_COVSRT
}{
   Numerical Recipes in Fortran routine called by SCULIB\_MRQMIN
}{
   \sstdescription{
      Expand in storage the covariance matrix COVAR, so as to take into account
      parameters that are being held fixed. (For the latter return zero
      covariances). Derived from COVSRT in Numerical Recipes in Fortran, p.669,
      with STATUS added.
   }
   \sstinvocation{
      CALL SCULIB\_COVSRT (COVAR, NPC, MA, IA, MFIT, STATUS)
   }
   \sstarguments{
      \sstsubsection{
         COVAR (  NPC, NPC ) = REAL (Given \& Returned)
      }{
         Covariance matrix
      }
      \sstsubsection{
         INTEGER NPC = INTEGER (Given)
      }{
         Size of covariance matrix
      }
      \sstsubsection{
         INTEGER MA = INTEGER (Given)
      }{
         Size of IA array. Total number of parameters.
      }
      \sstsubsection{
         INTEGER IA ( MA ) = INTEGER (Given)
      }{
         Ordering of parameters. Not a free parameter if containes 0.
      }
      \sstsubsection{
         INTEGER MFIT = INTEGER (Given)
      }{
         Number of required parameters
      }
      \sstsubsection{
         STATUS = INTEGER (Given \& Returned)
      }{
         Global status
      }
   }
   \sstnotes{
      Derived from the {\tt '}Numerical Recipes in Fortran{\tt '} COVSRT routine.
   }
   \sstdiytopic{
      Copyright
   }{
      Copyright \copyright 1993,1999 Particle Physics and Astronomy
      Research Council. All Rights Reserved.
   }
}
\sstroutine{
   SCULIB\_CROSSTALK
}{
   crosstalk measurements
}{
   \sstdescription{
      Experimental routine for calculating crosstalk between
      A-to-D cards. Averages signal for each bolometer and writes
      to a file. The file name comes from the FILE ADAM parameter.
   }
   \sstinvocation{
       CALL SCULIB\_CROSSTALK (SWITCH\_PER\_EXP, EXP\_PER\_INT,
         N\_INTEGRATIONS, N\_MEASUREMENTS, DEMOD\_POINTER, N\_BOLS,
         NUMPOS, DATA, BOL\_CHAN, BOL\_ADC, STATUS)
   }
   \sstarguments{
      \sstsubsection{
         SWITCH\_PER\_EXP         = INTEGER (Given)
      }{
         number of switches per exposure
      }
      \sstsubsection{
         EXP\_PER\_INT            = INTEGER (Given)
      }{
         number of exposures per integration
      }
      \sstsubsection{
         N\_INTEGRATIONS         = INTEGER (Given)
      }{
         number of integrations
      }
      \sstsubsection{
         N\_MEASUREMENTS         = INTEGER (Given)
      }{
         number of measurments
      }
      \sstsubsection{
         DEMOD\_POINTER (SWITCH\_PER\_EXP, EXP\_PER\_INT, N\_INTEGRATIONS,
      }{
         N\_MEASUREMENTS)      = INTEGER (Given)
             pointer to start of demodulated data for each switch in
             DATA array
      }
      \sstsubsection{
         N\_BOLS                 = INTEGER (Given)
      }{
         number of bolometers being measured
      }
      \sstsubsection{
         NUMPOS                 = INTEGER (Given)
      }{
         total number of positions measured
      }
      \sstsubsection{
         DATA (4, N\_BOLS, NUMPOS)
      }{
                          = REAL (Given)
         the demodulated data
      }
      \sstsubsection{
         BOL\_CHAN (N\_BOLS)      = INTEGER (Given)
      }{
         the channel number of the measured bolometers
      }
      \sstsubsection{
         BOL\_ADC (N\_BOLS)       = INTEGER (Given)
      }{
         the ADC numbers of the measured bolometers
      }
      \sstsubsection{
         STATUS                 = INTEGER (Given and returned)
      }{
         global status
      }
   }
   \sstdiytopic{
      Copyright
   }{
      Copyright \copyright 1995,1999 Particle Physics and Astronomy
      Research Council. All Rights Reserved.
   }
}
\sstroutine{
   SCULIB\_DAY
}{
   returns date and time as day number since 1st Jan
}{
   \sstdescription{
      Calculates the time since 1st January. The time is returned
      as a day number and fractional day.
   }
   \sstinvocation{
      DAY = SCULIB\_DAY ()
   }
   \sstarguments{
      \sstsubsection{
         None
      }{
      }
   }
   \sstreturnedvalue{
      \sstsubsection{
         SCULIB\_DAY = DOUBLE PRECISION
      }{
         date and time as a day number since 1st Jan
      }
   }
   \sstdiytopic{
      Copyright
   }{
      Copyright \copyright 1994,1999 Particle Physics and Astronomy
      Research Council. All Rights Reserved.
   }
}
\sstroutine{
   SCULIB\_DECODE\_ANGLE
}{
   convert angle string to double precision angle
}{
   \sstdescription{
      This routine converts an angle in dd:mm:ss.dd format into a
      double precision number in radians. It assumes that the input
      string specifies the angle in degrees. The process involves
      removing the {\tt '}:{\tt '} delimiters, then calling SLA\_DAFIN to perform
      the conversion. An error will be returned if there are less than 2 {\tt '}:{\tt '}
      delimiters in the string, or if the SLA routine errors.
   }
   \sstinvocation{
      CALL SCULIB\_DECODE\_ANGLE (STRING, ANGLE, STATUS)
   }
   \sstarguments{
      \sstsubsection{
         STRING         = CHARACTER$*$($*$) (Given)
      }{
         angle in xx:mm:ss.ddd format
      }
      \sstsubsection{
         ANGLE          = DOUBLE PRECISION (Returned)
      }{
         angle in radians
      }
      \sstsubsection{
         STATUS         = INTEGER (Given and returned)
      }{
         Global status
      }
   }
   \sstdiytopic{
      Copyright
   }{
      Copyright \copyright 1995,1996,1997,1998,1999 Particle Physics and Astronomy
      Research Council. All Rights Reserved.
   }
}
\sstroutine{
   SCULIB\_DECODE\_COMPONENT
}{
   This routine decodes the value of a single component in a SCUBA
   data-spec
}{
   \sstdescription{
      This routine decodes the value of a single component in a SCUBA
      data-spec. It is called by SCULIB\_DECODE\_SPEC. The component can be:-

      \sstitemlist{

         \sstitem
         [1]   {\tt "}$*${\tt "}                     select all data

         \sstitem
         [2]   {\tt "}$<$index$>${\tt "}               select data at position $<$index$>$

         \sstitem
         [3]   {\tt "}$<$index1$>$,$<$index2$>${\tt "}     select data at positions $<$index1$>$ and
                                   $<$index2$>$

         \sstitem
         [4]   {\tt "}$<$index1$>$:$<$index2$>${\tt "}     select all data in the range $<$index1$>$ to
                                   $<$index2$>$

      }
      or any combination of 2, 3 and 4 separated by commas. Example
      component values are:-

      \sstitemlist{

         \sstitem
          {\tt "}$*${\tt "}                     select all data

         \sstitem
          {\tt "}1,5{\tt "}                   select data at indices 1 and 5

         \sstitem
          {\tt "}5:10,16{\tt "}               select data at indices 5 through 10 and 16

      }
      Errors will occur if:-

      \sstitemlist{

         \sstitem
           You mix the $*$ format with selection by index.

         \sstitem
           Any index selected lies outside the range 1 to N.

         \sstitem
           In a selection range $<$index2$>$ is less than $<$index1$>$.

         \sstitem
           The component does not conform to the design syntax.

      }
      Output from the routine is in the form of a mask array SELECT.
      Data indices selected will be marked as 1 in SELECT, while indices
      not selected will be zeroes.
   }
   \sstinvocation{
      CALL SCULIB\_DECODE\_COMPONENT (COMPONENT, N, SELECT, STATUS)
   }
   \sstarguments{
      \sstsubsection{
         COMPONENT                        = CHARACTER$*$($*$) (Given)
      }{
         the component  to be decoded
      }
      \sstsubsection{
         N                                = INTEGER (Given)
      }{
         the number of items that can be selected
      }
      \sstsubsection{
         SELECT (N)                       = INTEGER (Returned)
      }{
         the items selected; 1 for selected, 0 for not selected
      }
      \sstsubsection{
         STATUS                           = INTEGER (Given and Returned)
      }{
         The global status.
      }
   }
   \sstdiytopic{
      Copyright
   }{
      Copyright \copyright 1995,1996,1997,1998,1999 Particle Physics and Astronomy
      Research Council. All Rights Reserved.
   }
}
\sstroutine{
   SCULIB\_DECODE\_FILTER
}{
   decode filter name into names and wavelengths of
   filters in front of each instrument section in use
}{
   \sstdescription{
      This routine decodes the filter name to give the name and central
      wavelength of the filter in front of each sub-instrument being used.

         The filter name should be in format $<$short$>$:$<$long$>$ where these
      are the names of the filters in front of the short and long-wave
      focal planes.

\texttt{
\begin{tabular}{crrcrl}
\multicolumn{6}{l}{$<$short$>$ can be:-}\\
'350'  &    short wavelength &=&  350.0 & name =& '350'\\
'450'  &                     & &  450.0 &       & '450'\\
'600'  &                     & &  600.0 &       & '600'\\
'750'  &                     & &  750.0 &       & '750'\\
'850'  &                     & &  850.0 &       & '850'\\
'PHOT' &                     & &   -1.0 &       & 'BLANK'\\
&&&&&\\
\multicolumn{6}{l}{$<$long$>$ can be:-}\\
'350'  &         long array  &=&  350.0  &name =& '350'\\
       &P1100, P1300, P2000  & &  350.0  &      & '350'\\
'450'  &         long array  &=&  450.0  &      & '450'\\
       &P1100, P1300, P2000  & &  450.0  &      & '450'\\
'600'  &         long array  &=&  600.0  &      & '600'\\
       &P1100, P1300, P2000  & &  600.0  &      & '600'\\
'750'  &         long array  &=&  750.0  &      & '750'\\
       &P1100, P1300, P2000  & &  750.0  &      & '750'\\
'850'  &         long array  &=&  850.0  &      & '850'\\
       &P1100, P1300, P2000  & &  850.0  &      & '850'\\
'PHOT' &         long array  &=&   -1.0  &      &  'BLANK'\\
       &              P1100  &=& 1100.0  &      &'1100'\\
       &              P1300  &=& 1300.0  &      &'1300'\\
       &              P2000  &=& 2000.0  &      &'2000'\\
\end{tabular}
}

      If the filter name does not fit the `pattern' described an error message
      will be output and the routine return with bad status. The routine is
      insensitive to the case of `PHOT'.
   }
   \sstinvocation{
       CALL SCULIB\_DECODE\_FILTER (FILTER, N\_SUB, SUB\_INSTRUMENT,
         SUB\_FILTER, WAVELENGTH, STATUS)
   }
   \sstarguments{
      \sstsubsection{
         FILTER                   = CHARACTER$*$($*$) (Given)
      }{
         name of filter
      }
      \sstsubsection{
         N\_SUB                    = INTEGER (Given)
      }{
         number of sub instruments being used
      }
      \sstsubsection{
         SUB\_INSTRUMENT (N\_SUB)   = CHARACTER$*$($*$) (Given)
      }{
         the names of the sub instruments being used
      }
      \sstsubsection{
         SUB\_FILTER (N\_SUB)       = CHARACTER$*$($*$) (Returned)
      }{
         the name of the filter in front of each sub-instrument
      }
      \sstsubsection{
         WAVELENGTH (N\_SUB)       = REAL (Returned)
      }{
         the wavelength in microns of the filters in front of the sub
         instruments being used. An inappropriate filter will have
         wavelength set to VAL\_\_BADR
      }
      \sstsubsection{
         STATUS                   = INTEGER (Given and returned)
      }{
         global status
      }
   }
   \sstnotes{
      The wavelength values are probably out of step with those currently
      in use by the real-time system.
   }
   \sstdiytopic{
      Copyright
   }{
      Copyright \copyright 1995,1996,1997,1998,1999 Particle Physics and Astronomy
      Research Council. All Rights Reserved.
   }
}
\sstroutine{
   SCULIB\_SPLIT\_DECODE\_REBIN\_LINE
}{
   Splits a string into a filename and parameters for REBIN
}{
   \sstdescription{
      This routine takes a string and splits it into bits.
      The format is assumed to be:

          FILESPEC WEIGHT SHIFT\_DX SHIFT\_DY

      The string is broken up from left to right.
      There must always be a weight if there is a shift
      There must always be a shift if there is a section. etc...
      Comment character is \#. A line contains a comment once a \#
      appears on the line. Only chars before the \# are valid.
   }
   \sstinvocation{
       CALL SCULIB\_DECODE\_REBIN\_LINE(LINE, N\_FOUND, NAME,
            WEIGHT, SHIFT\_DX, SHIFT\_DY, STATUS)
   }
   \sstarguments{
      \sstsubsection{
         LINE = CHAR (Given)
      }{
         Line read from file
      }
      \sstsubsection{
         N\_FOUND = INTEGER (Returned)
      }{
         Number of parameters found in LINE
      }
      \sstsubsection{
         NAME = CHARACTER (Returned)
      }{
         File name $+$ (optional) SCUBA section
      }
      \sstsubsection{
         WEIGHT = REAL (Returned)
      }{
         Relative weight given to filename
      }
      \sstsubsection{
         SHIFT\_DX = REAL (Returned)
      }{
         Shift in X
      }
      \sstsubsection{
         SHIFT\_DY = REAL (Returned)
      }{
         Shift in Y
      }
      \sstsubsection{
         STATUS = INTEGER (Given and Returned)
      }{
         Global Status value
      }
   }
   \sstdiytopic{
      Copyright
   }{
      Copyright \copyright 1995,1996,1997,1998,1999 Particle Physics and Astronomy
      Research Council. All Rights Reserved.
   }
}
\sstroutine{
   SCULIB\_DECODE\_SPEC
}{
   Decode SCUBA-style data specifications
}{
   \sstdescription{
      This routine decodes a SCUBA-style data specification. The data-spec
      will be of the form \{$<$component$>$;$<$component$>$;...\}, where components
      are one of the following:-

      \sstitemlist{

         \sstitem
         B$<$index\_spec$>$   - specifying bolometer indices

         \sstitem
         P$<$index\_spec$>$   -            position indices

         \sstitem
         S$<$index\_spec$>$   -            switch indices

         \sstitem
         E$<$index\_spec$>$   -            exposure indices

         \sstitem
         I$<$index\_spec$>$   -            integration indices

         \sstitem
         M$<$index\_spec$>$   -            measurement indices

      }
      and the $<$index\_spec$>$ is a list like, for example, 2,5:7,17 to select
      indices 2, 5 through 7 and 17. Alternatively, $<$index\_spec$>$ can be $*$
      which will select all data in that component coordinate.
      The data spec can be passed in as an array of data specs. The
      number is specified by N\_SPEC.
      By default all components in a dataset are selected. Thus the
      empty data-spec \{\} will return all components selected. Example
      data-specs are:-

      \sstitemlist{

         \sstitem
          {\tt "}\{\}               -         select all data

         \sstitem
          {\tt "}\{B7,12;P57\}{\tt "}     -          select data for bolometers 7 and 12 at
                                     measurement position 57

         \sstitem
          {\tt "}\{S2;E1;I3:M2\}{\tt "}   -          select data for all bolometers in
                                     switch 2 of exposure 1 of integration 3
                                     in measurement 2 of the observation

         \sstitem
          {\tt "}\{B29\}{\tt "}           -          select all data for bolometer 29

         \sstitem
          {\tt "}\{B29;E1\}{\tt "}        -          select data for bolometer 29 in the
                                     first exposure of each integration

      }
      The data-spec is case-insensitive and blanks are ignored.

      Errors will occur:-

      \sstitemlist{

         \sstitem
          If you attempt to select indices outside the dimensions input to
           the routine.

         \sstitem
          If you attempt to select data both by position Pxxx and by switch
           Sxxx, exposure Exxx, integration Ixxx or measurement Mxxx.

         \sstitem
          If you attempt to select by switch Sxxx when the SWITCH\_EXPECTED
           flag is input .FALSE.

      }
      Output consists of a flag POS\_SELECTED, to say whether or not the
      data were selected by position, and mask arrays that are set
      to 1 at the coordinate of selected data and 0 otherwise. Even if the
      data were not selected by position the POS\_S array, which is the mask
      for the position coordinate, will be set correctly. Conversely,
      however, the MEAS\_S, INT\_S, EXP\_S and SWITCH\_S arrays, which are
      masks for the measurement, integration, exposure and switch will not
      be set to sensible values if the data are position selected.
   }
   \sstinvocation{
       CALL SCULIB\_DECODE\_SPEC (SPEC, DEMOD\_POINTER,
            N\_SWITCHES,N\_EXPOSURES, N\_INTEGRATIONS, N\_MEASUREMENTS,
            N\_POS, N\_BOLS, SWITCH\_EXPECTED, POS\_SELECTED, POS\_S,
            SWITCH\_S, EXP\_S, INT\_S, MEAS\_S, BOL\_S, STATUS)
   }
   \sstarguments{
      \sstsubsection{
         N\_SPEC                           = INTEGER (Given)
      }{
         Number of specifications supplied in SPEC
      }
      \sstsubsection{
         SPEC( N\_SPEC )                   = CHARACTER$*$($*$) (Given)
      }{
         the specification to be decoded
      }
      \sstsubsection{
         DEMOD\_POINTER (N\_SWITCHES, N\_EXPOSURES, N\_INTEGRATIONS,  N\_MEASUREMENTS)
      }{
      }
      \sstsubsection{
         = INTEGER (Given)
      }{
         the pointer to the location in the main data array of the data
         for each switch of the observation
      }
      \sstsubsection{
         N\_SWITCHES                       = INTEGER (Given)
      }{
         the number of switches per exposure
      }
      \sstsubsection{
         N\_EXPOSURES                      = INTEGER (Given)
      }{
         the number of exposures per integration
      }
      \sstsubsection{
         N\_INTEGRATIONS                   = INTEGER (Given)
      }{
         the number of integrations per measurement
      }
      \sstsubsection{
         N\_MEASUREMENTS                   = INTEGER (Given)
      }{
         the number of measurements in the observation
      }
      \sstsubsection{
         N\_POS                            = INTEGER (Given)
      }{
         the number of positions measured in the observation
      }
      \sstsubsection{
         N\_BOLS                           = INTEGER (Given)
      }{
         the number of bolometers measured in the observation
      }
      \sstsubsection{
         SWITCH\_EXPECTED                  = LOGICAL (Given)
      }{
         .TRUE. if a switch component is allowed in the data-spec
      }
      \sstsubsection{
         POS\_SELECTED                     = LOGICAL (Returned)
      }{
         .TRUE. if the P component is used in the data-spec
      }
      \sstsubsection{
         POS\_S (N\_POS)                    = INTEGER (Returned)
      }{
         the position mask array; 1 for selected positions, 0 otherwise
      }
      \sstsubsection{
         SWITCH\_S (N\_SWITCHES)            = INTEGER (Returned)
      }{
         the switch mask array; 1 for selected switches, 0 otherwise
      }
      \sstsubsection{
         EXP\_S (N\_EXPOSURES)              = INTEGER (Returned)
      }{
         the exposure mask array; 1 for selected exposures, 0 otherwise
      }
      \sstsubsection{
         INT\_S (N\_INTEGRATIONS)           = INTEGER (Returned)
      }{
         the integration mask array; 1 for selected integrations, 0
         otherwise
      }
      \sstsubsection{
         MEAS\_S (N\_MEASUREMENTS)          = INTEGER (Returned)
      }{
         the measurement mask array; 1 for selected measurements, 0
         otherwise
      }
      \sstsubsection{
         BOL\_S (N\_BOLS)                   = INTEGER (Returned)
      }{
         the bolometer mask array; 1 for selected bolometers, 0 otherwise
      }
      \sstsubsection{
         STATUS                           = INTEGER (Given and Returned)
      }{
         The global status.
      }
   }
   \sstdiytopic{
      Copyright
   }{
      Copyright \copyright 1995,1996,1997,1998,1999 Particle Physics and Astronomy
      Research Council. All Rights Reserved.
   }
}
\sstroutine{
   SCULIB\_DIV\_CALIBRATOR
}{
   divides the calibrator signal into the chop signal
}{
   \sstdescription{
      This routine divides the calibrator signal into the chop signal and
      variance over an array of measurements. No division will occur if the
      measurement quality is bad or the square of the calibrator signal is zero.
   }
   \sstinvocation{
       CALL SCULIB\_DIV\_CALIBRATOR (NELM, DATA, VARIANCE, CALIBRATOR,
         QUALITY)
   }
   \sstarguments{
      \sstsubsection{
         NELM                = INTEGER (Given)
      }{
         Number of elements in each array
      }
      \sstsubsection{
         DATA (NELM)         = REAL (Given and returned)
      }{
         Demodulated data array
      }
      \sstsubsection{
         VARIANCE (NELM)     = REAL (Given and returned)
      }{
         Variance on DATA
      }
      \sstsubsection{
         CALIBRATOR (NELM)   = REAL (Given)
      }{
         Calibration array
      }
      \sstsubsection{
         QUALITY (NELM)      = BYTE (Given)
      }{
         Quality on DATA
      }
      \sstsubsection{
         STATUS = INTEGER (Given \& Returned)
      }{
         inherited status
      }
   }
   \sstdiytopic{
      Copyright
   }{
      Copyright \copyright 1995,1996,1997,1998,1999 Particle Physics and Astronomy
      Research Council. All Rights Reserved.
   }
}
\sstroutine{
   SCULIB\_DIV\_CALIBRATOR\_2
}{
   divides the mean of the calibrator signal into
   the chop signal
}{
   \sstdescription{
      This routine divides the calibrator signal into the chop signal and
      variance over an array of measurements. No division will occur if the
      measurement quality is bad or the square of the calibrator signal is zero.
   }
   \sstinvocation{
       CALL SCULIB\_DIV\_CALIBRATOR\_2 (N\_BOLS, N\_POS, DATA, VARIANCE,
         CALIBRATOR, QUALITY)
   }
   \sstarguments{
      \sstsubsection{
         N\_BOLS              = INTEGER (Given)
      }{
         Number of bolometers measured
      }
      \sstsubsection{
         N\_POS               = INTEGER (Given)
      }{
         Number of positions measured
      }
      \sstsubsection{
         DATA (N\_BOLS,N\_POS) = REAL (Given and returned)
      }{
         Demodulated data array
      }
      \sstsubsection{
         VARIANCE (N\_BOLS,N\_POS)
      }{
                    = REAL (Given and returned)
         Variance on DATA
      }
      \sstsubsection{
         CALIBRATOR (N\_BOLS,N\_POS)
      }{
                    = REAL (Given)
         Calibration array
      }
      \sstsubsection{
         QUALITY (N\_BOLS,N\_POS)
      }{
                    = BYTE (Given and returned)
         Quality on DATA
      }
      \sstsubsection{
         STATUS = INTEGER (Given \& Returned)
      }{
         inherited status
      }
   }
   \sstdiytopic{
      Copyright
   }{
      Copyright \copyright 1995,1996,1997,1998,1999 Particle Physics and Astronomy
      Research Council. All Rights Reserved.
   }
}
\sstroutine{
   SCULIB\_EXTRACT\_2DIM\_B
}{
   To extract the second dimension (at a given X) from a 2d array
}{
   \sstdescription{
      This routine extracts a second dimension from a 2D byte array.
   }
   \sstinvocation{
      CALL SCULIB\_EXTRACT\_2DIM\_B(X, NX, NY, TWOD, ONED, STATUS)
   }
   \sstarguments{
      \sstsubsection{
         X = INTEGER (Given)
      }{
         The X index in the 2D array to be extracted
      }
      \sstsubsection{
         NX = INTEGER (Given)
      }{
         Size of the array in x
      }
      \sstsubsection{
         NY = INTEGER (Given)
      }{
         Size of the array in y
      }
      \sstsubsection{
         TWOD = BYTE (Given)
      }{
         2D array input dataset
      }
      \sstsubsection{
         ONED = BYTE (Returned)
      }{
         1D strip extracted from 2D array
      }
      \sstsubsection{
         STATUS = INTEGER (Given \& Returned)
      }{
         Global status
      }
   }
   \sstnotes{
      This is the BYTE routine.
   }
   \sstdiytopic{
      Copyright
   }{
      Copyright \copyright 1995,1996,1997,1998,1999 Particle Physics and Astronomy
      Research Council. All Rights Reserved.
   }
}
\sstroutine{
   SCULIB\_EXTRACT\_2DIM\_D
}{
   To extract the second dimension (at a given X) from a 2d array
}{
   \sstdescription{
      This routine extracts a second dimension from a 2D DOUBLE array.
   }
   \sstinvocation{
      CALL SCULIB\_EXTRACT\_2DIM\_D(X, NX, NY, TWOD, ONED, STATUS)
   }
   \sstarguments{
      \sstsubsection{
         X = INTEGER (Given)
      }{
         The X index in the 2D array to be extracted
      }
      \sstsubsection{
         NX = INTEGER (Given)
      }{
         Size of the array in x
      }
      \sstsubsection{
         NY = INTEGER (Given)
      }{
         Size of the array in y
      }
      \sstsubsection{
         TWOD = DOUBLE (Given)
      }{
         2D array input dataset
      }
      \sstsubsection{
         ONED = DOUBLE (Returned)
      }{
         1D strip extracted from 2D array
      }
      \sstsubsection{
         STATUS = INTEGER (Given \& Returned)
      }{
         Global status
      }
   }
   \sstnotes{
      This is the DOUBLE PRECISION routine.
   }
   \sstdiytopic{
      Copyright
   }{
      Copyright \copyright 1995,1996,1997,1998,1999 Particle Physics and Astronomy
      Research Council. All Rights Reserved.
   }
}
\sstroutine{
   SCULIB\_EXTRACT\_2DIM\_R
}{
   To extract the second dimension (at a given X) from a 2d array
}{
   \sstdescription{
      This routine extracts a second dimension from a 2D REAL array.
   }
   \sstinvocation{
      CALLSCULIB\_EXTRACT\_2DIM\_R(X, NX, NY, TWOD, ONED, STATUS)
   }
   \sstarguments{
      \sstsubsection{
         X = INTEGER (Given)
      }{
         The X index in the 2D array to be extracted
      }
      \sstsubsection{
         NX = INTEGER (Given)
      }{
         Size of the array in x
      }
      \sstsubsection{
         NY = INTEGER (Given)
      }{
         Size of the array in y
      }
      \sstsubsection{
         TWOD = REAL (Given)
      }{
         2D array input dataset
      }
      \sstsubsection{
         ONED = REAL (Returned)
      }{
         1D strip extracted from 2D array
      }
      \sstsubsection{
         STATUS = INTEGER (Given \& Returned)
      }{
         Global status
      }
   }
   \sstnotes{
      This is the REAL routine.
   }
   \sstdiytopic{
      Copyright
   }{
      Copyright \copyright 1995,1996,1997,1998,1999 Particle Physics and Astronomy
      Research Council. All Rights Reserved.
   }
}
\sstroutine{
   SCULIB\_EXTRACT\_BOL
}{
   To extract a single bolometer from a data array
}{
   \sstdescription{
      This routine extracts a single bolometer from a SCUBA data array.
   }
   \sstinvocation{
       CALL SCULIB\_EXTRACT\_BOL(NBOL, N\_BOLS, N\_POS, SCUDATA,
            SCUQUAL, BOLDATA, BOLQUAL, STATUS)
   }
   \sstarguments{
      \sstsubsection{
         NBOL = INTEGER (Given)
      }{
         The bolometer number
      }
      \sstsubsection{
         N\_BOLS = \_INTEGER (Given)
      }{
         Total number of bolometers in data array
      }
      \sstsubsection{
         N\_POS = INTEGER (Given)
      }{
         Number of jiggle positions in data set
      }
      \sstsubsection{
         SCUDATA(N\_BOLS, N\_POS) = REAL (Given)
      }{
         The data
      }
      \sstsubsection{
         SCUQUAL(N\_BOLS, N\_POS) = BYTE (Given)
      }{
         The data quality
      }
      \sstsubsection{
         BOLDATA(N\_POS) = REAL (Returned)
      }{
         The specified bolometer data
      }
      \sstsubsection{
         BOLQUAL(N\_POS) = \_BYTE (Returned)
      }{
         The bolometer quality
      }
      \sstsubsection{
         STATUS = INTEGER (Given and Returned)
      }{
         Global Status value
      }
   }
   \sstdiytopic{
      Copyright
   }{
      Copyright \copyright 1995,1996,1997,1998,1999 Particle Physics and Astronomy
      Research Council. All Rights Reserved.
   }
   \sstimplementationstatus{
      Propogates quality
   }
}
\sstroutine{
   SCULIB\_FIND\_INT
}{
   Return the indices of the start and end of the specified integration
}{
   \sstdescription{
      This routine find that start and end indices of the specified
      integration. In general these can be found directly from
      the DEMOD\_POINTER array but with the warning that the last
      integration will not necessarily point to the end of the
      data array if multiple measurements are being taken. For this
      reason, put this knowledge in one place.
   }
   \sstinvocation{
       CALL SCULIB\_FIND\_INT(DEMOD\_POINTER, N\_SWITCHES,
            N\_EXPOSURES, N\_INTEGRATIONS, N\_MEASUREMENTS, N\_POS,
            INTEG, MEAS, I\_START, I\_END, STATUS)
   }
   \sstarguments{
      \sstsubsection{
         DEMOD\_POINTER ( N\_SWITCHES, N\_EXPOSURES, N\_INTEGRATIONS, N\_MEASUREMENTS ) = INTEGER (Given)
      }{
         an array of pointers to the start of each switch
      }
      \sstsubsection{
         N\_SWITCHES            = INTEGER (Given)
      }{
         the number of switches per exposure
      }
      \sstsubsection{
         N\_EXPOSURES           = INTEGER (Given)
      }{
         the number of exposures per integration
      }
      \sstsubsection{
         N\_INTEGRATIONS        = INTEGER (Given)
      }{
         the number of integrations per measurement
      }
      \sstsubsection{
         N\_MEASUREMENTS        = INTEGER (Given)
      }{
         the number of measurements in the observation
      }
      \sstsubsection{
         N\_POS                 = INTEGER (Given)
      }{
         the total number of positions measured in the observation
      }
      \sstsubsection{
         INTEG                 = INTEGER (Given)
      }{
         the integration number of the switch
      }
      \sstsubsection{
         MEAS                  = INTEGER (Given)
      }{
         the measurement number of the switch
      }
      \sstsubsection{
         I\_START               = INTEGER (Returned)
      }{
         pointer to the start of the integration
      }
      \sstsubsection{
         I\_END                 = INTEGER (Returned)
      }{
         pointer to the end of the integration
      }
      \sstsubsection{
         STATUS                = INTEGER (Given and returned)
      }{
         global status
      }
   }
   \sstdiytopic{
      Copyright
   }{
      Copyright \copyright 1999 Particle Physics and Astronomy
      Research Council. All Rights Reserved.
   }
}
\sstroutine{
   SCULIB\_FIND\_SWITCH
}{
   find the start and end indices of a switch in the
   demodulated data array
}{
   \sstdescription{
      This routine finds the start and end indices of the specified switch
      in a SCUBA observation. The array DEMOD\_POINTER contains the start
      indices, so this part is easy. However, the end index will be 1 less
      than the start index of the next, or the index of the last position
      measured if the specified switch is the last in the observation. This
      makes things a little complicated because consecutive switches within
      an exposure need not have been taken and stored consecutively,
      though exposures, integrations and measurements are ordered as one
      would expect.
   }
   \sstinvocation{
       CALL SCULIB\_FIND\_SWITCH (DEMOD\_POINTER, N\_SWITCHES, N\_EXPOSURES,
         N\_INTEGRATIONS, N\_MEASUREMENTS, N\_POS, SWITCH, EXP, INT, MEAS,
         S\_START, S\_END, STATUS)
   }
   \sstarguments{
      \sstsubsection{
         DEMOD\_POINTER (N\_SWITCHES, N\_EXPOSURES, N\_INTEGRATIONS, N\_MEASUREMENTS) = INTEGER (Given)
      }{
         an array of pointers to the start of each switch
      }
      \sstsubsection{
         N\_SWITCHES            = INTEGER (Given)
      }{
         the number of switches per exposure
      }
      \sstsubsection{
         N\_EXPOSURES           = INTEGER (Given)
      }{
         the number of exposures per integration
      }
      \sstsubsection{
         N\_INTEGRATIONS        = INTEGER (Given)
      }{
         the number of integrations per measurement
      }
      \sstsubsection{
         N\_MEASUREMENTS        = INTEGER (Given)
      }{
         the number of measurements in the observation
      }
      \sstsubsection{
         N\_POS                 = INTEGER (Given)
      }{
         the total number of positions measured in the observation
      }
      \sstsubsection{
         SWITCH                = INTEGER (Given)
      }{
         the index of the switch
      }
      \sstsubsection{
         EXP                   = INTEGER (Given)
      }{
         the exposure number of the switch
      }
      \sstsubsection{
         INT                   = INTEGER (Given)
      }{
         the integration number of the switch
      }
      \sstsubsection{
         MEAS                  = INTEGER (Given)
      }{
         the measurement number of the switch
      }
      \sstsubsection{
         S\_START               = INTEGER (Returned)
      }{
         pointer to the start of the switch
      }
      \sstsubsection{
         S\_END                 = INTEGER (Returned)
      }{
         pointer to the end of the switch
      }
      \sstsubsection{
         STATUS                = INTEGER (Given and returned)
      }{
         global status
      }
   }
   \sstdiytopic{
      Copyright
   }{
      Copyright \copyright 1995,1996,1997,1998,1999 Particle Physics and Astronomy
      Research Council. All Rights Reserved.
   }
}
\sstroutine{
   SCULIB\_FIT\_2D\_GAUSSIAN
}{
   fit a 2D Gaussian to an image
}{
   \sstdescription{

      This routine fits a 2D Gaussian to an image of a source. It takes as
      input data a map made of the source with pixels on a square grid.

 If status is good on entry, the routine will calculate the 0th ($U_0$) and
      1st order ($U_x$, $U_y$) moments of the image on the map. Map pixels
      with bad quality are ignored:

 The 0th order is calculated as
\begin{equation}
U_0 = \sum_{ij} f_{ij},
\end{equation}
where i,j are the pixel indices. The 1st order moment in X is calculated as
\begin{equation}
U_x = \sum_{ij} x f_{ij},
\end{equation}
where x is the x position of pixel i,j; and the 1st order moment in Y is
calculated as
\begin{equation}
U_y = \sum_{ij} y f_{ij},
\end{equation}
where y is the y position of pixel i,j.

      If the 0th order moment of the image was 0, i.e. there is no image on
      the map, then a warning message will be output and all image parameters
      set to 0 or bad quality. Otherwise, the x,y centre of the image is
      calculated from:-

\begin{equation}
X_\mathrm{cen} = \frac{U_x}{U_0}
\end{equation}
\begin{equation}
Y_\mathrm{cen} = \frac{U_y}{U_0}
\end{equation}

      With this information the 2nd order moments can be calculated about the
      centre of the image:-

\begin{equation}
          U_{xx} = \sum_{ij} (x - X_{cen})^2 f_{ij}
\end{equation}
\begin{equation}
          U_{yy} = \sum_{ij} (y - Y_{cen})^2  f_{ij}
\end{equation}

      If the image had a 2D Gaussian profile,
\begin{equation}
f = A \exp(-a_x^2 x^2) \exp(-a_y^2 y^2),
\end{equation}
      then it can be shown that for x:-

\begin{equation}
a_x^2 = \frac{U_0}{2 U_{xx}}
\end{equation}

      and the half-width at half-maximum is calculated from:-

\begin{equation}
X_\mathrm{HWHM} = \sqrt{\ln(2)}{a_x}
\end{equation}

         Lastly, the volume under the image and the variance on it are
      calculated. These will be the sum of all map pixels and their variances
      that lie within a radius of 2 $*$ max(HWHM\_x, HWHM\_y) from X\_CENTRE,
      Y\_CENTRE. If the HWHM in either axis is zero then a warning message will
      be output and PEAK\_QUAL set bad. Otherwise, PEAK will be set to
      volume / ($\pi$ $*$ HWHM\_x $*$ HWHM\_y).
   }
   \sstinvocation{
       CALL SCULIB\_FIT\_2D\_GAUSSIAN (IDIM, JDIM, MAP\_DATA,
         MAP\_VARIANCE, MAP\_QUALITY, X, Y, PEAK, PEAK\_VAR,
         PEAK\_QUAL, X\_CENTRE, Y\_CENTRE, X\_HWHM, Y\_HWHM, STATUS)
   }
   \sstarguments{
      \sstsubsection{
         IDIM                              = INTEGER (Given)
      }{
         `I{\tt '} dimension of map
      }
      \sstsubsection{
         JDIM                              = INTEGER (Given)
      }{
         `J{\tt '} dimension of map
      }
      \sstsubsection{
         MAP\_DATA (IDIM, JDIM)             = REAL (Given)
      }{
         map data values
      }
      \sstsubsection{
         MAP\_VARIANCE (IDIM, JDIM)         = REAL (Given)
      }{
         map variance
      }
      \sstsubsection{
         MAP\_QUALITY (IDIM, JDIM)          = REAL (Given)
      }{
         map quality
      }
      \sstsubsection{
         X (IDIM)                          = REAL (Given)
      }{
         x axis of map
      }
      \sstsubsection{
         Y (JDIM)                          = REAL (Given)
      }{
         y axis of map
      }
      \sstsubsection{
         PEAK                              = REAL (Returned)
      }{
         peak height of Gaussian
      }
      \sstsubsection{
         PEAK\_VAR                          = REAL (Returned)
      }{
         variance on PEAK
      }
      \sstsubsection{
         PEAK\_QUAL                         = INTEGER (Returned)
      }{
         quality of PEAK
      }
      \sstsubsection{
         X\_CENTRE                          = REAL (Returned)
      }{
         x of image centre
      }
      \sstsubsection{
         Y\_CENTRE                          = REAL (Returned)
      }{
         y of image centre
      }
      \sstsubsection{
         X\_HWHM                            = REAL (Returned)
      }{
         half-width half-max of image in x-axis
      }
      \sstsubsection{
         Y\_HWHM                            = REAL (Returned)
      }{
         half-width half-max of image in y-axis
      }
      \sstsubsection{
         STATUS                            = INTEGER (Given and returned)
      }{
         global status
      }
   }
   \sstdiytopic{
      Copyright
   }{
      Copyright \copyright 1994,1995,1996,1997 Particle Physics and Astronomy
      Research Council. All Rights Reserved.
   }
}
\sstroutine{
   SCULIB\_FIT\_2D\_PARABOLA
}{
   fit 2d parabola to data
}{
   \sstdescription{
      This routine performs a least-squares fit of a 2d parabola to a set
      of data. The form of the parabola is:-

\begin{equation}
      z = A0 + A1 * (x - X0)^2 + A1 * (y - Y0)^2
\end{equation}

      The apex of the parabola is at X0,Y0 and its value there is A0.

         If status is good on entry the routine will work through the data
      calculating the various sums required for the least-squares method.
      Data with bad quality or zero variance are ignored. If no `good{\tt '} data
      are found then an error will be reported and the routine will return
      with bad status.

         A matrix equation of the form $M * B = Z$ is constructed, where:-

\begin{equation}
M = \begin{pmatrix}
\sum \frac{1}{var}  & \sum \frac{x}{var}         & \sum \frac{y}{var}         &\sum \frac{x^2+y^2}{var}\\
\sum \frac{x}{var} & \sum \frac{x^2}{var}       & \sum \frac{xy}{var}        &\sum \frac{x^3+xy^2}{var}\\
\sum \frac{y}{var} & \sum \frac{xy}{var}        & \sum \frac{y^2}{var}       &\sum \frac{x^2y+y^3}{var}\\
\sum \frac{x^2+y^2}{var} & \sum \frac{x^3+xy^2}{var}& \sum \frac{x^2y+y^3}{var} &\sum \frac{(x^2+y^2)^2}{var}\\
\end{pmatrix}
\end{equation}

\begin{equation}
B = \begin{pmatrix}
A0 + A1 * X0^2 + A1 * Y0^2 \\
 -  2 * A1 * X0 \\
 -  2 * A1 * Y0 \\
    A1 \\
\end{pmatrix}
\end{equation}

\begin{equation}
Z = \begin{pmatrix}
 \sum \frac{z}{var} \\
 \sum \frac{xz}{var}\\
 \sum \frac{yz}{var}\\
 \sum \frac{(x^2 + y^2)z}{var}\\
\end{pmatrix}
\end{equation}

      where x, y, z and var are the x,y coords, value and variance of the
      measured points and the sums are over all
      valid points.

         SCULIB\_INVERT\_MATRIX is called to invert the matrix and
      SCULIB\_FIT\_MULT to multiply Z by the inverse. This yields the fitted
      value for B from which the other fit parameters are derived.
   }
   \sstinvocation{
       CALL SCULIB\_FIT\_2D\_PARABOLA (N, DATA, VARIANCE, QUALITY,
         X, Y, A0, A1, X0, Y0, Z\_PEAK, Z\_PEAK\_VAR, BADBIT, STATUS)
   }
   \sstarguments{
      \sstsubsection{
         N                      = INTEGER (Given)
      }{
         number of data points
      }
      \sstsubsection{
         DATA (N)               = REAL (Given)
      }{
         the data
      }
      \sstsubsection{
         VARIANCE (N)           = REAL (Given)
      }{
         variance on the data
      }
      \sstsubsection{
         QUALITY (N)            = BYTE (Given)
      }{
         the quality on the data
      }
      \sstsubsection{
         X (N)                  = REAL (Given)
      }{
         the x offsets of the data
      }
      \sstsubsection{
         Y (N)                  = REAL (Given)
      }{
         the y offsets of the data
      }
      \sstsubsection{
         A0                     = REAL (Returned)
      }{
         fit result
      }
      \sstsubsection{
         A1                     = REAL (Returned)
      }{
         fit result
      }
      \sstsubsection{
         X0                     = REAL (Returned)
      }{
         x coord of parabola apex
      }
      \sstsubsection{
         Y0                     = REAL (Returned)
      }{
         y coord of parabola apex
      }
      \sstsubsection{
         Z\_PEAK                 = REAL (Returned)
      }{
         the value of the parabola apex
      }
      \sstsubsection{
         Z\_PEAK\_VAR             = REAL (Returned)
      }{
         the variance on Z\_PEAK
      }
      \sstsubsection{
         BADBIT                 = BYTE (Given)
      }{
         bad bit mask
      }
      \sstsubsection{
         STATUS                 = INTEGER (Given and returned)
      }{
         global status
      }
   }
   \sstnotes{
      If the variances are not a true reflection of the errors
      on the data then very strange numbers can result.
   }
   \sstdiytopic{
      Copyright
   }{
      Copyright \copyright 1995,1996,1997,1998,1999 Particle Physics and Astronomy
      Research Council. All Rights Reserved.
   }
}
\sstroutine{
   SCULIB\_FIT\_D2XISQ\_DAJ2
}{
   calculate second differential of chi-squared
   with respect to a fit parameter
}{
   \sstdescription{
      If status is good on entry this routine will calculate the 2nd order
      differential of chi-squared with respect to A(J). It does this
      numerically by calling the supplied routine XISQ\_ROUTINE to calculate
      chi-squared at the current parameter values, then again at A(J)$+$DELTA\_AJ
      and A(J)-DELTA\_AJ, then using the formula:-

\begin{equation}
\frac{\partial^2\chi^2}{\partial AJ^2} = \frac{\chi^2(AJ+\delta AJ) -
2\chi^2(AJ) + \chi^2(AJ - \delta AJ)}{\delta AJ^2}
\end{equation}

      DELTA\_AJ is equal to the absolute value of 0.001 $*$ A(J) or, if this is
      zero, 0.001.
   }
   \sstinvocation{
       CALL SCULIB\_FIT\_D2XISQ\_DAJ2 (XISQ\_ROUTINE, N, A, J,
         D2XISQ\_DAJ2, STATUS)
   }
   \sstarguments{
      \sstsubsection{
         XISQ\_ROUTINE (XISQ, N, A, STATUS) = EXTERNAL ROUTINE (Given)
      }{
         name of routine that will calculate chi-squared
      }
      \sstsubsection{
         N                       = INTEGER (Given)
      }{
         the number of parameters in the fit
      }
      \sstsubsection{
         A(N)                    = DOUBLE PRECISION (Given and returned)
      }{
         the fit parameters
      }
      \sstsubsection{
         J                       = INTEGER (Given)
      }{
         the index of the parameter to be varied
      }
      \sstsubsection{
         D2XISQ\_DAJ2             = DOUBLE PRECISION (Returned)
      }{
         the second differential of chi-squared with respect to A(J)
      }
      \sstsubsection{
         STATUS                  = INTEGER (Given and returned)
      }{
         global status
      }
   }
   \sstdiytopic{
      Copyright
   }{
      Copyright \copyright 1995,1996,1997,1998,1999 Particle Physics and Astronomy
      Research Council. All Rights Reserved.
   }
}
\sstroutine{
   SCULIB\_FIT\_D2XISQ\_DAJK
}{
   calculate differential of chi-squared with respect
   to fit parameters J and K
}{
   \sstdescription{
      If status is good on entry this routine will calculate the differential
      of chi-squared with respect to A(J) and A(K). It does this numerically
      by calling the supplied routine XISQ\_ROUTINE to calculate chi-squared at
      [A(J)$+$DELTA\_AJ,A(K)$+$DELTA\_AK], [A(J)-DELTA\_AJ,A(K)$+$DELTA\_AK],
      [A(J)$+$DELTA\_AJ,A(K)-DELTA\_AK] and [A(J)-DELTA\_AJ,A(K)-DELTA\_AK].
      The result is then calculated using the formula:-

\begin{equation}
\begin{split}
\frac{\partial^2\chi^2}{\partial AJ\partial AK}=& \chi^2 (AJ+\delta AJ,AK+\delta AK) \\
                    &- \chi^2 (AJ-\delta AJ,AK+\delta AK) \\
                    &- \frac{\chi^2 (AJ-\delta AJ,AK+\delta AK) + \chi^2
(AJ-\delta AJ,AK-\delta AK)}{4 \delta AJ \delta AK}\\
\end{split}
\end{equation}

      DELTA\_AJ is equal to the absolute value of 0.001 $*$ A(J) or, if this is
      zero, 0.001. DELTA\_AK is calculated in a similar way.
   }
   \sstinvocation{
       CALL SCULIB\_FIT\_D2XISQ\_DAJK (XISQ\_ROUTINE, N, A, J, K,
         D2XISQ\_DAJK, STATUS)
   }
   \sstarguments{
      \sstsubsection{
         XISQ\_ROUTINE (XISQ, N, A, STATUS)  = EXTERNAL ROUTINE (Given)
      }{
         name of routine that will calculate xi-squared
      }
      \sstsubsection{
         N                       = INTEGER (Given)
      }{
         the number of parameters in the fit
      }
      \sstsubsection{
         A(N)                    = DOUBLE PRECISION (Given and returned)
      }{
         the fit parameters
      }
      \sstsubsection{
         J                       = INTEGER (Given)
      }{
         the index of the first parameter to be varied
      }
      \sstsubsection{
         K                       = INTEGER (Given)
      }{
         the index of the second parameter to be varied
      }
      \sstsubsection{
         D2XISQ\_DAJK             = DOUBLE PRECISION (Returned)
      }{
         the differential of chi-squared with respect to A(J) and A(K)
      }
      \sstsubsection{
         STATUS                  = INTEGER (Given and returned)
      }{
         global status
      }
   }
   \sstdiytopic{
      Copyright
   }{
      Copyright \copyright 1995,1996,1997,1998,1999 Particle Physics and Astronomy
      Research Council. All Rights Reserved.
   }
}
\sstroutine{
   SCULIB\_FIT\_DXISQ\_DAJ
}{
   calculate gradient in chi-squared with `a{\tt '}
}{
   \sstdescription{
      If status is good on entry this routine will calculate the gradient in
      the chi-squared of the fit with variation in fit parameter J. It does
      this numerically by calling the supplied routine XISQ\_ROUTINE to
      calculate chi-squared for (A(J)$+$DELTA\_AJ) and (A(J)-DELTA\_AJ), then
      using the equation:-

\begin{equation}
\frac{\partial\chi^2}{\partial AJ} = \frac{\chi^2 (A(J)+\delta AJ) - \chi^2 (A(J)-\delta AJ)}{2.0 \delta AJ}
\end{equation}

      DELTA\_AJ is equal to the absolute value of 0.001 $*$ A(J) or, if this is
      zero, 0.001.
   }
   \sstinvocation{
       CALL SCULIB\_FIT\_DXISQ\_DAJ (XISQ\_ROUTINE, N, A, J, DXISQ\_DAJ,
         STATUS)
   }
   \sstarguments{
      \sstsubsection{
         XISQ\_ROUTINE (XISQ, N, A, STATUS) = EXTERNAL ROUTINE (Given)
      }{
         name of routine that will calculate xi-squared
      }
      \sstsubsection{
         N                       = INTEGER (Given)
      }{
         the number of parameters in the fit
      }
      \sstsubsection{
         A (N)                   = DOUBLE PRECISION (Given and returned)
      }{
         the fit parameters
      }
      \sstsubsection{
         J                       = INTEGER (Given)
      }{
         the index of the parameter to be varied
      }
      \sstsubsection{
         DXISQ\_DAJ               = DOUBLE PRECISION (Returned)
      }{
         the gradient of chi-squared with A(J)
      }
      \sstsubsection{
         STATUS                  = INTEGER (Given and returned)
      }{
         global status
      }
   }
   \sstdiytopic{
      Copyright
   }{
      Copyright \copyright 1995,1996,1997,1998,1999 Particle Physics and Astronomy
      Research Council. All Rights Reserved.
   }
}
\sstroutine{
   SCULIB\_FIT\_FUNCTION
}{
   routine to fit a general function to data
}{
   \sstdescription{
      This routine fits a general function to a data-set by searching the
      chi-squared plane. Chi-squared at each test fit position will be
      evaluated by the external routine XISQ\_ROUTINE. The data to be fit
      are passed to the external routine through common; see
      SCULIB\_GAUSSIAN\_XISQ for an example. Each call to the routine will
      generate one iteration of the search. On exit the fit should have an
      improved chi-squared that should be checked for convergence (i.e. has
      chi-squared changed by less than CHICUT between 2 iterations) before
      the routine is called again to perform the next iteration. It is
      recommended that the routine be entered for the first iteration with
      LAMBDA set to 0.001. Thereafter, the routine will return the value
      of LAMBDA to be used in the next iteration.

        The search method used is that due to Marquardt as described in
      chapter 8 of `Data Reduction and Error Analysis{\tt '} by Bevington and
      Robinson. The code is an adaptation of the Pascal routine Marquardt
      listed in that book.

        If the routine is entered with LAMBDA = 0 then it will return the
      `error{\tt '} matrix for the given fit parameters.
   }
   \sstinvocation{
       CALL SCULIB\_FIT\_FUNCTION (XISQ\_ROUTINE, CHICUT, N, A,
         LAMBDA, ALPHA, BETA, IK, JK, DA, STATUS)
   }
   \sstarguments{
      \sstsubsection{
         XISQ\_ROUTINE (XISQ, N, A, STATUS) = EXTERNAL ROUTINE (Given)
      }{
         routine to calculate chi-squared of fit
      }
      \sstsubsection{
         CHICUT                  = DOUBLE PRECISION (Given)
      }{
         increase in chi-squared that will make routine search with
         larger LAMBDA
      }
      \sstsubsection{
         N                       = INTEGER (Given)
      }{
         number of fitted parameters
      }
      \sstsubsection{
         A (N)                   = DOUBLE PRECISION (Given and returned)
      }{
         the fitted parameters
      }
      \sstsubsection{
         LAMBDA                  = DOUBLE PRECISION (Given and returned)
      }{
         curvature factor for alpha matrix
      }
      \sstsubsection{
         ALPHA (N,N)             = DOUBLE PRECISION (Scratch)
      }{
         storage for alpha array
      }
      \sstsubsection{
         BETA (N)                = DOUBLE PRECISION (Scratch)
      }{
         storage for beta array
      }
      \sstsubsection{
         IK (N)                  = INTEGER (Scratch)
      }{
         used for inverting matrix
      }
      \sstsubsection{
         JK (N)                  = INTEGER (Scratch)
      }{
         used for inverting matrix
      }
      \sstsubsection{
         DA (N)                  = DOUBLE PRECISION (Scratch)
      }{
         increment in A
      }
      \sstsubsection{
         STATUS                  = INTEGER (Given and returned)
      }{
         global status
      }
   }
   \sstdiytopic{
      Copyright
   }{
      Copyright \copyright 1995,1996,1997,1998,1999 Particle Physics and Astronomy
      Research Council. All Rights Reserved.
   }
   \sstdiytopic{
      Method
   }{
      If status is good on entry the routine will enter a loop

      \sstitemlist{

         \sstitem
           SCULIB\_FIT\_MAKEBETA is called to calculate the `beta{\tt '} matrix

         \sstitem
           SCULIB\_FIT\_MAKEALPHA is called to calculate the `alpha{\tt '} matrix

         \sstitem
           The diagonal elements of the `alpha{\tt '} matrix are multiplied by
              1 $+$ LAMBDA following the method due to Marquardt

         \sstitem
           SCULIB\_INVERT\_MATRIX is called to invert the modified `alpha{\tt '}
              matrix

         \sstitem
           If LAMBDA = 0 this is all that{\tt '}s required because this input
              is given if the routine is required to calculate the `error{\tt '}
              matrix

         \sstitem
           Otherwise, SCULIB\_FIT\_MULT is called to multiply the inverted
              `alpha{\tt '} matrix by the `beta{\tt '} to give the next iteration of the
              fit for this LAMBDA

         \sstitem
           Chi-squared is calculated for the new fit and compared with
              that for the previous fit. If the chi-squared is greater than
              the previous value plus CHICUT then things have got worse.
              In this case the routine will increase LAMBDA by a factor of
              10, return the fit to the previous value and return to the
              start of the loop. If LAMBDA exceeds 1000 status is set bad,
              an error is reported and the loop is exited

         \sstitem
           However, if the chi-squared is an improvement then LAMBDA is
              decreased by a factor of 10 and the loop is exited

      }
      End of loop
   }
}
\sstroutine{
   SCULIB\_FIT\_MAKEALPHA
}{
   calculates alpha matrix for SCULIB\_FIT\_FUNCTION
}{
   \sstdescription{
      This routine calculates the elements of the alpha matrix for non-linear
      fitting as described in `Data Reduction and Error Analysis for the
      Physical Sciences{\tt '} by Bevington and Robinson, section 8.6 `the Marquardt
      Method{\tt '}.

         The routine does this by calling routines SCULIB\_FIT\_D2XISQ\_DAJ2
      and SCULIB\_FIT\_D2XISQ\_DAJK for each element required. The routine will
      only execute if entered with good status.
   }
   \sstinvocation{
      CALL SCULIB\_FIT\_MAKEALPHA (XISQ\_ROUTINE, N, A, ALPHA, STATUS)
   }
   \sstarguments{
      \sstsubsection{
         XISQ\_ROUTINE (XISQ, N, A, STATUS) = EXTERNAL ROUTINE (Given)
      }{
         routine to calculate chi-squared of current fit
      }
      \sstsubsection{
         N                         = INTEGER (Given)
      }{
         the number of fit parameters
      }
      \sstsubsection{
         A (N)                     = DOUBLE PRECISION (Given)
      }{
         the values of the fit parameters
      }
      \sstsubsection{
         ALPHA (N,N)               = DOUBLE PRECISION (Returned)
      }{
         the alpha matrix
      }
      \sstsubsection{
         STATUS                    = INTEGER (Given and returned)
      }{
         global status
      }
   }
   \sstdiytopic{
      Copyright
   }{
      Copyright \copyright 1995,1996,1997,1998,1999 Particle Physics and Astronomy
      Research Council. All Rights Reserved.
   }
}
\sstroutine{
   SCULIB\_FIT\_MAKEBETA
}{
   calculates beta matrix for SCULIB\_FIT\_FUNCTION
}{
   \sstdescription{
      This routine calculates the elements of the beta matrix for non-linear
      fitting as described in `Data Reduction and Error Analysis for the
      Physical Sciences{\tt '} by Bevington and Robinson, section 8.6 `the Marquardt
      Method{\tt '}.

         The routine does this by calling the routine SCULIB\_FIT\_DXISQ\_DAJ
      for each element required. The routine will only execute if entered
      with good status.
   }
   \sstinvocation{
      CALL SCULIB\_FIT\_MAKEBETA (XISQ\_ROUTINE, N, A, BETA, STATUS)
   }
   \sstarguments{
      \sstsubsection{
         XISQ\_ROUTINE (XISQ, N, A, STATUS) = EXTERNAL ROUTINE (Given)
      }{
         routine to calculate chi-squared of current fit
      }
      \sstsubsection{
         N                         = INTEGER (Given)
      }{
         the number of fit parameters
      }
      \sstsubsection{
         A (N)                     = DOUBLE PRECISION (Given)
      }{
         the values of the fit parameters
      }
      \sstsubsection{
         BETA (N)                  = DOUBLE PRECISION (Returned)
      }{
         the beta vector
      }
      \sstsubsection{
         STATUS                    = INTEGER (Given and returned)
      }{
         global status
      }
   }
   \sstdiytopic{
      Copyright
   }{
      Copyright \copyright 1995,1996,1997,1998,1999 Particle Physics and Astronomy
      Research Council. All Rights Reserved.
   }
}
\sstroutine{
   SCULIB\_FIT\_MULT
}{
   matrix vector multiplication routine
}{
   \sstdescription{
      If status is good on entry, this routine will multiply the N-vector
      BETA by the NxN matrix ALPHA to give the N-vector GAMMA. GAMMA should
      not be the same as BETA as this would mean BETA being overwritten
      during the multiplication process giving an incorrect answer.
   }
   \sstinvocation{
      CALL SCULIB\_FIT\_MULT (N, ALPHA, BETA, GAMMA, STATUS)
   }
   \sstarguments{
      \sstsubsection{
         N                      = INTEGER (Given)
      }{
         dimensions of square matrix
      }
      \sstsubsection{
         ALPHA (N,N)            = DOUBLE PRECISION (Given)
      }{
         the multiplying matrix
      }
      \sstsubsection{
         BETA (N)               = DOUBLE PRECISION (Given)
      }{
         the multiplied vector
      }
      \sstsubsection{
         GAMMA (N)              = DOUBLE PRECISION (Returned)
      }{
         the result vector (GAMMA and BETA should not be the same array)
      }
      \sstsubsection{
         STATUS                 = INTEGER (Given and returned)
      }{
         global status
      }
   }
   \sstdiytopic{
      Copyright
   }{
      Copyright \copyright 1995,1996,1997,1998,1999 Particle Physics and Astronomy
      Research Council. All Rights Reserved.
   }
}
\sstroutine{
   SCULIB\_FIT\_PLANE
}{
   fits a plane to some x, y, z data
}{
   \sstdescription{
      This routine does a least squares fit of data to a plane. The plane
      is described by an equation of form:-

\begin{equation}
           z = M_X * x  +  M_Y * y  +  C
\end{equation}
   }
   \sstinvocation{
      CALL SCULIB\_FIT\_PLANE (N, X, Y, Z, QUALITY, MX, MY, C, STATUS)
   }
   \sstarguments{
      \sstsubsection{
         N                      = INTEGER (Given)
      }{
         the number of data points
      }
      \sstsubsection{
         X (N)                  = REAL (Given)
      }{
         the x coords of the measured points
      }
      \sstsubsection{
         Y (N)                  = REAL (Given)
      }{
         the y coords of the measured points
      }
      \sstsubsection{
         Z (N)                  = REAL (Given)
      }{
         the measured values
      }
      \sstsubsection{
         QUALITY (N)            = INTEGER (Given)
      }{
         the quality on the data
      }
      \sstsubsection{
         MX                     = REAL (Returned)
      }{
         the slope of the fitted plane in the x-axis
      }
      \sstsubsection{
         MY                     = REAL (Returned)
      }{
         the slope of the fitted plane in the y-axis
      }
      \sstsubsection{
         C                      = REAL (Returned)
      }{
         the constant of the fitted plane
      }
      \sstsubsection{
         STATUS                 = INTEGER (Given and returned)
      }{
         global status
      }
   }
   \sstdiytopic{
      Method
   }{
      If status is good on entry the routine will loop through the
      input data calculating the sums required and the number of valid
      data points (i.e. those with good quality flags).

      If there are less than 3 valid data then
         an error message will be output and bad status returned
      else
         the plane coefficients will be calculated from the following
         formulae:-

\begin{equation}
M_X = \frac{\sum y^2(n_v\sum xz - \sum x \sum z) -
                (\sum y)^2 \sum xz  - n_v\sum xy \sum yz +
                \sum y(\sum xy \sum z  + \sum x \sum yz)
}{2\sum xy\sum x\sum y - n_v(\sum xy)^2 - (\sum x)^2\sum y^2
+ \sum x^2(n_v\sum y^2 + (\sum y)^2)}
\end{equation}

\begin{equation}
M_Y = \frac{n_v\sum yz - \sum y \sum z + M_X(\sum x \sum y - n_v\sum xy)}{
                          n_v\sum y^2 - \left(\sum y\right)^2}
\end{equation}


\begin{equation}
C = \frac{\sum z - M_X\sum x - M_Y\sum y}{n_v}
\end{equation}

         where $n_v$ is the number of valid data points. If the denominator of
         the expression for MX is 0 then an error message will be output and
         bad status returned.

      end if
   }
   \sstdiytopic{
      Copyright
   }{
      Copyright \copyright 1995,1996,1997,1998,1999 Particle Physics and Astronomy
      Research Council. All Rights Reserved.
   }
}
\sstroutine{
   SCULIB\_FIT\_SKYDIP
}{
   Fit the skydip data
}{
   \sstdescription{
      This routine fits a sub-instrument{\tt '}s measurements of the sky
      brightness at a range of airmasses to obtain the sky opacity, ETAtel
      and `b{\tt '} parameters. The fit is to the function:-

\begin{equation}
J_\mathrm{meas} = (1 - \eta_\mathrm{tel}) J_\mathrm{tel} +
                   \eta_\mathrm{tel} J_\mathrm{atm} -
                   b  \eta_\mathrm{tel} J_\mathrm{atm} e^{-A\tau},
\end{equation}

      Jtel is known and Jatm is related to Jamb by:-

 \begin{equation}
J_\mathrm{atm} = J_\mathrm{amb} \int_0^{40}\! A \left[k\exp\left(-\frac{h}{h_2}
\right)\exp\left[A k h_2 \left(\exp\left(-\frac{h}{h_2}\right)-1\right)\right]
\left(1-\frac{h}{h_1}\right)\right]\mathrm{d}h,
\end{equation}

      where h2 = scale height of absorption ( = 2km),
            h1 = coefficient to give 6.5K/km temperature drop in absorber,
             A = airmass,
             k = extinction.

      which relation has been fudged by Bill Duncan to:-

\begin{equation}
J_\mathrm{atm} = J_\mathrm{amb} X_\mathrm{g} \left[1-\exp\left(-A k h_2\right)\right],
\end{equation}
with $X_\mathrm{g}$ having the form
\begin{equation}
X_\mathrm{g} = 1 + \frac{h_2 T_\mathrm{lapse}}{T_\mathrm{amb}}\exp\left(-\frac{A \tau}{X_\mathrm{gconst}}\right)
\end{equation}

      with

            h1 = 2,
            h2 = -6.5 (note this h2 is defined differently to that in the
                       previous equation),
       $X_\mathrm{gconst}$ = 3.669383

      See `Calibration of mm and sub-mm Photometers by Skydipping{\tt '}, W.D.Duncan
      preprint and `Inversion of Sky Dips{\tt '}, SCU/WDD/31.1/1093 for further
      details.

         The fit can be made with ETAtel and/or b either fixed or
      varying. To allow one of these parameters to vary it should be
      input to the routine with a value below zero. If the input
      value is greater than zero then the routine will fix it at
      that for the fit.
   }
   \sstinvocation{
       CALL SCULIB\_FIT\_SKYDIP (CVAR, N\_MEASUREMENTS, AIRMASS, J\_MEASURED,
         J\_VARIANCE, SUB\_WAVELENGTH, SUB\_INSTRUMENT, SUB\_FILTER, T\_TEL,
         T\_AMB, ETA\_TEL\_IN, B\_IN, ETA\_TEL\_FIT, B\_FIT, TAUZ\_FIT,
         REXISQ, TAU\_ERROR, ETA\_ERROR, B\_ERROR, RESIDUAL, SIGMA, STATUS)
   }
   \sstarguments{
      \sstsubsection{
         CVAR                      = LOGICAL (Given)
      }{
         flag to govern whether to use a fixed variance (true)
         or the actual variance. The fixed variance is the mean
         of the actual variances.
      }
      \sstsubsection{
         N\_MEASUREMENTS            = INTEGER (Given)
      }{
         the number of SKYDIP measurements
      }
      \sstsubsection{
         AIRMASS (N\_MEASUREMENTS)  = REAL (Given)
      }{
         the airmasses at which the measurements were made
      }
      \sstsubsection{
         DATA (N\_MEASUREMENTS)     = REAL (Given)
      }{
         the measured sky brightness temperatures
      }
      \sstsubsection{
         VARIANCE (N\_MEASUREMENTS) = REAL (Given)
      }{
         the variance on DATA
      }
      \sstsubsection{
         SUB\_WAVELENGTH            = REAL (Given)
      }{
         the wavelength of the measurements
      }
      \sstsubsection{
         SUB\_INSTRUMENT            = CHARACTER$*$($*$) (Given)
      }{
         the name of the sub-instrument used
      }
      \sstsubsection{
         SUB\_FILTER                = CHARACTER$*$($*$) (Given)
      }{
         the name of the filter used
      }
      \sstsubsection{
         T\_TEL                     = REAL (Given)
      }{
         the telescope temperature (K)
      }
      \sstsubsection{
         T\_AMB                     = REAL (Given)
      }{
         the ambient temperature (K)
      }
      \sstsubsection{
         ETA\_TEL\_IN                = REAL (Given)
      }{
         if $>$= 0 then this will the ETAtel assumed in the fit.
         if $<$ 0 then ETAtel will be allowed to vary in the fit.
      }
      \sstsubsection{
         B\_IN                      = REAL (Given)
      }{
         if $>$=0 then this value of b will be assumed in the fit.
         if $<$ 0 then b will be allowed to vary in the fit.
      }
      \sstsubsection{
         ETA\_TEL\_FIT               = REAL (Returned)
      }{
         the result for ETAtel
      }
      \sstsubsection{
         B\_FIT                     = REAL (Returned)
      }{
         the result for b
      }
      \sstsubsection{
         TAUZ\_FIT                  = REAL (Returned)
      }{
         the fitted result for tauz
      }
      \sstsubsection{
         REXISQ                    = REAL (Returned)
      }{
         the reduced chi square of the fit
      }
      \sstsubsection{
         TAU\_ERROR                 = REAL (Returned)
      }{
         error in the tau
      }
      \sstsubsection{
         ETA\_ERROR                 = REAL (Returned)
      }{
         error in eta\_tel
      }
      \sstsubsection{
         B\_ERROR                   = REAL (Returned)
      }{
         error in B
      }
      \sstsubsection{
         RESIDUAL                     = DOUBLE (Returned)
      }{
        Absolute difference between the model and the fit.
      }
      \sstsubsection{
         SIGMA                     = DOUBLE (Returned)
      }{
         standard deviation of the difference between the
         model and the fit.
      }
      \sstsubsection{
         STATUS                    = INTEGER (Given and returned)
      }{
         Global status
      }
   }
}
\sstroutine{
   SCULIB\_FIX\_SCAN\_V10
}{
   Correct scan positions for version 1.0 data from SCUCD
}{
   \sstdescription{
      For data taken with version 1.0 of SCUCD the positions of
      the scan ends were incorrect. This was because of a bug
      in the software that calculated the scan ends concerning
      the use of arcseconds instead of radians.
      This routine attempts to recreate the bug so that the
      correct scan offsets can be calculated.
   }
   \sstinvocation{
       CALL SCULIB\_FIX\_SCAN\_V10(CENTRE\_COORDS, LAT\_OBS, LONG, LAT,
            MJD, RA\_OFF\_START, DEC\_OFF\_START, RA\_OFF\_END,
            DEC\_OFF\_END, RA\_NEW\_START, DEC\_NEW\_START,
            RA\_NEW\_END, DEC\_NEW\_END, STATUS)
   }
   \sstarguments{
      \sstsubsection{
         CENTRE\_COORDS = CHAR (Given)
      }{
         The coordinate system of the map centre
      }
      \sstsubsection{
         LAT\_OBS                = DOUBLE PRECISION (Given)
      }{
         Latitude of observatory (radians)
      }
      \sstsubsection{
         LONG = DOUBLE (Given)
      }{
         The longitude (ra) of the map centre (radians)
      }
      \sstsubsection{
         LAT = DOUBLE (Given)
      }{
         The latitude (dec) of the map centre (radians)
      }
      \sstsubsection{
         MJD = DOUBLE (Given)
      }{
         The modified Julian date at which the data was taken
      }
      \sstsubsection{
         RA\_OFF\_START = DOUBLE (Given)
      }{
         The RA start position of the scan as stored in the header (RD)
      }
      \sstsubsection{
         DEC\_OFF\_START = DOUBLE (Given)
      }{
         The Dec start position of the scan as stored in the header (RD)
      }
      \sstsubsection{
         RA\_OFF\_END = DOUBLE (Given)
      }{
         The RA end position of the scan as stored in the header (RD)
      }
      \sstsubsection{
         DEC\_OFF\_END = DOUBLE (Given)
      }{
         The DEC end position of the scan as stored in the header (RD)
      }
      \sstsubsection{
         RA\_NEW\_START = DOUBLE (Returned)
      }{
         The corrected RA start position of the scan  (RD)
      }
      \sstsubsection{
         DEC\_NEW\_START = DOUBLE (Returned)
      }{
         The corrected Dec start position of the scan  (RD)
      }
      \sstsubsection{
         RA\_NEW\_END = DOUBLE (Returned)
      }{
         The corrected RA end position of the scan  (RD)
      }
      \sstsubsection{
         RA\_NEW\_END = DOUBLE (Returned)
      }{
         The corrected Dec end position of the scan  (RD)
      }
      \sstsubsection{
         STATUS = INTEGER (Given \& Returned)
      }{
      }
   }
   \sstnotes{
      The scan ends are all stored as apparent RA/Decs.
   }
   \sstdiytopic{
      Copyright
   }{
      Copyright \copyright 1995,1996,1997,1998,1999 Particle Physics and Astronomy
      Research Council. All Rights Reserved.
   }
}
\sstroutine{
   SCULIB\_FLATFIELD\_DATA
}{
   flatfield the data in an array
}{
   \sstdescription{
      This routine multiplies the data in the input array by the flatfield
      values of the bolometers that took it. Output quality will be set
      bad if input quality was bad or the quality of the flatfield for that
      bolometer was bad.
   }
   \sstinvocation{
       CALL SCULIB\_FLATFIELD\_DATA (N\_BOL, N\_POS, N\_BEAM, BDATA, VARIANCE,
         QUALITY, BOL\_CHAN, BOL\_ADC, NUM\_CHAN, NUM\_ADC, BOL\_FLAT,
         BOL\_QUALITY, STATUS)
   }
   \sstarguments{
      \sstsubsection{
         N\_BOL                  = INTEGER (Given)
      }{
         number of bolometers measured
      }
      \sstsubsection{
         N\_POS                  = INTEGER (Given)
      }{
         number of positions at which the bolometers were measured
      }
      \sstsubsection{
         N\_BEAM                 = INTEGER (Given)
      }{
         number of beams data have been reduced into
      }
      \sstsubsection{
         BDATA (N\_BOL, N\_POS, N\_BEAM)
      }{
                          = REAL (Given and returned)
         measured data
      }
      \sstsubsection{
         VARIANCE (N\_BOL, N\_POS, N\_BEAM)
      }{
                          = REAL (Given and returned)
         variance on BDATA
      }
      \sstsubsection{
         QUALITY (N\_BOL, N\_POS, N\_BEAM)
      }{
                          = BYTE (Given and returned)
         quality on BDATA
      }
      \sstsubsection{
         BOL\_CHAN (N\_BOL)       = INTEGER (Given)
      }{
         the channel numbers of the measured bolometers
      }
      \sstsubsection{
         BOL\_ADC (N\_BOL)        = INTEGER (Given)
      }{
         the A/D numbers of the measured bolometers
      }
      \sstsubsection{
         NUM\_CHAN               = INTEGER (Given)
      }{
         number of channels per A/D
      }
      \sstsubsection{
         NUM\_ADC                = INTEGER (Given)
      }{
         number of A/D cards
      }
      \sstsubsection{
         BOL\_FLAT (NUM\_CHAN, NUM\_ADC)
      }{
                          = REAL (Given)
         the flatfield values
      }
      \sstsubsection{
         BOL\_QUALITY (NUM\_CHAN, NUM\_ADC)
      }{
                          = INTEGER (Given)
         quality on BOL\_FLAT
      }
      \sstsubsection{
         STATUS                 = INTEGER (Given and returned)
      }{
         global status
      }
   }
   \sstdiytopic{
      Copyright
   }{
      Copyright \copyright 1995,1996,1997,1998,1999 Particle Physics and Astronomy
      Research Council. All Rights Reserved.
   }
}
\sstroutine{
   SCULIB\_FLATFIELD\_SEQUENCE
}{
   get bolometer measurement sequence for FLATFIELD
}{
   \sstdescription{
      This routine works out the bolometer measurement sequence for a FLATFIELD
      observation. Only one of SCUBA{\tt '}s sub-instruments can be flat-fielded at
      a time. If the sub-instrument is one of the arrays then the measurement
      sequence will be ref-bol-ref-bol-......-bol-ref, where the `bol{\tt '}s are the
      bolometers to be measured and `ref{\tt '} the reference bolometer on the array.
      If the sub-instrument is one of the photometry pixels then only that
      bolometer will be measured.

         Whatever bolometer is the target of the measurement, data will be
      taken from all A/D channels.

         If status is good on entry SCULIB\_BOLSELECT is called to decode
      the bolometers to be measured and the sub-instruments involved. If
      the bolometers belong to more than one sub-instrument an error will
      be reported and bad status returned.

         If the sub-instrument is one of the arrays then the name of the
      reference bolometer will be read from parameters LONGREF\_BOL or
      SHORTREF\_BOL as appropriate. The number of measurements and their
      sequence is set as described above in N\_MEASUREMENTS, FLAT\_CHAN and
      FLAT\_ADC. If the sub-instrument is not one of the arrays then the
      reference bolometer will not be used and the bolometers(s) will be
      measured in sequence.

         For each measurement of a target bolometer data will be taken from
      all data channels; BOLS\_MEASURED is set to `ALL{\tt '} and SCULIB\_BOLSELECT
      called to decode this to channel and ADC numbers. A check is made
      that the reference bolometer, if used, is among those being measured.
      If not, an error message will be output and bad status returned.

         Lastly, the FLAT\_INDEX array is set so that it points to the
      position in the datablock of data from the target bolometer at each
      measurement.
   }
   \sstinvocation{
       CALL SCULIB\_FLATFIELD\_SEQUENCE (BOLOMETERS, NUM\_CHAN,
         NUM\_ADC, NUM\_SUB, BOL\_TYPE, BOL\_CALB, BOL\_DU3, BOL\_DU4,
         BOL\_QUAL, BOL\_ENABLED, BOLS\_MEASURED, N\_BOLS, BOL\_SELECT\_CHAN,
         BOL\_SELECT\_ADC, N\_SUBS, SUB\_INSTRUMENT, FLATREF\_CHAN,
         FLATREF\_ADC, N\_MEASUREMENTS, FLAT\_CHAN, FLAT\_ADC, FLAT\_INDEX,
         STATUS)
   }
   \sstarguments{
      \sstsubsection{
         BOLOMETERS                  = CHARACTER$*$($*$)
      }{
         the bolometers to be measured in this FLATFIELD
      }
      \sstsubsection{
         NUM\_CHAN                    = INTEGER (Given)
      }{
         the number of channels per A/D
      }
      \sstsubsection{
         NUM\_ADC                     = INTEGER (Given)
      }{
         the number of A/Ds
      }
      \sstsubsection{
         NUM\_SUB                     = INTEGER (Given)
      }{
         the number of sub-instruments in SCUBA
      }
      \sstsubsection{
         BOL\_TYPE (NUM\_CHAN,NUM\_ADC) = CHARACTER$*$($*$) (Given)
      }{
         the type of each bolometer
      }
      \sstsubsection{
         BOL\_CALB (NUM\_CHAN,NUM\_ADC) = REAL (Given)
      }{
         the flat-field factor for each bolometer
      }
      \sstsubsection{
         BOL\_DU3 (NUM\_CHAN,NUM\_ADC)  = REAL (Given)
      }{
         the dU3 coord of each bolometer
      }
      \sstsubsection{
         BOL\_DU4 (NUM\_CHAN,NUM\_ADC)  = REAL (Given)
      }{
         the dU4 coord of each bolometer
      }
      \sstsubsection{
         BOL\_QUAL (NUM\_CHAN,NUM\_ADC) = INTEGER (Given)
      }{
         the quality of each bolometer
      }
      \sstsubsection{
         BOL\_ENABLED (NUM\_CHAN,NUM\_ADC) = LOGICAL (Returned)
      }{
         .TRUE. if a bolometer is to be measured
      }
      \sstsubsection{
         BOLS\_MEASURED               = CHARACTER$*$($*$) (Returned)
      }{
         the bolometers to be measured for each FLATFIELD measurement
      }
      \sstsubsection{
         N\_BOLS                      = INTEGER (Returned)
      }{
         the number of bolometers to be measured at each measurement
      }
      \sstsubsection{
         BOL\_SELECT\_CHAN (NUM\_CHAN $*$ NUM\_ADC) = INTEGER (Returned)
      }{
         the channel numbers of the bolometers to be measured at each
         measurement
      }
      \sstsubsection{
         BOL\_SELECT\_ADC (NUM\_CHAN $*$ NUM\_ADC) = INTEGER (Returned)
      }{
         the A/D numbers of the bolometers to be measured at each measurement
      }
      \sstsubsection{
         N\_SUBS                      = INTEGER (Returned)
      }{
         the number of sub-instruments being used
      }
      \sstsubsection{
         SUB\_INSTRUMENT (NUM\_SUB)    = CHARACTER$*$($*$) (Returned)
      }{
         name of sub-instrument being used
      }
      \sstsubsection{
         FLATREF\_CHAN                = INTEGER (Returned)
      }{
         channel number of reference bolometer
      }
      \sstsubsection{
         FLATREF\_ADC                 = INTEGER (Returned)
      }{
         A/D number of reference bolometer
      }
      \sstsubsection{
         N\_MEASUREMENTS              = INTEGER (Returned)
      }{
         the number of measurements in this FLATFIELD observation
      }
      \sstsubsection{
         FLAT\_CHAN (2 $*$ NUM\_CHAN $*$ NUM\_ADC) = INTEGER (Returned)
      }{
         the channel number of the target bolometer in each measurement
      }
      \sstsubsection{
         FLAT\_ADC (2 $*$ NUM\_CHAN $*$ NUM\_ADC) = INTEGER (Returned)
      }{
         the A/D number of the target bolometer in each measurement
      }
      \sstsubsection{
         FLAT\_INDEX (2 $*$ NUM\_CHAN $*$ NUM\_ADC) = INTEGER (Returned)
      }{
         the index in the datablock of the target bolometer in each
         measurement
      }
      \sstsubsection{
         STATUS                      = INTEGER (Given and returned)
      }{
         global status
      }
   }
   \sstdiytopic{
      Copyright
   }{
      Copyright \copyright 1995,1996,1997,1998,1999 Particle Physics and Astronomy
      Research Council. All Rights Reserved.
   }
}
\sstroutine{
   SCULIB\_FREE
}{
   release virtual memory
}{
   \sstdescription{
      This routine frees virtual memory obtained by SCULIB\_MALLOC.

      If status is bad on entry the routine will return immediately.

      If START\_PTR is not equal to 0 then:

      \sstitemlist{

         \sstitem
          The sentinel integers above and below the used piece of memory
            will be checked against the values they were set to by
            SCULIB\_MALLOC.

         \sstitem
          PSX\_FREE will be called to free the virtual memory. If
            that{\tt '}s successful START\_PTR and END\_PTR will be set to
            the bad value.

         \sstitem
          If either of the sentinel integers was corrupted then an
            error will be reported and bad status returned.
      }
   }
   \sstinvocation{
      CALL SCULIB\_FREE (NAME, START\_PTR, END\_PTR, STATUS)
   }
   \sstarguments{
      \sstsubsection{
         NAME                = CHARACTER$*$($*$) (Given)
      }{
         name associated with VM
      }
      \sstsubsection{
         START\_PTR           = INTEGER (Given and returned)
      }{
         pointer to beginning of virtual memory
      }
      \sstsubsection{
         END\_PTR             = INTEGER (Given and returned)
      }{
         pointer to end of virtual memory
      }
      \sstsubsection{
         STATUS              = INTEGER (Given and returned)
      }{
         global status
      }
   }
   \sstdiytopic{
      Copyright
   }{
      Copyright \copyright 1995,1996,1997,1998,1999 Particle Physics and Astronomy
      Research Council. All Rights Reserved.
   }
}
\sstroutine{
   SCULIB\_GAUSS\_WTINIT
}{
   Generate a weighting function for rebinning
}{
   \sstdescription{
      A Gaussian weighting function
      The look up table is initialised such that one can access the
      weight directly given the square of the distance between the input
      and output pixel.
      The scale length is taken to be radius at half-width-half-max
      The gaussian should go out to a radius of at least 3 scale
      lengths ($<$0.2\% of peak). Governed by the RADIUS argument
      The function is simply taken as
        F(x) = exp(-X$*$$*$2)
      where the scale size (FWHM) is at X = SQRT(loge(2))
   }
   \sstinvocation{
      CALL SCULIB\_GAUSS\_WTINIT (WTFN, RADIUS, RES, STATUS)
   }
   \sstarguments{
      \sstsubsection{
         WTFN ( RADIUS $*$ RADIUS $*$ RES $*$ RES $+$ 1) = REAL (Returned)
      }{
         The weighting function generated by this routine.
         The index corresponds to the square of the distance from the
         centre in scale units.
      }
      \sstsubsection{
         RADIUS = INTEGER (Given)
      }{
         Size of the weighting function in scale units.
      }
      \sstsubsection{
         RES = INTEGER (Given)
      }{
         Number of points per scale lenghth.
      }
      \sstsubsection{
         STATUS = INTEGER (Given \& Returned)
      }{
         Global status
      }
   }
   \sstdiytopic{
      Copyright
   }{
      Copyright \copyright 1995,1996,1997,1998,1999 Particle Physics and Astronomy
      Research Council. All Rights Reserved.
   }
}
\sstroutine{
   SCULIB\_GAUSSIAN\_XISQ
}{
   calculate chi-squared of Gaussian fit
}{
   \sstdescription{
      If entered with good status this routine calculates the chi-squared
      between a dataset and a Gaussian function. Data points with bad quality
      will be ignored, as will points with zero variance. A warning will be
      issued if any data points with zero variance are encountered. If no
      valid data points are found then an error will be reported and bad
      status returned. The data are passed in via common.
   }
   \sstinvocation{
      CALL SCULIB\_GAUSSIAN\_XISQ (XISQ, N, FIT, STATUS)
   }
   \sstarguments{
      \sstsubsection{
         XISQ                   = DOUBLE PRECISION (Returned)
      }{
         the chi-squared of the current fit
      }
      \sstsubsection{
         N                      = INTEGER (Given)
      }{
         the number of parameters being fit, should be 6
      }
      \sstsubsection{
         FIT(N)                 = DOUBLE PRECISION (Given)
      }{
         the fit parameters:-
         \sstitemlist{

            \sstitem
               FIT(1) = peak height

            \sstitem
               FIT(2) = length of the `a{\tt '} axis of the sigma ellipse

            \sstitem
               FIT(3) = length of the `b{\tt '} axis of the sigma ellipse

            \sstitem
               FIT(4) = the angle between the `a{\tt '} axis and
                          the x axis ($+$ve anticlockwise, radians)

            \sstitem
               FIT(5) = the x coord of the centre

            \sstitem
               FIT(6) = the y coord of the centre
         }
      }
      \sstsubsection{
         STATUS                 = INTEGER (Given and returned)
      }{
         global status
      }
   }
   \sstdiytopic{
      Copyright
   }{
      Copyright \copyright 1995,1996,1997,1998,1999 Particle Physics and Astronomy
      Research Council. All Rights Reserved.
   }
}
\sstroutine{
   SCULIB\_GAUSSJ
}{
   Numerical Recipes in Fortran routine for solution of
   linear equations by Gauss-Jordan elimination
}{
   \sstdescription{
      Linear equation solution by Gauss-Jordan elimination. A(1:N,1:N) is an
      input matrix stored in an array of physical dimensions NP by NP.
      B(1:N,1:M) is an input matrix containing the M right-hand side vectors,
      stored in an array of physical dimensions NP by MP. On output,
      A(1:N,1:N) is replaced by its matrix inverse and B(1:N,1:M) is replaced
      by the corresponding set of solution vectors. Parameter NMAX is the
      largest anticipated value of N.
      Copied from GAUSSJ on p.30 of Numerical Recipes in Fortran, with
      STATUS added.
   }
   \sstinvocation{
      CALL SCULIB\_GAUSSJ (A, N, NP, B, M, MP, STATUS)
   }
   \sstarguments{
      \sstsubsection{
         A ( NP, NP ) = REAL (Given \& Returned)
      }{
         Input matrix. On exit, contains its inverse.
      }
      \sstsubsection{
         N = INTEGER (Given)
      }{
         Required size of A ($<$ NP)
      }
      \sstsubsection{
         NP = INTEGER (Given)
      }{
         Dimensions of input matrix.
      }
      \sstsubsection{
         B ( NP, MP) = REAL (Given \& Returned)
      }{
         Input matrix containing M right-hand side vectors. On exit, contains
         the solution vectors.
      }
      \sstsubsection{
         M = INTEGER (Given)
      }{
         Required size of B. ($<$ MP)
      }
      \sstsubsection{
         MP = INTEGER (Given)
      }{
         Size of second dimension of B.
      }
      \sstsubsection{
         STATUS = INTEGER (Given \& Returned)
      }{
         Global status.
      }
   }
   \sstdiytopic{
      Copyright
   }{
      Copyright \copyright 1995,1996,1997,1998,1999 Particle Physics and Astronomy
      Research Council. All Rights Reserved.
   }
}
\sstroutine{
   SCULIB\_GENSYCONFN
}{
   Generate a convolution function which will set to zero the
   spatial frequencies with no signal in Dual Beam maps
}{
   \sstdescription{
      This routine will return the symmetric convolution function
      required for setting to zero the spatial frequencies that have zero
      sensitivity in the raw chop-scan data. This function is not described
      in Emerson et al., 1979. A\&A 96, 92, which (I think) describes an
      early version of the NOD2 algorithm before all the wrinkles had been
      properly understood.

      This convolution function has to set to zero those points in the FT of
      the raw data at zero frequency and harmonics of 1/chop spacing. In
      the Fourier domain this means multiplying by a function which is 1
      at all points apart from those specified, where it is zero. This is
      the sum of 2 functions, the first being 1 everywhere, the second
      being -1 at the points to be zeroed. The convolution function required
      is the sum of the inverse FTs of these two functions. The inverse
      FT of the first function is simply a delta function at the origin,
      that of the second is a series of negative spikes with the first at the
      origin and the others separated by the chop spacing.

      The convolution function must cover the map even when the centre of
      the function is at the left or right hand end of the map. Hence the
      convolution function must be twice the length of the raw map.

      Since the raw data is not sampled such that the chop spacing is an
      integer number of samples, the actual convolution function must be
      rebinned onto the sample mesh by sinc interpolation.

      Otherwise the only tricky part of the algorithm is the need to normalise
      the two component functions such that in the Fourier domain the
      addition of the 2 functions does result in zeroes at the desired points.
      This is most easily achieved by concentrating on the zero spatial
      frequency which is just the sum of all the points in each of the 2
      functions. The first function is a delta function set to 1, so the
      sum of that function is 1. The sum of the second function is that of
      all the points in the convolution function from its centre to one end,
      a length that corresponds to the map size. This is the length to be
      used because, in the Fourier domain, the width of the sample to be
      set to zero is 1/map\_size.
   }
   \sstinvocation{
      CALL SCULIB\_GENSYCONFN( BSEP, PIXSEP, NPIX, NCFN, CONF, STATUS )
   }
   \sstarguments{
      \sstsubsection{
         BSEP = REAL (Given)
      }{
         The beam separation in arcseconds
      }
      \sstsubsection{
         PIXSEP = REAL (Given)
      }{
         The pixel separation in arcseconds
      }
      \sstsubsection{
         NPIX = INTEGER (Given)
      }{
         The number of pixels in the x direction
      }
      \sstsubsection{
         UNBAL = REAL (Given)
      }{
         The relative amplitudes of the right and left hand beams
         UNBAL=amp(lhb)/abs(amp(rhb))
      }
      \sstsubsection{
         NCFN = INTEGER (Returned)
      }{
         The length of the convolution array
      }
      \sstsubsection{
         CONF( $*$ ) = REAL (Returned)
      }{
         The convolution function
      }
      \sstsubsection{
         STATUS = INTEGER (Given and returned)
      }{
         Global status
      }
   }
   \sstdiytopic{
      Copyright
   }{
      Copyright \copyright 1995,1996,1997,1998,1999 Particle Physics and Astronomy
      Research Council. All Rights Reserved.
   }
}
\sstroutine{
   SCULIB\_GET\_BOL\_DESC
}{
   Get the bolometer description arrays and check their dimensions
}{
   \sstdescription{
      This routine obtains the bolometer description arrays. ie:
      BOL\_TYPE, BOL\_ADC, BOL\_CHAN, BOL\_DU3 and BOL\_DU4 arrays.
      The dimensions of these arrays are checked and an error returned
      if necessary.
   }
   \sstinvocation{
       CALL SCULIB\_GET\_BOL\_DESC(LOC, NUM\_CHAN, NUM\_ADC, BOL\_TYPE,
            BOL\_DU3, BOL\_DU4, BOL\_ADC, BOL\_CHAN, STATUS)
   }
   \sstarguments{
      \sstsubsection{
         LOC = \_CHARACTER (Given)
      }{
         Locator to the SCUBA extension
      }
      \sstsubsection{
         NUM\_CHAN = \_INTEGER (Given)
      }{
         Number of SCUBA channels per ADC
      }
      \sstsubsection{
         NUM\_ADC = \_INTEGER (Given)
      }{
         Number of ADC cards
      }
      \sstsubsection{
         N\_BOL = \_INTEGER (Given)
      }{
         Number of bolometers
      }
      \sstsubsection{
         BOL\_TYPE = \_CHARACTER$*$20() (Returned)
      }{
         Bolometer types
      }
      \sstsubsection{
         BOL\_DU3  = \_REAL() (Returned)
      }{
         du3 Nasmyth coordinated of bolometers
      }
      \sstsubsection{
         BOL\_DU4  = \_REAL() (Returned)
      }{
         du4 Nasmyth coordinated of bolometers
      }
      \sstsubsection{
         BOL\_ADC  = \_INTEGER() (Returned)
      }{
         A/D numbers of bolometers measured
      }
      \sstsubsection{
         BOL\_CHAN  = \_INTEGER() (Returned)
      }{
         Channel numbers of bolometers measured
      }
      \sstsubsection{
         STATUS = \_INTEGER (Given \& Returned)
      }{
         Global status
      }
   }
   \sstnotes{
      This routine does not annul the locator.
   }
   \sstdiytopic{
      Prior Requirements
   }{
      The locator to the structure must already be available.
   }
   \sstdiytopic{
      Copyright
   }{
      Copyright \copyright 1995,1996,1997,1998,1999 Particle Physics and Astronomy
      Research Council. All Rights Reserved.
   }
}
\sstroutine{
   SCULIB\_GET\_DEM\_PNTR
}{
   Get the DEM\_PNTR array and its dimensions
}{
   \sstdescription{
      This routine obtains the dem\_pntr array and returns the
      dimensions in terms of switches, exposures, inegrations and
      measurements.
   }
   \sstinvocation{
       CALL SCULIB\_GET\_DEM\_PNTR(ACTDIM, LOC, DEM\_PNTR\_PTR,
            N\_SWITCHES, N\_EXPOSURES, N\_INTEGRATIONS, N\_MEASUREMENTS,
            STATUS)
   }
   \sstarguments{
      \sstsubsection{
         ACTDIM = \_INTEGER (Given)
      }{
         Number of expected dimensions
      }
      \sstsubsection{
         LOC = \_CHARACTER (Given)
      }{
         Locator to the SCUBA extension
      }
      \sstsubsection{
         DEM\_PNTR\_PTR = \_INTEGER (Returned)
      }{
         Pointer to location of DEM\_PTR array in memory
      }
      \sstsubsection{
         N\_SWITCHES = \_INTEGER (Returned)
      }{
         Number of switches indicated by DEM\_PNTR
      }
      \sstsubsection{
         N\_EXPOSURES = \_INTEGER (Returned)
      }{
         Number of exposures indicated by DEM\_PNTR
      }
      \sstsubsection{
         N\_INTEGRATIONS = \_INTEGER (Returned)
      }{
         Number of integrations indicated by DEM\_PNTR
      }
      \sstsubsection{
         N\_MEASUREMENTS = \_INTEGER (Returned)
      }{
         Number of measurements indicated by DEM\_PNTR
      }
      \sstsubsection{
         STATUS = \_INTEGER (Given \& Returned)
      }{
         Global status
      }
   }
   \sstnotes{
      This routine does not annul the locator.
   }
   \sstdiytopic{
      Prior Requirements
   }{
      The locator to the structure must already be available.
   }
   \sstdiytopic{
      Copyright
   }{
      Copyright \copyright 1995,1996,1997,1998,1999 Particle Physics and Astronomy
      Research Council. All Rights Reserved.
   }
}
\sstroutine{
   SCULIB\_GET\_FILENAME
}{
   Find the filename associated with a parameter
}{
   \sstdescription{
      This routine finds the name of the filename associated with
      a NDF parameter.
   }
   \sstinvocation{
      CALL SCULIB\_GET\_FILENAME(PARAM, FILENAME, STATUS)
   }
   \sstarguments{
      \sstsubsection{
         PARAM = CHAR (Given)
      }{
         Name of the parameter
      }
      \sstsubsection{
         FILENAME = CHAR (Returned)
      }{
         Name of the file associated with PARAM
      }
      \sstsubsection{
         STATUS = INTEGER (Given \& Returned)
      }{
         Global status
      }
   }
   \sstdiytopic{
      Copyright
   }{
      Copyright \copyright 1995,1996,1997,1998,1999 Particle Physics and Astronomy
      Research Council. All Rights Reserved.
   }
   \sstimplementationstatus{
      Uses SUBPAR
   }
}
\sstroutine{
   SCULIB\_GET\_FITS\_C
}{
   get the value of specified FITS character keyword
}{
   \sstdescription{
      This routine will get the value of a specified FITS character keyword
      held in the FITS extension of an NDF file. The FITS extension must
      have been read into the input array FITS before this routine is called.

         The routine assumes that each line in the FITS array will contain
      a string with format:-

       {\tt "}KEYWORD= {\tt '}VALUE{\tt '}           / this is a comment{\tt "}

      The string is extracted from between the quotes and the {\tt '}/{\tt '} is not
      required.

      It will search the input array for a line containing the required
      keyword and return VALUE. If the keyword is not found and error will
      be reported and bad status returned. If the keyword is found but the
      line does not conform to the above format an error will be reported
      and bad status returned.
   }
   \sstinvocation{
      CALL SCULIB\_GET\_FITS\_C (MAX\_FITS, N\_FITS, FITS, NAME,
         VALUE, STATUS)
   }
   \sstarguments{
      \sstsubsection{
         MAX\_FITS             = INTEGER (Given)
      }{
         the maximum number of items in the FITS array
      }
      \sstsubsection{
         N\_FITS               = INTEGER (Given)
      }{
         the actual number of items in the FITS array
      }
      \sstsubsection{
         FITS (MAX\_FITS)      = CHARACTER$*$($*$) (Given)
      }{
         array containing the FITS items
      }
      \sstsubsection{
         NAME                 = CHARACTER$*$($*$) (Given)
      }{
         the name of the FITS keyword whose value is required
      }
      \sstsubsection{
         VALUE                = CHARACTER$*$($*$) (Returned)
      }{
         the value of the FITS keyword
      }
      \sstsubsection{
         STATUS               = INTEGER (Given and returned)
      }{
         global status
      }
   }
   \sstdiytopic{
      Copyright
   }{
      Copyright \copyright 1995,1996,1997,1998,1999 Particle Physics and Astronomy
      Research Council. All Rights Reserved.
   }
}
\sstroutine{
   SCULIB\_GET\_FITS\_D
}{
   get the value of specified FITS double keyword
}{
   \sstdescription{
      This routine will get the value of a specified FITS double precision
      keyword held in the FITS extension of an NDF file. The FITS extension must
      have been read into the input array FITS before this routine is called.

         The routine assumes that each line in the FITS array will contain
      a string with format:-

       {\tt "}KEYWORD=VALUE           / this is a comment{\tt "}

      It will search the input array for a line containing the required
      keyword and return VALUE. If the keyword is not found an error will
      be reported and bad status returned. If the keyword is found but the
      line does not conform to the above format an error will be reported
      and bad status returned.
   }
   \sstinvocation{
       CALL SCULIB\_GET\_FITS\_D (MAX\_FITS, N\_FITS, FITS, NAME, VALUE,
         STATUS)
   }
   \sstarguments{
      \sstsubsection{
         MAX\_FITS             = INTEGER (Given)
      }{
         the maximum number of items in the FITS array
      }
      \sstsubsection{
         N\_FITS               = INTEGER (Given)
      }{
         the actual number of items in the FITS array
      }
      \sstsubsection{
         FITS (MAX\_FITS)      = CHARACTER$*$($*$) (Given)
      }{
         array containing the FITS items
      }
      \sstsubsection{
         NAME                 = CHARACTER$*$($*$) (Given)
      }{
         the name of the FITS keyword whose value is required
      }
      \sstsubsection{
         VALUE                = DOUBLE PRECISION (Returned)
      }{
         the value of the FITS keyword
      }
      \sstsubsection{
         STATUS               = INTEGER (Given and returned)
      }{
         global status
      }
   }
   \sstdiytopic{
      Copyright
   }{
      Copyright \copyright 1995,1996,1997,1998,1999 Particle Physics and Astronomy
      Research Council. All Rights Reserved.
   }
}
\sstroutine{
   SCULIB\_GET\_FITS\_I
}{
   get the value of specified FITS integer keyword
}{
   \sstdescription{
      This routine will get the value of a specified FITS integer keyword
      held in the FITS extension of an NDF file. The FITS extension must
      have been read into the input array FITS before this routine is called.

         The routine assumes that each line in the FITS array will contain
      a string with format:-

       {\tt "}KEYWORD=INTEGER\_VALUE           / this is a comment{\tt "}

      It will search the input array for a line containing the required
      keyword and return VALUE. If the keyword is not found an error will
      be reported and bad status returned. If the keyword is found but the
      line does not conform to the above format an error will be reported
      and bad status returned.
   }
   \sstinvocation{
       CALL SCULIB\_GET\_FITS\_I (MAX\_FITS, N\_FITS, FITS, NAME, VALUE,
         STATUS)
   }
   \sstarguments{
      \sstsubsection{
         MAX\_FITS             = INTEGER (Given)
      }{
         the maximum number of items in the FITS array
      }
      \sstsubsection{
         N\_FITS               = INTEGER (Given)
      }{
         the actual number of items in the FITS array
      }
      \sstsubsection{
         FITS (MAX\_FITS)      = CHARACTER$*$($*$) (Given)
      }{
         array containing the FITS items
      }
      \sstsubsection{
         NAME                 = CHARACTER$*$($*$) (Given)
      }{
         the name of the FITS keyword whose value is required
      }
      \sstsubsection{
         VALUE                = INTEGER (Returned)
      }{
         the value of the FITS keyword
      }
      \sstsubsection{
         STATUS               = INTEGER (Given and returned)
      }{
         global status
      }
   }
   \sstdiytopic{
      Copyright
   }{
      Copyright \copyright 1995,1996,1997,1998,1999 Particle Physics and Astronomy
      Research Council. All Rights Reserved.
   }
}
\sstroutine{
   SCULIB\_GET\_FITS\_L
}{
   get the value of specified FITS logical keyword
}{
   \sstdescription{
      This routine will get the value of a specified FITS logical keyword
      held in the FITS extension of an NDF file. The FITS extension must
      have been read into the input array FITS before this routine is called.

         The routine assumes that each line in the FITS array will contain
      a string with format:-

       {\tt "}KEYWORD=LOGICAL           / this is a comment{\tt "}

      It will search the input array for a line containing the required
      keyword and return VALUE. If the keyword is not found an error will
      be reported and bad status returned. If the keyword is found but the
      line does not conform to the above format an error will be reported
      and bad status returned.
   }
   \sstinvocation{
       CALL SCULIB\_GET\_FITS\_L (MAX\_FITS, N\_FITS, FITS, NAME, VALUE,
         STATUS)
   }
   \sstarguments{
      \sstsubsection{
         MAX\_FITS             = INTEGER (Given)
      }{
         the maximum number of items in the FITS array
      }
      \sstsubsection{
         N\_FITS               = INTEGER (Given)
      }{
         the actual number of items in the FITS array
      }
      \sstsubsection{
         FITS (MAX\_FITS)      = CHARACTER$*$($*$) (Given)
      }{
         array containing the FITS items
      }
      \sstsubsection{
         NAME                 = CHARACTER$*$($*$) (Given)
      }{
         the name of the FITS keyword whose value is required
      }
      \sstsubsection{
         VALUE                = LOGICAL (Returned)
      }{
         the value of the FITS keyword
      }
      \sstsubsection{
         STATUS               = INTEGER (Given and returned)
      }{
         global status
      }
   }
   \sstdiytopic{
      Copyright
   }{
      Copyright \copyright 1995,1996,1997,1998,1999 Particle Physics and Astronomy
      Research Council. All Rights Reserved.
   }
}
\sstroutine{
   SCULIB\_GET\_FITS\_R
}{
   get the value of specified FITS real keyword
}{
   \sstdescription{
      This routine will get the value of a specified FITS real keyword
      held in the FITS extension of an NDF file. The FITS extension must
      have been read into the input array FITS before this routine is called.

         The routine assumes that each line in the FITS array will contain
      a string with format:-

       {\tt "}KEYWORD=VALUE           / this is a comment{\tt "}

      It will search the input array for a line containing the required
      keyword and return VALUE. If the keyword is not found an error will
      be reported and bad status returned. If the keyword is found but the
      line does not conform to the above format an error will be reported
      and bad status returned.
   }
   \sstinvocation{
       CALL SCULIB\_GET\_FITS\_R (MAX\_FITS, N\_FITS, FITS, NAME, VALUE,
         STATUS)
   }
   \sstarguments{
      \sstsubsection{
         MAX\_FITS             = INTEGER (Given)
      }{
         the maximum number of items in the FITS array
      }
      \sstsubsection{
         N\_FITS               = INTEGER (Given)
      }{
         the actual number of items in the FITS array
      }
      \sstsubsection{
         FITS (MAX\_FITS)      = CHARACTER$*$($*$) (Given)
      }{
         array containing the FITS items
      }
      \sstsubsection{
         NAME                 = CHARACTER$*$($*$) (Given)
      }{
         the name of the FITS keyword whose value is required
      }
      \sstsubsection{
         VALUE                = REAL (Returned)
      }{
         the value of the FITS keyword
      }
      \sstsubsection{
         STATUS               = INTEGER (Given and returned)
      }{
         global status
      }
   }
   \sstdiytopic{
      Copyright
   }{
      Copyright \copyright 1995,1996,1997,1998,1999 Particle Physics and Astronomy
      Research Council. All Rights Reserved.
   }
}
\sstroutine{
   SCULIB\_GET\_JIGGLE
}{
   Get the jiggle parameters
}{
   \sstdescription{
      This routine obtains the parameters of the JIGGLE pattern:
      JIGL\_X, JIGL\_Y and some FITS parameters.
      The dimensions of these arrays are checked and an error returned
      if necessary.
   }
   \sstinvocation{
       CALL SCULIB\_GET\_JIGGLE(LOC, MAX\_JIGGLE, N\_FITS,
            FITS, JIGGLE\_COUNT, JIGGLE\_REPEAT, JIGGLE\_P\_SWITCH,
            SAMPLE\_PA, SAMPLE\_COORDS, JIGGLE\_X, JIGGLE\_Y,
            STATUS)
   }
   \sstarguments{
      \sstsubsection{
         LOC = \_CHARACTER (Given)
      }{
         Locator to the SCUBA extension
      }
      \sstsubsection{
         MAX\_JIGGLE = \_INTEGER (Given)
      }{
         Maximum number of jiggle positions
      }
      \sstsubsection{
         N\_FITS = \_INTEGER (Given)
      }{
         Size of FITS array
      }
      \sstsubsection{
         FITS  = \_CHARACTER (Given)
      }{
         FITS values
      }
      \sstsubsection{
         JIGGLE\_COUNT = \_INTEGER (Returned)
      }{
         Actual size of jiggle pattern
      }
      \sstsubsection{
         JIGGLE\_REPEAT = \_INTEGER (Returned)
      }{
         Number of times jiggle pattern is repeated in a switch
      }
      \sstsubsection{
         JIGGLE\_P\_SWITCH = \_INTEGER (Returned)
      }{
         Number of jiggles per switch
      }
      \sstsubsection{
         SAMPLE\_PA = \_REAL
      }{
         position angle of sample x axis relative to x axis of SAMPLE\_COORDS
      }
      \sstsubsection{
         SAMPLE\_COORDS = \_CHARACTER
      }{
         coordinate system of sample offsets
      }
      \sstsubsection{
         JIGGLE\_X = \_REAL (Returned)
      }{
         X jiggle positions
      }
      \sstsubsection{
         JIGGLE\_Y = \_REAL (Returned)
      }{
         Y jiggle positions
      }
      \sstsubsection{
         STATUS = \_INTEGER (Given \& Returned)
      }{
         Global status
      }
   }
   \sstnotes{
      This routine does not annul the locator.
   }
   \sstdiytopic{
      Prior Requirements
   }{
      The locator to the structure must already be available.
   }
   \sstdiytopic{
      Copyright
   }{
      Copyright \copyright 1995,1996,1997,1998,1999 Particle Physics and Astronomy
      Research Council. All Rights Reserved.
   }
}
\sstroutine{
   SCULIB\_GET\_LST\_STRT
}{
   Get the LST\_STRT array and check its dimensions
}{
   \sstdescription{
      This routine obtains the lst\_strt array and compares the
      dimensions with the number of exposures, inegrations and
      measurements. The number of switches is returned.
   }
   \sstinvocation{
       CALL SCULIB\_GET\_LST\_STRT(LOC, LST\_STRT\_PTR,
            N\_SWITCHES, N\_EXPOSURES, N\_INTEGRATIONS, N\_MEASUREMENTS,
            STATUS)
   }
   \sstarguments{
      \sstsubsection{
         LOC = \_CHARACTER (Given)
      }{
         Locator to the SCUBA extension
      }
      \sstsubsection{
         LST\_STRT\_PTR = \_INTEGER (Returned)
      }{
         Pointer to location of DEM\_PTR array in memory
      }
      \sstsubsection{
         N\_SWITCHES = \_INTEGER (Returned)
      }{
         Number of switches indicated by LST\_STRT
      }
      \sstsubsection{
         N\_EXPOSURES = \_INTEGER (Given)
      }{
         Actual number of exposures
      }
      \sstsubsection{
         N\_INTEGRATIONS = \_INTEGER (Given)
      }{
         Actual number of integrations
      }
      \sstsubsection{
         N\_MEASUREMENTS = \_INTEGER (Given)
      }{
         Actual number of measurements
      }
      \sstsubsection{
         STATUS = \_INTEGER (Given \& Returned)
      }{
         Global status
      }
   }
   \sstnotes{
      This routine does not annul the locator.
      The array must be unmapped before finishing the program.
   }
   \sstdiytopic{
      Prior Requirements
   }{
      The locator to the structure must already be available.
   }
   \sstdiytopic{
      Copyright
   }{
      Copyright \copyright 1995,1996,1997,1998,1999 Particle Physics and Astronomy
      Research Council. All Rights Reserved.
   }
}
\sstroutine{
   SCULIB\_GET\_MJD
}{
   Obtain the Modified Julian date from the UT stored in FITS
}{
   \sstdescription{
      This routine retrieves the UTSTART and UTDATE values from the FITS
      array and calculates the MJD from these. It can corrects for the
      time taken before beginning the observation by comparing the
      LST stored in the header (written when the observation begins and
      before the telescope gets on source) and the reference LST supplied to
      the routine (which is assumed to be the LST when the data acquisition
      starts).

   }
   \sstinvocation{
      CALL SCULIB\_GET\_MJD(N\_FITS, FITS, MJD, STATUS)
   }
   \sstarguments{
      \sstsubsection{
         N\_FITS = \_INTEGER (Given)
      }{
         Size of FITS array
      }
      \sstsubsection{
         FITS  = \_CHARACTER (Given)
      }{
         FITS values
      }
      \sstsubsection{
         LST\_REF = DOUBLE PRECISION (Given)
      }{
        Reference LST. This is the LST at which data acquisition
        begins. The MJD should be calculated for the start of
        data taking rather than the time the telescope begain its
        slew. This number is normally read from the LST\_STRT array
        and should be in radians. If a negative value or a zero
        value is given it is assumed that no correction should be applied.
        This does mean that LST\_REF can never be exactly zero but that
        is thought be a remarkably unlikely event.
      }
      \sstsubsection{
         MJD = DOUBLE PRECISION (Returned)
      }{
         Modified Julian date of observation
      }
      \sstsubsection{
         EPOCH = \_REAL (Returned)
      }{
         Epoch of observation
      }
      \sstsubsection{
        STARTUP\_TIME = REAL (Returned)
      }{
        The time taken, in seconds, for the observation to begin.
        If the time stored in the STSTART is larger than LST\_REF
        the assumption is that a day boundary has been crossed since
        the LST\_REF should always be more recent than STSTART (the
        data can not be taken before the observation starts). A startup
        time longer than 5 minutes or so should be treated with
        suspicion. If the startup time is longer than 10 minutes it is
        not used for the MJD calculation since this should never
        be the case. Set to zero if LST\_REF was negative.
      }
      \sstsubsection{
         STATUS = \_INTEGER (Given \& Returned)
      }{
         Global status
      }
   }
   \sstnotes{
     Uses the UTDATE and UTSTART and STSTART FITS headers.
     These must be present.
      \sstitemlist{

         \sstitem
           The date must be stored in UTDATE in format YYYY:M:D

         \sstitem
           The time must be stored in UTSTART in format HH:MM:SS.SS

        \sstitem
     The MJD is calculated for the start of data taking rather
     than simply using the values stored in UTSTART. This is because
     there is always a delay between writing the FITS headers and
     taking data.
      }
   }
   \sstdiytopic{
      Copyright
   }{
      Copyright \copyright 1995-2000 Particle Physics and Astronomy
      Research Council. All Rights Reserved.
   }
}
\sstroutine{
   SCULIB\_GET\_RASTER
}{
   Get the raster parameters
}{
   \sstdescription{
      This routine obtains the parameters of the SCAN/MAP raster:
      RA1, RA2, DEC1 and DEC2.
      The dimensions of these arrays are checked and an error returned
      if necessary.
   }
   \sstinvocation{
       CALL SCULIB\_GET\_RASTER(LOC, N\_SWITCHES, N\_EXPOSURES,
            N\_INTEGRATIONS, N\_MEASUREMENTS, RA1\_PTR, RA2\_PTR,
            DEC1\_PTR, DEC2\_PTR, STATUS)
   }
   \sstarguments{
      \sstsubsection{
         LOC = \_CHARACTER (Given)
      }{
         Locator to the SCUBA extension
      }
      \sstsubsection{
         N\_SWITCHES = \_INTEGER (Given)
      }{
         Actual number of switches
      }
      \sstsubsection{
         N\_EXPOSURES = \_INTEGER (Given)
      }{
         Actual number of exposures
      }
      \sstsubsection{
         N\_INTEGRATIONS = \_INTEGER (Given)
      }{
         Actual number of integrations
      }
      \sstsubsection{
         N\_MEASUREMENTS = \_INTEGER (Given)
      }{
         Actual number of measurements
      }
      \sstsubsection{
         RA1\_PTR = \_INTEGER (Given)
      }{
         Pointer to mapped RA1 array
      }
      \sstsubsection{
         RA2\_PTR = \_INTEGER (Given)
      }{
         Pointer to mapped RA2 array
      }
      \sstsubsection{
         DEC1\_PTR = \_INTEGER (Given)
      }{
         Pointer to mapped DEC1 array
      }
      \sstsubsection{
         DEC2\_PTR = \_INTEGER (Given)
      }{
         Pointer to mapped DEC2 array
      }
      \sstsubsection{
         STATUS = \_INTEGER (Given \& Returned)
      }{
         Global status
      }
   }
   \sstnotes{
      This routine does not annul the locator.
   }
   \sstdiytopic{
      Prior Requirements
   }{
      The locator to the structure must already be available.
   }
   \sstdiytopic{
      Copyright
   }{
      Copyright \copyright 1995,1996,1997,1998,1999 Particle Physics and Astronomy
      Research Council. All Rights Reserved.
   }
}
\sstroutine{
   SCULIB\_GET\_SUB\_BOLS
}{
   copy bolometers belonging to a particular sub-
   instrument from the input data array to the output
}{
   \sstdescription{
      This routine extracts data for bolometers belonging to a particular
      sub-instrument from an input data array which may contain data for
      several sub-instruments. The input array IN\_POINTER points to the
      indices in the first dimension of the input data array that contain
      the data of interest. IN\_POINTER should have been calculated by an
      earlier call to SCULIB\_CALC\_SUB\_BOLS.
   }
   \sstinvocation{
       CALL SCULIB\_GET\_SUB\_BOLS (N\_BOL\_IN, N\_POS, N\_BEAM, IN\_DATA,
         IN\_VARIANCE, IN\_QUALITY, N\_BOL\_OUT, IN\_POINTER, OUT\_DATA,
         OUT\_VARIANCE, OUT\_QUALITY, STATUS)
   }
   \sstarguments{
      \sstsubsection{
         N\_BOL\_IN                     = INTEGER (Given)
      }{
         number of bolometers in input array
      }
      \sstsubsection{
         N\_POS                        = INTEGER (Given)
      }{
         number of positions measured in input array
      }
      \sstsubsection{
         N\_BEAM                       = INTEGER (Given)
      }{
         the number of beams in the input array
      }
      \sstsubsection{
         IN\_DATA (N\_BOL\_IN,N\_POS,N\_BEAM)
      }{
                                = REAL (Given)
         input data array
      }
      \sstsubsection{
         IN\_VARIANCE (N\_BOL\_IN,N\_POS,N\_BEAM)
      }{
                                = REAL (Given)
         variance on IN\_DATA
      }
      \sstsubsection{
         IN\_QUALITY (N\_BOL\_IN,N\_POS,N\_BEAM)
      }{
                                = BYTE (Given)
         quality on IN\_DATA
      }
      \sstsubsection{
         N\_BOL\_OUT                    = INTEGER (Given)
      }{
         number of bolometers in output array
      }
      \sstsubsection{
         IN\_POINTER (N\_BOL\_OUT)       = INTEGER (Given)
      }{
         pointers from bolometers in output array to their indices in input
      }
      \sstsubsection{
         OUT\_DATA (N\_BOL\_OUT,N\_POS,N\_BEAM)
      }{
                                = REAL (Returned)
         output data array
      }
      \sstsubsection{
         OUT\_VARIANCE (N\_BOL\_OUT,N\_POS,N\_BEAM)
      }{
                                = REAL (Returned)
         variance on OUT\_DATA
      }
      \sstsubsection{
         OUT\_QUALITY (N\_BOL\_OUT,N\_POS,N\_BEAM)
      }{
                                = BYTE (Returned)
         quality on OUT\_DATA
      }
      \sstsubsection{
         STATUS                       = INTEGER (Given and returned)
      }{
         global status
      }
   }
   \sstdiytopic{
      Copyright
   }{
      Copyright \copyright 1995,1996,1997,1998,1999 Particle Physics and Astronomy
      Research Council. All Rights Reserved.
   }
}
\sstroutine{
   SCULIB\_GET\_SUB\_INST
}{
   Ask for the specific sub-instrument
}{
   \sstdescription{
      This routine finds which sub-instruments are present in the
      file. If there is only one this is returned, otherwise
      the user is given a choice.
   }
   \sstinvocation{
       CALL SCULIB\_GET\_SUB\_INST(PACKAGE, N\_FITS, FITS, PARAM,
            SUB\_POINTER, WAVE, INST, FILT, STATUS)
   }
   \sstarguments{
      \sstsubsection{
         PACKAGE = \_CHAR (Given)
      }{
         Name of software package for informational message
      }
      \sstsubsection{
         N\_FITS = \_INTEGER (Given)
      }{
         Number of FITS items
      }
      \sstsubsection{
         FITS(N\_FITS) = \_CHAR$*$80 (Given \& Returned)
      }{
         The FITS array
      }
      \sstsubsection{
         PARAM = \_CHAR (Given)
      }{
         Name of the ADAM parameter used to ask for sub instrument
      }
      \sstsubsection{
         N\_SUB = \_INTEGER (Returned)
      }{
         Number of sub instruments in file
      }
      \sstsubsection{
         SUB\_POINTER = \_INTEGER (Returned)
      }{
         Index of selected sub instrument
      }
      \sstsubsection{
         WAVE = \_REAL (Returned)
      }{
         Wavelength (microns) of selected sub instrument
      }
      \sstsubsection{
         INST = \_CHAR (Returned)
      }{
         Name of selected sub instrument
      }
      \sstsubsection{
         FILT = \_CHAR (Returned)
      }{
         Name of selected filter
      }
      \sstsubsection{
         STATUS = \_INTEGER (Given \& Returned)
      }{
         Global status
      }
   }
   \sstnotes{
      I actually need to include REDS\_SYS (for SCUBA\_\_MAX\_SUB) but I am
      doing it by hand at the moment.
   }
   \sstdiytopic{
      Copyright
   }{
      Copyright \copyright 1995,1996,1997,1998,1999 Particle Physics and Astronomy
      Research Council. All Rights Reserved.
   }
}
\sstroutine{
   SCULIB\_INSERT\_BOL
}{
   To insert a single bolometer into a data array
}{
   \sstdescription{
      This routine inserts a single bolometer into a SCUBA data array.
   }
   \sstinvocation{
       CALL SCULIB\_INSERT\_BOL(NBOL, N\_BOLS, N\_POS, BOLDATA,
            BOLQUAL, SCUDATA, SCUQUAL, STATUS)
   }
   \sstarguments{
      \sstsubsection{
         NBOL = INTEGER (Given)
      }{
         The bolometer number
      }
      \sstsubsection{
         N\_BOLS = \_INTEGER (Given)
      }{
         Total number of bolometers in data array
      }
      \sstsubsection{
         N\_POS = INTEGER (Given)
      }{
         Number of jiggle positions in data set
      }
      \sstsubsection{
         BOLDATA(N\_POS) = REAL (Given)
      }{
         The specified bolometer data
      }
      \sstsubsection{
         BOLQUAL(N\_POS) = \_BYTE (Given)
      }{
         The bolometer quality
      }
      \sstsubsection{
         SCUDATA(N\_BOLS, N\_POS) = REAL (Given \& Returned)
      }{
         The data
      }
      \sstsubsection{
         SCUQUAL(N\_BOLS, N\_POS) = BYTE (Given \&  Returned)
      }{
         The data quality
      }
      \sstsubsection{
         STATUS = INTEGER (Given and Returned)
      }{
         Global Status value
      }
   }
   \sstdiytopic{
      Copyright
   }{
      Copyright \copyright 1995,1996,1997,1998,1999 Particle Physics and Astronomy
      Research Council. All Rights Reserved.
   }
   \sstimplementationstatus{
      Propogates quality
   }
}
\sstroutine{
   SCULIB\_INTEGRATE\_PHOTOM\_JIGGLE
}{
   integrate the jiggle map made by a
   bolometer during a PHOTOM observation
}{
   \sstdescription{
      This routine just sums the data for the specified bolometer over the
      jiggle pattern.

        After checking status on entry the routine checks that the bolometer
      to be integrated is among those that were measured. All being well it
      then loops through the jiggle pattern summing the valid data and variance
      measured at each position for that bolometer. If no valid data were
      obtained then the result quality will be set bad.
   }
   \sstinvocation{
       CALL SCULIB\_INTEGRATE\_PHOTOM\_JIGGLE (BOL, N\_BOLS, J\_COUNT,
         DATA, VARIANCE, QUALITY, RESULT\_D, RESULT\_V, RESULT\_Q, STATUS)
   }
   \sstarguments{
      \sstsubsection{
         BOL                          = INTEGER (Given)
      }{
         the index of the bolometer whose data is to be integrated
      }
      \sstsubsection{
         N\_BOLS                       = INTEGER (Given)
      }{
         the number of bolometers being measured
      }
      \sstsubsection{
         J\_COUNT                      = INTEGER (Given)
      }{
         the number of jiggles in the pattern
      }
      \sstsubsection{
         DATA (N\_BOLS, J\_COUNT)       = REAL (Given)
      }{
         the measured data
      }
      \sstsubsection{
         VARIANCE (N\_BOLS, J\_COUNT)   = REAL (Given)
      }{
         the variance
      }
      \sstsubsection{
         QUALITY (N\_BOLS, J\_COUNT)    = INTEGER (Given)
      }{
         the quality
      }
      \sstsubsection{
         RESULT\_D                     = REAL (Returned)
      }{
         the sum of the input data
      }
      \sstsubsection{
         RESULT\_V                     = REAL (Returned)
      }{
         the sum of the input variances
      }
      \sstsubsection{
         RESULT\_Q                     = INTEGER (Returned)
      }{
         the quality on the sum
      }
      \sstsubsection{
         STATUS                       = INTEGER (Given and returned)
      }{
         global status
      }
   }
   \sstdiytopic{
      Copyright
   }{
      Copyright \copyright 1995,1996,1997,1998,1999 Particle Physics and Astronomy
      Research Council. All Rights Reserved.
   }
}
\sstroutine{
   SCULIB\_INVERT\_MATRIX
}{
   invert a square matrix
}{
   \sstdescription{
      This routine inverts a matrix by a method opaque to casual inspection
      but which has the advantage that it doesn{\tt '}t require any more space
      than the input matrix itself provides. The method is described in
      `Data Reduction and Error Analysis for the Physical Sciences{\tt '} by
      Bevington and Robinson, and the code is adapted from the Pascal
      version listed there.
   }
   \sstinvocation{
      CALL SCULIB\_INVERT\_MATRIX (M, ARRAY, DET, IK, JK, STATUS)
   }
   \sstarguments{
      \sstsubsection{
         M                      = INTEGER (Given)
      }{
         the dimensions of the matrix
      }
      \sstsubsection{
         ARRAY (M,M)            = DOUBLE PRECISION (Given and returned)
      }{
         the matrix to be inverted
      }
      \sstsubsection{
         DET                    = DOUBLE PRECISION (Returned)
      }{
         the determinant of the matrix
      }
      \sstsubsection{
         IK (M)                 = INTEGER (Scratch)
      }{
         i indices of largest array elements
      }
      \sstsubsection{
         JK (M)                 = INTEGER (Scratch)
      }{
         j indices of largest array elements
      }
      \sstsubsection{
         STATUS                 = INTEGER (Given and returned)
      }{
         global status
      }
   }
   \sstdiytopic{
      Copyright
   }{
      Copyright \copyright 1995,1996,1997,1998,1999 Particle Physics and Astronomy
      Research Council. All Rights Reserved.
   }
}
\sstroutine{
   SCULIB\_J\_THEORETICAL
}{
   calculate the theoretical sky brightness temperature
}{
   \sstdescription{
      This routine calculates the theoretical sky brightness temperature
      according to Bill Duncan{\tt '}s model:-


\begin{equation}
J_\mathrm{meas} = (1 - \eta_\mathrm{tel}) J_\mathrm{tel} +
                   \eta_\mathrm{tel} J_\mathrm{atm} -
                   b  \eta_\mathrm{tel} J_\mathrm{atm} e^{-A\tau},
\end{equation}

      Jtel is known and Jatm is related to Jamb by:-

 \begin{equation}
J_\mathrm{atm} = J_\mathrm{amb} \int_0^{40}\! A \left[k\exp\left(-\frac{h}{h_2}
\right)\exp\left[A k h_2 \left(\exp\left(-\frac{h}{h_2}\right)-1\right)\right]
\left(1-\frac{h}{h_1}\right)\right]\mathrm{d}h,
\end{equation}

      where h2 = scale height of absorption ( = 2km),
            h1 = coefficient to give 6.5K/km temperature drop in absorber,
             A = airmass,
             k = extinction.

      which relation has been fudged by Bill Duncan to:-

\begin{equation}
J_\mathrm{atm} = J_\mathrm{amb} X_\mathrm{g} \left[1-\exp\left(-A k h_2\right)\right],
\end{equation}
with $X_\mathrm{g}$ having the form
\begin{equation}
X_\mathrm{g} = 1 + \frac{h_2 T_\mathrm{lapse}}{T_\mathrm{amb}}\exp\left(-\frac{A \tau}{X_\mathrm{gconst}}\right)
\end{equation}

      with

            h1 = 2,
            h2 = -6.5 (note this h2 is defined differently to that in the
                       previous equation),
       $X_\mathrm{gconst}$ = 3.669383

      where Jatm is the effective brightness temperature of the sky.

      The routine will return immediately if entered with bad status.
      Errors will be reported and bad status returned if:-
      \sstitemlist{

         \sstitem
           TAUZ is less than 0

         \sstitem
           AIRMASS is less than 0

         \sstitem
           WAVELENGTH is less than or equal to 0

         \sstitem
           T\_AMB is less than or equal to 0

         \sstitem
           ETA\_TEL is outside the range 0 to 1

         \sstitem
           B is outside the range 0 to 1
      }
   }
   \sstinvocation{
       CALL SCULIB\_J\_THEORETICAL (TAUZ, AIRMASS, T\_TEL, T\_AMB,
         WAVELENGTH, ETA\_TEL, B, J\_THEORETICAL, STATUS)
   }
   \sstarguments{
      \sstsubsection{
         TAUZ                = REAL (Given)
      }{
         zenith sky optical depth
      }
      \sstsubsection{
         AIRMASS             = REAL (Given)
      }{
         airmass
      }
      \sstsubsection{
         T\_TEL               = REAL (Given)
      }{
         temperature of the telescope (K)
      }
      \sstsubsection{
         T\_AMB               = REAL (Given)
      }{
         temperature of ambient air (K)
      }
      \sstsubsection{
         WAVELENGTH          = REAL (Given)
      }{
         wavelength of interest (microns)
      }
      \sstsubsection{
         ETA\_TEL             = REAL (Given)
      }{
         telescope transmission
      }
      \sstsubsection{
         B                   = REAL (Given)
      }{
         bandwidth factor
      }
      \sstsubsection{
         J\_THEORETICAL       = REAL (Returned)
      }{
         the theoretical sky brightness temperature
      }
      \sstsubsection{
         STATUS              = INTEGER (Given and returned)
      }{
         global status
      }
   }
   \sstdiytopic{
      Copyright
   }{
      Copyright \copyright 1995,1996,1997,1998,1999 Particle Physics and Astronomy
      Research Council. All Rights Reserved.
   }
}
\sstroutine{
   SCULIB\_JNU
}{
   Calculate the Rayleigh-Jeans corrected brightness temperature
}{
   \sstdescription{
      This function calculates the Rayleigh-Jeans corrected brightness
      temperature of radiation at frequency NU ($\nu$) and temperature T.

\begin{equation}
\jmath_\nu = \frac{X T}{\exp(X) - 1}
\end{equation}
where
\begin{equation}
X = \frac{h\nu}{kT}
\end{equation}

      If the absolute value of X is less than 1e-4 JNU = T, or if the absolute
      value of X is greater than 20 JNU = 0.
   }
   \sstinvocation{
      JNU = SCULIB\_JNU (NU, T, STATUS)
   }
   \sstarguments{
      \sstsubsection{
         NU                = REAL (Given)
      }{
         frequency (Hz)
      }
      \sstsubsection{
         T                 = REAL (Given)
      }{
         temperature (K)
      }
      \sstsubsection{
         STATUS            = INTEGER (Given and returned)
      }{
         global status
      }
   }
   \sstreturnedvalue{
      \sstsubsection{
         SCULIB\_JNU = REAL
      }{
         Rayleigh-Jeans corrected brightness temperature
      }
   }
   \sstdiytopic{
      Copyright
   }{
      Copyright \copyright 1995,1996,1997,1998,1999 Particle Physics and Astronomy
      Research Council. All Rights Reserved.
   }
}
\sstroutine{
   SCULIB\_LINEAR\_WTINIT
}{
   Set up weighting function for linear rebinning
}{
   \sstdescription{
      This is a FORTRAN version of the C transputer code.
      Here is the C description:
      Initialise the array containing the weighting function for
      linear interpolation.
      The weighting function is a cone, going from 1.0 at the centre to
      0.0 at the edge. The look-up table represents a radial slice through the
      cone, tabulated to enable a look-up in terms of the square of the
      radius.
   }
   \sstinvocation{
      CALL SCULIB\_LINEAR\_WTINIT(WTFN, RES, STATUS)
   }
   \sstarguments{
      \sstsubsection{
         WTFN ( RES $*$ RES ) = REAL (Returned)
      }{
         The weighting function generated by this routine.
         The index corresponds to the square of the distance from the
         centre in scale units.
      }
      \sstsubsection{
         RES = INTEGER (Given)
      }{
         Number of points per scale lenghth.
      }
      \sstsubsection{
         STATUS = INTEGER (Given \& Returned)
      }{
         Global status
      }
   }
   \sstdiytopic{
      Copyright
   }{
      Copyright \copyright 1995,1996,1997,1998,1999 Particle Physics and Astronomy
      Research Council. All Rights Reserved.
   }
}
\sstroutine{
   SCULIB\_LST
}{
   returns LST in radians
}{
   \sstdescription{
      Calculates the current LST in radians.
   }
   \sstinvocation{
      LST = SCULIB\_LST
   }
   \sstarguments{
      \sstsubsection{
         None
      }{
      }
   }
   \sstreturnedvalue{
      \sstsubsection{
         LST = DOUBLE PRECISION
      }{
         LST in radians
      }
   }
   \sstdiytopic{
      Copyright
   }{
      Copyright \copyright 1995,1996,1997,1998,1999 Particle Physics and Astronomy
      Research Council. All Rights Reserved.
   }
}
\sstroutine{
   SCULIB\_MALLOC
}{
   get virtual memory
}{
   \sstdescription{
      This routine gets SIZE bytes of virtual memory.

      If status is bad on entry the routine will return immediately.

      If START\_PTR is not equal to 0 then
      \sstitemlist{

         \sstitem
           It{\tt '}s possible that VM may have already been allocated to
            START\_PTR so issue an error message and set status bad.

      }
      else

      \sstitemlist{

         \sstitem
          Call PSX\_MALLOC to allocate the required memory plus space for
            2 integers at either end.

         \sstitem
          If status is good then
               Set START\_PTR and END\_PTR to point to the first and last
               bytes in the section of memory to be used and
               call SCULIB\_CFILLI to set the sentinel integers at either
               end of the block to 37. SCULIB\_FREE will check these for
               corruption when the time comes to free the VM.

      }
      end if
   }
   \sstinvocation{
      CALL SCULIB\_MALLOC (SIZE, START\_PTR, END\_PTR, STATUS)
   }
   \sstarguments{
      \sstsubsection{
         SIZE                = INTEGER (Given)
      }{
         number of bytes required
      }
      \sstsubsection{
         START\_PTR           = INTEGER (Given and returned)
      }{
         pointer to beginning of virtual memory
      }
      \sstsubsection{
         END\_PTR             = INTEGER (Returned)
      }{
         pointer to end of virtual memory
      }
      \sstsubsection{
         STATUS              = INTEGER (Given and returned)
      }{
         global status
      }
   }
   \sstdiytopic{
      Copyright
   }{
      Copyright \copyright 1995,1996,1997,1998,1999 Particle Physics and Astronomy
      Research Council. All Rights Reserved.
   }
}
\sstroutine{
   SCULIB\_MAP\_ALLAN\_VARIANCE
}{
   incorporate latest set of MAP demodulated data
   into Allan variance
}{
   \sstdescription{
      This routine incorporates a data slice into a run of data and
      updates the Allan variance of the run.
         The Allan variance is calculated for a range of simulated integration
      times for which artificial samples are calculated from the input data.
      The Allan variance for a particular integration time is calculated from:-

\begin{equation}
variance = \sum_{i=1\rightarrow n-1}\frac{(sample(i) - sample(i+1))^2}{2n}
\end{equation}

      where n is the number of artificial samples available.
   }
   \sstinvocation{
       CALL SCULIB\_MAP\_ALLAN\_VARIANCE (DEMOD, N\_BOLS, N\_SAMPLES,
         N\_ALLAN\_BOL, ALLAN\_BOL, KMAX, SUM, N\_SUM, ARTIFICIAL,
         N\_ARTIFICIAL, ALLAN\_VARIANCE, ALLAN\_QUALITY, STATUS)
   }
   \sstarguments{
      \sstsubsection{
         DEMOD (4, N\_BOLS, N\_SAMPLES)    = REAL (Given)
      }{
         the demodulated data (data,-,-,quality)
      }
      \sstsubsection{
         N\_BOLS                          = INTEGER (Given)
      }{
         the number of bolometers taking data
      }
      \sstsubsection{
         N\_SAMPLES                       = INTEGER (Given)
      }{
         the number of samples taken for each bolometer
      }
      \sstsubsection{
         N\_ALLAN\_BOL                     = INTEGER (Given)
      }{
         the number of bolometers whose data is to be averaged together
         at each sample position
      }
      \sstsubsection{
         ALLAN\_BOL (N\_ALLAN\_BOL)         = INTEGER (Given)
      }{
         the positions in the demodulated data array of the bolometers
         whose measurements are to be used for the calculation of Allan
         variance
      }
      \sstsubsection{
         KMAX                            = INTEGER (Given)
      }{
         the size of the Allan variance array
      }
      \sstsubsection{
         SUM (KMAX)                      = REAL (Given and returned)
      }{
         the accumulator for current artificial sample of length K real
         samples
      }
      \sstsubsection{
         N\_SUM (KMAX)                    = INTEGER (Given and returned)
      }{
         the number of samples currently in the artificial sample
         accumulator
      }
      \sstsubsection{
         ARTIFICIAL (KMAX)               = REAL (Given and returned)
      }{
         the last artificial sample of length K samples that was
         calculated
      }
      \sstsubsection{
         N\_ARTIFICIAL (KMAX)             = INTEGER (Given and returned)
      }{
         the number of artificial samples of length K samples that
         have been calculated
      }
      \sstsubsection{
         ALLAN\_VARIANCE (KMAX)           = REAL (Given and returned)
      }{
         the Allan variance
      }
      \sstsubsection{
         ALLAN\_QUALITY (KMAX)            = INTEGER (Returned)
      }{
         quality on the Allan variance
      }
      \sstsubsection{
         STATUS                          = INTEGER (Given and returned)
      }{
         global status
      }
   }
   \sstdiytopic{
      Method
   }{
      If status on entry is good the routine will:-

      loop through the data samples to be incorporated -

         loop through the bolometers whose data is to be averaged into
         this sample -

            if the bolometer was measured -
               if the measurement had good quality -

                  add the measurement to the average for this sample

               end if
            end if

         end of loop

         if any bolometers had good data for this sample -

            calculate the sample average

            loop through the Allan variances -

               add the sample to the buffer used to calculate the latest
               artificial sample at the simulated integration time of this
               variance

               if all the data has been obtained for the latest artificial
               sample -

                  calculate the artificial sample

                  if there is only one artificial sample for this simulated
                  integration time we can{\tt '}t calculate the Allan variance, so
                  set its quality to bad

                  if there are 2 artifical samples then the Allan variance
                  can be calculated from -

\begin{equation}
variance = \frac{(artificial\_2 - artificial\_1)^2}{2 \times 2}
\end{equation}

                  if there are more than 2 artificial samples

                     the sum of the squares of the differences between the
                     previous samples is recovered from the current value of the
                     Allan variance

                     the square of the difference between the current artificial
                     sample and the previous one is added to the sum

                     the Allan variance is re-calculated

                  end if

                  the current artifical sample is stored to be used as the
                  `previous{\tt '} sample next time round

                  the variables used to calculate the artificial samples at
                  this simulated integration time are reset

               end if

            end of loop through Allan variances

         end if

      end of loop through samples in this dataslice
   }
   \sstdiytopic{
      Copyright
   }{
      Copyright \copyright 1995,1996,1997,1998,1999 Particle Physics and Astronomy
      Research Council. All Rights Reserved.
   }
}
\sstroutine{
   SCULIB\_MASK\_DATA
}{
   Set a data array from a SCUBA section
}{
   \sstdescription{
      From an array of SCUBA sections, this routine creates a mask
      which is then used to modify the given data array
   }
   \sstinvocation{
      CALL SCULIB\_MASK\_DATA( STATUS )
   }
   \sstarguments{
      \sstsubsection{
         USE\_SECT   = LOGICAL (Given)
      }{
         Are we using the section (TRUE) or inverse (FALSE)
      }
      \sstsubsection{
         TYPE = CHAR (Given)
      }{
         Type of data. Allowed values are:
           REAL - data is real array
           BYTE - data is byte array
           BIT  - affect single bit of byte array
      }
      \sstsubsection{
         N\_SPEC                           = INTEGER (Given)
      }{
         Number of specifications supplied in SPEC
      }
      \sstsubsection{
         SPEC( N\_SPEC )                   = CHARACTER$*$($*$) (Given)
      }{
         the specification to be decoded
      }
      \sstsubsection{
         DEMOD\_POINTER (N\_SWITCHES, N\_EXPOSURES, N\_INTEGRATIONS,  N\_MEASUREMENTS)
      }{
      }
      \sstsubsection{
         = INTEGER (Given)
      }{
         the pointer to the location in the main data array of the data
         for each switch of the observation
      }
      \sstsubsection{
         N\_SWITCHES                       = INTEGER (Given)
      }{
         the number of switches per exposure
      }
      \sstsubsection{
         N\_EXPOSURES                      = INTEGER (Given)
      }{
         the number of exposures per integration
      }
      \sstsubsection{
         N\_INTEGRATIONS                   = INTEGER (Given)
      }{
         the number of integrations per measurement
      }
      \sstsubsection{
         N\_MEASUREMENTS                   = INTEGER (Given)
      }{
         the number of measurements in the observation
      }
      \sstsubsection{
         N\_POS                            = INTEGER (Given)
      }{
         the number of positions measured in the observation
      }
      \sstsubsection{
         N\_BOLS                           = INTEGER (Given)
      }{
         the number of bolometers measured in the observation
      }
      \sstsubsection{
         N\_BEAM = INTEGER (Given)
      }{
         Number of beams in DATA\_PTR
      }
      \sstsubsection{
         SWITCH\_EXPECTED                  = LOGICAL (Given)
      }{
         .TRUE. if a switch component is allowed in the data-spec
      }
      \sstsubsection{
         VALUE   = REAL (Given)
      }{
         If TYPE is {\tt '}REAL{\tt '} then use this value for masked data
      }
      \sstsubsection{
         BVALUE   = BYTE (Given)
      }{
         If TYPE is {\tt '}BYTE{\tt '} then use this value for masked data
      }
      \sstsubsection{
         BITNUM   = INTEGER (Given)
      }{
         If TYPE is {\tt '}BIT{\tt '} then affect this bit position
      }
      \sstsubsection{
         BIT\_STATE = LOGICAL (Given)
      }{
         If TYPE is {\tt '}BIT{\tt '} then clear the bit (FALSE) or set it (TRUE)
      }
      \sstsubsection{
         DATA\_PTR = INTEGER (Given \& Returned)
      }{
         Pointer to data array that is to be masked
      }
      \sstsubsection{
         STATUS                           = INTEGER (Given and Returned)
      }{
         The global status.
      }
   }
   \sstdiytopic{
      Copyright
   }{
      Copyright \copyright 1995,1996,1997,1998,1999 Particle Physics and Astronomy
      Research Council. All Rights Reserved.
   }
}
\sstroutine{
   SCULIB\_MJD\_TO\_DATEOBS
}{
   Convert a modified julian date to the correct FITS DATE-OBS format
}{
   \sstdescription{
      Given a modified Julian date  generate
      a correctly formatted FITS DATE-OBS string.
   }
   \sstinvocation{
      CALL SCULIB\_MJD\_TO\_DATEOBS ( MJD, DATE\_OBS, STATUS )
   }
   \sstarguments{
      \sstsubsection{
         MJD = DOUBLE PRECISION (Given)
      }{
         Modified Julian date (Defined as JD - 2400000.5)
      }
      \sstsubsection{
         DATE\_OBS = CHARACTER $*$ ($*$) (Returned)
      }{
         FITS DATE-OBS string. Should have at least 24
         characters for YYYY-MM-DDThh:mm:ss.sssZ format
      }
      \sstsubsection{
         STATUS = INTEGER (Given \& Returned)
      }{
         Global Status
      }
   }
   \sstnotes{
      The format used for the DATE-OBS keyword depends on the value of the
      keyword. For DATE-OBS $<$ 1999.0, use the old {\tt "}dd/mm/yy{\tt "} format.
      Otherwise, use the new {\tt "}ccyy-mm-ddThh:mm:ss[.ssss]Z{\tt "} format.
   }
   \sstdiytopic{
      Copyright
   }{
      Copyright (C) 2000 Particle Physics and Astronomy
      Research Council. All Rights Reserved.
   }
}

\sstroutine{
   SCULIB\_MRQCOF
}{
   called by SCULIB\_MRQMIN
}{
   \sstdescription{
      Utility routine called by SCULIB\_MRQMIN, copied from Numerical
      Recipes in Fortran, p. 681. Used by SCULIB\_MRQMIN to evaluate the
      linearised fitting matrix ALPHA, and vector BETA, and calculate
      chi-squared.
         Copied from MRQCOF on p.681 of Numerical Recipes in Fortran, with
      STATUS added.
   }
   \sstinvocation{
       CALL SCULIB\_MRQCOF (X, Y, SIG, NDATA, A, IA, MA, ALPHA,
         BETA, NALP, CHISQ, FUNCS, STATUS)
   }
   \sstarguments{
      \sstsubsection{
         X( NDATA ) = REAL (Given)
      }{
         X Data points
      }
      \sstsubsection{
         Y( NDATA ) = REAL (Given)
      }{
         Y data points
      }
      \sstsubsection{
         SIG ( NDATA ) = REAL (Given)
      }{
      }
      \sstsubsection{
         NDATA = INTEGER (Given)
      }{
         Number of data points in arrays
      }
      \sstsubsection{
         A( MA ) = REAL (Given \& Returned)
      }{
         The input parameter values. Contains the best-fit values on exit.
      }
      \sstsubsection{
         IA( MA ) = INTEGER (Given)
      }{
         Array specifying which components should be fitted. A zero
         indicates the parameter should not be fitted -- the parameters
         are kept at their input values. Non-zero indicates the parameter
         is free to be modified.
      }
      \sstsubsection{
         MA = INTEGER (Given)
      }{
         Number of coefficients/parameters for function.
      }
      \sstsubsection{
         ALPHA( NALP, NALP ) = REAL (Returned)
      }{
      }
      \sstsubsection{
         BETA( MA ) = REAL (Returned)
      }{
      }
      \sstsubsection{
         NALP = INTEGER  (Given)
      }{
         Size of ALPHA.
      }
      \sstsubsection{
         CHISQ = REAL (Returned)
      }{
         ChiSq of fit.
      }
      \sstsubsection{
         FUNCS = REAL FUNCTION (Given)
      }{
         The function to be fitted. Must take arguments of X, A, YFIT,
         DYDA and MA, where YFIT is the fitting function and DYDA the
         the derivatives with respect to parameters A at X.
      }
      \sstsubsection{
         STATUS = INTEGER (Given \& Returned)
      }{
         Global Status
      }
   }
   \sstdiytopic{
      Copyright
   }{
      Copyright \copyright 1995,1996,1997,1998,1999 Particle Physics and Astronomy
      Research Council. All Rights Reserved.
   }
}
\sstroutine{
   SCULIB\_MRQMIN
}{
   Levenberg-Marquardt method non-linear least-squares
   fit from Numerical Recipes in FORTRAN
}{
   \sstdescription{
      Levenberg-Marquardt method, attempting to reduce the value of
      chi-squared of a fit between a set of data points X(1:NDATA),
      Y(1:NDATA), and a non-linear function dependent on MA coefficients
      A(1:MA). The input array IA(1:MA) indicates by non-zero components
      those components of A that should be fitted for, and by 0 entries
      those components that should be held fixed at their input values.
      The program returns current best-fit values for the parameters
      A(1:MA), and chi-squared = CHISQ. The arrays COVAR (1:NCA,1:NCA),
      ALPHA (1:NCA,1:NCA) with physical dimension NCA ($>$= the number of
      fitted parameters) are used as working space during most iterations.
      Supply a subroutine FUNCS (X, A, YFIT, DYDA, MA) that evaluates
      the fitting function YFIT, and its derivatives DYDA with respect to the
      the fitting paramters A at X. On the first call provide an initial
      guess for the parameters A, and set ALAMDA $<$ 0 for initialisation
      (which then sets ALAMDA = 0.001). If a step succeeds CHISQ becomes
      smaller and ALAMDA decreases by a factor of 10. If a step fails
      ALAMDA grows by a factor of 10. You must call this routine repeatedly
      until convergence is achieved. Then, make one final call with
      ALAMDA = 0, so that COVAR (1:MA,1:MA) returns the covariance matrix,
      and ALPHA the curvature matrix. Parameters held fixed will return
      zero covariance.
        Copied from MRQMIN on p.680 of Numerical Recipes in Fortran, with
      STATUS added.
   }
   \sstinvocation{
       CALL SCULIB\_MRQMIN (X, Y, SIG, NDATA, A, IA, MA, COVAR,
         ALPHA, NCA, CHISQ, FUNCS, ALAMDA, STATUS)
   }
   \sstarguments{
      \sstsubsection{
         X( NDATA ) = REAL (Given)
      }{
         X Data points
      }
      \sstsubsection{
         Y( NDATA ) = REAL (Given)
      }{
         Y data points
      }
      \sstsubsection{
         SIG ( NDATA ) = REAL (Given)
      }{
      }
      \sstsubsection{
         NDATA = INTEGER (Given)
      }{
         Number of data points in arrays
      }
      \sstsubsection{
         A( MA ) = REAL (Given \& Returned)
      }{
         The input parameter values. Contains the best-fit values on exit.
      }
      \sstsubsection{
         IA( MA ) = INTEGER (Given)
      }{
         Array specifying which components should be fitted. A zero
         indicates the parameter should not be fitted -- the parameters
         are kept at their input values. Non-zero indicates the parameter
         is free to be modified.
      }
      \sstsubsection{
         MA = INTEGER (Given)
      }{
         Number of coefficients/parameters for function.
      }
      \sstsubsection{
         COVAR ( NCA, NCA ) = REAL (Returned)
      }{
         The covariance matrix. Fixed parameters will return zero covariance.
      }
      \sstsubsection{
         ALPHA ( NCA, NCA ) = REAL (Returned)
      }{
         The curvature matrix.
      }
      \sstsubsection{
         NCA = INTEGER (Given)
      }{
         Size of output arrays. Must be greater than the number of fitted
         parameters
      }
      \sstsubsection{
         CHISQ = REAL (Returned)
      }{
         ChiSq of fit.
      }
      \sstsubsection{
         FUNCS = REAL FUNCTION (Given)
      }{
         The function to be fitted. Must take arguments of X, A, YFIT,
         DYDA and MA, where YFIT is the fitting function and DYDA the
         the derivatives with respect to parameters A at X.
      }
      \sstsubsection{
         ALAMDA = REAL (Given \& Returned)
      }{
         The stepping factor for the iteration. Should be set to $<$0
         initially.
      }
      \sstsubsection{
         STATUS = INTEGER (Given \& Returned)
      }{
         Global Status
      }
   }
   \sstdiytopic{
      Copyright
   }{
      Copyright \copyright 1995,1996,1997,1998,1999 Particle Physics and Astronomy
      Research Council. All Rights Reserved.
   }
}
\sstroutine{
   SCULIB\_MULCAD
}{
   multiply double precision array by a constant double
}{
   \sstdescription{
      multiplies a double precision array by a double constant
   }
   \sstinvocation{
      CALL SCULIB\_MULCAD (N, IN, DVAL, OUT, STATUS)
   }
   \sstarguments{
      \sstsubsection{
         N              = INTEGER (Given)
      }{
         number of array elements
      }
      \sstsubsection{
         IN (N)         = DOUBLE (Given)
      }{
         array to be multiplied
      }
      \sstsubsection{
         DVAL           = DOUBLE (Given)
      }{
         multiplication factor
      }
      \sstsubsection{
         OUT (N)        = DOUBLE (Returned)
      }{
         output array (can be same as input)
      }
      \sstsubsection{
         STATUS         = INTEGER (Given \& Returned)
      }{
         Global status
      }
   }
   \sstdiytopic{
      Copyright
   }{
      Copyright \copyright 1995,1996,1997,1998,1999 Particle Physics and Astronomy
      Research Council. All Rights Reserved.
   }
}
\sstroutine{
   SCULIB\_MULCAR
}{
   multiply real array by a constant
}{
   \sstdescription{
      multiplies a real array by a real constant
   }
   \sstinvocation{
      CALL SCULIB\_MULCAR (N, IN, RVAL, OUT)
   }
   \sstarguments{
      \sstsubsection{
         N              = INTEGER (Given)
      }{
         number of array elements
      }
      \sstsubsection{
         IN (N)         = REAL (Given)
      }{
         array to be multiplied
      }
      \sstsubsection{
         RVAL           = REAL (Given)
      }{
         multiplication factor
      }
      \sstsubsection{
         OUT (N)        = REAL (Returned)
      }{
         output array (can be same as input)
      }
   }
   \sstdiytopic{
      Copyright
   }{
      Copyright \copyright 1995,1996,1997,1998,1999 Particle Physics and Astronomy
      Research Council. All Rights Reserved.
   }
}
\sstroutine{
   SCULIB\_MULTARE
}{
   multiplies 2 real arrays with optional error and quality
   handling
}{
   \sstdescription{
      Multiplies two floating point arrays.  The arrays
      may have any dimensions; they are treated here as
      linear in order to generate more efficient code.
   }
   \sstinvocation{
       CALL SCULIB\_MULTARE (NELM, ARRAY1, ARRAY2, ARRAY3, Q1DATA,
         Q2DATA, Q3DATA, E1DATA, E2DATA, E3DATA, QUAL, FLAGS, FBAD,
         VARIANCE)
   }
   \sstarguments{
      \sstsubsection{
         NELM                = INTEGER (Given)
      }{
         Number of elements in each array
      }
      \sstsubsection{
         ARRAY1 (NELM)       = REAL (Given)
      }{
         Input array
      }
      \sstsubsection{
         ARRAY2 (NELM)       = REAL (Given)
      }{
         Second input array.
      }
      \sstsubsection{
         ARRAY3 (NELM)       = REAL (Returned)
      }{
         Result array.  ARRAY3=ARRAY1$*$ARRAY2
      }
      \sstsubsection{
         Q1DATA (NELM)       = INTEGER (Given)
      }{
         Quality array for first input array
      }
      \sstsubsection{
         Q2DATA (NELM)       = INTEGER (Given)
      }{
         Quality array for second input array
      }
      \sstsubsection{
         Q3DATA (NELM)       = INTEGER (Returned)
      }{
         Quality array for output array
      }
      \sstsubsection{
         E1DATA (NELM)       = REAL (Given)
      }{
         Variance array for input array
      }
      \sstsubsection{
         E2DATA (NELM)       = REAL (Given)
      }{
         Variance array for second input array
      }
      \sstsubsection{
         E3DATA (NELM)       = REAL (Returned)
      }{
         Variance array for output array
      }
      \sstsubsection{
         QUAL                = LOGICAL (Given)
      }{
         True if input has quality information
      }
      \sstsubsection{
         FLAGS               = LOGICAL (Given)
      }{
         True if input has flagged data values
      }
      \sstsubsection{
         FBAD                = REAL (Given)
      }{
         Flag value
      }
      \sstsubsection{
         VARIANCE            = LOGICAL (Given)
      }{
         True if both input arrays have variance arrays
      }
   }
   \sstnotes{
      \sstitemlist{

         \sstitem
         Any of the arrays may be the same.

         \sstitem
         Consider using VEC\_MULR (SUN/39) instead.
      }
   }
   \sstdiytopic{
      Copyright
   }{
      Copyright \copyright 1995,1996,1997,1998,1999 Particle Physics and Astronomy
      Research Council. All Rights Reserved.
   }
}
\sstroutine{
   SCULIB\_NFILLI
}{
   fill an integer array with its indices
}{
   \sstdescription{
      fills an integer array with its indices
   }
   \sstinvocation{
      CALL SCULIB\_NFILLI (N, ARRAY)
   }
   \sstarguments{
      \sstsubsection{
         N            = INTEGER (Given)
      }{
         number of array elements
      }
      \sstsubsection{
         ARRAY (N)    = INTEGER (Returned)
      }{
         array to be filled
      }
   }
   \sstdiytopic{
      Copyright
   }{
      Copyright \copyright 1995,1996,1997,1998,1999 Particle Physics and Astronomy
      Research Council. All Rights Reserved.
   }
   \sstimplementationstatus{
      No status checking
   }
}
\sstroutine{
   SCULIB\_NFILLR
}{
   fill a real array with its indices
}{
   \sstdescription{
      fills a real array with its indices
   }
   \sstinvocation{
      CALL SCULIB\_NFILLR (N, ARRAY)
   }
   \sstarguments{
      \sstsubsection{
         N            = INTEGER (Given)
      }{
         number of array elements
      }
      \sstsubsection{
         ARRAY (N)    = REAL (Returned)
      }{
         array to be filled
      }
   }
   \sstdiytopic{
      Copyright
   }{
      Copyright \copyright 1995,1996,1997,1998,1999 Particle Physics and Astronomy
      Research Council. All Rights Reserved.
   }
   \sstimplementationstatus{
      No status checking
   }
}
\sstroutine{
   SCULIB\_NOISE\_MEAN
}{
   Calculate statistics of NOISE data from raw demodulated data
}{
   \sstdescription{
      This routine calculates the mean and variance of a set of integrations
      in a NOISE measurement.
         If, only one integration was taken then the demodulated chop signal,
      chop variance, calibrator signal, calibrator variance and quality will
      be returned in CHOP\_DATA, CHOP\_VARIANCE, CAL\_DATA, CAL\_VARIANCE and
      QUALITY as they were obtained from the transputers.
         If more than 1 integration was obtained then:-

           If, for a given bolometer, no valid data was obtained for any of the
           integrations then the QUALITY for that bolometer will be returned
           as 1.

           If, for that bolometer, one valid integration was obtained then the
           demodulated chop signal and variance will be returned in CHOP\_DATA
           and CHOP\_VARIANCE, the calibrator signal and variance will be
           returned in CAL\_DATA and CAL\_VARIANCE, and QUALITY will be set
           to 0.

           If more than one valid integration was obtained then the means
           of the chop and calibrator signals will be returned in CHOP\_DATA
           and CAL\_DATA, and the variances of the individual samples in
           CHOP\_VARIANCE and  CAL\_VARIANCE, and QUALITY will be set to 0.
   }
   \sstinvocation{
       SCULIB\_NOISE\_MEAN ( N\_INTEGRATIONS, N\_BOLS, DEMOD,
            CHOP\_DATA, CHOP\_VARIANCE, CAL\_DATA, CAL\_VARIANCE,
            QUALITY, WORKSPACE, STATUS)
   }
   \sstarguments{
      \sstsubsection{
         N\_INTEGRATIONS                    = INTEGER (Given)
      }{
         number of integrations taken in measurement
      }
      \sstsubsection{
         N\_BOLS                            = INTEGER (Given)
      }{
         number of bolometers measured
      }
      \sstsubsection{
         DEMOD (5, N\_BOLS, N\_INTEGRATIONS) = REAL (Given)
      }{
         the demodulated data; chop, chop variance, calibrator, cal
         variance, quality
      }
      \sstsubsection{
         CHOP\_DATA (N\_BOLS)                = REAL (Returned)
      }{
         the mean of the demodulated chop signal over the integrations
      }
      \sstsubsection{
         CHOP\_VARIANCE (N\_BOLS)            = REAL (Returned)
      }{
         the variance on CHOP\_DATA, calculated from the dispersion
         about the mean, or the demodulated variance if there was only
         one integration
      }
      \sstsubsection{
         CAL\_DATA (N\_BOLS)                 = REAL (Returned)
      }{
         the mean of the demodulated calibrator signal over the integrations
      }
      \sstsubsection{
         CAL\_VARIANCE (N\_BOLS)             = REAL (Returned)
      }{
         the variance on CAL\_DATA, calculated from the dispersion about the
         mean, or the demodulated variance if there was only one integration
      }
      \sstsubsection{
         QUALITY (N\_BOLS)                  = BYTE (Returned)
      }{
         the quality on the returned data
      }
      \sstsubsection{
         WORKSPACE (N\_BOLS)                = INTEGER (Scratch)
      }{
         some scratch space for keeping track of the number of
         valid observations to be included in the averaging
         for each bolometer
      }
      \sstsubsection{
         STATUS = INTEGER (Given \& Returned)
      }{
         inherited status
      }
   }
   \sstdiytopic{
      Copyright
   }{
      Copyright \copyright 1995,1996,1997,1998,1999 Particle Physics and Astronomy
      Research Council. All Rights Reserved.
   }
}
\sstroutine{
   SCULIB\_PAR\_GET0?
}{
   Wrapper for the standard PAR\_GET0 routines
}{
   \sstdescription{
      This routine provides a wrapper for PAR\_GET0? so that
      STATUS can be checked and an informative error message
      added without adding large numbers of lines to other routines.
   }
   \sstinvocation{
      CALL SCULIB\_PAR\_GET0?(PARAM, VALUE, STATUS)
   }
   \sstarguments{
      \sstsubsection{
         PARAM  = CHARACTER (Given)
      }{
         Name of requested parameter
      }
      \sstsubsection{
         VALUE  = TYPE (Returned)
      }{
         Parameter value (bad if bad status from PAR\_GET)
      }
      \sstsubsection{
         STATUS = INTEGER (Given \& Returned)
      }{
         Global status value
      }
   }
   \sstdiytopic{
      Copyright
   }{
      Copyright \copyright 1995,1996,1997,1998,1999 Particle Physics and Astronomy
      Research Council. All Rights Reserved.
   }
   \sstimplementationstatus{
      \sstitemlist{

         \sstitem
         All return types are present as different subroutines
      }
   }
}
\sstroutine{
   SCULIB\_PHOTOM\_BOLSELECT
}{
   select photometers for a PHOTOM observation
}{
   \sstdescription{
      This routine selects the bolometers to be used in a PHOTOM observation.
      If status is good on entry the routine will start by setting to zero
      the arrays describing which bolometers have been selected. Then it will
      call SCULIB\_BOLSELECT to get the chan/ADC numbers of the bolometers
      directly specified in BOLOMETERS.

        The observer can specify between 1 and 3 bolometers by name. Otherwise
      an error will be reported and the routine will return with bad status.
      Likewise, if any of these directly selected bolometers have bad quality
      or have identical Nasmyth offsets.

      Now the routine branches on the number of directly selected bolometers:-

      \sstitemlist{

         \sstitem
          3 bolometers specified. In this case the observer must want to chop
          between 3 bolometers on an array.
            All 3 bolometers should belong to a single array, if not the routine
          will error and return with bad status. The routine will arbitrarily
          assign the 3 bolometers to the `left{\tt '}, {\tt '}middle{\tt '} and `right{\tt '} projected
          positions in the order they were selected. The routine now checks that
          the 3 bolometers do lie in roughly a straight line on the sky (allowing
          for distortion) and that the left and right bolometers lie roughly the
          same distance on either side of the middle. Errors will occur and the
          routine return with bad status if these conditions are not met.
            If all is OK the routine will now store the channel/ADC numbers of
          the directly selected bolometers, set CHOP\_COORDS to Nasmyth, the
          bolometer spacing to half the distance between the left and right,
          the chopper position angle, and the instrument `centre{\tt '} to the coords of
          the middle bolometer.
            Next, the routine will search the bolometers belonging to the other
          array. If 3 can be found that match the positions of the 3 direct
          bolometers and have good quality then the routine will store the
          channel/ADC numbers of those bolometers as well.
            Lastly, the routine will set the bolometers to actually be measured
          to all those belonging to the array(s) containing the bolometers
          already selected.

         \sstitem
          2 bolometers specified. In this case the observer must want to observe
          a source by chopping between the 2 named bolometers.
            The routine will arbitrarily call the first of the 2 bolometers
          the `middle{\tt '} projected bolometer, the second `right{\tt '}. Then it will
          store the channel/ADC numbers of these directly selected bolometers and
          set the types of sub-instruments that they imply. If any of the
          sub-instruments are arrays the routine will search for bolometer(s)
          at the same Nasmyth offset in the other array and select them too.
            Lastly, the routine will set the bolometers to actually be measured
          to all those in the selected sub-instruments. CHOP\_COORDS is set to
          Nasmyth, the bolometer spacing to the distance between the
          directly selected bolometers, the chopper position angle is set, and
          the instrument `centre{\tt '} to the Nasmyth coords of the middle bolometer.

         \sstitem
          1 bolometer specified. In this case the observer must want to observe
          a source without chopping between different bolometers.
            The routine will arbitrarily call the primary selected bolometer
          the `middle{\tt '} projected bolometer. `Left{\tt '} and `right{\tt '} bolometers will
          not be assigned.
            The routine will now store the channel/ADC numbers of the directly
          selected bolometer and set the type of sub-instrument that is implied.
          If the sub-instrument is one of the arrays the routine will search
          for a bolometer at the same Nasmyth offset in the other array and
          select it too.
            Lastly, the routine will set the bolometers to actually be measured
          to all those in the selected sub-instruments. The instrument `centre{\tt '}
          is set to the Nasmyth offset of the primary selected bolometer.
          CHOP\_COORDS, the bolometer spacing and chopper position angle are
          all set to bad values so that the calling routine knows these have to
          be read explicitly from the observation definition file.

      }
      Finally, SCULIB\_BOLSELECT is called to select for measurement all the
      bolometers belonging to the sub-instruments involved in this PHOTOM
      observation. An array is then calculated holding the index in the array
      of ALL the photometers to be measured of each projected bolometer
      in each sub-instrument that was directly selected.
   }
   \sstinvocation{
       CALL  SCULIB\_PHOTOM\_BOLSELECT (BOLOMETERS, BOL\_TYPE,
          BOL\_CALIB, BOL\_DU3, BOL\_DU4, BOL\_QUAL, BOL\_ENABLED, NUM\_CHAN,
          NUM\_ADC, BOL\_SELECT\_CHAN, BOL\_SELECT\_ADC, N\_BOLS, MAX\_SUB,
          SUB\_INSTRMNT, N\_SUB, CENTRE\_DU3, CENTRE\_DU4, CHOP\_COORDS,
          BOL\_SPACING, CHOP\_PA, N\_BOL\_SUB, BOLS\_MEASURED,
          PHOT\_BEAM\_CHAN, PHOT\_BEAM\_ADC, PHOT\_BEAM\_BOL, STATUS)
   }
   \sstarguments{
      \sstsubsection{
         BOLOMETERS                  = CHARACTER$*$($*$) (Given)
      }{
         list of bolometer selections
      }
      \sstsubsection{
         BOL\_TYPE (NUM\_CHAN, NUM\_ADC)
      }{
                               = CHARACTER$*$($*$) (Given)
         type of bolometer
      }
      \sstsubsection{
         BOL\_CALIB (NUM\_CHAN, NUM\_ADC)
      }{
                               = REAL (Given)
         target calibrator values for bolometers
      }
      \sstsubsection{
         BOL\_DU3 (NUM\_CHAN, NUM\_ADC) = REAL (Given)
      }{
         Nasmyth DU3 offset of bolometer from field centre
      }
      \sstsubsection{
         BOL\_DU4 (NUM\_CHAN, NUM\_ADC) = REAL (Given)
      }{
         Nasmyth DU4 offset of bolometer from field centre
      }
      \sstsubsection{
         BOL\_QUAL (NUM\_CHAN, NUM\_ADC)
      }{
                               = INTEGER (Given)
         quality of bolometers
      }
      \sstsubsection{
         BOL\_ENABLED (NUM\_CHAN, NUM\_ADC)
      }{
                               = LOGICAL (Returned)
         .TRUE. if bolometer was selected
      }
      \sstsubsection{
         NUM\_CHAN                    = INTEGER (Given)
      }{
         number of channels per A/D
      }
      \sstsubsection{
         NUM\_ADC                     = INTEGER (Given)
      }{
         number of A/D cards
      }
      \sstsubsection{
         BOL\_SELECT\_CHAN (NUM\_CHAN $*$ NUM\_ADC)
      }{
                               = INTEGER (Returned)
         channel numbers of selected bolometers
      }
      \sstsubsection{
         BOL\_SELECT\_ADC (NUM\_CHAN $*$ NUM\_ADC)
      }{
                               = INTEGER (Returned)
         A/D card numbers of selected bolometers
      }
      \sstsubsection{
         N\_BOLS                      = INTEGER (Returned)
      }{
         total number of bolometers selected
      }
      \sstsubsection{
         MAX\_SUB                     = INTEGER (Given)
      }{
         maximum number of sub instruments
      }
      \sstsubsection{
         SUB\_INSTRMNT (MAX\_SUB)      = CHARACTER$*$($*$) (Returned)
      }{
         names of sub instrument sections to be used
      }
      \sstsubsection{
         N\_SUB                       = INTEGER (Returned)
      }{
         the number of sub instruments to be used
      }
      \sstsubsection{
         CENTRE\_DU3                  = REAL (Returned)
      }{
         Nasmyth DU3 offset from instrument centre to which telescope
         is to pointed
      }
      \sstsubsection{
         CENTRE\_DU4                  = REAL (Returned)
      }{
         Nasmyth DU4 offset from instrument centre to which telescope
         is to pointed
      }
      \sstsubsection{
         CHOP\_COORDS                 = CHARACTER$*$($*$) (Returned)
      }{
         Chopper coordinate system required
      }
      \sstsubsection{
         BOL\_SPACING                 = REAL (Returned)
      }{
         Spacing between bolometers (arcsec)
      }
      \sstsubsection{
         CHOP\_PA                     = REAL (Returned)
      }{
         Chop position angle required (degrees)
      }
      \sstsubsection{
         N\_BOL\_SUB (MAX\_SUB)         = INTEGER (Returned)
      }{
         Number of bolometers selected in each sub instrument
      }
      \sstsubsection{
         BOLS\_MEASURED               = CHARACTER$*$($*$) (Returned)
      }{
         Bolometers to be measured by transputer system
      }
      \sstsubsection{
         PHOT\_BEAM\_CHAN (3, MAX\_SUB) = INTEGER (Returned)
      }{
         Channel numbers of selected bolometers projected left, middle,
         right on sky for each sub-instrument (I index = 1 for left, 2
         for middle and 3 for right)
      }
      \sstsubsection{
         PHOT\_BEAM\_ADC (3, MAX\_SUB)  = INTEGER (Returned)
      }{
         ADC numbers of selected bolometers projected left, middle,
         right on sky for each sub-instrument
      }
      \sstsubsection{
         PHOT\_BEAM\_BOL (3, MAX\_SUB)  = INTEGER (Returned)
      }{
         The index in the array of selected bolometers of those projected
         left, middle, right on sky for each sub-instrument.
      }
      \sstsubsection{
         STATUS                      = INTEGER (Given and returned)
      }{
         global status
      }
   }
   \sstdiytopic{
      Copyright
   }{
      Copyright \copyright 1995,1996,1997,1998,1999 Particle Physics and Astronomy
      Research Council. All Rights Reserved.
   }
}
\sstroutine{
   SCULIB\_POWER2
}{
   Calculate next highest power of 2
}{
   \sstdescription{
      Given a number, returns that number or the next highest
      number that is a power of 2, and returns the power itself.
      If N$<$1, GEN\_POWER2 is returned as -1, and NP2 is returned as 0.
   }
   \sstinvocation{
      I = SCULIB\_POWER2 (N, NP2)
   }
   \sstarguments{
      \sstsubsection{
         N              = INTEGER (Given)
      }{
         The number in question.
      }
      \sstsubsection{
         NP2            = INTEGER (Returned)
      }{
         The next number that is a power of 2. Returns 0 if N is less
         than 1.
      }
   }
   \sstreturnedvalue{
      \sstsubsection{
         SCULIB\_POWER2  = INTEGER (Returned)
      }{
         The power of 2.  ie NP2=2$*$$*$SCULIB\_POWER2.
         Returns -1 if N is less than 1.
      }
   }
   \sstdiytopic{
      Copyright
   }{
      Copyright \copyright 1992,1999 Particle Physics and Astronomy
      Research Council. All Rights Reserved.
   }
}
\sstroutine{
   SCULIB\_PUT\_FITS\_C
}{
   write a string to a FITS item specified
}{
   \sstdescription{
      This routine  writes a FITS character item into a character array
      ready to be written out to the .MORE.FITS structure in an NDF file.
      An error will be reported and bad status returned if the name of the
      FITS item is blank or more than 8 characters long. An error will also
      occur if the FITS character array is full.
      The FITS array must contain a final entry of simply `END' unless
      N\_FITS is 0. The END entry is automatically moved to the last field
      as entries are added to the FITS array.
  }
   \sstinvocation{
       CALL SCULIB\_PUT\_FITS\_C (MAX\_FITS, N\_FITS, FITS, NAME, VALUE,
         COMMENT, STATUS)
   }
   \sstarguments{
      \sstsubsection{
         MAX\_FITS                       = INTEGER (Given)
      }{
         the size of the FITS array
      }
      \sstsubsection{
         N\_FITS                         = INTEGER (Given and returned)
      }{
         the number of FITS items currently in the FITS array.
           This should include an `END' entry. If N\_FITS is 1 it will
           be assumed to be an array with just an `END'. If N\_FITS=0
           on entry the array will be empty and N\_FITS will be changed to 2
           on exit (the entry and the END card).
      }
      \sstsubsection{
         FITS (MAX\_FITS)                = CHARACTER$*$80 (Given and returned)
      }{
         array containing the FITS items
      }
      \sstsubsection{
         NAME                           = CHARACTER$*$($*$) (Given)
      }{
         the name of the FITS item to be written
      }
      \sstsubsection{
         VALUE                          = CHARACTER$*$($*$) (Given)
      }{
         the value to be assigned to the item
      }
      \sstsubsection{
         COMMENT                        = CHARACTER$*$($*$) (Given)
      }{
         a comment applying to the FITS item
      }
      \sstsubsection{
         STATUS                         = INTEGER (Given and returned)
      }{
         global status
      }
   }
   \sstdiytopic{
      Copyright
   }{
      Copyright \copyright 1995-2000 Particle Physics and Astronomy
      Research Council. All Rights Reserved.
   }
}
\sstroutine{
   SCULIB\_PUT\_FITS\_D
}{
   write a DOUBLE to a FITS item specified
}{
   \sstdescription{
      This routine  writes a FITS double precision item into a character array
      ready to be written out to the .MORE.FITS structure in an NDF file.
      An error will be reported and bad status returned if the name of the
      FITS item is blank or more than 8 characters long. An error will also
      occur if the FITS character array is full.
      The FITS array must contain a final entry of simply `END' unless
      N\_FITS is 0. The END entry is automatically moved to the last field
      as entries are added to the FITS array.
   }
   \sstinvocation{
       CALL SCULIB\_PUT\_FITS\_D (MAX\_FITS, N\_FITS, FITS, NAME, VALUE,
         COMMENT, STATUS)
   }
   \sstarguments{
      \sstsubsection{
         MAX\_FITS                       = INTEGER (Given)
      }{
         the size of the FITS array
      }
      \sstsubsection{
         N\_FITS                         = INTEGER (Given and returned)
      }{
         The number of FITS items currently in the FITS array.
           This should include an `END' entry. If N\_FITS is 1 it will
           be assumed to be an array with just an `END'. If N\_FITS=0
           on entry the array will be empty and N\_FITS will be changed to 2
           on exit (the entry and the END card).
      }
      \sstsubsection{
         FITS (MAX\_FITS)                = CHARACTER$*$80 (Given and returned)
      }{
         array containing the FITS items
      }
      \sstsubsection{
         NAME                           = CHARACTER$*$($*$) (Given)
      }{
         the name of the FITS item to be written
      }
      \sstsubsection{
         VALUE                          = DOUBLE PRECISION (Given)
      }{
         the value to be assigned to the item
      }
      \sstsubsection{
         COMMENT                        = CHARACTER$*$($*$) (Given)
      }{
         a comment applying to the FITS item
      }
      \sstsubsection{
         STATUS                         = INTEGER (Given and returned)
      }{
         global status
      }
   }
   \sstdiytopic{
      Copyright
   }{
      Copyright \copyright 1995,1996,1997,1998,1999 Particle Physics and Astronomy
      Research Council. All Rights Reserved.
   }
}
\sstroutine{
   SCULIB\_PUT\_FITS\_I
}{
   write an integer to a FITS item specified
}{
   \sstdescription{
      This routine writes a FITS integer item into a character array
      ready to be written out to the .MORE.FITS structure in an NDF file.
      An error will be reported and bad status returned if the name of the
      FITS item is blank or more than 8 characters long. An error will also
      occur if the FITS character array is full.
      The FITS array must contain a final entry of simply `END' unless
      N\_FITS is 0. The END entry is automatically moved to the last field
      as entries are added to the FITS array.
   }
   \sstinvocation{
       CALL SCULIB\_PUT\_FITS\_I (MAX\_FITS, N\_FITS, FITS, NAME, VALUE,
         COMMENT, STATUS)
   }
   \sstarguments{
      \sstsubsection{
         MAX\_FITS                       = INTEGER (Given)
      }{
         the size of the FITS array
      }
      \sstsubsection{
         N\_FITS                         = INTEGER (Given and returned)
      }{
         The number of FITS items currently in the FITS array.
           This should include an `END' entry. If N\_FITS is 1 it will
           be assumed to be an array with just an `END'. If N\_FITS=0
           on entry the array will be empty and N\_FITS will be changed to 2
           on exit (the entry and the END card).
      }
      \sstsubsection{
         FITS (MAX\_FITS)                = CHARACTER$*$80 (Given and returned)
      }{
         array containing the FITS items
      }
      \sstsubsection{
         NAME                           = CHARACTER$*$($*$) (Given)
      }{
         the name of the FITS item to be written
      }
      \sstsubsection{
         VALUE                          = INTEGER (Given)
      }{
         the value to be assigned to the item
      }
      \sstsubsection{
         COMMENT                        = CHARACTER$*$($*$) (Given)
      }{
         a comment applying to the FITS item
      }
      \sstsubsection{
         STATUS                         = INTEGER (Given and returned)
      }{
         global status
      }
   }
   \sstdiytopic{
      Copyright
   }{
      Copyright \copyright 1995,1996,1997,1998,1999 Particle Physics and Astronomy
      Research Council. All Rights Reserved.
   }
}
\sstroutine{
   SCULIB\_RAD2STRING
}{
   Translate an angle or time in radians to a nicely formatted string
}{
   \sstdescription{
      Converts an angle or time in radians to a string that can be
      used for display. The sexagesimal format is colon separated.
      For time output, the output is format SHH:MM:SS.NDP.
      For angle output, the format used is SDD:MM:SS.NDP. Note
      that the sign is used for the first character but this can
      be either blank or a {\tt '}-{\tt '}.
   }
   \sstinvocation{
      CALL SCULIB\_RAD2STRING( ANGLE, NDP, ISTIME, RESLT, STATUS )
   }
   \sstarguments{
      \sstsubsection{
         ANGLE = DOUBLE PRECISION (Given)
      }{
         Angle to be translated (radians)
      }
      \sstsubsection{
         NDP = INTEGER (Given)
      }{
         Number of decimal places to use for the seconds
         Can not be greater than 6.
      }
      \sstsubsection{
         ISTIME = LOGICAL (Given)
      }{
         If true the angle will be assumed to be a time and will
         be converted to HH:MM:SS.NDP format. If false it will be
         assumed to be an angle and converted to SDDD:MM:SS.NDP.
      }
      \sstsubsection{
         RESLT = CHARACTER (Returned)
      }{
         String containing the translated angle. Colon separated.
         Should be at least 11$+$NDP characters long.
      }
      \sstsubsection{
         STATUS = INTEGER (Given \& Returned)
      }{
         Global Status
      }
   }
   \sstnotes{
      Essentially a wrapper around SLA\_DR2AF and SLA\_DR2TF
   }
   \sstdiytopic{
      Copyright
   }{
      Copyright (C) 2000 Particle Physics and Astronomy Research Council.
      All Rights Reserved.
   }
}


\sstroutine{
   SCULIB\_RANGED
}{
   Finds the maximum and minimum values in a double precision array
}{
   \sstdescription{
      Find the maximum and minimum values in a double precision array.
   }
   \sstinvocation{
      CALL SCULIB\_RANGED( ARRAY, IST, IEN, VMAX, VMIN, STATUS)
   }
   \sstarguments{
      \sstsubsection{
         ARRAY( IEN ) = DOUBLE PRECISION (Given)
      }{
         Array containing the values to be checked.
      }
      \sstsubsection{
         IST = INTEGER (Given)
      }{
         The first element of ARRAY to be examined.
      }
      \sstsubsection{
         IEN = INTEGER (Given)
      }{
         The last element of ARRAY to be examined.
      }
      \sstsubsection{
         VMAX = DOUBLE PRECISION (Returned)
      }{
         The maximum value of those examined
      }
      \sstsubsection{
         VMIN = DOUBLE PRECISION (Returned)
      }{
         The minumum value of those examined
      }
      \sstsubsection{
         STATUS = INTEGER (Given \& Returned)
      }{
         Global status
      }
   }
   \sstdiytopic{
      Copyright
   }{
      Copyright \copyright 1995,1996,1997,1998,1999 Particle Physics and Astronomy
      Research Council. All Rights Reserved.
   }
}
\sstroutine{
   SCULIB\_READ\_JIGGLE
}{
   read a jiggle pattern
}{
   \sstdescription{
      This routine reads in a jiggle pattern and sets some variables associated
      with it; the number of jiggles in the pattern, the number of jiggles to
      be measured in each exposure, the number of times the pattern will be
      repeated in each exposure, the maximum offsets in the pattern.

         After checking status on entry, SCULIB\_READ\_NUMBERS is called to read
      in the jiggle offsets from the file named in JIGGLE\_NAME, an error will
      be reported and bad status returned if no offsets are read. The maximum
      and minimum x and y offsets in the pattern are calculated.

         Next, the variables governing the way the jiggle pattern will be
      divided among the exposures making up each integration is worked out.
      An error will be reported and bad status returned if JIGGLE\_P\_SWITCH
      is less than or equal to zero. Otherwise, the number of exposures
      required to execute the whole jiggle pattern, each exposure containing
      JIGGLE\_P\_SWITCH jiggles (except perhaps the last which will hold less),
      is calculated. If JIGGLE\_P\_SWITCH is larger than the number of jiggles
      in the pattern then the jiggle pattern will be repeated an integer
      number of times during each switch of the exposure and JIGGLE\_P\_SWITCH
      will be reset accordingly.
   }
   \sstinvocation{
       CALL SCULIB\_READ\_JIGGLE (JIGGLE\_NAME, MAX\_JIGGLE,
         JIGGLE\_P\_SWITCH, JIGGLE\_REPEAT, EXP\_PER\_INT, JIGGLE\_X, JIGGLE\_Y,
         JIGGLE\_COUNT, JIGGLE\_X\_MAX, JIGGLE\_X\_MIN, JIGGLE\_Y\_MAX,
         JIGGLE\_Y\_MIN, STATUS)
   }
   \sstarguments{
      \sstsubsection{
         JIGGLE\_NAME            = CHARACTER$*$($*$) (Given)
      }{
         the name of the file containing the jiggle pattern
      }
      \sstsubsection{
         MAX\_JIGGLE             = INTEGER (Given)
      }{
         the maximum number of jiggles that can be read
      }
      \sstsubsection{
         JIGGLE\_P\_SWITCH        = INTEGER (Given and returned)
      }{
         the requested (given) and actual (returned) number of jiggles
         that will be performed in each exposure of the integration
      }
      \sstsubsection{
         JIGGLE\_REPEAT          = INTEGER (Returned)
      }{
         the number of times the jiggle pattern wil be repeated in each
         exposure
      }
      \sstsubsection{
         EXP\_PER\_INT            = INTEGER (Returned)
      }{
         the number of exposures required per integration
      }
      \sstsubsection{
         JIGGLE\_X (MAX\_JIGGLE)  = REAL (Returned)
      }{
         the x jiggle offsets
      }
      \sstsubsection{
         JIGGLE\_Y (MAX\_JIGGLE)  = REAL (Returned)
      }{
         the y jiggle offsets
      }
      \sstsubsection{
         JIGGLE\_COUNT           = INTEGER (Returned)
      }{
         the number of jiggles in the pattern
      }
      \sstsubsection{
         JIGGLE\_X\_MAX           = REAL (Returned)
      }{
         the maximum jiggle offset in the x axis
      }
      \sstsubsection{
         JIGGLE\_X\_MIN
      }{
         the minimum jiggle offset in the x axis
      }
      \sstsubsection{
         JIGGLE\_Y\_MAX           = REAL (Returned)
      }{
         the maximum jiggle offset in the y axis
      }
      \sstsubsection{
         JIGGLE\_Y\_MIN
      }{
         the minimum jiggle offset in the y axis
      }
      \sstsubsection{
         STATUS =       INTEGER (Given and returned)
      }{
         global status
      }
   }
   \sstdiytopic{
      Copyright
   }{
      Copyright \copyright 1995,1996,1997,1998,1999 Particle Physics and Astronomy
      Research Council. All Rights Reserved.
   }
}
\sstroutine{
   SCULIB\_READ\_NUMBERS
}{
   routine to read numbers from an ASCII file
}{
   \sstdescription{
      This routine reads a file of ASCII numbers. The name of the file
      is specified by FILENAME, and the routine assumes that the numbers
      are in NARRAY columns each containing SIZE numbers. The maximum
      value for NARRAY is 3. The read operation itself is free format.
      An error will be returned if there is trouble reading from the
      file. A warning will be given if there are more numbers in the
      file than can be read into the arrays.
   }
   \sstinvocation{
       CALL SCULIB\_READ\_NUMBERS (FILENAME, NARRAY, SIZE,
          ARRAY1, ARRAY2, ARRAY3, LENGTH, STATUS)
   }
   \sstarguments{
      \sstsubsection{
         FILENAME         = CHARACTER$*$($*$) (Given)
      }{
         Name of file containing numbers
      }
      \sstsubsection{
         NARRAY           = INTEGER (Given)
      }{
         Number of columns in file ($<$=3)
      }
      \sstsubsection{
         SIZE             = INTEGER (Given)
      }{
         Dimension of arrays
      }
      \sstsubsection{
         ARRAY1 (SIZE)    = REAL (Returned)
      }{
         Output array to contain first column of numbers
      }
      \sstsubsection{
         ARRAY2 (SIZE)    = REAL (Returned)
      }{
         Output array to contain second column of numbers (if present)
      }
      \sstsubsection{
         ARRAY3 (SIZE)    = REAL (Returned)
      }{
         Output array to contain third column of numbers (if present)
      }
      \sstsubsection{
         LENGTH           = INTEGER (Returned)
      }{
         Number of items read into output arrays
      }
      \sstsubsection{
         STATUS = INTEGER (Given and returned)
      }{
         Global status
      }
   }
   \sstdiytopic{
      Copyright
   }{
      Copyright \copyright 1995,1996,1997,1998,1999 Particle Physics and Astronomy
      Research Council. All Rights Reserved.
   }
}
\sstroutine{
   SCULIB\_READ\_SKY
}{
   read sky parameters from a named file
}{
   \sstdescription{
      This routine reads sky and telescope parameters for each SCUBA
      sub-instrument/filter combination from an ASCII file named in FILE.
      The format of each line in the file that contains sky information is
      assumed to be:-

         FIT1 $<$qual$>$ $<$sub$>$ $<$filter$>$ $<$eta\_tel$>$ $<$b$>$ $<$tauz$>$ $<$date$>$ $<$day$>$ $<$run$>$

      The values returned for a given filter/sub-instrument will be from the
      good quality entry ($<$qual$>$ = 0) with the highest associated $<$day$>$ number
      in the file.

      Errors will be reported and bad status returned if;

      \sstitemlist{

         \sstitem
           there is an error opening or reading from the file

         \sstitem
           $<$qual$>$ does not convert to an integer

         \sstitem
           $<$sub$>$ is not one of SHORT, LONG, P1100, P1300, P2000

         \sstitem
           $<$eta\_tel$>$ does not convert to a real

         \sstitem
           $<$b$>$ does not convert to a real

         \sstitem
           $<$tauz$>$ does not convert to a real

         \sstitem
           $<$run$>$ does not convert to an integer

         \sstitem
           More than MAX\_SKY combinations are read.

      }
      Blank lines are ignored as are those parts of lines following
      a ! character (comments) and lines not beginning with the word FIT1.

      This routine reads a file written by SCUDR\_END\_SKYDIP, so these 2
      routines should be changed together when necessary.

      If the file name is {\tt '}NULL{\tt '} (case insensitive) then the sky parameters
      will be reset to null values.
   }
   \sstinvocation{
       CALL SCULIB\_READ\_SKY (FILE, MAX\_SKY, N\_SKY, SUB,
         FILTER, ETA\_TEL, B, TAUZ, DATEM, DAY, RUN, STATUS)
   }
   \sstarguments{
      \sstsubsection{
         FILE                              = CHARACTER$*$($*$) (Given)
      }{
         the name of the file
      }
      \sstsubsection{
         MAX\_SKY                           = INTEGER (Given)
      }{
         the maximum number of sub-instrument/filter combinations
      }
      \sstsubsection{
         N\_SKY                             = INTEGER (Returned)
      }{
         the number of sub-instrument/filter combinations read
      }
      \sstsubsection{
         SUB (MAX\_SKY)                     = CHARACTER$*$($*$) (Returned)
      }{
         the name of the sub-instrument
      }
      \sstsubsection{
         FILTER (MAX\_SKY)                  = CHARACTER$*$($*$) (Returned)
      }{
         the name of the filter
      }
      \sstsubsection{
         ETA\_TEL (MAX\_SKY)                 = REAL (Returned)
      }{
         telescope transmission
      }
      \sstsubsection{
         B (MAX\_SKY)                       = REAL (Returned)
      }{
         atmospheric bandwidth b factor
      }
      \sstsubsection{
         TAUZ (MAX\_SKY)                    = REAL (Returned)
      }{
         zenith atmospheric optical depth
      }
      \sstsubsection{
         DATEM (MAX\_SKY)                   = CHARACTER$*$($*$) (Returned)
      }{
         the date on which a measurement was made
      }
      \sstsubsection{
         DAY (MAX\_SKY)                     = DOUBLE PRECISION (Returned)
      }{
         the date and time at which the measurement was made, measured
         as a day number from 1st Jan
      }
      \sstsubsection{
         RUN (MAX\_SKY)                     = INTEGER (Returned)
      }{
         the run number of a measurement
      }
      \sstsubsection{
         STATUS                            = INTEGER (Given and returned)
      }{
         global status
      }
   }
   \sstdiytopic{
      Copyright
   }{
      Copyright \copyright 1995,1996,1997,1998,1999 Particle Physics and Astronomy
      Research Council. All Rights Reserved.
   }
}
\sstroutine{
   SCULIB\_READ\_TAUZ
}{
   read sky zenith optical depths from a named file
}{
   \sstdescription{
      This routine reads filter names and associated sky zenith opacities
      from an ASCII file named in FILE. The format of each line in the file
      that contains sky information is assumed to be:-

       FIT1 $<$qual$>$ $<$sub\_inst$>$ $<$filter$>$ $<$eta\_tel$>$ $<$b$>$ $<$tauz$>$ $<$date$>$ $<$day$>$ $<$run$>$

      Errors will be reported and bad status returned if;

      \sstitemlist{

         \sstitem
           there is an error opening or reading the file

         \sstitem
           $<$qual$>$ is not an integer

         \sstitem
           $<$sub-inst$>$ is not one of SHORT, LONG, P1100, P1300, P2000

         \sstitem
           $<$eta\_tel$>$ does not convert to a real

         \sstitem
           $<$b$>$ does not convert to a real

         \sstitem
           $<$tauz$>$ does not convert to a real

         \sstitem
           $<$day$>$ does not convert to a double

         \sstitem
           $<$run$>$ does not convert to an integer

      }
      The value of tauz returned for a given filter will be the good
      quality entry ($<$qual$>$ = 0) with the highest associated $<$day$>$ number
      in the file.

      Blank lines are ignored as are those parts of lines following
      a ! character (comments) and any line whose first line is not
      FIT1.

      This routine reads a file written by SCUDR\_END\_SKYDIP, so these 2
      routines should be changed together when necessary.
   }
   \sstinvocation{
       CALL SCULIB\_READ\_TAUZ (FILE, MAX\_FILT, N\_FILT, FILTER,
         TAUZ, DATEM, DAY, RUN, STATUS)
   }
   \sstarguments{
      \sstsubsection{
         FILE                              = CHARACTER$*$($*$) (Given)
      }{
         the name of the file
      }
      \sstsubsection{
         MAX\_FILT                          = INTEGER (Given)
      }{
         the maximum number of filters
      }
      \sstsubsection{
         N\_FILT                            = INTEGER (Returned)
      }{
         the number of filters read
      }
      \sstsubsection{
         FILTER (MAX\_FILT)                 = CHARACTER$*$($*$) (Returned)
      }{
         the names of the filters
      }
      \sstsubsection{
         TAUZ (MAX\_FILT)                   = REAL (Returned)
      }{
         zenith atmospheric optical depth for that filter
      }
      \sstsubsection{
         DATEM (MAX\_FILT)                  = CHARACTER$*$($*$) (Returned)
      }{
         the date on which the measurement was made
      }
      \sstsubsection{
         DAY (MAX\_FILT)                    = DOUBLE PRECISION (Returned)
      }{
         the date and time at which the measurement was made, measured
         as a day number from 1st Jan
      }
      \sstsubsection{
         RUN (MAX\_FILT)                    = INTEGER (Returned)
      }{
         the run number of the measurement
      }
      \sstsubsection{
         STATUS                            = INTEGER (Given and returned)
      }{
         global status
      }
   }
   \sstdiytopic{
      Copyright
   }{
      Copyright \copyright 1995,1996,1997,1998,1999 Particle Physics and Astronomy
      Research Council. All Rights Reserved.
   }
}
\sstroutine{
   SCULIB\_READBOLS
}{
   read bolometer data from a name file
}{
   \sstdescription{
      This routine reads bolometer data from an ASCII file named in
      FILE. The format of the file is assumed to be:-

\texttt{
       proc SET\_BOLS \\
       \{ flat field data file written by SCUDR \\
       \{ Tues Apr 20 18:09:20 1993\\
       \{ observation run number 6\\
       \{ long-wave array\\
         SETBOL name type du3 du4 calib theta a b qual day run ref\\
            "     ditto for other long-wave array bolometers    "\\
       \{ short-wave array
         SETBOL name type du3 du4 calib theta a b qual day run ref\\
            "     ditto for other short-wave array bolometers   "\\
       \{ P1100\\
         SETBOL name type du3 du4 calib theta a b qual day run ref\\
       \{ P1300\\
         SETBOL name type du3 du4 calib theta a b qual day run ref\\
       \{ P2000\\
         SETBOL name type du3 du4 calib theta a b qual day run ref\\
       end proc\\
}

      where:-

$<$name$>$ is the name of a bolometer (e.g. a16), case insensitive\\
$<$type$>$ is the type of the bolometer, e.g. SHORT, LONG, P1100, P1300,
          P2000\\
$<$du3$>$ is the Nasmyth DU3 coordinate of the bolometer\\
$<$du4$>$ is the Nasmyth DU4 coordinate of the bolometer\\
$<$calib$>$ is the flat-field value of the bolometer\\
$<$theta$>$ is the angle between x axis an `a{\tt '} axis of fitted ellipse\\
$<$a$>$ is the semi-length of the `a{\tt '} axis\\
$<$b$>$ is the semi-length of the `b{\tt '} axis\\
$<$qual$>$ is the quality of the bolometer; 0=good, 1=bad\\
$<$day$>$ is the date when the information was measured, in days from 1 Jan\\
$<$run$>$ is the number of the run when the information was measured\\
$<$ref$>$ is the name of the reference bolometer\\

      Values for bolometers that are not specified in the file will be set
      to default values; {\tt '}BAD{\tt '} for strings, 0 for all other values except
      BOL\_QUAL which will be 1.

      Errors will be reported and bad status returned if:-

      \sstitemlist{

         \sstitem
          $<$name$>$ is not a valid bolometer name

         \sstitem
          $<$du3$>$, $<$du4$>$, $<$calib$>$, $<$theta$>$, $<$a$>$, $<$b$>$ do not convert to reals

         \sstitem
          $<$qual$>$, $<$run$>$ do not convert to integers

         \sstitem
          $<$day$>$ does not convert to double

         \sstitem
          An attempt is made to set values for a bolometer twice

      }
      Each line in the file will be converted to upper case. Characters to the
      right of a \{ character in a line will be treated as comments and ignored.
   }
   \sstinvocation{
       CALL SCULIB\_READBOLS (FILE, NUM\_CHAN, NUM\_ADC, BOL\_TYPE,
         BOL\_DU3, BOL\_DU4, BOL\_CALIB, BOL\_THETA, BOL\_A, BOL\_B,
         BOL\_QUAL, BOL\_DAY, BOL\_RUN, BOL\_REF, STATUS)
   }
   \sstarguments{
      \sstsubsection{
         FILE                              = CHARACTER$*$($*$) (Given)
      }{
         the name of the file
      }
      \sstsubsection{
         NUM\_CHAN                          = INTEGER (Given)
      }{
         number of channels per ADC
      }
      \sstsubsection{
         NUM\_ADC                           = INTEGER (Given)
      }{
         number of ADC cards
      }
      \sstsubsection{
         BOL\_TYPE (NUM\_CHAN, NUM\_ADC)      = CHARACTER$*$($*$) (Returned)
      }{
         the type of bolometer; SHORT, LONG, P1100, etc.
      }
      \sstsubsection{
         BOL\_DU3 (NUM\_CHAN, NUM\_ADC)       = REAL (Returned)
      }{
         the DU3 Nasmyth coordinate of the bolometer
      }
      \sstsubsection{
         BOL\_DU4 (NUM\_CHAN, NUM\_ADC)       = REAL (Returned)
      }{
         the DU4 Nasmyth coordinate of the bolometer
      }
      \sstsubsection{
         BOL\_CALIB (NUM\_CHAN, NUM\_ADC)     = REAL (Returned)
      }{
         the target calibration of the bolometer
      }
      \sstsubsection{
         BOL\_THETA (NUM\_CHAN, NUM\_ADC)     = REAL (Returned)
      }{
         angle between x axis and `a{\tt '} axis of fitted ellipse (radians)
      }
      \sstsubsection{
         BOL\_A (NUM\_CHAN, NUM\_ADC)         = REAL (Returned)
      }{
         half length of `a{\tt '} axis of fitted ellipse
      }
      \sstsubsection{
         BOL\_B (NUM\_CHAN, NUM\_ADC)         = REAL (Returned)
      }{
         half length of `b{\tt '} axis of fitted ellipse
      }
      \sstsubsection{
         BOL\_QUAL (NUM\_CHAN, NUM\_ADC)      = INTEGER (Returned)
      }{
         the bolometer quality
      }
      \sstsubsection{
         BOL\_DAY (NUM\_CHAN, NUM\_ADC)       = DOUBLE PRECISION (Returned)
      }{
         the date when the bolometer was measured
      }
      \sstsubsection{
         BOL\_RUN (NUM\_CHAN, NUM\_ADC)       = INTEGER (Returned)
      }{
         the run when the bolometer was measured
      }
      \sstsubsection{
         BOL\_REF (NUM\_CHAN, NUM\_ADC)       = CHARACTER$*$($*$) (Returned)
      }{
         the name of the reference bolometer used
      }
      \sstsubsection{
         STATUS                            = INTEGER (Given and returned)
      }{
         global status
      }
   }
   \sstdiytopic{
      Copyright
   }{
      Copyright \copyright 1995,1996,1997,1998,1999 Particle Physics and Astronomy
      Research Council. All Rights Reserved.
   }
}
\sstroutine{
   SCULIB\_REDUCE\_SWITCH
}{
   reduce the demodulated data from the switches of
   an exposure into the exposure result
}{
   \sstdescription{
      This routine reduces the switches in an exposure to give the
      exposure result, and returns the weight to be given to the
      specified projected bolometer position when combining different
      projected bolometers into the result for a sub-instrument.
        Exposure results are calculated assuming that:-

      For a SQUARE chop-function -\\
       switch 1 has the object in the R beam of the `middle{\tt '} projected
       bolometer and in the L beam of the `right{\tt '} bolometer.\\
       switch 2 has the object in the L beam of the `middle{\tt '} bolometer
       and in R beam of the `left{\tt '} bolometer.

      For a TRIPOS (3-position) chop-function -\\
       switch 1 has the object in the M beam of the `middle{\tt '} bolometer,
       in the L beam of the right and the R beam of the `left{\tt '}.\\
       switch 2 has the object in the R beam of the `middle{\tt '} bolometer
       and in the M beam of the `right{\tt '}.\\
       switch 3 has the object in the L beam of the `middle{\tt '} bolometer
       and in the M beam of the `left{\tt '}.

      The reduction method depends on the number of switches taken per
      exposure, the chop function used and the projected beam that the
      bolometers are assumed to be in.

      SWITCH\_PER\_EXP = 1\\
         CHOP\_FUN = SQUARE (SCUBAWAVE or RAMPWAVE)\\
            left bolometer = BAD (was never looking at source), weight =  0.0\\
          middle bolometer = switch 1, weight = 1.0\\
           right bolometer = - switch 1, weight = 1.0\\
         CHOP\_FUN = TRIPOS\\
            left bolometer = - 2 $*$ switch 1, weight = 0.5\\
          middle bolometer = switch 1, weight = 1.0\\
           right bolometer = - 2 $*$ switch 1, weight = 0.5\\

      SWITCH\_PER\_EXP = 2\\
         CHOP\_FUN = SQUARE\\
            left bolometer = switch 2 - switch 1, weight = 0.5\\
          middle bolometer = (switch 1 - switch 2) / 2, weight = 1.0\\
           right bolometer = switch 2 - switch 1, weight = 0.5\\
         CHOP\_FUN = TRIPOS\\
            left bolometer = - 2 $*$ (switch 1 - switch 2), weight = 0.5\\
          middle bolometer = 2/3 $*$ (switch 1 - switch 2), weight = 1.5\\
           right bolometer = -2/3 $*$ (switch 1 - switch 2), weight = 1.5\\

      SWITCH\_PER\_EXP = 3\\
         CHOP\_FUN = SQUARE\\
            error\\
         CHOP\_FUN = TRIPOS\\
            left bolometer = -1/2 $*$ (2 $*$ switch 1 - (switch 2 $+$ switch 3)),
                             weight = 2/3\\
          middle bolometer = 1/3 $*$ (2 $*$ switch 1 - (switch 2 $+$ switch 3)),
                             weight = 1.0\\
           right bolometer = -1/2 $*$ (2 $*$ switch 1 - (switch 2 $+$ switch 3)),
                             weight = 2/3

      Any other combinations of parameters will give rise to an error report
      and the routine will return with bad status.
   }
   \sstinvocation{
       CALL SCULIB\_REDUCE\_SWITCH (CHOP\_FUN, SWITCH\_PER\_EXP,
         N\_DATA,
         SWITCH\_1\_DATA, SWITCH\_1\_VARIANCE, SWITCH\_1\_QUALITY,
         SWITCH\_2\_DATA, SWITCH\_2\_VARIANCE, SWITCH\_2\_QUALITY,
         SWITCH\_3\_DATA, SWITCH\_3\_VARIANCE, SWITCH\_3\_QUALITY,
         BEAM,
         EXP\_DATA, EXP\_VARIANCE, EXP\_QUALITY, WEIGHT, STATUS)
   }
   \sstarguments{
      \sstsubsection{
         CHOP\_FUN                              = CHARACTER$*$($*$) (Given)
      }{
         the chop function used
      }
      \sstsubsection{
         SWITCH\_PER\_EXP                        = INTEGER (Given)
      }{
         the number of switches per exposure
      }
      \sstsubsection{
         N\_DATA                                = INTEGER (Given)
      }{
         the number of measurements
      }
      \sstsubsection{
         SWITCH\_1\_DATA (N\_DATA)                = REAL (Given)
      }{
         data for switch 1
      }
      \sstsubsection{
         SWITCH\_1\_VARIANCE (N\_DATA)            = REAL (Given)
      }{
         variance for switch 1
      }
      \sstsubsection{
         SWITCH\_1\_QUALITY (N\_DATA)             = BYTE (Given)
      }{
         quality for switch 1
      }
      \sstsubsection{
         SWITCH\_2\_DATA (N\_DATA)                = REAL (Given)
      }{
         likewise for switch 2
      }
      \sstsubsection{
         SWITCH\_2\_VARIANCE (N\_DATA)            = REAL (Given)
      }{
      }
      \sstsubsection{
         SWITCH\_2\_QUALITY (N\_DATA)             = BYTE (Given)
      }{
      }
      \sstsubsection{
         SWITCH\_3\_DATA (N\_DATA)                = REAL (Given)
      }{
         and switch 3
      }
      \sstsubsection{
         SWITCH\_3\_VARIANCE (N\_DATA)            = REAL (Given)
      }{
      }
      \sstsubsection{
         SWITCH\_3\_QUALITY (N\_DATA)             = BYTE (Given)
      }{
      }
      \sstsubsection{
         BEAM                                  = INTEGER (Given)
      }{
         the projected beam assumed for the bolometers; 1 = LEFT
         2 = MIDDLE, 3 = RIGHT
      }
      \sstsubsection{
         EXP\_DATA (N\_DATA)                     = REAL (Returned)
      }{
         the exposure result
      }
      \sstsubsection{
         EXP\_VARIANCE (N\_DATA)                 = REAL (Returned)
      }{
         the variance on the result
      }
      \sstsubsection{
         EXP\_QUALITY (N\_DATA)                  = BYTE (Returned)
      }{
         the quality on the result
      }
      \sstsubsection{
         WEIGHT                                = REAL (Returned)
      }{
         the relative weight to assign to this projected-beam compared
         to the others when adding them together to give a final result
      }
      \sstsubsection{
         STATUS                                = INTEGER (Given and returned)
      }{
         global status
      }
   }
   \sstdiytopic{
      Copyright
   }{
      Copyright \copyright 1995,1996,1997,1998,1999 Particle Physics and Astronomy
      Research Council. All Rights Reserved.
   }
}
\sstroutine{
   SCULIB\_REMOVE\_DEMOD\_INT
}{
   remove demodulated data for one sub-instrument
   in an integration from a coadded result
}{
   \sstdescription{
      This routine removes the demodulated data for a specified sub-instrument
      and integration from the coadded measurement.

        After checking status on entry the routine cycles through the
      bolometers for which data has been taken. If the bolometer belongs
      to the sub-instrument whose data is to be removed, the routine will
      then cycle through the jiggle pattern removing the integration data
      from the coadd. If this reduces the number of coadded integrations
      to 0 then all the coadd numbers are set to 0 apart from quality,
      which is set to 1. If the number of coadded integrations is reduced
      to 1 then the coadd variance can no longer be estimated from the
      spread of data about the mean and will be set to 0.
   }
   \sstinvocation{
       CALL SCULIB\_REMOVE\_DEMOD\_INT (REMOVE\_TYPE, NUM\_CHAN,
         NUM\_ADC, BOL\_TYPE, N\_BOLS, BOL\_CHAN, BOL\_ADC, J\_COUNT,
         INT\_DATA, INT\_QUALITY, COADD\_DATA, COADD\_VARIANCE,
         COADD\_QUALITY, COADD\_NUMBER, STATUS)
   }
   \sstarguments{
      \sstsubsection{
         REMOVE\_TYPE                   = CHARACTER$*$($*$) (Given)
      }{
         the type of the sub-instrument whose data is to be removed
      }
      \sstsubsection{
         NUM\_CHAN                      = INTEGER (Given)
      }{
         the number of channels per ADC in SCUBA
      }
      \sstsubsection{
         NUM\_ADC                       = INTEGER (Given)
      }{
         the number of ADCs in SCUBA
      }
      \sstsubsection{
         BOL\_TYPE (NUM\_CHAN, NUM\_ADC)  = CHARACTER$*$($*$) (Given)
      }{
         the types of the bolometers in SCUBA
      }
      \sstsubsection{
         N\_BOLS                        = INTEGER (Given)
      }{
         the number of bolometers being measured
      }
      \sstsubsection{
         BOL\_CHAN (N\_BOLS)             = INTEGER (Given)
      }{
         the channel numbers of the bolometers being measured
      }
      \sstsubsection{
         BOL\_ADC (N\_BOLS)              = INTEGER (Given)
      }{
         the ADC numbers of the bolometers being measured
      }
      \sstsubsection{
         J\_COUNT                       = INTEGER (Given)
      }{
         the number of jiggle positions in the pattern
      }
      \sstsubsection{
         INT\_DATA (N\_BOLS, J\_COUNT)    = REAL (Given)
      }{
         the integration data
      }
      \sstsubsection{
         INT\_QUALITY (N\_BOLS, J\_COUNT) = INTEGER (Given)
      }{
         quality on the integration data
      }
      \sstsubsection{
         COADD\_DATA (N\_BOLS, J\_COUNT)  = REAL (Given and returned)
      }{
         the mean of the coadded data
      }
      \sstsubsection{
         COADD\_VARIANCE (N\_BOLS, J\_COUNT)
      }{
                                 = REAL (Given and returned)
         the variance on the mean of the coadded data
      }
      \sstsubsection{
         COADD\_QUALITY (N\_BOLS, J\_COUNT)
      }{
                                 = INTEGER (Given and returned)
         the quality on the mean of the coadded data
      }
      \sstsubsection{
         COADD\_NUMBER (N\_BOLS, J\_COUNT)= INTEGER (Given and returned)
      }{
         the number of data coadded
      }
      \sstsubsection{
         STATUS                        = INTEGER (Given and returned)
      }{
         the global status
      }
   }
   \sstdiytopic{
      Copyright
   }{
      Copyright \copyright 1995,1996,1997,1998,1999 Particle Physics and Astronomy
      Research Council. All Rights Reserved.
   }
}
\sstroutine{
   SCULIB\_REMOVE\_LINEAR\_BASELINE
}{
   Remove linear baseline from each exposure
}{
   \sstdescription{
      This routine takes a data array. It then removes a linear baseline
      from each scan. The ends of the scan are the size of the chop throw
      (should be no source at the ends of a scan)
   }
   \sstinvocation{
       CALL SCULIB\_REMOVE\_LINEAR\_BASELINE(DORLB, N\_EXPOSURES,
            N\_INTEGRATIONS, N\_MEASUREMENTS, DEM\_PNTR, N\_BOL, N\_POS,
            IN\_DATA, IN\_QUALITY, SAMPLE\_DX, CHOP\_THROW, OUT\_DATA,
            OUT\_QUALITY, BADBIT, STATUS)
   }
   \sstarguments{
      \sstsubsection{
         DORLB                       = LOGICAL (Given)
      }{
         control whether we are subtracting the baseline (TRUE)
         or storing the basline (FALSE)
      }
      \sstsubsection{
         N\_EXPOSURES                 = INTEGER (Given)
      }{
         maximum number of exposures per integration
      }
      \sstsubsection{
         N\_INTEGRATIONS              = INTEGER (Given)
      }{
         number of integrations in the observation
      }
      \sstsubsection{
         N\_MEASUREMENTS              = INTEGER (Given)
      }{
         number of measurements in the observation
      }
      \sstsubsection{
         DEMOD\_POINTER (N\_EXPOSURES, N\_INTEGRATIONS, N\_MEASUREMENTS) = INTEGER (Given)
      }{
         array pointing to start and finish of scans in IN\_DATA
      }
      \sstsubsection{
         N\_BOL                       = INTEGER (Given)
      }{
         the number of bolometers for which data was taken
      }
      \sstsubsection{
         N\_POS                       = INTEGER (Given)
      }{
         the number of positions measured in the scan
      }
      \sstsubsection{
         IN\_DATA (N\_BOL, N\_POS)      = REAL (Given)
      }{
         the measured data
      }
      \sstsubsection{
         IN\_VARIANCE (N\_BOL, N\_POS)      = REAL (Given)
      }{
         the measured variance
      }
      \sstsubsection{
         IN\_QUALITY (N\_BOL, N\_POS)   = BYTE (Given)
      }{
         the quality on IN\_DATA
      }
      \sstsubsection{
         NSTART                      = INTEGER (Given)
      }{
         Number of points used for fit at start of scan
      }
      \sstsubsection{
         NEND                        = INTEGER (Given)
      }{
         Number of points from end of scan
      }
      \sstsubsection{
         OUT\_DATA (N\_BOL, N\_POS)     = REAL (Returned)
      }{
         the data with baseline removed
      }
      \sstsubsection{
         OUT\_VARIANCE (N\_BOL, N\_POS)      = REAL (Given)
      }{
         the output variance
      }
      \sstsubsection{
         OUT\_QUALITY (N\_BOL, N\_POS)  = BYTE (Returned)
      }{
         the quality on OUT\_DATA
      }
      \sstsubsection{
         BADBIT                      = BYTE (Given)
      }{
         bad bit mask
      }
      \sstsubsection{
         STATUS = INTEGER (Given and Returned)
      }{
         Global Status value
      }
   }
   \sstdiytopic{
      Copyright
   }{
      Copyright \copyright 1995,1996,1997,1998,1999 Particle Physics and Astronomy
      Research Council. All Rights Reserved.
   }
}
\sstroutine{
   SCULIB\_REMOVE\_OPACITY
}{
   remove sky opacity from demodulated data
}{
   \sstdescription{
      This routine corrects demodulated data for the effect of sky opacity.
      All data is corrected assuming the same airmass of observation, but
      the zenith sky opacity used will depend which SCUBA sub-instrument each
      bolometer belongs to.

        After checking status on entry, the routine will loop through the
      bolometers measured. For each bolometer it will then ascertain the
      appropriate zenith sky opacity from the parent sub-instrument and
      wavelength of observation. The total sky optical depth will then
      be calculated and, if this is in the range 0 to 20, the flux correction
      factor derived. The routine will then loop through the positions
      measured by this bolometer, correcting the fluxes and variances.
      If the bolometer data quality was bad, or if sky optical depth lay
      outside the above range, the data quality will be set bad and no
      correction applied.
   }
   \sstinvocation{
       CALL SCULIB\_REMOVE\_OPACITY (N\_BOLS, N\_POS,
         BOL\_SELECT\_CHAN, BOL\_SELECT\_ADC, NUM\_CHAN, NUM\_ADC, BOL\_TYPE,
         N\_SUB, SUB\_INSTRUMENT, TAUZ, AIRMASS, EXP\_DATA, EXP\_VARIANCE,
         EXP\_QUALITY, STATUS)
   }
   \sstarguments{
      \sstsubsection{
         N\_BOLS                        = INTEGER (Given)
      }{
         the number of bolometers measured
      }
      \sstsubsection{
         N\_POS                         = INTEGER (Given)
      }{
         the number of positions they were measured at
      }
      \sstsubsection{
         BOL\_SELECT\_CHAN (N\_BOLS)      = INTEGER (Given)
      }{
         the channel numbers of the measured bolometers
      }
      \sstsubsection{
         BOL\_SELECT\_ADC (N\_BOLS)       = INTEGER (Given)
      }{
         the ADC numbers of the measured bolometers
      }
      \sstsubsection{
         NUM\_CHAN                      = INTEGER (Given)
      }{
         the number of channels per ADC
      }
      \sstsubsection{
         NUM\_ADC                       = INTEGER (Given)
      }{
         the number of ADCs
      }
      \sstsubsection{
         BOL\_TYPE (NUM\_CHAN, NUM\_ADC)  = CHARACTER$*$($*$) (Given)
      }{
         the type of the bolometer on each channel
      }
      \sstsubsection{
         N\_SUB                         = INTEGER (Given)
      }{
         the number of SCUBA sub-instruments being measured
      }
      \sstsubsection{
         SUB\_INSTRUMENT (N\_SUB)        = CHARACTER$*$($*$)
      }{
         the names of the sub-instruments being measured
      }
      \sstsubsection{
         TAUZ (NSUB)                   = REAL (Given)
      }{
         the zenith sky opacity for each sub-instrument
      }
      \sstsubsection{
         AIRMASS                       = REAL (Given)
      }{
         the airmass at which the observations were made
      }
      \sstsubsection{
         EXP\_DATA (N\_BOLS, N\_POS)      = REAL (Given and returned)
      }{
         the bolometer data
      }
      \sstsubsection{
         EXP\_VARIANCE (N\_BOLS, N\_POS)  = REAL (Given and returned)
      }{
         the variance on the data
      }
      \sstsubsection{
         EXP\_QUALITY (N\_BOLS, N\_POS)   = REAL (Given and returned)
      }{
         the quality on the data
      }
      \sstsubsection{
         STATUS                        = INTEGER (Given and returned)
      }{
         global status
      }
   }
   \sstdiytopic{
      Copyright
   }{
      Copyright \copyright 1995,1996,1997,1998,1999 Particle Physics and Astronomy
      Research Council. All Rights Reserved.
   }
}
\sstroutine{
   SCULIB\_REM\_SKY
}{
   To remove sky background from SCUBA data
}{
   \sstdescription{
      This routine goes through each jiggle position in turn, finding
      the sky background level by using the selected SKY bolometers.
      This level is removed from each bolometer.
      The background level is determined by the median of the SKY data
      for each jiggle.
      If requested, via ADD\_BACK, the mean of the level that was
      removed is added back onto the data so that flux is conserved
      in the image.
   }
   \sstinvocation{
       CALL SCULIB\_REM\_SKY(MODE, ADD\_BACK, N\_BOLS, N\_POS, SCUDATA,
            SCUVAR, SCUQUAL, CLIP, N\_SKYBOLS, SKYBOLS, BADBIT, STATUS)
   }
   \sstarguments{
      \sstsubsection{
         MODE = CHARACTER $*$ ( $*$ ) (Given)
      }{
         Clip mode
      }
      \sstsubsection{
         ADD\_BACK = LOGICAL (Given)
      }{
         Add the mean removed level back onto the final data set. (if TRUE)
      }
      \sstsubsection{
         N\_BOLS = INTEGER (Given)
      }{
         Number of bolometers in data set
      }
      \sstsubsection{
         N\_POS = INTEGER (Given)
      }{
         Number of jiggle positions in data set
      }
      \sstsubsection{
         SCUDATA(N\_BOLS, N\_POS) = REAL (Given \& Returned)
      }{
         The data
      }
      \sstsubsection{
         SCUVAR(N\_BOLS, N\_POS) = REAL (Given \& Returned)
      }{
         The variance on the data
      }
      \sstsubsection{
         SCUQUAL(N\_BOLS, N\_POS) = BYTE (Given \& Returned)
      }{
         The data quality
      }
      \sstsubsection{
         CLIP = \_REAL (Given)
      }{
         Iterative sigma clipping level
      }
      \sstsubsection{
         N\_SKYBOLS = INTEGER (Given)
      }{
         Number of sky bolometers
      }
      \sstsubsection{
         SKYBOLS (N\_SKYBOLS) = INTEGER (Given)
      }{
         List of sky bolometers
      }
      \sstsubsection{
         BADBIT = BYTE (Given)
      }{
         Bad bit mask
      }
      \sstsubsection{
         STATUS = INTEGER (Given and Returned)
      }{
         Global Status value
      }
   }
   \sstdiytopic{
      Copyright
   }{
      Copyright \copyright 1995,1996,1997,1998,1999 Particle Physics and Astronomy
      Research Council. All Rights Reserved.
   }
}
\sstroutine{
   SCULIB\_REWRITE\_FITS\_C
}{
   rewrite the value of specified FITS character keyword
}{
   \sstdescription{
      This routine will overwrite the value of a specified FITS character
      keyword held in the FITS extension of an NDF file. The FITS extension
      must have been read into the input array FITS before this routine is
      called and be written out again afterwards for the change to take effect.

         The routine assumes that each line in the FITS array will contain
      a string with format:-

       {\tt "}KEYWORD={\tt '}VALUE{\tt '}           / this is a comment{\tt "}

      It will search the input array for a line containing the required
      keyword and replace VALUE. If the keyword is not found an error will
      be reported and bad status returned. If the keyword is found but the
      line does not conform to the above format an error will be reported
      and bad status returned.
   }
   \sstinvocation{
       CALL SCULIB\_REWRITE\_FITS\_C (MAX\_FITS, N\_FITS, FITS, NAME,
         VALUE, STATUS)
   }
   \sstarguments{
      \sstsubsection{
         MAX\_FITS             = INTEGER (Given)
      }{
         the maximum number of items in the FITS array
      }
      \sstsubsection{
         N\_FITS               = INTEGER (Given)
      }{
         the actual number of items in the FITS array
      }
      \sstsubsection{
         FITS (MAX\_FITS)      = CHARACTER$*$($*$) (Given and returned)
      }{
         array containing the FITS items
      }
      \sstsubsection{
         NAME                 = CHARACTER$*$($*$) (Given)
      }{
         the name of the FITS keyword whose value is to be changed
      }
      \sstsubsection{
         VALUE                = CHARACTER$*$($*$) (Given)
      }{
         the value that the FITS keyword is to have
      }
      \sstsubsection{
         STATUS               = INTEGER (Given and returned)
      }{
         global status
      }
   }
   \sstdiytopic{
      Copyright
   }{
      Copyright \copyright 1995,1996,1997,1998,1999 Particle Physics and Astronomy
      Research Council. All Rights Reserved.
   }
}
\sstroutine{
   SCULIB\_REWRITE\_FITS\_I
}{
   rewrite the value of specified FITS integer keyword
}{
   \sstdescription{
      This routine will overwrite the value of a specified FITS integer
      keyword held in the FITS extension of an NDF file. The FITS extension
      must have been read into the input array FITS before this routine is
      called and be written out again afterwards for the change to take effect.

         The routine assumes that each line in the FITS array will contain
      a string with format:-

       {\tt "}KEYWORD=VALUE           / this is a comment{\tt "}

      It will search the input array for a line containing the required
      keyword and replace VALUE. If the keyword is not found an error will
      be reported and bad status returned. If the keyword is found but the
      line does not conform to the above format an error will be reported
      and bad status returned.
   }
   \sstinvocation{
       CALL SCULIB\_REWRITE\_FITS\_I (MAX\_FITS, N\_FITS, FITS, NAME,
         VALUE, STATUS)
   }
   \sstarguments{
      \sstsubsection{
         MAX\_FITS             = INTEGER (Given)
      }{
         the maximum number of items in the FITS array
      }
      \sstsubsection{
         N\_FITS               = INTEGER (Given)
      }{
         the actual number of items in the FITS array
      }
      \sstsubsection{
         FITS (MAX\_FITS)      = CHARACTER$*$($*$) (Given and returned)
      }{
         array containing the FITS items
      }
      \sstsubsection{
         NAME                 = CHARACTER$*$($*$) (Given)
      }{
         the name of the FITS keyword whose value is to be rewritten
      }
      \sstsubsection{
         VALUE                = INTEGER (Given)
      }{
         the value that the FITS keyword is to have
      }
      \sstsubsection{
         STATUS               = INTEGER (Given and returned)
      }{
         global status
      }
   }
   \sstdiytopic{
      Copyright
   }{
      Copyright \copyright 1995,1996,1997,1998,1999 Particle Physics and Astronomy
      Research Council. All Rights Reserved.
   }
}
\sstroutine{
   SCULIB\_REWRITE\_FITS\_R
}{
   rewrite the value of specified FITS real keyword
}{
   \sstdescription{
      This routine will overwrite the value of a specified FITS real
      keyword held in the FITS extension of an NDF file. The FITS extension
      must have been read into the input array FITS before this routine is
      called and be written out again afterwards for the change to take effect.

         The routine assumes that each line in the FITS array will contain
      a string with format:-

       {\tt "}KEYWORD=VALUE           / this is a comment{\tt "}

      It will search the input array for a line containing the required
      keyword and replace VALUE. If the keyword is not found an error will
      be reported and bad status returned. If the keyword is found but the
      line does not conform to the above format an error will be reported
      and bad status returned.
   }
   \sstinvocation{
       CALL SCULIB\_REWRITE\_FITS\_R (MAX\_FITS, N\_FITS, FITS, NAME,
         VALUE, STATUS)
   }
   \sstarguments{
      \sstsubsection{
         MAX\_FITS             = INTEGER (Given)
      }{
         the maximum number of items in the FITS array
      }
      \sstsubsection{
         N\_FITS               = INTEGER (Given)
      }{
         the actual number of items in the FITS array
      }
      \sstsubsection{
         FITS (MAX\_FITS)      = CHARACTER$*$($*$) (Given and returned)
      }{
         array containing the FITS items
      }
      \sstsubsection{
         NAME                 = CHARACTER$*$($*$) (Given)
      }{
         the name of the FITS keyword whose value is to be rewritten
      }
      \sstsubsection{
         VALUE                = REAL (Given)
      }{
         the value that the FITS keyword is to have
      }
      \sstsubsection{
         STATUS               = INTEGER (Given and returned)
      }{
         global status
      }
   }
   \sstdiytopic{
      Copyright
   }{
      Copyright \copyright 1995,1996,1997,1998,1999 Particle Physics and Astronomy
      Research Council. All Rights Reserved.
   }
}
\sstroutine{
   SCULIB\_SCAN\_2\_RD
}{
   Calculate the apparent RA/Dec of a scan
}{
   \sstdescription{
      The start and end of each scan is calculated in different
      coordinate systems dependent on the coordinate system of the
      tracking centre.
      This routine calculates the apparent RA/DEC centre for
      the supplied long and Lat using knowledge of the transputer system.
   }
   \sstinvocation{
       CALL SCULIB\_SCAN\_2\_RD(VERSION, CENTRE\_COORDS, RA\_CEN, DEC\_CEN,
            LONG, LAT, LST, MJD, LAT\_OBS, RA\_APP, DEC\_APP,
            STATUS )
   }
   \sstarguments{
      \sstsubsection{
         VERSION = REAL (Given)
      }{
         Version number of the file. This governs whether we need to
         even run this subroutine.
      }
      \sstsubsection{
         CENTRE\_COORDS = CHAR (Given)
      }{
         Centre coordinates of tracking centre
      }
      \sstsubsection{
         RA\_CEN = DOUBLE (Given)
      }{
         Apparent RA of map centre
      }
      \sstsubsection{
         DEC\_CEN = DOUBLE (Given)
      }{
         Apparent dec of map centre
      }
      \sstsubsection{
         LONG = DOUBLE (Given)
      }{
         Longitude of array at LST
      }
      \sstsubsection{
         LAT = DOUBLE (Given)
      }{
         Latitude of array at LST
      }
      \sstsubsection{
         LST = DOUBLE (Given)
      }{
         Local sidereal time (radians)
      }
      \sstsubsection{
         MJD = DOUBLE (Given)
      }{
         Modified Julian data of observation (should be the MJD of the
         time for which lst = LST).
      }
      \sstsubsection{
         LAT\_OBS                = DOUBLE PRECISION (Given)
      }{
         Latitude of observatory in radians.
      }
      \sstsubsection{
         RA\_APP                 = DOUBLE PRECISION (Returned)
      }{
         Apparent RA of point at date (radians)
      }
      \sstsubsection{
         DEC\_APP                = DOUBLE PRECISION (Returned)
      }{
         Apparent Dec
      }
      \sstsubsection{
         STATUS                 = INTEGER (Given and returned)
      }{
         Global status
      }
   }
   \sstnotes{
      Before November 1997 the situation is a bit tricky:
      \sstitemlist{

         \sstitem
           [1] Scan ends were assumed to be RD by the transputers

         \sstitem
           [2] The telescope assumed RJ offsets

      }
      Therefore I need to do the following to recreate what actually
      happened:
      \sstitemlist{

         \sstitem
           [1] Calculate tangent plane offsets from the RD centre

         \sstitem
           [2] Add these offsets onto the actual RJ tracking centre

         \sstitem
           [3] Convert back into RD

      }
      Post November 1997 the offsets really are RD even when the telescope
      goes to RJ so this routine should not be called. Version 1.0 data
      will not be modified.
   }
   \sstdiytopic{
      Copyright
   }{
      Copyright \copyright 1995,1996,1997,1998,1999 Particle Physics and Astronomy
      Research Council. All Rights Reserved.
   }
}
\sstroutine{
   SCULIB\_SCAN\_APPARENT\_TP\_2\_AZNA
}{
   Calculate NAsmyth and AZel coordinates for SCAN/MAP
}{
   \sstdescription{
      This routine calculates the NA/AZ offsets from tangent plane
      offsets for SCAN map data.
   }
   \sstinvocation{
       CALL SCULIB\_SCAN\_APPARENT\_TP\_2\_AZNA(OUT\_COORDS, N\_POS,
            N\_BOL, ELEVATION, PAR\_ANGLE,
            BOL\_DEC, BOL\_RA, STATUS)
   }
   \sstarguments{
      \sstsubsection{
         OUT\_COORDS = \_CHAR (Given)
      }{
         Output coordinate system
      }
      \sstsubsection{
         N\_POS = \_INTEGER (Given)
      }{
         Number of `samples{\tt '} taken
      }
      \sstsubsection{
         N\_BOL = \_INTEGER (Given)
      }{
         Number of bolometers in the input data
      }
      \sstsubsection{
         ELEVATION ( N\_POS ) = \_DOUBLE (Given)
      }{
         Elevation of each exposure
      }
      \sstsubsection{
         PAR\_ANGLE ( N\_POS ) = \_DOUBLE (Given)
      }{
         Parallactic angle of each exposure
      }
      \sstsubsection{
         RA\_CEN = \_DOUBLE (Given)
      }{
         apparent RA of output map centre (radians) Only used for JIGGLE data.
      }
      \sstsubsection{
         DEC\_CEN = \_DOUBLE (Given)
      }{
         apparent Dec of output map centre (radians) Only used for JIGGLE data.
      }
      \sstsubsection{
         BOL\_DEC(N\_BOL, N\_POS) = \_DOUBLE (Returned)
      }{
         Apparent DEC of bolometers for each measurement for MJD\_STANDARD
      }
      \sstsubsection{
         BOL\_RA(N\_BOL, N\_POS) = \_DOUBLE (Returned)
      }{
         Apparent RA of bolometers for each measurement for MJD\_STANDARD
      }
      \sstsubsection{
         STATUS = \_INTEGER (Given \& Returned)
      }{
         Global status
      }
   }
   \sstnotes{
      This routine does not annul the locator.
   }
   \sstdiytopic{
      Prior Requirements
   }{
      The locator to the structure must already be available.
   }
   \sstdiytopic{
      Copyright
   }{
      Copyright \copyright 1995,1996,1997,1998,1999 Particle Physics and Astronomy
      Research Council. All Rights Reserved.
   }
}
\sstroutine{
   SCULIB\_SEARCH\_DATADIR
}{
   Open an NDF using the parameter system whilst searching DATADIR
}{
   \sstdescription{
      This routine reads a parameter and attempt to open an NDF.
      If it fails it searches the directory DATADIR and then opens
      an NDF there. If that fails it asks again.
      This routine also recognises the SCUBA\_PREFIX environment variable.
      When a number is given it is expanded to 4 digites (leading zeroes)
      and prepended with SCUBA\_PREFIX and {\tt '}\_dem\_{\tt '} before opening the file.
      The priority is:
      \sstitemlist{

         \sstitem
              File specified in current directory

         \sstitem
              Check if filename was a number and look in current directory
            then either

         \sstitem
              Look for full filename in DATADIR
            or

         \sstitem
              Look for numbered name in DATADIR
         Only one of the last two will happen since I will know from the status
         return of CHR\_CTOI whether a number was given.
      }
   }
   \sstinvocation{
      CALL SCULIB\_SEARCH\_DATADIR(PACKAGE, PARAM, INDF, STATUS)
   }
   \sstarguments{
      \sstsubsection{
         PACKAGE = CHAR (Given)
      }{
         String to identify the software system when the information
         messages appear.
      }
      \sstsubsection{
         PARAM = CHAR (Given)
      }{
         Name of parameter associated with NDF
      }
      \sstsubsection{
         INDF = INTEGER (Returned)
      }{
         NDF identifier of successfully opened file
      }
      \sstsubsection{
         STATUS = INTEGER (Given and Returned)
      }{
         Global Status value
      }
   }
   \sstdiytopic{
      Copyright
   }{
      Copyright \copyright 1995,1996,1997,1998,1999 Particle Physics and Astronomy
      Research Council. All Rights Reserved.
   }
   \sstimplementationstatus{
      \sstitemlist{

         \sstitem
         Uses some non PSX calls for accessing Current working directory
           These are GETCWD and CHDIR. This may be a portability issue.

         \sstitem
         Checks for PAR\_\_ABORT or PAR\_\_NULL status from parameter.
      }
   }
}
\sstroutine{
   SCULIB\_SET\_DATA
}{
   set data to a real value given a byte mask
}{
   \sstdescription{
      This routine uses a byte mask (N\_BOLS $*$ N\_POS) to set a data
      value in the output. A mask value of 1 indicates that a value
      should be changed. This routine does not distinguish {\tt '}beams{\tt '}
   }
   \sstinvocation{
       CALL SCULIB\_SET\_DATA (USE\_THIS, N\_BOLS, N\_POS,
            N\_BEAM, MASK, VALUE, IN\_DATA, STATUS)
   }
   \sstarguments{
      \sstsubsection{
         USE\_THIS   = LOGICAL (Given)
      }{
         Describes whether the mask should be set to the value (TRUE)
         or whether the rest of the data should be set (FALSE)
      }
      \sstsubsection{
         N\_BOLS                  = INTEGER (Given)
      }{
         number of bolometers measured
      }
      \sstsubsection{
         N\_POS                   = INTEGER (Given)
      }{
         number of positions measured
      }
      \sstsubsection{
         N\_BEAM                  = INTEGER (Given)
      }{
         number of beams used
      }
      \sstsubsection{
         MASK( N\_BOLS, N\_POS)    = BYTE (Given)
      }{
         input mask
      }
      \sstsubsection{
         VALUE               = REAL (Given)
      }{
         The value of the masked/unmasked data
      }
      \sstsubsection{
         IN\_DATA (N\_BOLS, N\_POS, N\_BEAM) = REAL (Given and returned)
      }{
         the data to be masked
      }
      \sstsubsection{
         STATUS                  = INTEGER (Given and returned)
      }{
         global status
      }
   }
   \sstdiytopic{
      Copyright
   }{
      Copyright \copyright 1995,1996,1997,1998,1999 Particle Physics and Astronomy
      Research Council. All Rights Reserved.
   }
}
\sstroutine{
   SCULIB\_SET\_DATA\_BIT
}{
   set a bit in data given a byte mask
}{
   \sstdescription{
      This routine uses a byte mask (N\_BOLS $*$ N\_POS) to set a data
      bit in the output. A mask value of 1 indicates that a value
      should be changed. This routine does not distinguish {\tt '}beams{\tt '}
      A bit can be set or unset.
   }
   \sstinvocation{
       CALL SCULIB\_SET\_DATA\_BIT(USE\_THIS, N\_BOLS, N\_POS, N\_BEAM,
           MASK, BITNUM, STATE, IN\_DATA, STATUS)
   }
   \sstarguments{
      \sstsubsection{
         USE\_THIS   = LOGICAL (Given)
      }{
         Describes whether the mask should be set to the value (TRUE)
         or whether the rest of the data should be set (FALSE)
      }
      \sstsubsection{
         N\_BOLS                  = INTEGER (Given)
      }{
         number of bolometers measured
      }
      \sstsubsection{
         N\_POS                   = INTEGER (Given)
      }{
         number of positions measured
      }
      \sstsubsection{
         N\_BEAM                  = INTEGER (Given)
      }{
         number of beams used
      }
      \sstsubsection{
         MASK( N\_BOLS, N\_POS)    = BYTE (Given)
      }{
         input mask
      }
      \sstsubsection{
         BITNUM                  = INTEGER (Given)
      }{
         Bit number to affect
      }
      \sstsubsection{
         STATE                   = LOGICAL (Given)
      }{
         Set the bits if true. Otherwise clear them
      }
      \sstsubsection{
         IN\_DATA (N\_BOLS, N\_POS, N\_BEAM) = BYTE (Given and returned)
      }{
         the data to be masked
      }
      \sstsubsection{
         STATUS                  = INTEGER (Given and returned)
      }{
         global status
      }
   }
   \sstdiytopic{
      Copyright
   }{
      Copyright \copyright 1995,1996,1997,1998,1999 Particle Physics and Astronomy
      Research Council. All Rights Reserved.
   }
}
\sstroutine{
   SCULIB\_SET\_DATA\_UB
}{
   set data to a byte value given a byte mask
}{
   \sstdescription{
      This routine uses a byte mask (N\_BOLS $*$ N\_POS) to set a data
      value in the output. A mask value of 1 indicates that a value
      should be changed. This routine does not distinguish {\tt '}beams{\tt '}
   }
   \sstinvocation{
       CALL SCULIB\_SET\_DATA\_UB(USE\_THIS, N\_BOLS, N\_POS, N\_BEAM,
         MASK, BVALUE, IN\_DATA, STATUS)
   }
   \sstarguments{
      \sstsubsection{
         USE\_THIS   = LOGICAL (Given)
      }{
         Describes whether the mask should be set to the value (TRUE)
         or whether the rest of the data should be set (FALSE)
      }
      \sstsubsection{
         N\_BOLS                  = INTEGER (Given)
      }{
         number of bolometers measured
      }
      \sstsubsection{
         N\_POS                   = INTEGER (Given)
      }{
         number of positions measured
      }
      \sstsubsection{
         N\_BEAM                  = INTEGER (Given)
      }{
         number of beams used
      }
      \sstsubsection{
         MASK( N\_BOLS, N\_POS)    = BYTE (Given)
      }{
         input mask
      }
      \sstsubsection{
         BVALUE                  = BYTE (Given)
      }{
         The value to be given to the masked/unmasked data
      }
      \sstsubsection{
         IN\_DATA (N\_BOLS, N\_POS, N\_BEAM) = BYTE (Given and returned)
      }{
         the data to be masked
      }
      \sstsubsection{
         STATUS                  = INTEGER (Given and returned)
      }{
         global status
      }
   }
   \sstdiytopic{
      Copyright
   }{
      Copyright \copyright 1995,1996,1997,1998,1999 Particle Physics and Astronomy
      Research Council. All Rights Reserved.
   }
}
\sstroutine{
   SCULIB\_SET\_QUAL
}{
   set quality bits in a subset of a quality array
}{
   \sstdescription{
      Set the bits in the quality array as specified by a mask.
      Can be used to set or unset bits as well as using the mask
      or the inverse of the mask. This is used to mask scuba sections.
   }
   \sstinvocation{
       CALL SCULIB\_SET\_QUAL (USE\_SECT, QUALITY, N\_BOLS, N\_POS, N\_BEAM,
         BOL\_S, POS\_S, BIT\_POS, BIT\_SWITCH, STATUS)
   }
   \sstarguments{
      \sstsubsection{
         USE\_SECT                = LOGICAL (Given)
      }{
         am I changing SECTION or .NOT. SECTION
      }
      \sstsubsection{
         QUALITY (N\_BOLS, N\_POS, N\_BEAM) = BYTE (Given and returned)
      }{
         the quality array
      }
      \sstsubsection{
         N\_BOLS                  = INTEGER (Given)
      }{
         number of bolometers measured
      }
      \sstsubsection{
         N\_POS                   = INTEGER (Given)
      }{
         number of positions measured
      }
      \sstsubsection{
         N\_BEAM                  = INTEGER (Given)
      }{
         number of beams used
      }
      \sstsubsection{
         BOL\_S (N\_BOLS)          = INTEGER (Given)
      }{
         array containing 1s for bolometers whose quality is to be
         changed
      }
      \sstsubsection{
         POS\_S (N\_POS)           = INTEGER (Given)
      }{
         array containing 1s for positions whose quality is to be
         changed
      }
      \sstsubsection{
         BIT\_POS                 = INTEGER (Given)
      }{
         position of bit ot be set (0 - 7)
      }
      \sstsubsection{
         BIT\_SWITCH               = LOGICAL (Given)
      }{
         If .TRUE. we set the bit using SCULIB\_BITON. If .FALSE.
         we unset the bit with SCULIB\_BITOFF
      }
      \sstsubsection{
         STATUS                  = INTEGER (Given and returned)
      }{
         global status
      }
   }
   \sstdiytopic{
      Copyright
   }{
      Copyright \copyright 1995,1996,1997,1998,1999 Particle Physics and Astronomy
      Research Council. All Rights Reserved.
   }
}
\sstroutine{
   SCULIB\_SET\_QUALITY
}{
   set quality bits in a subset of a quality array
}{
   \sstdescription{
      set quality bits in a subset of a quality array
   }
   \sstinvocation{
        CALL  SCULIB\_SET\_QUALITY (N\_BOLS, N\_POS, N\_BEAM, QUALITY,
         START\_BOL, END\_BOL, START\_POS, END\_POS, START\_BEAM, END\_BEAM,
         BIT\_POS, BIT\_VALUE, STATUS)
   }
   \sstarguments{
      \sstsubsection{
         N\_BOLS                  = INTEGER (Given)
      }{
         number of bolometers measured
      }
      \sstsubsection{
         N\_POS                   = INTEGER (Given)
      }{
         number of positions measured
      }
      \sstsubsection{
         N\_BEAM                  = INTEGER (Given)
      }{
         number of beams used
      }
      \sstsubsection{
         QUALITY (N\_BOLS, N\_POS, N\_BEAM) = BYTE (Given and returned)
      }{
         the quality array
      }
      \sstsubsection{
         START\_BOL               = INTEGER (Given)
      }{
         index of first bolometer to be set bad
      }
      \sstsubsection{
         END\_BOL                 = INTEGER (Given)
      }{
         index of last bolometer to be set bad
      }
      \sstsubsection{
         START\_POS               = INTEGER (Given)
      }{
         index of first position to be set bad
      }
      \sstsubsection{
         END\_POS                 = INTEGER (Given)
      }{
         index of last position to be set bad
      }
      \sstsubsection{
         START\_BEAM              = INTEGER (Given)
      }{
         index of first beam  to be set bad
      }
      \sstsubsection{
         END\_BEAM                = INTEGER (Given)
      }{
         index of last beam  to be set bad
      }
      \sstsubsection{
         BIT\_POS                 = INTEGER (Given)
      }{
         position of bit ot be set (0 - 7)
      }
      \sstsubsection{
         BIT\_VALUE               = INTEGER (Given)
      }{
         value to which bit is to be set (0 if zero, 1 otherwise)
      }
      \sstsubsection{
         STATUS                  = INTEGER (Given and returned)
      }{
         global status
      }
   }
   \sstdiytopic{
      Copyright
   }{
      Copyright \copyright 1995,1996,1997,1998,1999 Particle Physics and Astronomy
      Research Council. All Rights Reserved.
   }
}
\sstroutine{
   SCULIB\_SET\_USER
}{
   set the USER array used by SCULIB\_SKYFUNC\_1
}{
   \sstdescription{
      E04UPF is a NAG routine that is used to fit a theoretical
      sky-dip curve to the measured data by varying ETA\_TEL, B and TAU.
      SCULIB\_SKYFUNC\_1 is called by E04UPF to calculate the M sub-functions
      fi(x) at the current solution vector X, and/or their Jacobian
      matrix (see NAG manual for further info). That routine obtains the
      data needed for its calculations via the USER array passed into it,
      and this routine fills USER with the necessary numbers:-
      \sstitemlist{

         \sstitem
          USER (1) = J\_TEL

         \sstitem
          USER (2) = J\_ATM

         \sstitem
          USER (3) = not used

         \sstitem
          USER (4:M$+$3)      = the measured airmasses

         \sstitem
          USER (M$+$4:2M$+$3)   = the measured sky temperatures

         \sstitem
          USER (2M$+$4:3M$+$3)  = the errors on the measured sky temperatures
      }
   }
   \sstinvocation{
       CALL SCULIB\_SET\_USER (J\_TEL, J\_ATM, N\_MEASUREMENTS,
         AIRMASS, J\_MEAS\_D, J\_MEAS\_V, USER)
   }
   \sstarguments{
      \sstsubsection{
         J\_TEL                     = REAL (Given)
      }{
         the temperature of the telescope
      }
      \sstsubsection{
         J\_ATM                     = REAL (Given)
      }{
         the temperature of the atmosphere
      }
      \sstsubsection{
         N\_MEASUREMENTS            = INTEGER (Given)
      }{
         the number of measurements taken in the sky-dip
      }
      \sstsubsection{
         AIRMASS (N\_MEASUREMENTS)  = REAL (Given)
      }{
         the airmass at each measurement
      }
      \sstsubsection{
         J\_MEAS\_D (N\_MEASUREMENTS) = REAL (Given)
      }{
         the brightness temperature of the atmosphere at each airmass
      }
      \sstsubsection{
         J\_MEAS\_V (N\_MEASUREMENTS) = REAL (Given)
      }{
         the variance on J\_MEAS\_D
      }
      \sstsubsection{
         USER (3 $*$ N\_MEASUREMENTS $+$ 3) = DOUBLE PRECISION (Returned)
      }{
         the USER array required
      }
   }
   \sstdiytopic{
      Copyright
   }{
      Copyright \copyright 1995,1996,1997,1998,1999 Particle Physics and Astronomy
      Research Council. All Rights Reserved.
   }
}
\sstroutine{
   SCULIB\_SINC
}{
   Calculate SIN(PI$*$X)/(PI$*$X)
}{
   \sstdescription{
      Calculate sinc function.
   }
   \sstinvocation{
      RESULT = SCULIB\_SINC( X )
   }
   \sstarguments{
      \sstsubsection{
         X = REAL (Given)
      }{
         The argument to the sinc function - In turns
      }
   }
   \sstreturnedvalue{
      \sstsubsection{
         SCULIB\_SINC = REAL
      }{
         SINC(PI$*$X)
      }
   }
   \sstdiytopic{
      Copyright
   }{
      Copyright \copyright 1995,1996,1997,1998,1999 Particle Physics and Astronomy
      Research Council. All Rights Reserved.
   }
   \sstimplementationstatus{
      \sstitemlist{

         \sstitem
         No status checking

         \sstitem
         No bad value checking
      }
   }
}
\sstroutine{
   SCULIB\_SKYCON\_1
}{
   routine to calculate non-linear constraints for
   NAG non-linear fitting routine E04UPF when fitting
   ETA\_TEL, B and TAU in sky-dip analysis
}{
   \sstdescription{
      E04UPF is a NAG routine used to fit a theoretical SKYDIP
      curve to the measured data by varying ETA\_TEL, B and TAU. This
      routine is called by E04UPF to calculate the 1 non-linear constraint
      in the problem and/or it{\tt '}s Jacobian at the current solution vector
      X. See the NAG manual if you need to know more about E04UPF.

      The solution vector is composed as follows:-
      \sstitemlist{

         \sstitem
              x(1) = ETA\_TEL

         \sstitem
              x(2) = B

         \sstitem
              x(3) = TAU

      }
      The non-linear constraint limits ETA\_TEL $*$ B. So, evaluation
      of the constraint gives:-

           C(1) = ETA\_TEL $*$ B

      and its Jacobian has components :-

\begin{equation}
\frac{\partial C}{\partial x(1)} = B
\end{equation}
\begin{equation}
\frac{\partial C}{\partial x(2)} = \eta_\mathrm{tel}
\end{equation}
\begin{equation}
\frac{\partial C}{\partial x(2)} = 0.0
\end{equation}

   }
   \sstinvocation{
       CALL SCULIB\_SKYCON\_1 (MODE, NCNLN, N, LDCJ, NEEDC, X, C,
         CJAC, NSTATE, IUSER, USER)
   }
   \sstarguments{
      \sstsubsection{
         See NAG manual description of CONFUN parameter in E04UPF.
      }{
      }
   }
   \sstdiytopic{
      Copyright
   }{
      Copyright \copyright 1993-1999 Particle Physics and Astronomy
      Research Council. All Rights Reserved.
   }
}
\sstroutine{
   SCULIB\_SKYDIP\_ALLAN\_VARIANCE
}{
   incorporate latest set of SKYDIP data
   samples into Allan variance
}{
   \sstdescription{
      This routine incorporates a data slice into a run of data and updates
      the Allan variance of the run.

         The Allan variance is calculated for a range of simulated integration
      times for which artificial samples are calculated from the input data.
      The Allan variance for a particular integration time is calculated from:-


\begin{equation}
variance = \sum_{i=1\rightarrow n-1}\frac{(sample(i) - sample(i+1))^2}{2n}
\end{equation}
      where n is the number of artificial samples available.
   }
   \sstinvocation{
       CALL SCULIB\_SKYDIP\_ALLAN\_VARIANCE (SUB\_INSTRUMENT, NUM\_CHAN,
         NUM\_ADC, BOL\_TYPE, N\_BOLS, BOL\_CHAN, BOL\_ADC, N\_SAMPLES,
         SAMPLE, SAMPLE\_QUALITY, KMAX, SUM, N\_SUM, ARTIFICIAL,
         N\_ARTIFICIAL, ALLAN\_VARIANCE, ALLAN\_QUALITY, STATUS)
   }
   \sstarguments{
      \sstsubsection{
         SUB\_INSTRUMENT                  = CHARACTER$*$($*$) (Given)
      }{
         the name of the sub-instrument for which the Allan variance
         is to be calculated
      }
      \sstsubsection{
         NUM\_CHAN                        = INTEGER (Given)
      }{
         number of channels per A/D card
      }
      \sstsubsection{
         NUM\_ADC                         = INTEGER (Given)
      }{
         number of A/D cards
      }
      \sstsubsection{
         BOL\_TYPE (NUM\_CHAN, NUM\_ADC)    = CHARACTER$*$($*$) (Given)
      }{
         the types of the bolometers
      }
      \sstsubsection{
         N\_BOLS                          = INTEGER (Given)
      }{
         the number of bolometers making measurements
      }
      \sstsubsection{
         BOL\_CHAN (N\_BOLS)               = INTEGER (Given)
      }{
         channel numbers of selected bolometers
      }
      \sstsubsection{
         BOL\_ADC (N\_BOLS)                = INTEGER (Given)
      }{
         ADC numbers of selected bolometers
      }
      \sstsubsection{
         N\_SAMPLES                       = INTEGER (Given)
      }{
         the number of samples
      }
      \sstsubsection{
         SAMPLE (N\_BOLS, N\_SAMPLES)      = REAL (Given)
      }{
         the samples
      }
      \sstsubsection{
         SAMPLE\_QUALITY (N\_BOLS, N\_SAMPLES) = INTEGER (Given)
      }{
         quality on the samples
      }
      \sstsubsection{
         KMAX                            = INTEGER (Given)
      }{
         the size of the Allan variance array
      }
      \sstsubsection{
         SUM (KMAX)                      = REAL (Given and returned)
      }{
         the accumulator for current artificial sample of length K real
         samples
      }
      \sstsubsection{
         N\_SUM (KMAX)                    = INTEGER (Given and returned)
      }{
         the number of samples currently in the artificial sample
         accumulator
      }
      \sstsubsection{
         ARTIFICIAL (KMAX)               = REAL (Given and returned)
      }{
         the last artificial sample of length K samples that was
         calculated
      }
      \sstsubsection{
         N\_ARTIFICIAL (KMAX)             = INTEGER (Given and returned)
      }{
         the number of artificial samples of length K samples that
         have been calculated
      }
      \sstsubsection{
         ALLAN\_VARIANCE (KMAX)           = REAL (Given and returned)
      }{
         the Allan variance
      }
      \sstsubsection{
         ALLAN\_QUALITY (KMAX)            = INTEGER (Returned)
      }{
         quality on the Allan variance
      }
      \sstsubsection{
         STATUS                          = INTEGER (Given and returned)
      }{
         global status
      }
   }
   \sstdiytopic{
      Copyright
   }{
      Copyright \copyright 1995,1996,1997,1998,1999 Particle Physics and Astronomy
      Research Council. All Rights Reserved.
   }
   \sstdiytopic{
      Method
   }{
      If status on entry is good the routine will:-

      loop through the data samples to be incorporated -

         loop through the bolometers for which data is available and calculate
         an average value for the sample from those bolometers belonging to
         the specified sub-instrument

         if the quality of the average sample is good -

            loop through the Allan variances -

               add the sample to the buffer used to calculate the latest
               artificial sample at the simulated integration time of this
               variance

               if all the data has been obtained for the latest artificial
               sample -

                  calculate the artificial sample

                  if there is only one artificial sample for this simulated
                  integration time we can{\tt '}t calculate the Allan variance, so
                  set its quality to bad

                  if there are 2 artifical samples then the Allan variance
                  can be calculated from -
\begin{equation}
variance = \frac{(artificial\_2 - artificial\_1)^2}{2 \times 2}
\end{equation}

                  if there are more than 2 artificial samples

      \sstitemlist{

         \sstitem
                        the sum of the squares of the differences between the
                        previous samples is recovered from the current value of the
                        Allan variance

         \sstitem
                        the square of the difference between the current artificial
                        sample and the previous one is added to the sum

         \sstitem
                        the Allan variance is re-calculated

      }
                  end if

                  the current artifical sample is stored to be used as the
                  `previous{\tt '} sample next time round

                  the variables used to calculate the artificial samples at
                  this simulated integration time are reset

               end if

            end of loop through Allan variances

         end if

      end of loop through samples in this dataslice
   }
}
\sstroutine{
   SCULIB\_SKYDIP\_BOLS
}{
   returns bolometers to be measured in a SKYDIP
}{
   \sstdescription{
      This routine returns a string containing the names of the bolometers
      to be measured in a SKYDIP observation.

         Only those sub-instruments that are looking out through a suitable
      filter will be measured. Thus, if the filter in front of the SHORT
      array is {\tt '}350{\tt '} or {\tt '}450{\tt '} the returned string will contain the word
      SHORT\_DC. Other filter/sub/bolometer combinations are as follows:-
      \sstitemlist{

         \sstitem
           600 or 750 or 850 in front of the LONG array gives LONG\_DC

         \sstitem
           1100 in front of the P1100 photometer gives        P1100\_DC

         \sstitem
           1300 in front of the P1300 photometer gives        P1300\_DC

         \sstitem
           2000 in front of the P2000 photometer gives        P2000\_DC

      }
      The returned BOLOMETERS string will contain the words for each
      sub-instrument to be measured separated by commas.
   }
   \sstinvocation{
      CALL SCULIB\_SKYDIP\_BOLS (FILTER, BOLOMETERS, STATUS)
   }
   \sstarguments{
      \sstsubsection{
         FILTER                    = CHARACTER$*$($*$) (Given)
      }{
         the name of the filter combination in use
      }
      \sstsubsection{
         BOLOMETERS                = CHARACTER$*$($*$) (Returned)
      }{
         the names of the bolometers to be measured
      }
      \sstsubsection{
         STATUS                    = INTEGER (Given and returned)
      }{
         global status
      }
   }
   \sstdiytopic{
      Copyright
   }{
      Copyright \copyright 1995,1996,1997,1998,1999 Particle Physics and Astronomy
      Research Council. All Rights Reserved.
   }
}
\sstroutine{
   SCULIB\_SKYDIP\_TEMPERATURES
}{
   derive sky temperatures from chopper-wheel data
}{
   \sstdescription{
      This routine derives sky brightness temperatures from data containing
      the measured signals from cold and ambient loads and sky. The sky
      temperature is calculated for each sample by linear interpolation between
      the signals from the cold and ambient loads.

         In addition, mean values of the sky temperature for each bolometer
      are derived from the samples taken. If more than one sample contributes
      to the mean then the variance is calculated from the spread of samples
      about the mean.
   }
   \sstinvocation{
       CALL SCULIB\_SKYDIP\_TEMPERATURES (T\_COLD, T\_HOT, N\_SUB,
      :  SUB\_INSTRUMENT, WAVELENGTH, NUM\_CHAN, NUM\_ADC, BOL\_TYPE,
      :  N\_BOLS, BOL\_CHAN, BOL\_ADC, N\_SAMPLES\_IN, RAW\_DATA,
      :  N\_SAMPLES\_OUT, J\_SKY, J\_SKY\_VARIANCE, J\_SKY\_QUALITY, J\_SKY\_AV,
      :  J\_SKY\_AV\_VARIANCE, J\_SKY\_AV\_QUALITY, STATUS)
   }
   \sstarguments{
      \sstsubsection{
         T\_COLD (N\_SUB)             = REAL (Given)
      }{
         the temperature of the cold load for each sub-instrument
      }
      \sstsubsection{
         T\_HOT (N\_SUB)              = REAL (Given)
      }{
         the temperature of the hot load for each sub-instrument
      }
      \sstsubsection{
         N\_SUB                      = INTEGER (Given)
      }{
         the number of sub-instruments being used
      }
      \sstsubsection{
         SUB\_INSTRUMENT (N\_SUB)     = CHARACTER$*$($*$) (Given)
      }{
         the names of the sub-instruments being used
      }
      \sstsubsection{
         WAVELENGTH (N\_SUB)         = REAL (Given)
      }{
         the wavelengths of the filters in front of the sub-instruments
         (microns)
      }
      \sstsubsection{
         NUM\_CHAN                   = INTEGER (Given)
      }{
         number of A/D channels per A/D card
      }
      \sstsubsection{
         NUM\_ADC                    = INTEGER (Given)
      }{
         number of A/D cards
      }
      \sstsubsection{
         BOL\_TYPE (NUM\_CHAN, NUM\_ADC) = CHARACTER$*$($*$) (Given)
      }{
         the types of the bolometers
      }
      \sstsubsection{
         N\_BOLS                     = INTEGER (Given)
      }{
         the number of bolometers making measurements
      }
      \sstsubsection{
         BOL\_CHAN (N\_BOLS)          = INTEGER (Given)
      }{
         channel numbers of selected bolometers
      }
      \sstsubsection{
         BOL\_ADC (N\_BOLS)           = INTEGER (Given)
      }{
         ADC numbers of selected bolometers
      }
      \sstsubsection{
         N\_SAMPLES\_IN               = INTEGER (Given)
      }{
         the number of samples taken by each bolometer
      }
      \sstsubsection{
         RAW\_DATA (3, N\_BOLS, N\_SAMPLES\_IN) = REAL (Given)
      }{
         the measured signals from the cold and ambient loads and sky
         \sstitemlist{

            \sstitem
              RAW\_DATA (1,i,j) = ambient load signal

            \sstitem
              RAW\_DATA (2,i,j) = sky signal

            \sstitem
              RAW\_DATA (3,i,j) = cold load signal
         }
      }
      \sstsubsection{
         N\_SAMPLES\_OUT              = INTEGER (Given)
      }{
         the number of samples requested for each bolometer
      }
      \sstsubsection{
         J\_SKY (N\_BOLS, N\_SAMPLES\_OUT) = REAL (Returned)
      }{
         the brightness temperature of the sky
      }
      \sstsubsection{
         J\_SKY\_VARIANCE (N\_BOLS, N\_SAMPLES\_OUT) = REAL (Returned)
      }{
         the variance on J\_SKY
      }
      \sstsubsection{
         J\_SKY\_QUALITY (N\_BOLS, N\_SAMPLES\_OUT) = BYTE (Returned)
      }{
         the quality on J\_SKY
      }
      \sstsubsection{
         J\_SKY\_AV (N\_BOLS)          = REAL (Returned)
      }{
         the average of J\_SKY for each bolometer
      }
      \sstsubsection{
         J\_SKY\_AV\_VARIANCE (N\_BOLS) = REAL (Returned)
      }{
         the variance on the average of J\_SKY
      }
      \sstsubsection{
         J\_SKY\_AV\_QUALITY (N\_BOLS)  = BYTE (Returned)
      }{
         the quality of the average of J\_SKY
      }
      \sstsubsection{
         STATUS                     = INTEGER (Given and returned)
      }{
         Global status
      }
   }
   \sstdiytopic{
      Copyright
   }{
      Copyright (C) 1995,1996,1997,1998,1999 Particle Physics and Astronomy
      Research Council. All Rights Reserved.
   }
   \sstimplementationstatus{
      \sstitemlist{

         \sstitem
         Should be converted to use double precision internally since sometimes
           can return a negative temperature.
      }
   }
}

\sstroutine{
   SCULIB\_SKYDIP\_VAR
}{
   calculate variance of skydip data points
}{
   \sstdescription{
      If entered with good status this routine calculates the variance
      between a SKYDIP dataset and the theoretical curve generated by
      the the sky and telescope parameters in FIT.  Additionally, the residual
      of the fit is also returned (absolute value of measured data mins
      fitted values). Data points with bad
      quality will be ignored. If no valid data points are found then an
      error will be reported and bad status returned. The data are passed
      in via common.

      The FIT vector is composed as follows:-
      \sstitemlist{

         \sstitem
               x(1) = ETA\_TEL

         \sstitem
               x(2) = B

         \sstitem
               x(3) = TAU

      }
      For a given airmass:-

\begin{equation}
J_\mathrm{theory} = (1-\eta_\mathrm{tel}) J_\mathrm{tel} +
\eta_\mathrm{tel}J_\mathrm{atm} (1-B \exp(-\tau Airmass))
\end{equation}

\begin{equation}
            J_\mathrm{atm} = J_\mathrm{amb} X_G
\end{equation}

\begin{equation}
              X_G = 1 + \frac{h1\times h2}{J_\mathrm{amb}} \exp (-\tau \times
Airmass(i) / X\_Gconst)
\end{equation}

   }
   \sstinvocation{
      CALL SCULIB\_SKYDIP\_VAR (RESIDUAL, VAR, N, FIT, NDEG, STATUS)
   }
   \sstarguments{
      \sstsubsection{
        RESIDUAL = DOUBLE PRECISION (Returned)
      }{
        residual of the fit (absolute difference between the
            model and the data)
      }
      \sstsubsection{
         VAR                     = DOUBLE PRECISION (Returned)
      }{
         the chi-squared value of the current fit
      }
      \sstsubsection{
         N                        = INTEGER (Given)
      }{
         the number of parameters being fit, should be 3
      }
      \sstsubsection{
         FIT (N)                  = DOUBLE PRECISION (Given)
      }{
         the fit parameters:-
         \sstitemlist{

            \sstitem
                                FIT (1) = ETA\_TEL

            \sstitem
                                FIT (2) = B

            \sstitem
                                FIT (3) = TAUZ
         }
      }
      \sstsubsection{
         NDEG                     = INTEGER (Given)
      }{
         the number of degrees of freedom
      }
      \sstsubsection{
         STATUS                   = INTEGER (Given and returned)
      }{
         global status
      }
   }
   \sstdiytopic{
      Copyright
   }{
      Copyright \copyright 1995-2000 Particle Physics and Astronomy
      Research Council. All Rights Reserved.
   }
}
\sstroutine{
   SCULIB\_SKYDIP\_XISQ
}{
   calculate chi-square of sky fit
}{
   \sstdescription{
      If entered with good status this routine calculates the chi-squared
      between a SKYDIP dataset and the theoretical curve generated by
      the the sky and telescope parameters in FIT. Data points with bad
      quality will be ignored. If no valid data points are found then an
      error will be reported and bad status returned. The data are passed
      in via common.

      The FIT vector is composed as follows:-
      \sstitemlist{

         \sstitem
               x(1) = ETA\_TEL

         \sstitem
               x(2) = B

         \sstitem
               x(3) = TAU

      }
      For a given airmass:-
\begin{equation}
J_\mathrm{theory} = (1-\eta_\mathrm{tel}) J_\mathrm{tel} +
\eta_\mathrm{tel}J_\mathrm{atm} (1-B \exp(-\tau Airmass))
\end{equation}

\begin{equation}
            J_\mathrm{atm} = J_\mathrm{amb} X_G
\end{equation}

\begin{equation}
              X_G = 1 + \frac{h1\times h2}{J_\mathrm{amb}} \exp (-\tau \times
Airmass(i) / X\_Gconst)
\end{equation}

   }
   \sstinvocation{
      CALL SCULIB\_SKYDIP\_XISQ (XISQ, N, FIT, STATUS)
   }
   \sstarguments{
      \sstsubsection{
         XISQ                     = DOUBLE PRECISION (Returned)
      }{
         the chi-squared value of the current fit
      }
      \sstsubsection{
         N                        = INTEGER (Given)
      }{
         the number of parameters being fit, should be 3
      }
      \sstsubsection{
         FIT (N)                  = DOUBLE PRECISION (Given)
      }{
         the fit parameters:-
         \sstitemlist{

            \sstitem
                               FIT (1) = ETA\_TEL

            \sstitem
                               FIT (2) = B

            \sstitem
                               FIT (3) = TAUZ
         }
      }
      \sstsubsection{
         STATUS                   = INTEGER (Given and returned)
      }{
         global status
      }
   }
   \sstdiytopic{
      Copyright
   }{
      Copyright \copyright 1995,1996,1997,1998,1999 Particle Physics and Astronomy
      Research Council. All Rights Reserved.
   }
}
\sstroutine{
   SCULIB\_SKYFUNC
}{
   Return the value of the skydip function
}{
   \sstdescription{
      This routine returns the value of the skydip function.
      The parameters are passed as follows:
      \sstitemlist{

         \sstitem
             P(1) = ETA\_TEL

         \sstitem
             P(2) = B

         \sstitem
             P(3) = TAU

         \sstitem
             P(4) = J\_AMB

         \sstitem
             P(5) = J\_TEL

      }
      This subroutine returns the theoretical value of the skydip
      for the given input parameters.

\begin{equation}
J_\mathrm{theory} = (1-\eta_\mathrm{tel}) J_\mathrm{tel} +
\eta_\mathrm{tel}J_\mathrm{atm} (1-B \exp(-\tau Airmass))
\end{equation}

\begin{equation}
            J_\mathrm{atm} = J_\mathrm{amb} X_G
\end{equation}

\begin{equation}
              X_G = 1 + \frac{h1\times h2}{J_\mathrm{amb}} \exp (-\tau \times
Airmass(i) / X\_Gconst)
\end{equation}

   }
   \sstinvocation{
      CALL SCULIB\_SKYFUNC(F, X, P, M)
   }
   \sstarguments{
      \sstsubsection{
         F = REAL (Returned)
      }{
         The value of the skydip function
      }
      \sstsubsection{
         X = REAL (Given)
      }{
         The airmass used to calculate the function
      }
      \sstsubsection{
         P = REAL (Given)
      }{
         The parameter values
      }
      \sstsubsection{
         M = INTEGER (Given)
      }{
         The number of parameters
      }
   }
   \sstdiytopic{
      Copyright
   }{
      Copyright \copyright 1995,1996,1997,1998,1999 Particle Physics and Astronomy
      Research Council. All Rights Reserved.
   }
}
\sstroutine{
   SCULIB\_SKYFUNC\_1
}{
   calculate f(x) and its Jacobian for NAG non-linear
   fitting routine E04UPF when fitting ETA\_TEL, B
   and TAU in SKYDIP analysis
}{
   \sstdescription{
      E04UPF is a NAG routine that is used to fit a theoretical
      sky-dip curve to the measured data by varying ETA\_TEL, B and TAU.
      This routine is called by E04UPF to calculate the M sub-functions
      fi(x) at the current solution vector X, and/or their Jacobian
      matrix. See the NAG manual if you need to know more about E04UPF.

      The solution vector is composed as follows:-
      \sstitemlist{

         \sstitem
               x(1) = ETA\_TEL

         \sstitem
               x(2) = B

         \sstitem
               x(3) = TAU

      }
      For each of the M measured airmasses the subfunction f(x) is:-

\begin{equation}
             f(x) = \frac{J_\mathrm{theory} (Airmass(i)) - J_\mathrm{meas} (Airmass(i))}{J_\mathrm{noise} (Airmass(i))}
\end{equation}

      where:-

\begin{equation}
J_\mathrm{theory} = (1-\eta_\mathrm{tel}) J_\mathrm{tel} +
\eta_\mathrm{tel}J_\mathrm{atm} (1-B \exp(-\tau Airmass))
\end{equation}

\begin{equation}
            J_\mathrm{atm} = J_\mathrm{amb} X_G
\end{equation}

\begin{equation}
              X_G = 1 + \frac{h1\times h2}{J_\mathrm{amb}} \exp (-\tau \times
Airmass(i) / X\_Gconst)
\end{equation}


      The corresponding Jacobian has component (i,j):-


\begin{equation}
\frac{\partial f(x)}{\partial x(j)} = \frac{\partial}{\partial x(j)}J_\mathrm{theory}(Airmass(i)) \times\frac{1}{J_\mathrm{noise}(i)}
\end{equation}

leading to:

\begin{equation}
\frac{\partial f(i)}{\partial x(1)} = -\frac{J_\mathrm{tel}+J_\mathrm{atm} - B
J_\mathrm{tel} \exp(-\tau A(i))}{J_\mathrm{noise}(i)}
\end{equation}

\begin{equation}
\frac{\partial f(i)}{\partial x(2)} = - \frac{\eta_\mathrm{tel} J_\mathrm{atm}
\exp(-\tau A(i))}{J_\mathrm{noise}(i)}
\end{equation}

\begin{equation}
\frac{\partial f(i)}{\partial x(3)} = \frac{\eta_\mathrm{tel} \frac{\partial
J_\mathrm{atm}}{\partial\tau} - \eta_\mathrm{tel} B \exp(-\tau A(i)) \left(\frac{\partial
J_\mathrm{atm}}{\partial\tau} - J_\mathrm{atm} A(i)\right)}{J_\mathrm{noise}(i)}
\end{equation}

where

\begin{equation}
\frac{\partial J_\mathrm{atm}}{\partial\tau} = J_\mathrm{amb} \frac{\partial
X_G}{\partial\tau}
\end{equation}

and

\begin{equation}
\frac{\partial X_G}{\partial\tau} = - \frac{H_1 H_2 A(i)}{J_\mathrm{amb}
X_\mathrm{Gconst}} \exp(-\tau A(i) / X_\mathrm{Gconst})
\end{equation}

      The array USER is used to get the necessary ancillary information into
      the routine:-
      \sstitemlist{

         \sstitem
          USER (1) = J\_TEL

         \sstitem
          USER (2) = J\_AMB

         \sstitem
          USER (3) = not used

         \sstitem
          USER (4:M$+$3)      = the measured airmasses

         \sstitem
          USER (M$+$4:2M$+$3)   = the measured sky temperatures

         \sstitem
          USER (2M$+$4:3M$+$3)  = the errors on the measured sky temperatures
      }
   }
   \sstinvocation{
       CALL SCULIB\_SKYFUNC\_1 (MODE, M, N, LDFJ, X, F, FJAC, NSTATE,
           IUSER, USER)
   }
   \sstarguments{
      \sstsubsection{
         See NAG manual description of OBJFUN parameter in E04UPF.
      }{
      }
   }
   \sstdiytopic{
      Copyright
   }{
      Copyright \copyright 1995,1996,1997,1998,1999 Particle Physics and Astronomy
      Research Council. All Rights Reserved.
   }
}
\sstroutine{
   SCULIB\_SKYFUNC\_2
}{
   calculate f(x) and its Jacobian for NAG non-linear
   fitting routine E04UPF when fitting B and TAU in
   sky-dip analysis
}{
   \sstdescription{
      E04UPF is a NAG routine that is used to fit a theoretical
      sky-dip curve to the measured data by varying B and TAU.
      This routine is called by E04UPF to calculate the M sub-functions
      f(x) at the current solution vector X, and/or the Jacobian
      matrix. Please see the NAG manual if you need to know more about
      E04UPF.

      The solution vector is composed as follows:-
      \sstitemlist{

         \sstitem
              x(1) = B

         \sstitem
              x(2) = TAU

      }
      For each of the M measured airmasses the subfunction f(x) is:-


\begin{equation}
             f(x) = \frac{J_\mathrm{theory} (Airmass(i)) - J_\mathrm{meas} (Airmass(i))}{J_\mathrm{noise} (Airmass(i))}
\end{equation}

      where:-

\begin{equation}
J_\mathrm{theory} = (1-\eta_\mathrm{tel}) J_\mathrm{tel} +
\eta_\mathrm{tel}J_\mathrm{atm} (1-B \exp(-\tau Airmass))
\end{equation}

      The corresponding Jacobian has component (i,j):-

\begin{equation}
\frac{\partial f(x)}{\partial x(j)} = \frac{\partial}{\partial x(j)}J_\mathrm{theory}(Airmass(i)) \times\frac{1}{J_\mathrm{noise}(i)}
\end{equation}

leading to:

\begin{equation}
\frac{\partial f(i)}{\partial x(1)} = - \eta_\mathrm{tel} J_\mathrm{atm}
\exp(-\tau A(i))
\end{equation}

\begin{equation}
\frac{\partial f(i)}{\partial x(2)} = A(i) B \eta_\mathrm{tel} J_\mathrm{atm}
\exp(-\tau A(i))
\end{equation}

      The array USER is used to get the necessary ancillary information into
      the routine:-
      \sstitemlist{

         \sstitem
          USER (1) = J\_TEL

         \sstitem
          USER (2) = J\_ATM

         \sstitem
          USER (3) = ETA\_TEL

         \sstitem
          USER (4:M$+$3)     = the measured airmasses

         \sstitem
          USER (M$+$4:2M$+$3)  = the measured sky temperatures

         \sstitem
          USER (2M$+$4:3M$+$3) = the errors on the measured sky temperatures
      }
   }
   \sstinvocation{
       CALL SCULIB\_SKYFUNC\_2 (MODE, M, N, LDFJ, X, F, FJAC, NSTATE,
         IUSER, USER)
   }
   \sstarguments{
      \sstsubsection{
         See NAG manual description of OBJFUN parameter in E04UPF.
      }{
      }
   }
   \sstdiytopic{
      Copyright
   }{
      Copyright \copyright 1995,1996,1997,1998,1999 Particle Physics and Astronomy
      Research Council. All Rights Reserved.
   }
}
\sstroutine{
   SCULIB\_SKYFUNCD
}{
   Return the value of the partial derivatives of the skydip function
}{
   \sstdescription{
      This routine returns the value of the skydip function.
      The parameters are passed as follows:
      \sstitemlist{

         \sstitem
              P(1) = ETA\_TEL

         \sstitem
              P(2) = B

         \sstitem
              P(3) = TAU

         \sstitem
              P(4) = J\_AMB

         \sstitem
              P(5) = J\_TEL

      }
      The  theoretical value of the skydip
      for the given input parameters is given by:

\begin{equation}
f(x) = \frac{J_\mathrm{theory} (Airmass) - J_\mathrm{meas} (Airmass)}{J_\mathrm{noise} (Airmass)}
\end{equation}

      where:-

\begin{equation}
J_\mathrm{theory} = (1-\eta_\mathrm{tel}) J_\mathrm{tel} +
\eta_\mathrm{tel}J_\mathrm{atm} (1-B \exp(-\tau Airmass))
\end{equation}

\begin{equation}
            J_\mathrm{atm} = J_\mathrm{amb} X_G
\end{equation}

\begin{equation}
X_G = 1 + \frac{h1\times h2}{J_\mathrm{amb}} \exp (-\tau \times
Airmass / X\_Gconst)
\end{equation}

      The partial derivatives are:

\begin{equation}
\frac{\partial f(i)}{\partial\eta_\mathrm{tel}} = DR(1) =
-J_\mathrm{tel}+J_\mathrm{atm} - B J_\mathrm{tel} \exp(-\tau A)
\end{equation}

\begin{equation}
\frac{\partial f(i)}{\partial B} = DR(2) = - \eta_\mathrm{tel} J_\mathrm{atm}
\exp(-\tau A)
\end{equation}

\begin{equation}
\frac{\partial X_G}{\partial\tau} = - \frac{H_1 H_2 A}{J_\mathrm{amb}
X_\mathrm{Gconst}} \exp(-\tau A / X_\mathrm{Gconst})
\end{equation}

\begin{equation}
\frac{\partial J_\mathrm{atm}}{\partial\tau} = J_\mathrm{amb} \frac{\partial
X_G}{\partial\tau}
\end{equation}

\begin{equation}
\frac{\partial f}{\partial \tau} = DR(3) = \eta_\mathrm{tel} \frac{\partial
J_\mathrm{atm}}{\partial\tau} - \eta_\mathrm{tel} B \exp(-\tau A) \left(\frac{\partial
J_\mathrm{atm}}{\partial\tau} - J_\mathrm{atm} A\right)
\end{equation}

\begin{equation}
\frac{\partial f}{\partial J_\mathrm{amb}} = \frac{\partial f}{\partial
J_\mathrm{tel}} = 0.0 \quad(\text{fixed parameters})
\end{equation}

   }
   \sstinvocation{
      CALL SCULIB\_SKYFUNCD(X, P, DR, M)
   }
   \sstarguments{
      \sstsubsection{
         X = REAL (Given)
      }{
         The airmass used to calculate the function
      }
      \sstsubsection{
         P = REAL (Given)
      }{
         The parameter values
      }
      \sstsubsection{
         M = INTEGER (Given)
      }{
         The number of parameters
      }
      \sstsubsection{
         DR = REAL (Returned)
      }{
         The values of the partial derivatives
      }
   }
   \sstdiytopic{
      Copyright
   }{
      Copyright \copyright 1995,1996,1997,1998,1999 Particle Physics and Astronomy
      Research Council. All Rights Reserved.
   }
}
\sstroutine{
   SCULIB\_SPLINE\_PDA\_IDBVIP
}{
   Fit a surface to an irregular grid using the PDA function PDA\_IDBVIP
}{
   \sstdescription{
      This routine provides a wrapper for the PDA\_IDBVIP spline interpolation
      routine. The quirks of the algorithm are dealt with here so that
      SCULIB\_SPLINE\_REGRID does not have to know anything about the
      interpolation routine.
   }
   \sstinvocation{
       CALL SCULIB\_SPLINE\_PDA\_IDBVIP (NDP, X\_IN, Y\_IN, DATA\_IN,
         N\_PTS, X\_OUT, Y\_OUT, DATA\_OUT, STATUS)
   }
   \sstarguments{
      \sstsubsection{
         NDP = INTEGER (Given)
      }{
         Number of data points on irregular grid
      }
      \sstsubsection{
         X\_IN ( NDP ) = REAL (Given)
      }{
         X coordinates of input data
      }
      \sstsubsection{
         Y\_IN ( NDP ) = REAL (Given)
      }{
         Y coordinates of input data
      }
      \sstsubsection{
         DATA\_IN ( NDP ) = REAL (Given)
      }{
         Data values for each X,Y
      }
      \sstsubsection{
         N\_PTS = INTEGER (Given)
      }{
         Number of points in output grid
      }
      \sstsubsection{
         X\_OUT ( N\_PTS ) = REAL (Given)
      }{
         X coordinates of output points
      }
      \sstsubsection{
         Y\_OUT ( N\_PTS ) = REAL (Given)
      }{
         Y coordinates of output points
      }
      \sstsubsection{
         DATA\_OUT ( N\_PTS ) = REAL (Returned)
      }{
         Output data values
      }
      \sstsubsection{
         STATUS = INTEGER (Given \& Returned)
      }{
         Global status
      }
   }
   \sstdiytopic{
      Copyright
   }{
      Copyright \copyright 1995,1996,1997,1998,1999 Particle Physics and Astronomy
      Research Council. All Rights Reserved.
   }
}
\sstroutine{
   SCULIB\_SPLINE\_PDA\_IDSFFT
}{
   Fit a surface to an irregular grid using the PDA function PDA\_IDSFFT
}{
   \sstdescription{
      This routine provides a wrapper for the PDA\_IDSFFT spline interpolation
      routine. The quirks of the algorithm are dealt with here so that
      SCULIB\_SPLINE\_REGRID does not have to know anything about the
      interpolation routine.
   }
   \sstinvocation{
       CALL SCULIB\_SPLINE\_PDA\_IDSFFT (NDP, X\_IN, Y\_IN, DATA\_IN,
         NX\_OUT, NY\_OUT, X\_OUT, Y\_OUT, DATA\_OUT, STATUS)
   }
   \sstarguments{
      \sstsubsection{
         NDP = INTEGER (Given)
      }{
         Number of data points on irregular grid
      }
      \sstsubsection{
         X\_IN ( NDP ) = REAL (Given)
      }{
         X coordinates of input data
      }
      \sstsubsection{
         Y\_IN ( NDP ) = REAL (Given)
      }{
         Y coordinates of input data
      }
      \sstsubsection{
         DATA\_IN ( NDP ) = REAL (Given)
      }{
         Data values for each X,Y
      }
      \sstsubsection{
         NX\_OUT = INTEGER (Given)
      }{
         Number of pixels in X direction
      }
      \sstsubsection{
         NY\_OUT = INTEGER (Given)
      }{
         Number of pixels in Y direction
      }
      \sstsubsection{
         X\_OUT ( NX\_OUT ) = REAL (Given)
      }{
         X coordinates of output points
      }
      \sstsubsection{
         Y\_OUT ( NY\_OUT ) = REAL (Given)
      }{
         Y coordinates of output points
      }
      \sstsubsection{
         DATA\_OUT ( NX\_OUT, NY\_OUT ) = REAL (Returned)
      }{
         Output data values
      }
      \sstsubsection{
         STATUS = INTEGER (Given \& Returned)
      }{
         Global status
      }
   }
   \sstdiytopic{
      Copyright
   }{
      Copyright \copyright 1995,1996,1997,1998,1999 Particle Physics and Astronomy
      Research Council. All Rights Reserved.
   }
}
\sstroutine{
   SCULIB\_SPLINE\_PDA\_SURFIT
}{
   Fit a surface to an irregular grid using the PDA function PDA\_SURFIT
}{
   \sstdescription{
      This routine provides a wrapper for the PDA\_SURFIT spline interpolation
      routine. The quirks of the algorithm are dealt with here so that
      SCULIB\_SPLINE\_REGRID does not have to know anything about the interpolation
      routine.
   }
   \sstinvocation{
       CALL SCULIB\_SPLINE\_PDA\_SURFIT (NDP, X\_IN, Y\_IN, DATA\_IN,
          VARIANCE\_IN, NX\_OUT, NY\_OUT, X\_OUT, Y\_OUT, DATA\_OUT, STATUS)
   }
   \sstarguments{
      \sstsubsection{
         NDP = INTEGER (Given)
      }{
         Number of data points on irregular grid
      }
      \sstsubsection{
         S = REAL (Given)
      }{
         Smoothness factor
      }
      \sstsubsection{
         X\_IN ( NDP ) = REAL (Given)
      }{
         X coordinates of input data
      }
      \sstsubsection{
         Y\_IN ( NDP ) = REAL (Given)
      }{
         Y coordinates of input data
      }
      \sstsubsection{
         DATA\_IN ( NDP ) = REAL (Given)
      }{
         Data values for each X,Y
      }
      \sstsubsection{
         VARIANCE\_IN ( NDP ) = REAL (Given)
      }{
         Variance for each data value
      }
      \sstsubsection{
         NX\_OUT = INTEGER (Given)
      }{
         Number of pixels in X direction
      }
      \sstsubsection{
         NY\_OUT = INTEGER (Given)
      }{
         Number of pixels in Y direction
      }
      \sstsubsection{
         X\_OUT ( NX\_OUT ) = REAL (Given)
      }{
         X coordinates of output points
      }
      \sstsubsection{
         Y\_OUT ( NY\_OUT ) = REAL (Given)
      }{
         Y coordinates of output points
      }
      \sstsubsection{
         DATA\_OUT ( NX\_OUT, NY\_OUT ) = REAL (Returned)
      }{
         Output data values
      }
      \sstsubsection{
         STATUS = INTEGER (Given \& Returned)
      }{
         Global status
      }
   }
   \sstdiytopic{
      Copyright
   }{
      Copyright \copyright 1995,1996,1997,1998,1999 Particle Physics and Astronomy
      Research Council. All Rights Reserved.
   }
}
\sstroutine{
   SCULIB\_SPLINE\_REGRID
}{
   Regrid supplied data onto a rectangular grid using a spline
}{
   \sstdescription{
      This routine takes data with a variance array and x y positions
      and regrids it onto a rectangular grid using a spline interpolation
      algorithm (from PDA). Each integration is regridded separately and
      added into a mean output image.
   }
   \sstinvocation{
       CALL SCULIB\_SPLINE\_REGRID( METHOD, SFACTOR, MAX\_MAPS, N\_MAPS,
            N\_BOLS, N\_INTS, DIAMETER, WAVELENGTH, PXSIZE, NX\_OUT, NY\_OUT,
            I\_CENTRE, J\_CENTRE, WEIGHT, INT\_PTR, DATA\_PTR, VAR\_PTR,
            XPOS\_PTR,YPOS\_PTR, OUT\_DATA, OUT\_VARIANCE,
            OUT\_QUALITY, OUT\_WEIGHT, STATUS )
   }
   \sstarguments{
      \sstsubsection{
         METHOD = CHARACTER (Given)
      }{
         Spline method to use.
      }
      \sstsubsection{
         SFACTOR = REAL (Given)
      }{
         Smoothing factor to be used for PDA\_SURFIT
      }
      \sstsubsection{
         MAX\_MAPS = INTEGER (Given)
      }{
         Max allowed number of input maps (eg FILES in rebin)
      }
      \sstsubsection{
         N\_MAPS = INTEGER (Given)
      }{
         Number of input maps
      }
      \sstsubsection{
         N\_BOLS( N\_MAPS ) = INTEGER (Given)
      }{
         Number of bolometers in each map
      }
      \sstsubsection{
         MAX\_INTS = INTEGER (Given)
      }{
         Maximum number of separate integrations
      }
      \sstsubsection{
         N\_INTS( N\_MAPS ) = INTEGER (Given)
      }{
         Number of integrations in each input map
      }
      \sstsubsection{
         DIAMETER = REAL (Given)
      }{
         Diameter of dish in metres
      }
      \sstsubsection{
         WAVELENGTH = REAL (Given)
      }{
         Wavelength of observation
      }
      \sstsubsection{
         PXSIZE = REAL (Given)
      }{
         Pixel size of output map (radians)
      }
      \sstsubsection{
         NX\_OUT = INTEGER (Given)
      }{
         Number of X pixels in output map
      }
      \sstsubsection{
         NY\_OUT = INTEGER (Given)
      }{
         Number of Y pixels in output map
      }
      \sstsubsection{
         I\_CENTRE = INTEGER (Given)
      }{
         Reference pixel in X
      }
      \sstsubsection{
         J\_CENTRE = INTEGER (Given)
      }{
         Reference pixel in Y
      }
      \sstsubsection{
         WEIGHT( N\_MAPS ) = REAL (Given)
      }{
         Relative weight of each input map
      }
      \sstsubsection{
         INT\_LIST( MAX\_MAPS, MAX\_INTS $+$ 1 ) = INTEGER (Given)
      }{
         Position of each integration in data (cf DEM\_PNTR for integrations)
      }
      \sstsubsection{
         DATA\_PTR( N\_MAPS ) = INTEGER (Given)
      }{
         Array of pointers to input data
      }
      \sstsubsection{
         VAR\_PTR( N\_MAPS ) = INTEGER (Given)
      }{
         Array of pointers to input data variance
      }
      \sstsubsection{
         XPOS\_PTR( N\_MAPS ) = INTEGER (Given)
      }{
         Array of pointers to input data X positions
      }
      \sstsubsection{
         YPOS\_PTR( N\_MAPS ) = INTEGER (Given)
      }{
         Array of pointers to input data Y positions
      }
      \sstsubsection{
         OUT\_DATA( NX\_OUT, NY\_OUT ) = REAL (Returned)
      }{
         Regular gridded data
      }
      \sstsubsection{
         OUT\_VARIANCE( NX\_OUT, NY\_OUT ) = REAL (Returned)
      }{
         Variance of output data
      }
      \sstsubsection{
         OUT\_DATA( NX\_OUT, NY\_OUT ) = BYTE (Returned)
      }{
         Quality of output map
      }
      \sstsubsection{
         OUT\_WEIGHT( NX\_OUT, NY\_OUT) = REAL (Returned)
      }{
         Total weight of each pixel
      }
      \sstsubsection{
         STATUS = \_INTEGER (Given \& Returned)
      }{
         Global status
      }
   }
   \sstdiytopic{
      Copyright
   }{
      Copyright \copyright 1995,1996,1997,1998,1999 Particle Physics and Astronomy
      Research Council. All Rights Reserved.
   }
}
\sstroutine{
   SCULIB\_SPLINE\_REGRID\_1
}{
   Calculate the areas of the output map that contain good data points
}{
   \sstdescription{
      Calculate the areas of the output map that will be affected by
      points in the input data array. This is used to generate a mask
      so that the rebinned output map only contains data points close
      to input data. Similar to SCULIB\_WTFN\_REGRID\_1.
   }
   \sstinvocation{
       CALL SCULIB\_SPLINE\_REGRID\_1 (DIAMETER, WAVELENGTH,
            IN\_DATA, X, Y, NPIX,
            PIXSPACE, NI, NJ, ICEN, JCEN, QUALITY, STATUS)
   }
   \sstarguments{
      \sstsubsection{
         DIAMETER   = REAL (Given)
      }{
         Diameter of telescope
      }
      \sstsubsection{
         WAVELENGTH = REAL (Given)
      }{
         Wavelength of the observation
      }
      \sstsubsection{
         IN\_DATA (NPIX)                   = REAL (Given)
      }{
         The input data values. (Used to define quality)
      }
      \sstsubsection{
         X( NPIX ) = REAL (Given)
      }{
         The x coordinates of the input pixels
      }
      \sstsubsection{
         Y( NPIX ) = REAL (Given)
      }{
         The y coordinates of the input pixels
      }
      \sstsubsection{
         NPIX = INTEGER (Given)
      }{
         the number of input pixels
      }
      \sstsubsection{
         PIXSPACE = DOUBLE PRECISION (Given)
      }{
         the pixel spacing of the output pixels
      }
      \sstsubsection{
         NI = INTEGER (Given)
      }{
         The number of output pixels in the x direction
      }
      \sstsubsection{
         NJ = INTEGER (Given)
      }{
         The number of output pixels in the y direction
      }
      \sstsubsection{
         ICEN = INTEGER (Given)
      }{
         the x index of the centre of the output array
      }
      \sstsubsection{
         JCEN = INTEGER (Given)
      }{
         the y index of the centre of the output array
      }
      \sstsubsection{
         SCRATCH (NI, NJ) = UBYTE (Returned)
      }{
         Quality mask
      }
      \sstsubsection{
         STATUS = INTEGER (Given and Returned)
      }{
         The global status.
      }
   }
   \sstdiytopic{
      Copyright
   }{
      Copyright \copyright 1995,1996,1997,1998,1999 Particle Physics and Astronomy
      Research Council. All Rights Reserved.
   }
}
\sstroutine{
   SCULIB\_SPLIT\_FILE\_SPEC
}{
   Splits a string into a filename and SCUBA section
}{
   \sstdescription{
      This routine takes a string and returns the filename (ie text
      before a \{) and the data specifications. Each specification
      is indicated by a \{\}. More than one spec can be returned.
      If the string contains a {\tt '}!{\tt '} then we will use the invers of the
      specified section.
   }
   \sstinvocation{
       CALL SCULIB\_SPLIT\_FILE\_SPEC( NAME, MAX\_SPEC, FILENAME, NSPEC,
            DATA\_SPEC, STATUS)
   }
   \sstarguments{
      \sstsubsection{
         NAME = CHAR (Given)
      }{
         Input string
      }
      \sstsubsection{
         MAX\_SPEC = INTEGER (Given)
      }{
         Max number of specifications allowed
      }
      \sstsubsection{
         FILENAME = CHAR (Returned)
      }{
         The name of the file (text before first spec)
      }
      \sstsubsection{
         NSPEC = INTEGER (Returned)
      }{
         Number of specs found/returned (can be 0)
      }
      \sstsubsection{
         DATA\_SPEC(NSPEC) = CHAR (Returned)
      }{
         The actual specification for the SCUBA section
      }
      \sstsubsection{
         USE\_SECTION = LOGICAL (Returned)
      }{
         Set to true by default and false if a {\tt '}-{\tt '} is found
         at the end of the section (ie after last closing brace) --
         a {\tt '}-{\tt '} indicates that we are using the inverse of the section.
      }
      \sstsubsection{
         STATUS = INTEGER (Given and Returned)
      }{
         Global Status value
      }
   }
   \sstdiytopic{
      Copyright
   }{
      Copyright \copyright 1995,1996,1997,1998,1999 Particle Physics and Astronomy
      Research Council. All Rights Reserved.
   }
}
\sstroutine{
   SCULIB\_SQROOTR
}{
   take the square root of a real array
}{
   \sstdescription{
      Puts square root of real input array into output array. If FLAGGED and
      QUALITY are both false then negative input numbers give a zero output.
      If FLAGGED is true then input bad values, or negative input numbers,
      give a bad output. If QUALITY is true then input values with bad
      quality, or negative input numbers, lead to zero output data and bad
      output quality.
   }
   \sstinvocation{
       CALL SCULIB\_SQROOTR (N, IN\_DATA, IN\_QUAL, ROOT, ROOT\_QUAL, QUALITY,
          FLAGGED)
   }
   \sstarguments{
      \sstsubsection{
         N                       = INTEGER (Given)
      }{
         the number of array elements
      }
      \sstsubsection{
         IN\_DATA (N)             = REAL (Given)
      }{
         the array whose root is to be taken
      }
      \sstsubsection{
         IN\_QUAL (N)             = INTEGER (Given)
      }{
         the quality on the input data
      }
      \sstsubsection{
         ROOT (N)                = REAL (Returned)
      }{
         the array to hold the square roots (may be the same as IN\_DATA)
      }
      \sstsubsection{
         ROOT\_QUAL (N)           = INTEGER (Returned)
      }{
         the quality on the output root
      }
      \sstsubsection{
         QUALITY                 = LOGICAL (Given)
      }{
         .TRUE. if input quality array exists
      }
      \sstsubsection{
         FLAGGED                 = LOGICAL (Given)
      }{
         .TRUE. if input data has flagged values
      }
   }
   \sstnotes{
      Consider using VEC\_SQRTR (SUN/39) instead.
   }
   \sstdiytopic{
      Copyright
   }{
      Copyright \copyright 1995,1996,1997,1998,1999 Particle Physics and Astronomy
      Research Council. All Rights Reserved.
   }
   \sstimplementationstatus{
      Uses INTEGER rather than BYTE Quality.
   }
}
\sstroutine{
   SCULIB\_STANDARD\_APPARENT
}{
   convert apparent RA,Decs from one date to another
}{
   \sstdescription{
      This routine converts a list of apparent RA,Decs for one date into
      apparent coordinates at another. To do this it calls SLA\_MAPPA and
      SLA\_AMPQK to convert the input coordinates into mean coordinates for
      a J2000.0 equinox, then SLA\_MAPPA and SLA\_MAPQKZ to convert the mean
      coordinates into apparent on the output date.
   }
   \sstinvocation{
       CALL SCULIB\_STANDARD\_APPARENT (N, RA\_APP, DEC\_APP, IN\_MJD, OUT\_MJD,
         STATUS)
   }
   \sstarguments{
      \sstsubsection{
         N                      = INTEGER (Given)
      }{
         number of positions to be converted
      }
      \sstsubsection{
         RA\_APP (N)             = DOUBLE PRECISION (Given and returned)
      }{
         apparent RA (radians)
      }
      \sstsubsection{
         DEC\_APP (N)            = DOUBLE PRECISION (Given and returned)
      }{
         apparent Dec (radians)
      }
      \sstsubsection{
         IN\_MJD                 = DOUBLE PRECISION (Given)
      }{
         date of input coordinates (modified Julian day)
      }
      \sstsubsection{
         OUT\_MJD                = DOUBLE PRECISION (Given)
      }{
         date of output coordinates (modified Julian day)
      }
      \sstsubsection{
         STATUS                 = INTEGER (Given and returned)
      }{
         global status
      }
   }
   \sstdiytopic{
      Copyright
   }{
      Copyright \copyright 1995,1996,1997,1998,1999 Particle Physics and Astronomy
      Research Council. All Rights Reserved.
   }
}
\sstroutine{
   SCULIB\_STATD
}{
   To calculate mean and standard deviation of a DOUBLE array
}{
   \sstdescription{
      This routine calculates the mean and standard deviation of
      a double array. If a positive value of CLIP is given, the data
      set will be clipped at CLIP sigma repeatedly until all
      points are within CLIP sigma. Mean and standard deviation are those
      values derived after iterative clipping.
   }
   \sstinvocation{
       CALL SCULIB\_STATD(N\_POS, SCUDATA, SCUQUAL, BADBIT,
            NGOOD, MEAN, MEDIAN, SUM, SUMSQ, STDEV, QSORT, STATUS)
   }
   \sstarguments{
      \sstsubsection{
         N\_POS = \_INTEGER (Given)
      }{
         Number of jiggle positions in data set
      }
      \sstsubsection{
         CLIP = \_REAL (Given)
      }{
         Number of sigma to clip iteratively
      }
      \sstsubsection{
         SCUDATA(N\_POS) = \_DOUBLE (Given)
      }{
         The data
      }
      \sstsubsection{
         SCUQUAL(N\_POS) = \_BYTE (Given)
      }{
         The data quality
      }
      \sstsubsection{
         BADBIT = \_BYTE (Given)
      }{
         Bad bit mask
      }
      \sstsubsection{
         NGOOD = \_INTEGER (Returned)
      }{
         Number of good pixels in mean
      }
      \sstsubsection{
         MEAN = \_DOUBLE (Returned)
      }{
         Mean value
      }
      \sstsubsection{
         MEDIAN = \_DOUBLE (Returned)
      }{
         Median of good data
      }
      \sstsubsection{
         SUM = \_DOUBLE (Returned)
      }{
         Data sum
      }
      \sstsubsection{
         SUMSQ = \_DOUBLE (Returned)
      }{
         Sum of squares
      }
      \sstsubsection{
         STDEV = \_DOUBLE (Returned)
      }{
         Standard deviation
      }
      \sstsubsection{
         QSORT = \_DOUBLE (Returned)
      }{
         Sorted data set
      }
      \sstsubsection{
         STATUS = INTEGER (Given and Returned)
      }{
         Global Status value
      }
   }
   \sstdiytopic{
      Copyright
   }{
      Copyright \copyright 1995,1996,1997,1998,1999 Particle Physics and Astronomy
      Research Council. All Rights Reserved.
   }
   \sstimplementationstatus{
      \sstitemlist{

         \sstitem
         Deals with bad pixels

         \sstitem
         Uses PDA sort routine
      }
   }
}
\sstroutine{
   SCULIB\_STATR
}{
   To calculate mean and standard deviation of a REAL array
}{
   \sstdescription{
      This routine calculates the mean and standard deviation of
      a real array. If a positive value of CLIP is given, the data
      set will be clipped at CLIP sigma repeatedly until all
      points are within CLIP sigma. Mean and standard deviation are those
      values derived after iterative clipping.
   }
   \sstinvocation{
       CALL SCULIB\_STATR(N\_POS, SCUDATA, SCUQUAL, BADBIT,
            NGOOD, MEAN, MEDIAN, SUM, SUMSQ, STDEV, QSORT, STATUS)
   }
   \sstarguments{
      \sstsubsection{
         N\_POS = \_INTEGER (Given)
      }{
         Number of jiggle positions in data set
      }
      \sstsubsection{
         CLIP = \_REAL (Given)
      }{
         Number of sigma to clip iteratively
      }
      \sstsubsection{
         SCUDATA(N\_POS) = \_REAL (Given)
      }{
         The data
      }
      \sstsubsection{
         SCUQUAL(N\_POS) = \_BYTE (Given)
      }{
         The data quality
      }
      \sstsubsection{
         BADBIT = \_BYTE (Given)
      }{
         Bad bit mask
      }
      \sstsubsection{
         NGOOD = \_INTEGER (Returned)
      }{
         Number of good pixels in mean
      }
      \sstsubsection{
         MEAN = \_DOUBLE (Returned)
      }{
         Mean value
      }
      \sstsubsection{
         MEDIAN = \_DOUBLE (Returned)
      }{
         Median of good data
      }
      \sstsubsection{
         SUM = \_DOUBLE (Returned)
      }{
         Data sum
      }
      \sstsubsection{
         SUMSQ = \_DOUBLE (Returned)
      }{
         Sum of squares
      }
      \sstsubsection{
         STDEV = \_DOUBLE (Returned)
      }{
         Standard deviation
      }
      \sstsubsection{
         QSORT = \_REAL (Returned)
      }{
         Sorted data set
      }
      \sstsubsection{
         STATUS = INTEGER (Given and Returned)
      }{
         Global Status value
      }
   }
   \sstdiytopic{
      Copyright
   }{
      Copyright \copyright 1995,1996,1997,1998,1999 Particle Physics and Astronomy
      Research Council. All Rights Reserved.
   }
   \sstimplementationstatus{
      \sstitemlist{

         \sstitem
         Deals with bad pixels

         \sstitem
         Uses PDA sort routine
      }
   }
}
\sstroutine{
   SCULIB\_SUB\_TAUZ
}{
   get tauz appropriate for each sub-instrument
}{
   \sstdescription{
      This routine finds the zenith sky optical depth appropriate for the
      filter in front of each sub-instrument being used. If it cannot find
      an optical depth for the required filter the routine will output a
      warning message and return an optical depth of 0.0.

      The times and dates at which the optical depths were measured are also
      returned. If the optical depth for the required filter cannot be found
      then the date will be that at which the routine is run.
   }
   \sstinvocation{
       CALL SCULIB\_GET\_TAUZ (N\_SUBS, SUB\_FILTER, N\_SKY, SKY\_FILTER,
         SKY\_TAUZ, SKY\_DAY, SUB\_TAUZ, SUB\_DAY, STATUS)
   }
   \sstarguments{
      \sstsubsection{
         N\_SUBS                         = INTEGER (Given)
      }{
         number of sub-instruments
      }
      \sstsubsection{
         SUB\_FILTER (N\_SUBS)            = CHARACTER$*$($*$) (Given)
      }{
         names of filters in front of sub-instruments
      }
      \sstsubsection{
         N\_SKY                          = INTEGER (Given)
      }{
         number of filters for which data is available
      }
      \sstsubsection{
         SKY\_FILTER (N\_SKY)             = CHARACTER$*$($*$) (Given)
      }{
         filters for which data is available
      }
      \sstsubsection{
         SKY\_TAUZ (N\_SKY)               = REAL (Given)
      }{
         tauz for each possible filter
      }
      \sstsubsection{
         SKY\_DAY (N\_SKY)                = DOUBLE PRECISION (Given)
      }{
         the date and time of the tauz measurement specified as day
         number since 1st Jan.
      }
      \sstsubsection{
         SUB\_TAUZ (N\_SUBS)              = REAL (Returned)
      }{
         tauz appropriate to sub-instruments
      }
      \sstsubsection{
         SUB\_DAY (N\_SUBS)               = DOUBLE PRECISION (Returned)
      }{
         the date and time of the tauz measurement
      }
      \sstsubsection{
         STATUS                         = INTEGER (Given and returned)
      }{
         global status
      }
   }
   \sstdiytopic{
      Copyright
   }{
      Copyright \copyright 1995,1996,1997,1998,1999 Particle Physics and Astronomy
      Research Council. All Rights Reserved.
   }
}
\sstroutine{
   SCULIB\_SUBARE
}{
   subtract one real array from another into a third
}{
   \sstdescription{
      Subtracts two real arrays. Note that any of the arrays may be the same.
   }
   \sstinvocation{
       CALL SCULIB\_SUBARE (N, ARRAY1, ARRAY2, ARRAY3, Q1DATA,
          Q2DATA, Q3DATA, V1DATA, V2DATA, V3DATA, QUALITY, FLAGGED,
          VARIANCE)
   }
   \sstarguments{
      \sstsubsection{
         N            = INTEGER (Given)
      }{
         Number of elements in each array
      }
      \sstsubsection{
         ARRAY1 (N)   = REAL (Given)
      }{
         Input array
      }
      \sstsubsection{
         ARRAY2 (N)   = REAL (Given)
      }{
         Second input array
      }
      \sstsubsection{
         ARRAY3 (N)   = REAL (Returned)
      }{
         Result array.  ARRAY3 = ARRAY1 - ARRAY2
      }
      \sstsubsection{
         Q1DATA (N)   = INTEGER (Given)
      }{
         Quality array for first input array
      }
      \sstsubsection{
         Q2DATA (N)   = INTEGER (Given)
      }{
         Quality array for second input array
      }
      \sstsubsection{
         Q3DATA (N)   = INTEGER (Returned)
      }{
         Quality array for output array
      }
      \sstsubsection{
         V1DATA (N)   = REAL (Given)
      }{
         Variance array for first input array
      }
      \sstsubsection{
         V2DATA (N)   = REAL (Given)
      }{
         Variance array for second input array
      }
      \sstsubsection{
         V3DATA (N)   = REAL (Returned)
      }{
         Variance array for output array
      }
      \sstsubsection{
         QUALITY      = LOGICAL (Given)
      }{
         True if input has quality information
      }
      \sstsubsection{
         FLAGGED      = LOGICAL (Given)
      }{
         True if input has flagged data values
      }
      \sstsubsection{
         VARIANCE     = LOGICAL (Given)
      }{
         True if both input arrays have variance arrays
      }
   }
   \sstnotes{
      Consider using VEC\_SUBR (SUN/39) instead
   }
   \sstdiytopic{
      Copyright
   }{
      Copyright \copyright 1995,1996,1997,1998,1999 Particle Physics and Astronomy
      Research Council. All Rights Reserved.
   }
}
\sstroutine{
   SCULIB\_SUMAD
}{
   sum the elements of a double precision array
}{
   \sstdescription{
      This routine adds up the NELM elements of the double precision ARRAY
      and puts the result in SUM.
   }
   \sstinvocation{
      CALL SCULIB\_SUMAD (NELM, ARRAY, SUM)
   }
   \sstarguments{
      \sstsubsection{
         NELM               = INTEGER (Given)
      }{
         the number of elements in the array
      }
      \sstsubsection{
         ARRAY (NELM)       = DOUBLE PRECISION (Given)
      }{
         the array whose elements are to be summed
      }
      \sstsubsection{
         SUM                = DOUBLE PRECISION (Returned)
      }{
         the sum
      }
   }
   \sstdiytopic{
      Copyright
   }{
      Copyright \copyright 1995,1996,1997,1998,1999 Particle Physics and Astronomy
      Research Council. All Rights Reserved.
   }
   \sstimplementationstatus{
      No status checking
   }
}
\sstroutine{
   SCULIB\_TIDY\_LINE
}{
   remove tabs and comments from a line
}{
   \sstdescription{
      This routine tidies up a character string and returns its length
      ignoring trailing blanks. If the returned length is 0 the returned
      string will be {\tt '} {\tt '}, otherwise it will be equal to the input string
      truncated at the used length.

      The tidying involves:-
      \sstitemlist{

         \sstitem
          calling CHR\_CLEAN to remove non-printable characters.

         \sstitem
          replacing all HT (tab) characters by spaces.

         \sstitem
          truncating the string at the character before a COMCHAR character, if
            present (the characters after COMCHAR are assumed to be comments).
      }
   }
   \sstinvocation{
      CALL SCULIB\_TIDY\_LINE (COMCHAR, LINE, LENGTH)
   }
   \sstarguments{
      \sstsubsection{
         COMCHAR    = CHARACTER$*$1 (Given)
      }{
         the character at the beginning of a comment
      }
      \sstsubsection{
         LINE       = CHARACTER$*$($*$) (Given and returned)
      }{
         the line to be tidied
      }
      \sstsubsection{
         LENGTH     = INTEGER (Returned)
      }{
         the length of the string, ignoring trailing blanks
      }
   }
   \sstdiytopic{
      Copyright
   }{
      Copyright \copyright 1995,1996,1997,1998,1999 Particle Physics and Astronomy
      Research Council. All Rights Reserved.
   }
}
\sstroutine{
   SCULIB\_UNPACK
}{
   routine to unpack compressed resampled map
}{
   \sstdescription{
      This routine unpacks a compressed SCUBA resampled data array into
      a 2-d image. It does this by cycling up the j slices of the image
      and -

         if the packed data has pixels for that slice
      \sstitemlist{

         \sstitem
              fill the image slice data (data sum, variance sum and weight sum)
                 with zeroes and bad quality before the packed slice starts.

         \sstitem
              copy the data and quality of the packed slice into the appropriate
                 part of the image slice

         \sstitem
              fill the image slice with zeroes and bad quality after the packed
                 slice ends.

      }
         else
      \sstitemlist{

         \sstitem
              fill the image slice with zeroes and bad quality

      }
         endif
   }
   \sstinvocation{
       CALL SCULIB\_UNPACK (NBS\_DIM1, NBS\_DIM2, NBS\_DIM3,
          RESNBS, ISTART, NPIX, POINTER, SUB, RES\_DIM1, RES\_DIM2,
          RES\_DS, RES\_VS, RES\_WS, RES\_QS)
   }
   \sstarguments{
      \sstsubsection{
         NBS\_DIM1                            = INTEGER (Given)
      }{
         first dimension of the RESNBS array
      }
      \sstsubsection{
         NBS\_DIM2                            = INTEGER (Given)
      }{
         second dimension of the RESNBS array
      }
      \sstsubsection{
         NBS\_DIM3                            = INTEGER (Given)
      }{
         third dimension of the RESNBS array
      }
      \sstsubsection{
         RESNBS (NBS\_DIM1, NBS\_DIM2, NBS\_DIM3)) = REAL (Given)
      }{
         the noticeboard `resampled data{\tt '}
      }
      \sstsubsection{
         ISTART (RES\_DIM2)                   = INTEGER (Given)
      }{
         the I index of the start of each packed slice
      }
      \sstsubsection{
         NPIX (RES\_DIM2)                     = INTEGER (Given)
      }{
         the number of pixels in each packed slice
      }
      \sstsubsection{
         POINTER (RES\_DIM2)                  = INTEGER (Given)
      }{
         the pointer to the start of each packed slice
      }
      \sstsubsection{
         SUB                                 = INTEGER (Given)
      }{
         the number of the sub-instrument map to be unpacked
      }
      \sstsubsection{
         RES\_DIM1                            = INTEGER (Given)
      }{
         first dimension of resampled image array
      }
      \sstsubsection{
         RES\_DIM2                            = INTEGER (Given)
      }{
         second dimension of resampled image array
      }
      \sstsubsection{
         RES\_DS (RES\_DIM1, RES\_DIM2)         = REAL (Returned)
      }{
         the unpacked data sum
      }
      \sstsubsection{
         RES\_VS (RES\_DIM1, RES\_DIM2)         = REAL (Returned)
      }{
         the unpacked variance sum
      }
      \sstsubsection{
         RES\_WS (RES\_DIM1, RES\_DIM2)         = REAL (Returned)
      }{
         the unpacked weight sum
      }
      \sstsubsection{
         RES\_QS (RES\_DIM1, RES\_DIM2)         = INTEGER (Returned)
      }{
         the unpacked quality
      }
   }
   \sstdiytopic{
      Copyright
   }{
      Copyright \copyright 1993-1999 Particle Physics and Astronomy
      Research Council. All Rights Reserved.
   }
}
\sstroutine{
   SCULIB\_UNPACK\_CHOPSCAN
}{
   routine to unpack compressed chop-scan data
}{
   \sstdescription{
      Used by the on-line system to unpack notice-board entries
      and retrieve demodulated data.
   }
   \sstinvocation{
       CALL SCULIB\_UNPACK\_CHOPSCAN (RESNBS, N\_POINTS, ISTART,
          NPIX, POINTER, RESDATA\_450, RESWT\_450, RESDATA\_850, RESWT\_850,
          RESDIM1, RESDIM2, BAD)
   }
   \sstarguments{
      \sstsubsection{
         RESNBS (4, N\_POINTS)                = REAL (Given)
      }{
         the `resampled data{\tt '} noticeboard
      }
      \sstsubsection{
         N\_POINTS                            = INTEGER (Given)
      }{
         the 2nd dimension of the RESNBS array
      }
      \sstsubsection{
         ISTART (RESDIM2)                    = INTEGER (Given)
      }{
         the I index of the start of the packed slice at index J
      }
      \sstsubsection{
         NPIX (RESDIM2)                      = INTEGER (Given)
      }{
         the number of pixels in the packed slice for row J
      }
      \sstsubsection{
         POINTER (RESDIM2)                   = INTEGER (Given)
      }{
         the pointer to the start of the slice in the packed array
      }
      \sstsubsection{
         RESDATA\_450 (RESDIM1, RESDIM2)      = REAL (Returned)
      }{
         the unpacked 450 data
      }
      \sstsubsection{
         RESWT\_450 (RESDIM1, RESDIM2)        = REAL (Returned)
      }{
         450 weights
      }
      \sstsubsection{
         RESDATA\_850 (RESDIM1, RESDIM2)      = REAL (Returned)
      }{
         the unpacked 850 image
      }
      \sstsubsection{
         RESWT\_850 (RESDIM1, RESDIM2)        = REAL (Returned)
      }{
         850 weights
      }
      \sstsubsection{
         RESDIM1                             = INTEGER (Given)
      }{
         first dimension of resampled image array
      }
      \sstsubsection{
         RESDIM2                             = INTEGER (Given)
      }{
         second dimension
      }
      \sstsubsection{
         BAD                                 = REAL (Given)
      }{
         value signalling bad pixel
      }
   }
   \sstdiytopic{
      Copyright
   }{
      Copyright \copyright 1995,1996,1997,1998,1999 Particle Physics and Astronomy
      Research Council. All Rights Reserved.
   }
}
\sstroutine{
   SCULIB\_UNPACK\_JIGGLE
}{
   unpack demodulated data onto 2-d map
}{
   \sstdescription{
      This routine unpacks the demodulated data from a switch onto a
      rectangular 2-d map.

         If status is good on entry the data and variance of the output map
      will be initialised to `bad{\tt '} values, the quality to 1. If there are
      any data to unpack the routine will then attempt to do so. The method
      used depends on whether the switch covered part/all of the entire jiggle
      pattern or contains data for several repeats of the jiggle pattern.

         In the first case, the routine will check that the indices of the
      section of jiggle pattern covered by this switch lie within the bounds
      of the full jiggle pattern. If not, an error will be reported and bad
      status returned. Otherwise, the datablock will be unpacked into the
      map as specified by the I\_JIGGLE, J\_JIGGLE arrays. Each map point
      just has its data, variance and quality copied from the demodulated
      data.

         When the data covers several repeats of the jiggle pattern, the
      routine will check that the jiggle index of the first point in the
      datablock is 1 and that the number of data points is an integer multiple
      of the size of the jiggle. If not, an error will be reported and bad
      status returned. Otherwise, the datablock will be unpacked into the
      map as specified by I\_JIGGLE, J\_JIGGLE. The unpacking differs from
      the first case in that demodulated data with bad quality will be
      ignored. Points in the unpacked map will be set to the average of the
      values contributed by the separate jiggles, variances will be calculated
      from the dispersion of the values about the mean if there was more than
      1, or set to the demodulated variance otherwise.
   }
   \sstinvocation{
       CALL SCULIB\_UNPACK\_JIGGLE (N\_JIGGLES, N\_BOLS, DEMOD,
         J\_START, JIGGLE\_COUNT, I\_JIGGLE, J\_JIGGLE, IDIM, JDIM,
         MAP\_DATA, MAP\_VARIANCE, MAP\_QUALITY, J\_END, STATUS)
   }
   \sstarguments{
      \sstsubsection{
         N\_JIGGLES                         = INTEGER (Given)
      }{
         number of jiggles in datablock
      }
      \sstsubsection{
         N\_BOLS                            = INTEGER (Given)
      }{
         number of bolometers measured
      }
      \sstsubsection{
         DEMOD (4, N\_BOLS, N\_JIGGLES)      = REAL (Given)
      }{
         demodulated data
      }
      \sstsubsection{
         J\_START                           = INTEGER (Given)
      }{
         index of first jiggle in datablock in the overall jiggle pattern
      }
      \sstsubsection{
         JIGGLE\_COUNT                      = INTEGER (Given)
      }{
         number of jiggles in jiggle pattern
      }
      \sstsubsection{
         I\_JIGGLE (JIGGLE\_COUNT)           = INTEGER (Given)
      }{
         i index in output map of each jiggle position in pattern
      }
      \sstsubsection{
         J\_JIGGLE (JIGGLE\_COUNT)           = INTEGER (Given)
      }{
         j index in output map of each jiggle position in pattern
      }
      \sstsubsection{
         IDIM                              = INTEGER (Given)
      }{
         i dimension of output map
      }
      \sstsubsection{
         JDIM                              = INTEGER (Given)
      }{
         j dimension of output map
      }
      \sstsubsection{
         MAP\_DATA (IDIM, JDIM, N\_BOLS)     = REAL (Returned)
      }{
         output map data
      }
      \sstsubsection{
         MAP\_VARIANCE (IDIM, JDIM, N\_BOLS) = REAL (Returned)
      }{
         output map variance
      }
      \sstsubsection{
         MAP\_QUALITY (IDIM, JDIM, N\_BOLS)  = INTEGER (Returned)
      }{
         output map quality
      }
      \sstsubsection{
         J\_END                             = INTEGER (Returned)
      }{
         index of last jiggle in datablock in the overall jiggle pattern
      }
      \sstsubsection{
         STATUS                            = INTEGER (Given and returned)
      }{
         global status
      }
   }
   \sstdiytopic{
      Copyright
   }{
      Copyright \copyright 1993-1999 Particle Physics and Astronomy
      Research Council. All Rights Reserved.
   }
}
\sstroutine{
   SCULIB\_UNPACK\_JIGGLE\_SEPARATES
}{
   unpack jiggle data arrays onto 2-d map
}{
   \sstdescription{
      This routine unpacks data from a jiggle observation into a
      rectangular 2-d map.

         If status is good on entry the data and variance of the output map
      will be initialised to `bad{\tt '} values, the quality to 1. If there are
      any data to unpack the routine will then attempt to do so. The method
      used depends on whether the switch covered part/all of the entire jiggle
      pattern or contains data for several repeats of the jiggle pattern.

         In the first case, the routine will check that the indices of the
      section of jiggle pattern covered by this switch lie within the bounds
      of the full jiggle pattern. If not, an error will be reported and bad
      status returned. Otherwise, the data will be unpacked into the
      map as specified by the I\_JIGGLE, J\_JIGGLE arrays.

         When the data cover several repeats of the jiggle pattern, the
      routine will check that the jiggle index of the first point in the
      datablock is 1 and that the number of data points is an integer multiple
      of the size of the jiggle. If not, an error will be reported and bad
      status returned. Otherwise, the data will be unpacked into the
      map as specified by I\_JIGGLE, J\_JIGGLE. The unpacking differs from
      the first case in that points in the unpacked map will be set to the
      average of the good quality values contributed by the separate jiggles,
      variances will be calculated from the dispersion of the values about
      the mean if there was more than 1, or set to the input variance
      otherwise.
   }
   \sstinvocation{
       CALL SCULIB\_UNPACK\_JIGGLE\_SEPARATES (N\_JIGGLES, N\_BOLS, J\_DATA,
         J\_VARIANCE, J\_QUALITY, J\_START, JIGGLE\_COUNT, I\_JIGGLE,
         J\_JIGGLE, IDIM, JDIM, MAP\_DATA, MAP\_VARIANCE, MAP\_QUALITY, MAP\_NUMPTS,
         J\_END, BADBIT, STATUS)
   }
   \sstarguments{
      \sstsubsection{
         N\_JIGGLES                         = INTEGER (Given)
      }{
         number of jiggles in datablock
      }
      \sstsubsection{
         N\_BOLS                            = INTEGER (Given)
      }{
         number of bolometers measured
      }
      \sstsubsection{
         J\_DATA (N\_BOLS, N\_JIGGLES)        = REAL (Given)
      }{
         data for each jiggle measured
      }
      \sstsubsection{
         J\_VARIANCE (N\_BOLS, N\_JIGGLES)    = REAL (Given)
      }{
         variance on J\_DATA
      }
      \sstsubsection{
         J\_QUALITY (N\_BOLS, N\_JIGGLES)     = BYTE (Given)
      }{
         quality on J\_DATA
      }
      \sstsubsection{
         J\_START                           = INTEGER (Given)
      }{
         index of first jiggle in datablock in the overall jiggle pattern
      }
      \sstsubsection{
         JIGGLE\_COUNT                      = INTEGER (Given)
      }{
         number of jiggles in jiggle pattern
      }
      \sstsubsection{
         I\_JIGGLE (JIGGLE\_COUNT)           = INTEGER (Given)
      }{
         i index in output map of each jiggle position in pattern
      }
      \sstsubsection{
         J\_JIGGLE (JIGGLE\_COUNT)           = INTEGER (Given)
      }{
         j index in output map of each jiggle position in pattern
      }
      \sstsubsection{
         IDIM                              = INTEGER (Given)
      }{
         i dimension of output map
      }
      \sstsubsection{
         JDIM                              = INTEGER (Given)
      }{
         j dimension of output map
      }
      \sstsubsection{
         MAP\_DATA (IDIM, JDIM, N\_BOLS)     = REAL (Returned)
      }{
         output map data
      }
      \sstsubsection{
         MAP\_VARIANCE (IDIM, JDIM, N\_BOLS) = REAL (Returned)
      }{
         output map variance
      }
      \sstsubsection{
         MAP\_NUMPTS (IDIM, JDIM, N\_BOLS)   = INTEGER (Returned)
      }{
         output map quality
      }
      \sstsubsection{
         MAP\_QUALITY (IDIM, JDIM, N\_BOLS)  = BYTE (Returned)
      }{
         output map quality
      }
      \sstsubsection{
         J\_END                             = INTEGER (Returned)
      }{
         index of last jiggle in datablock in the overall jiggle pattern
      }
      \sstsubsection{
         BADBIT                            = BYTE  (Given)
      }{
         bad bit mask
      }
      \sstsubsection{
         STATUS                            = INTEGER (Given and returned)
      }{
         global status
      }
   }
   \sstdiytopic{
      Copyright
   }{
      Copyright \copyright 1996,1997,1998,1999 Particle Physics and Astronomy
      Research Council. All Rights Reserved.
   }
}
\sstroutine{
   SCULIB\_UT1
}{
   returns UT1 as a modified Julian day
}{
   \sstdescription{
      Return the current modified Julian day calculated for Hawaiian
      local time.
   }
   \sstinvocation{
      UT1 = SCULIB\_UT1
   }
   \sstarguments{
      \sstsubsection{
         None
      }{
      }
   }
   \sstreturnedvalue{
      \sstsubsection{
         SCULIB\_UT1 = DOUBLE PRECISION
      }{
         UT1 expressed as a modified Julian day (JD - 2400000.5)
      }
   }
   \sstnotes{
      \sstitemlist{

         \sstitem
         Depends on system clock for local time and date

         \sstitem
         Uses the IDATE function to retrieve the current date. This
         only returns a 2-digit year (not Y2K compliant) but is not
         a problem since the output is immediately passed to the SLALIB
         routine SLA\_CALDJ which uses windowing to select the correct
         year.
      }
   }
   \sstdiytopic{
      Copyright
   }{
      Copyright \copyright 1995,1996,1997,1998,1999 Particle Physics and Astronomy
      Research Council. All Rights Reserved.
   }
   \sstbugs{
      The time zone is hard-wired for Hawaii (10).
   }
}
\sstroutine{
   SCULIB\_WTFN\_REGRID
}{
   Regrid supplied data onto a rectangular grid
}{
   \sstdescription{
      This routine takes data with a variance array and x y positions
      and regrids it onto a rectangular grid using a weight function
      interpolation.

      This is an image space implementation of fourier techniques.
      In Fourier terms the technique could be described
      as follows:-
      \sstitemlist{

         \sstitem
         [1] Do a 2-d discrete Fourier transform of the data points. The result
         is a repeating pattern made up of copies of the transform of
         the map as a continuous function, each copied displaced from its
         neighbours by 1/dx, where dx is the sample spacing of the input points
         (assumed equal in x and y). Different copies of the {\tt '}continuous map{\tt '}
         transforms will overlap ({\tt '}alias{\tt '}) if there is any power in the
         map at frequencies greater than 1/(2dx). It is not possible to unravel
         aliased spectra and this constraint leads to the Nyquist sampling
         criterion.

         \sstitem
         [2] We want to derive the {\tt '}continuous map{\tt '} so that map values on the new
         grid mesh can be derived. Do this by multiplying the transform of the
         data by a function that has zero value beyond a radius of 0.5/dx from the
         origin. This will get rid of all the repeats in the pattern and leave
         just the transform of the {\tt '}continuous map{\tt '}.

         \sstitem
         [3] Do an inverse FT on the remaining transform for the points where you
         wish the resampled points to be (note: FFTs implicitly assume that the
         data being transformed DO repeat ad infinitum so we{\tt '}d have to be careful
         when using them to do this).

      }
      The analogue of these process steps in image space is as follows:
      \sstitemlist{

         \sstitem
         [1] Nothing.

         \sstitem
         [2] Convolve the input data with the FT of the function used to isolate
         the {\tt '}continuous map{\tt '} transform.

         \sstitem
         [3] Nothing.

      }
      If the method is done properly, the rebinned map is in fact the map on
      the new sample mesh that has the same FT as the continuous function
      going through the original sample points.

      Convolution Functions:-
      \sstitemlist{

         \sstitem
         [Bessel:]
         For good data and with no time constraint on reduction the best
         convolution function would be one whose FT is a flat-topped cylinder in
         frequency space, centred on the origin and of a radius such that
         frequencies to which the telescope is not sensitive are set to zero.
         Unfortunately, this function is a Bessel function, which extends to
         infinity along both axes and has significant power out to a large radius.
         To work correctly this would require an infinite map and infinite
         computing time. However, a truncated Bessel function should work well on
         a large map, except near the edges. Edge effects can be removed by
         pretending that the map extends further - of course, this only works if
         you know what data the pretend map area should contain, i.e. zeros.
         Another problem with a Bessel function arises from the fact that it does
         truncate the FT of the map sharply. If the data are good then there
         should be nothing but noise power at the truncation radius and the
         truncation of the FT should have no serious effect. However, if the data
         has spikes (power at all frequencies in the FT) or suffers from seeing
         effects such that the data as measured DO have power beyond the
         truncation radius, then this will cause ringing in the rebinned map.

         \sstitem
         [Gaussian:]
         In fact, any function that is finite in frequency space will have
         infinite extent in image space (I think). As such they will all
         drag in some power from the aliased versions of the map transform
         and all suffer from edge effects and large compute times. Some are
         worse than others, however. For example, a Gaussian can have most
         of its power concentrated over a much smaller footprint than a
         Bessel function, so the convolution calculation will be much
         faster. It is also more robust in the presence of spikes and
         seeing problems because it does not truncate the map FT as sharply
         as the Bessel function - such effects give rise to smoother
         defects in the rebinned map though the defects will still BE
         there.
      }
   }
   \sstinvocation{
       CALL SCULIB\_WTFN\_REGRID(USEGRD, N\_MAPS, N\_PTS, WTFNRAD, WTFNRES,
            WEIGHTSIZE, SCALE, DIAMETER, WAVELENGTH, PXSIZE, NX\_OUT, NY\_OUT,
            I\_CENTRE, J\_CENTRE, WTFN, WEIGHT, BOLWT, N\_BOL, MAX\_BOLS,
            DATA\_PTR, VAR\_PTR,
            XPOS\_PTR,YPOS\_PTR, OUT\_DATA, OUT\_VARIANCE,
            OUT\_QUALITY, CONV\_WEIGHT,
            STATUS)
   }
   \sstarguments{
      \sstsubsection{
         USEGRD = LOGICAL (Given)
      }{
         Use guard ring for final stage of rebin
      }
      \sstsubsection{
         N\_MAPS = INTEGER (Given)
      }{
         Number of data files read in
      }
      \sstsubsection{
         N\_PTS (N\_MAPS) = INTEGER (Given)
      }{
         Number of points in each input dataset
      }
      \sstsubsection{
         WTFNRAD = INTEGER (Given)
      }{
         Radius of largest weight function in scale lengths
      }
      \sstsubsection{
         WTFNRES = INTEGER (Given)
      }{
         Number of values per scale length in the supplied weight function
      }
      \sstsubsection{
         WEIGHTSIZE = INTEGER (Given)
      }{
         Radius of supplied weighting function
      }
      \sstsubsection{
         SCALE = REAL (Given)
      }{
         Size of a scale length in the same units as PXSIZE
      }
      \sstsubsection{
         DIAMETER = REAL (Given)
      }{
         Diameter of telescope in metres
      }
      \sstsubsection{
         WAVELENGTH = REAL (Given)
      }{
         Wavelength of map in microns
      }
      \sstsubsection{
         PXSIZE = REAL (Given)
      }{
         Pixel size in radians
      }
      \sstsubsection{
         NX\_OUT = INTEGER (Given)
      }{
         Number of pixels in X direction of output map
      }
      \sstsubsection{
         NY\_OUT = INTEGER (Given)
      }{
         Number of pixels in Y direction of output map
      }
      \sstsubsection{
         I\_CENTRE = INTEGER (Given)
      }{
         X reference pixel
      }
      \sstsubsection{
         J\_CENTRE = INTEGER (Given)
      }{
         Y reference pixel
      }
      \sstsubsection{
         WTFN() = REAL (Given)
      }{
         The weighting function
      }
      \sstsubsection{
         WEIGHT(N\_MAPS) = REAL (Given)
      }{
         The weight of each input dataset
      }
      \sstsubsection{
         BOLWT (MAX\_BOLS, N\_MAPS) = REAL (Given)
      }{
         Relative weight for each bolometer
      }
      \sstsubsection{
         N\_BOL (N\_MAPS) = INTEGER (Given)
      }{
         Number of bolometers from each map
      }
      \sstsubsection{
         MAX\_BOLS = INTEGER (Given)
      }{
         Max number of bolometers allowable in BOLWT
      }
      \sstsubsection{
         DATA\_PTR(N\_MAPS) = INTEGER (Given)
      }{
         Array of pointers to REAL input data
      }
      \sstsubsection{
         VAR\_PTR(N\_MAPS) = INTEGER (Given)
      }{
         Array of pointers to REAL input variance
      }
      \sstsubsection{
         XPOS\_PTR(N\_MAPS) = INTEGER (Given)
      }{
         Array of pointers to X bolometer positions (DOUBLE)
      }
      \sstsubsection{
         YPOS\_PTR(N\_MAPS) = INTEGER (Given)
      }{
         Array of pointers to Y bolometer positions (DOUBLE)
      }
      \sstsubsection{
         OUT\_DATA ( NX\_OUT, NY\_OUT ) = REAL (Returned)
      }{
         Output map
      }
      \sstsubsection{
         OUT\_VARIANCE ( NX\_OUT, NY\_OUT ) = REAL (Returned)
      }{
         Output variance
      }
      \sstsubsection{
         OUT\_QUALITY ( NX\_OUT, NY\_OUT ) = BYTE (Returned)
      }{
         Output quality
      }
      \sstsubsection{
         CONV\_WEIGHT ( NX\_OUT, NY\_OUT ) = REAL (Returned)
      }{
         Contribution to each pixel
      }
      \sstsubsection{
         STATUS = \_INTEGER (Given \& Returned)
      }{
         Global status
      }
   }
   \sstdiytopic{
      Copyright
   }{
      Copyright \copyright 1997,1998,1999 Particle Physics and Astronomy
      Research Council. All Rights Reserved.
   }
}
\sstroutine{
   SCULIB\_WTFN\_REGRID\_1
}{
   Calculate the coverage of the data and total weight per cell
}{
   \sstdescription{
      This routine performs two tasks:
    \sstitemlist{
      \sstitem
      [1] Determine which output pixels contain data (effectively
         detecting the edge of the useful area of the output image.

      \sstitem
      [2] For each output pixel containing data the weight of the
         map is added to the total weight. This TOTAL\_WEIGHT is then
         used to modify the behaviour when two images overlap (the
         TOTAL\_WEIGHT will be 2.0 (say) in the overlap region and
         1 elsewhere.
    }
   }
   \sstinvocation{
       CALL SCULIB\_WTFN\_REGRID\_1 (DIAMETER, WAVELENGTH, WEIGHT,
            IN\_DATA, X, Y, NPIX,
            PIXSPACE, NI, NJ, ICEN, JCEN, AV\_WEIGHT, SCRATCH, STATUS)
   }
   \sstarguments{
      \sstsubsection{
         DIAMETER = REAL (Given)
      }{
         Diameter of telescope
      }
      \sstsubsection{
         WAVELENGTH = REAL (Given)
      }{
         Wavelength of the observation
      }
      \sstsubsection{
         WEIGHT = REAL (Given)
      }{
         The weight of the input dataset.
      }
      \sstsubsection{
         IN\_DATA (NPIX)                   = REAL (Given)
      }{
         The input data values. (Used to define quality)
      }
      \sstsubsection{
         X( NPIX ) = DOUBLE PRECISION (Given)
      }{
         The x coordinates of the input pixels
      }
      \sstsubsection{
         Y( NPIX ) = DOUBLE PRECISION (Given)
      }{
         The y coordinates of the input pixels
      }
      \sstsubsection{
         NPIX = INTEGER (Given)
      }{
         the number of input pixels
      }
      \sstsubsection{
         PIXSPACE = DOUBLE PRECISION (Given)
      }{
         the pixel spacing of the output pixels
      }
      \sstsubsection{
         NI = INTEGER (Given)
      }{
         The number of output pixels in the x direction
      }
      \sstsubsection{
         NJ = INTEGER (Given)
      }{
         The number of output pixels in the y direction
      }
      \sstsubsection{
         ICEN = INTEGER (Given)
      }{
         the x index of the centre of the output array
      }
      \sstsubsection{
         JCEN = INTEGER (Given)
      }{
         the y index of the centre of the output array
      }
      \sstsubsection{
         AV\_WEIGHT (NI, NJ) = REAL (Given and Returned)
      }{
         Given as a workspace array
      }
      \sstsubsection{
         SCRATCH (NI, NJ) = INTEGER (Returned)
      }{
         Temporary workspace.
      }
      \sstsubsection{
         STATUS = INTEGER (Given and Returned)
      }{
         The global status.
      }
   }
   \sstdiytopic{
      Copyright
   }{
      Copyright \copyright 1995,1996,1997,1998,1999 Particle Physics and Astronomy
      Research Council. All Rights Reserved.
   }
}
\sstroutine{
   SCULIB\_WTFN\_REGRID\_2
}{
   Perform convolution
}{
   \sstdescription{
      This routine convolves the input data with a weighting function
      onto a regularly spaced output grid.
      The weighting function is stored in a lookup table indexed
      by the square of the distance and is parametrized by RES,
      WEIGHTSIZE, SCLSZ and WTFN.

      This is an image space implementation of fourier techniques.
      In Fourier terms the technique could be described
      as follows:-
      \sstitemlist{

         \sstitem
         [1] Do a 2-d discrete Fourier transform of the data points. The result
         is a repeating pattern made up of copies of the transform of
         the map as a continuous function, each copied displaced from its
         neighbours by 1/dx, where dx is the sample spacing of the input points
         (assumed equal in x and y). Different copies of the {\tt '}continuous map{\tt '}
         transforms will overlap ({\tt '}alias{\tt '}) if there is any power in the
         map at frequencies greater than 1/(2dx). It is not possible to unravel
         aliased spectra and this constraint leads to the Nyquist sampling
         criterion.

         \sstitem
         [2] We want to derive the {\tt '}continuous map{\tt '} so that map values on the new
         grid mesh can be derived. Do this by multiplying the transform of the
         data by a function that has zero value beyond a radius of 0.5/dx from the
         origin. This will get rid of all the repeats in the pattern and leave
         just the transform of the {\tt '}continuous map{\tt '}.

         \sstitem
         [3] Do an inverse FT on the remaining transform for the points where you
         wish the resampled points to be (note: FFTs implicitly assume that the
         data being transformed DO repeat ad infinitum so we{\tt '}d have to be careful
         when using them to do this).

      }
      The analogue of these process steps in image space is as follows:
      \sstitemlist{

         \sstitem
         [1] Nothing.

         \sstitem
         [2] Convolve the input data with the FT of the function used to isolate
         the `continuous map{\tt '} transform.

         \sstitem
         [3] Nothing.

      }
      If the method is done properly, the rebinned map is in fact the map on
      the new sample mesh that has the same FT as the continuous function
      going through the original sample points.

      Convolution Functions:-
      \sstitemlist{

         \sstitem
         [Bessel]
         For good data and with no time constraint on reduction the best
         convolution function would be one whose FT is a flat-topped cylinder in
         frequency space, centred on the origin and of a radius such that
         frequencies to which the telescope is not sensitive are set to zero.
         Unfortunately, this function is a Bessel function, which extends to
         infinity along both axes and has significant power out to a large radius.
         To work correctly this would require an infinite map and infinite
         computing time. However, a truncated Bessel function should work well on
         a large map, except near the edges. Edge effects can be removed by
         pretending that the map extends further - of course, this only works if
         you know what data the pretend map area should contain, i.e. zeros.
         Another problem with a Bessel function arises from the fact that it does
         truncate the FT of the map sharply. If the data are good then there
         should be nothing but noise power at the truncation radius and the
         truncation of the FT should have no serious effect. However, if the data
         has spikes (power at all frequencies in the FT) or suffers from seeing
         effects such that the data as measured DO have power beyond the
         truncation radius, then this will cause ringing in the rebinned map.

         \sstitem
         [Gaussian]
         In fact, any function that is finite in frequency space will have
         infinite extent in image space (I think). As such they will all
         drag in some power from the aliased versions of the map transform
         and all suffer from edge effects and large compute times. Some are
         worse than others, however. For example, a Gaussian can have most
         of its power concentrated over a much smaller footprint than a
         Bessel function, so the convolution calculation will be much
         faster. It is also more robust in the presence of spikes and
         seeing problems because it does not truncate the map FT as sharply
         as the Bessel function - such effects give rise to smoother
         defects in the rebinned map though the defects will still BE
         there.
      }
   }
   \sstinvocation{
       CALL SCULIB\_WTFN\_REGRID\_2 (RES, IN\_DATA, IN\_VARIANCE,
         WEIGHT, USEVARWT, VARWT, X, Y, NPIX, PIXSPACE, NI, NJ, ICEN, JCEN,
         TOTAL\_WEIGHT, WAVELENGTH, CONV\_DATA\_SUM, CONV\_VARIANCE\_SUM,
         CONV\_WEIGHT, WEIGHTSIZE, SCLSZ, WTFN, STATUS)
   }
   \sstarguments{
      \sstsubsection{
         RES                            = INTEGER (Given)
      }{
         number of resolution elements per scale length
      }
      \sstsubsection{
         IN\_DATA (NPIX)                   = REAL (Given)
      }{
         The input data values.
      }
      \sstsubsection{
         IN\_VARIANCE (NPIX)               = REAL (Given)
      }{
         Variance on IN\_DATA.
      }
      \sstsubsection{
         WEIGHT                           = REAL (Given)
      }{
         The weight of this dataset.
      }
      \sstsubsection{
         USEVARWT                         = LOGICAL (Given)
      }{
         Are we using a different weight for each point
      }
      \sstsubsection{
         VARWT (NPIX)                     = REAL (Given)
      }{
         The different weight for each pixel
      }
      \sstsubsection{
         X (NPIX)                         = DOUBLE PRECISION (Given)
      }{
         The x coordinates of the input pixels.
      }
      \sstsubsection{
         Y (NPIX)                         = DOUBLE PRECISION (Given)
      }{
         The y coordinates of the input pixels.
      }
      \sstsubsection{
         NPIX                             = INTEGER (Given)
      }{
         the number of input pixels.
      }
      \sstsubsection{
         PIXSPACE                         = REAL (Given)
      }{
         the pixel spacing of the output map.
      }
      \sstsubsection{
         NI                               = INTEGER (Given)
      }{
         The number of output pixels in the x direction.
      }
      \sstsubsection{
         NJ                               = INTEGER (Given)
      }{
         The number of output pixels in the y direction.
      }
      \sstsubsection{
         ICEN                             = INTEGER (Given)
      }{
         the x index of the centre of the output array.
      }
      \sstsubsection{
         JCEN                             = INTEGER (Given)
      }{
         the y index of the centre of the output array.
      }
      \sstsubsection{
         TOTAL\_WEIGHT (NI,NJ)             = REAL (Given)
      }{
         the `total weight{\tt '} of each output pixel.
      }
      \sstsubsection{
         WAVELENGTH                       = REAL (Given)
      }{
         the wavelength at which the maps were made (microns).
      }
      \sstsubsection{
         CONV\_DATA\_SUM (NI,NJ)            = REAL (Given and returned)
      }{
         the convolution sum for each output pixel.
      }
      \sstsubsection{
         CONV\_VARIANCE\_SUM (NI,NJ)        = REAL (Given and returned)
      }{
         the variance convolution sum for each output pixel.
      }
      \sstsubsection{
         CONV\_WEIGHT (NI,NJ)              = REAL (Given and returned)
      }{
         the convolution weight for each output pixel.
      }
      \sstsubsection{
         WEIGHTSIZE                       = INTEGER (Given)
      }{
         radius of weight function in scale units (SCUIP\_\_FILTRAD for BESSEL, 1 for LINEAR)
      }
      \sstsubsection{
         SCLSZ                            = REAL
      }{
         1 scale length in the same units as PIXSPACE
      }
      \sstsubsection{
         WTFN (RES $*$ RES $*$ WEIGHTSIZE $*$ WEIGHTSIZE) = REAL (Given)
      }{
         convolution weighting function
      }
      \sstsubsection{
         STATUS                           = INTEGER (Given and Returned)
      }{
         The global status.
      }
   }
   \sstdiytopic{
      Copyright
   }{
      Copyright \copyright 1995,1996,1997,1998,1999 Particle Physics and Astronomy
      Research Council. All Rights Reserved.
   }
}
\sstroutine{
   SCULIB\_WTFN\_REGRID\_3
}{
   Sets up the {\tt '}guard ring{\tt '} of bolometers outside the data
}{
   \sstdescription{
      Takes the regridded image and adds in data of value {\tt '}zero{\tt '}
      for any points that are within range of the weighting function
      (ie if TOTAL\_WEIGHT is greater than 0) but that do not already
      contain any data (ie TOTAL\_WEIGHT is set to something becuase
      regrid\_1 determined that the pixel was close enough to be
      influenced by the convolution function, BUT no data was actually
      found to lie in that pixel). This has the effect of enforcing
      zero flux density at the edges of the map or where no data
      was present (eg surrounding a bad pixel).
      This is only really necessary for the bessel function regrid since
      that function is extremely sensitive to edge effects.
      Once the {\tt '}guard ring{\tt '} has been processed, the output image is
      normalised by dividing by the convolution weight (ie the actual
      weight used for each pixel). Any point that has zero convolution
      weight is set to bad.
   }
   \sstinvocation{
       CALL SCULIB\_WTFN\_REGRID\_3 (METHOD,RES,PIXSPACE, NI, NJ, ICEN, JCEN,
         TOT\_WEIGHT\_IN, WAVELENGTH, CONV\_DATA\_SUM, CONV\_VARIANCE\_SUM,
         CONV\_QUALITY\_SUM, CONV\_WEIGHT, WEIGHTSIZE, SCLSZ, WTFN, STATUS)
   }
   \sstarguments{
      \sstsubsection{
         USEGUARD                       = LOGICAL (Given)
      }{
         Logical to determine whether the gueard ring should be used.
      }
      \sstsubsection{
         RES                            = INTEGER (Given)
      }{
         number of resolution elements per scale length
      }
      \sstsubsection{
         PIXSPACE                       = REAL (Given)
      }{
         the pixel spacing of the output pixels (radians).
      }
      \sstsubsection{
         NI                             = INTEGER (Given)
      }{
         The number of output pixels in the x direction.
      }
      \sstsubsection{
         NJ                             = INTEGER (Given)
      }{
         The number of output pixels in the y direction.
      }
      \sstsubsection{
         ICEN                           = INTEGER (Given)
      }{
         the x index of the centre of the output array.
      }
      \sstsubsection{
         JCEN                           = INTEGER (Given)
      }{
         the y index of the centre of the output array.
      }
      \sstsubsection{
         TOT\_WEIGHT\_IN (NI, NJ)         = REAL (Given)
      }{
         the `total weight{\tt '} of each output pixel.
      }
      \sstsubsection{
         WAVELENGTH                     = REAL (Given)
      }{
         the wavelength at which the maps was made (microns).
      }
      \sstsubsection{
         CONV\_DATA\_SUM (NI,NJ)          = REAL (Given and returned)
      }{
         the convolution sum for each output pixel.
      }
      \sstsubsection{
         CONV\_VARIANCE\_SUM (NI,NJ)      = REAL (Given and returned)
      }{
         the convolved variance sum for each output pixel.
      }
      \sstsubsection{
         CONV\_QUALITY\_SUM (NI,NJ)       = BYTE (Given)
      }{
         the quality on the convolution sum for each output pixel.
      }
      \sstsubsection{
         CONV\_WEIGHT (NI,NJ)            = REAL (Given and returned)
      }{
         the convolution weight for each output pixel.
      }
      \sstsubsection{
         WEIGHTSIZE                       = INTEGER (Given)
      }{
         radius of weight function in scale units (SCUIP\_\_FILTRAD for BESSEL, 1 for LINEAR)
      }
      \sstsubsection{
         SCLSZ                          = REAL (Given)
      }{
         1 scale length in the same units as pixspace
      }
      \sstsubsection{
         WTFN (RES $*$ RES $*$ WEIGHTSIZE $*$ WEIGHTSIZE) = REAL (Given)
      }{
         convolution weighting function
      }
      \sstsubsection{
         STATUS = INTEGER (Given and Returned)
      }{
         The global status.
      }
   }
   \sstdiytopic{
      Copyright
   }{
      Copyright \copyright 1995,1996,1997,1998,1999 Particle Physics and Astronomy
      Research Council. All Rights Reserved.
   }
}




% ? End of main text
\end{document}
